\begin{solution}{1.01}
\textbf{Method 1.} Take the point of projection as the origin of the coordinate system, and
take the horizontal and vertical directions as those of the $x$- and $y$-axes respectively.
The trajectory corresponding to the ball being thrown over the wall with the minimum initial
speed must be the parabola tangent to the top of the wall. When the ball reaches the top of
the wall,

\begin{align}
d &= v_0\cos\theta\, t, \tag{1}\\
h &= v_0\sin\theta\, t - \frac{1}{2}gt^2. \tag{2}
\end{align}

Eliminating $t$ from (1) and (2), we obtain

\begin{equation}
h = d\tan\theta - \frac{gd^2}{2v_0^2\cos^2\theta}. \tag{3}
\end{equation}

Equation (3) can be written as

\begin{equation}
v_0^2 = \frac{gd^2}{2\cos\theta\,(d\sin\theta - h\cos\theta)}. \tag{4}
\end{equation}

From (4) it is seen that to find the minimum of $v_0^2$ we need only find the maximum of the
denominator on the right-hand side. This denominator can be transformed into
\begin{equation*}
2\cos\theta\,(d\sin\theta-h\cos\theta)
= 2\cos\theta\sqrt{d^2+h^2}\,\sin(\theta-\alpha)
= \sqrt{d^2+h^2}\left[\sin(2\theta-\alpha)-\sin\alpha\right],
\end{equation*}
where $\alpha$ is the angle between the line joining the point of projection to the top of the
wall and the horizontal, as shown in the figure. From the above it is seen that when $\theta$
satisfies $2\theta-\alpha=\pi/2$, i.e. $\theta=\pi/4+\alpha/2$, the denominator of (4) is
maximum, that is, the value of $v_0^2$ is minimum. The minimum value of $v_0^2$ is
\begin{equation*}
v_{0\min}^2
= \frac{gd^2}{\sqrt{d^2+h^2}\,(1-\sin\alpha)}
= \frac{gd^2}{\sqrt{d^2+h^2}-h}
= g\left(\sqrt{d^2+h^2}+h\right),
\end{equation*}
i.e. when the angle of projection is $\theta=\dfrac{\pi}{4}+\dfrac{1}{2}\arctan\dfrac{h}{d}$,
the initial speed with which the ball is thrown over the wall is minimum, its minimum value
being $v_{0\min}=\sqrt{g\left(\sqrt{d^2+h^2}+h\right)}$.

\begin{figbox}{1.01}
\begin{tikzpicture}[scale=0.95]
  \draw[thickline] (-0.5,0) -- (6.7,0);
  \draw[axis] (0,0) -- (6.5,0) node[right] {$x$};
  \draw[axis] (0,0) -- (0,3.4) node[above] {$y$};
  \node[below left] at (0,0) {$O$};
  % wall
  \draw[thickline,fill=black!12] (4,0) rectangle (4.16,2);
  % trajectory
  \draw[dashedline] plot[domain=0:5.79,samples=90] (\x,{1.618*\x-0.2795*\x*\x});
  % line of sight to the top of the wall
  \draw[thinline,dashed] (0,0) -- (4,2);
  \draw (1.35,0) arc (0:26.57:1.35);
  \node at (1.62,0.20) {$\alpha$};
  \draw (0.72,0) arc (0:58.3:0.72);
  \node at (0.55,0.50) {$\theta$};
  % initial velocity
  \draw[force] (0,0) -- (1.15,1.86) node[above right,inner sep=1pt] {$v_0$};
  % velocity triangle at the top of the wall
  \draw[force,dashed] (4,2) -- (4.60,2.97) node[above right,inner sep=1pt] {$v_0$};
  \draw[force] (4.60,2.97) -- (4.60,1.63) node[midway,right] {$gt$};
  \draw[force] (4,2) -- (4.60,1.63) node[midway,below right,inner sep=1pt] {$v$};
  % dimensions
  \draw[dim] (0,-0.55) -- (4.16,-0.55) node[midway,below] {$d$};
  \draw[dim] (3.70,0) -- (3.70,2) node[midway,left] {$h$};
\end{tikzpicture}
\end{figbox}

\textbf{Method 2.} When the ball is thrown to the top of the wall, the relation between the
velocity $\vec{v}$ of the ball and the initial velocity $\vec{v}_0$ is
\begin{equation*}
\vec{v} = \vec{v}_0 + \vec{g}t.
\end{equation*}
A vector triangle satisfying this equation can be drawn, as shown in the figure. The area of
this triangle is
\begin{equation}
S = \frac{1}{2}gtv_0\cos\theta = \frac{1}{2}gd. \tag{5}
\end{equation}

The above shows that however large the initial speed $v_0$ and however large the angle of
projection $\theta$ with which the ball is thrown, as long as the ball passes tangent to the top
of the wall, the area of the velocity vector triangle at the point of tangency is constant.
Since $v=\sqrt{v_0^2-2gh}$, it is seen that the minimum $v_0$ corresponds to the minimum $v$.
From geometry, only when the directions of $\vec{v}$ and $\vec{v}_0$ are mutually perpendicular
are the required $v_0$ and $v$ minimum; at this instant
\begin{equation}
S = \frac{1}{2}v_0v = \frac{1}{2}v_0\sqrt{v_0^2-2gh}. \tag{6}
\end{equation}

From (5) and (6) we obtain
\begin{equation*}
v_0\sqrt{v_0^2-2gh} = gd,
\end{equation*}
i.e.
\begin{equation*}
(v_0^2)^2 - 2gh(v_0^2) - g^2d^2 = 0,
\end{equation*}
which can be solved to give
\begin{equation*}
v_{0\min}^2 = g\left(\sqrt{d^2+h^2}+h\right).
\end{equation*}

From the velocity vector triangle,
\begin{equation*}
v_0\cos\theta = v\sin\theta = \sqrt{v_0^2-2gh}\,\sin\theta.
\end{equation*}
Substituting $v_0^2=g\left(\sqrt{d^2+h^2}+h\right)$ into the above, we obtain
\begin{equation*}
\tan\theta = \frac{h+\sqrt{d^2+h^2}}{d},
\end{equation*}
that is,
\begin{equation*}
\theta = \frac{\pi}{4}+\frac{\alpha}{2}
= \frac{\pi}{4}+\frac{1}{2}\arctan\frac{h}{d}.
\end{equation*}

\begin{commentary}
What is the shape of the parabola corresponding to throwing the ball over the wall with the
minimum initial speed? Is its highest point at the top of the wall, or in front of the wall, or
behind it? From the second method of this problem it is known that at the top of the wall the
velocity $\vec{v}$ of the ball should be perpendicular to $\vec{v}_0$; hence the highest point
must be in front of the wall, and not in either of the other two positions.

The second method combines the physical problem with a geometric one and analyses the problem
from another point of view, thereby making the solution more concise.
\end{commentary}
\end{solution}

\begin{solution}{1.02}
The parabola corresponding to throwing the ball over the two walls $A$ and $B$ with the minimum
initial speed must be tangent to the top of wall $A$. Let the velocity of the ball when it
reaches the top of wall $A$ be $v$, making an angle $\alpha$ with the horizontal. This is
equivalent to the ball being projected from the top of wall $A$ with speed $v$ and angle of
projection $\alpha$. As long as it can be thrown to the top of wall $B$, that is, pass tangent
to the top of wall $B$, it can be thrown over the two walls.

From the range formula for projectile motion,
\begin{equation}
\frac{v^2\sin 2\alpha}{g} = d. \tag{1}
\end{equation}
From this it is seen that when $\alpha=\pi/4$ the required $v$ is minimum:
\begin{equation}
v_{\min} = \sqrt{gd}. \tag{2}
\end{equation}

From (2) and the conservation of mechanical energy, the minimum initial speed $v_{0\min}$
required to throw the ball over the two walls is
\begin{equation}
v_{0\min} = \sqrt{v_{\min}^2+2gh} = \sqrt{g(d+2h)}. \tag{3}
\end{equation}

\begin{figbox}{1.02}
\begin{tikzpicture}[scale=1.0]
  \draw[thickline] (-0.4,0) -- (4.6,0);
  \draw[thickline,fill=black!12] (1.2,0) rectangle (1.32,2);
  \draw[thickline,fill=black!12] (3.0,0) rectangle (3.12,2);
  \node[above left] at (1.32,2.02) {$A$};
  \node[above left] at (3.12,2.02) {$B$};
  \draw[dashedline] plot[domain=0:4.2,samples=110] (\x,{2.3333*\x-0.5556*\x*\x});
  \draw[force] (0,0) -- (0.44,1.03) node[above right,inner sep=1pt] {$v_0$};
  \draw (0.60,0) arc (0:66.8:0.60);
  \node at (0.36,0.46) {$\theta$};
  \draw[force] (1.2,2) -- (1.62,2.42) node[above right,inner sep=1pt] {$v$};
  \draw[dashedline] (1.32,2) -- (2.10,2);
  \draw (1.70,2) arc (0:45:0.50);
  \node at (1.83,2.14) {$\alpha$};
  \draw[dim] (0,-0.55) -- (1.2,-0.55) node[midway,below] {$l$};
  \draw[dim] (1.32,-0.55) -- (3.0,-0.55) node[midway,below] {$d$};
  \draw[dim] (0.95,0) -- (0.95,2) node[midway,left] {$h$};
\end{tikzpicture}
\end{figbox}

Let the distance from the point of projection to wall $A$ be $l$ and the angle of projection be
$\theta$; then
\begin{equation*}
v_0\cos\theta = v\cos\alpha,
\end{equation*}
hence
\begin{equation}
\cos\theta = \frac{v}{v_0}\cos\alpha. \tag{4}
\end{equation}

Substituting $v_0=v_{0\min}=\sqrt{g(d+2h)}$, $v=v_{\min}=\sqrt{gd}$ and $\alpha=\pi/4$ into (4),
we obtain
\begin{equation*}
\theta = \arccos\sqrt{\frac{d}{2(d+2h)}}.
\end{equation*}

In addition,
\begin{align}
v_0\cos\theta\, t &= l, \tag{5}\\
v_0\sin\theta - gt &= v\sin\alpha. \tag{6}
\end{align}

From (5) and (6), eliminating $t$ and substituting (2), (3), (4) and $\alpha=\pi/4$, we obtain
\begin{equation*}
l = \frac{1}{2}\left(\sqrt{d(d+4h)}-d\right).
\end{equation*}

\begin{commentary}
This problem looks very similar to Problem 1.01, but because the point of projection is not
restricted, the method of solution must be completely different. The key to solving this problem
is to discuss separately the motion of the ball between the two wall tops and its motion on the
two sides of the walls, and then to solve the problem from the connection between them.
\end{commentary}
\end{solution}

\begin{solution}{1.03}
Take point $O$ as the origin of the coordinate system, take the direction parallel to the
shelter as the direction of the $x$-axis, and the direction perpendicular to the shelter as the
direction of the $y$-axis, as shown in the figure. In this coordinate system the initial
velocity and the acceleration of the shell in the $y$ direction are respectively
\begin{equation*}
v_{0y} = v_0\sin(\theta-\alpha),\qquad a_y = -g\cos\alpha.
\end{equation*}

If the trajectory of the shell is tangent to the shelter, then the $y$-coordinate of the point
of tangency $A$ is
\begin{equation*}
y_A = h = l\sin\alpha = \frac{v_{0y}^2}{-2a_y}
= \frac{v_0^2\sin^2(\theta-\alpha)}{2g\cos\alpha},
\end{equation*}
that is,
\begin{equation}
v_0^2\sin^2(\theta-\alpha) = gl\sin 2\alpha. \tag{1}
\end{equation}

\begin{figbox}{1.03}
\begin{tikzpicture}[scale=1.0]
  \draw[thickline] (-3.4,0) -- (6.5,0);
  % shelter
  \draw[thickline] (-3,0) -- (4.5,3.50);
  \node[left] at (-3,0) {$P$};
  % trajectory
  \draw[dashedline] plot[domain=0:6.19,samples=110] (\x,{1.732*\x-0.28*\x*\x});
  \draw[force] (0,0) -- (0.75,1.30) node[above right,inner sep=1pt] {$v_0$};
  \fill (2.27,2.46) circle (1.4pt) node[above left,inner sep=1pt] {$A$};
  % axes through O
  \draw[axis] (0,0) -- (2.45,1.14) node[right] {$x$};
  \draw[axis] (0,0) -- (-1.10,2.35) node[above] {$y$};
  \fill (0,0) circle (1.4pt) node[below] {$O$};
  % angle alpha at P
  \draw (0.55,0) arc (0:25:0.55);
  \node at (0.85,0.18) {$\alpha$};
  % angle theta at O
  \draw[dashedline] (0,0) -- (1.45,0);
  \draw (0.92,0) arc (0:60:0.92);
  \node at (1.12,0.44) {$\theta$};
  % perpendicular distance h
  \draw[dashedline] (0,0) -- (-0.54,1.15);
  \node[left] at (-0.10,0.72) {$h$};
  \draw[dim] (-3,-0.55) -- (0,-0.55) node[midway,below] {$l$};
\end{tikzpicture}
\end{figbox}

From (1) the following discussion can be made.

(1) If $v_0<\sqrt{gl\sin 2\alpha}$, the trajectory of the shell cannot be tangent to the shelter,
so it is sufficient to take $\theta=\pi/4$, and the shell can then obtain the maximum range
\begin{equation*}
L_{\max} = \frac{v_0^2}{g}.
\end{equation*}

(2) If $v_0\geqslant\sqrt{gl\sin 2\alpha}$, it is known from (1) that when the firing angle of
the shell takes the following value, its trajectory is tangent to the shelter:
\begin{equation}
\theta = \alpha+\arcsin\frac{\sqrt{gl\sin 2\alpha}}{v_0}. \tag{2}
\end{equation}

(i) If the $\theta$ given by (2) satisfies $\theta\geqslant\pi/4$, i.e.
$v_0\leqslant\sqrt{\dfrac{2gl\sin 2\alpha}{1-\sin 2\alpha}}$, then it is sufficient to take
$\theta=\pi/4$, and the shell can obtain the maximum range
\begin{equation*}
L_{\max} = \frac{v_0^2}{g}.
\end{equation*}

(ii) If $\theta<\pi/4$, i.e.
$v_0>\sqrt{\dfrac{2gl\sin 2\alpha}{1-\sin 2\alpha}}$, then the firing angle of the shell is
taken as
\begin{equation*}
\theta = \alpha+\arcsin\frac{\sqrt{gl\sin 2\alpha}}{v_0},
\end{equation*}
and the shell can obtain the maximum range
\begin{equation*}
L_{\max}
= \frac{v_0^2\sin 2\left(\alpha+\arcsin\dfrac{\sqrt{gl\sin 2\alpha}}{v_0}\right)}{g}.
\end{equation*}

\begin{commentary}
This problem asks for the maximum range obtained by a shell fired with a fixed initial speed.
Without any restriction, and without considering the influence of air resistance, one need only
fire the shell at an elevation angle $\theta=\pi/4$ to obtain the maximum range. Because of the
existence of the shelter in this problem, it is extremely necessary to discuss whether the
shell will collide with the shelter during its flight, and to find the corresponding conditions.
Moreover, when finding the condition for the trajectory of the shell to be tangent to the
shelter, choosing the two directions parallel and perpendicular to this slope as the directions
of the $x$- and $y$-axes of the coordinate system makes the solution more concise.
\end{commentary}
\end{solution}

\begin{solution}{1.04}
Take the position of the lamp as the origin of the coordinate system, with the coordinate axes
as shown in the figure. If there were no restriction by the ceiling, the fragments of the burst
lamp would fly out with initial speed $v_0$ and elevation angle $\theta$, and the distance from
their landing point to the point $O_2$ on the ground directly below the lamp is denoted by $R$.

From the formulas of projectile motion,
\begin{align}
v_0\cos\theta\, t &= R, \tag{1}\\
v_0\sin\theta\, t - \frac{1}{2}gt^2 &= -h_2. \tag{2}
\end{align}

Squaring both sides of (1) and (2) respectively and then adding the two equations, we obtain
\begin{equation}
R^2+h_2^2 = v_0^2t^2-gv_0\sin\theta\,t^3+\frac{1}{4}g^2t^4. \tag{3}
\end{equation}

Substituting (2) into (3), we obtain
\begin{align}
R^2 &= (v_0^2+gh_2)t^2-\frac{1}{4}g^2t^4-h_2^2\notag\\
&= \left(\frac{v_0^2}{g}+h_2\right)^2-h_2^2
-\left[\frac{1}{2}gt^2-\left(\frac{v_0^2}{g}+h_2\right)\right]^2. \tag{4}
\end{align}

From (4) it is seen that when $t=\dfrac{1}{g}\sqrt{2(v_0^2+gh_2)}$, $R$ is maximum, and its
value is
\begin{equation}
R_{\max} = \frac{v_0}{g}\sqrt{v_0^2+2gh_2}. \tag{5}
\end{equation}

Substituting (5) and $t=\dfrac{1}{g}\sqrt{2(v_0^2+gh_2)}$ into (1), we obtain the angle of
projection $\theta_0$ corresponding to $R_{\max}$:
\begin{equation}
\cos\theta_0 = \sqrt{\frac{v_0^2+2gh_2}{2(v_0^2+gh_2)}}. \tag{6}
\end{equation}

\begin{figbox}{1.04}
\begin{tikzpicture}[scale=1.0]
  % ceiling and floor
  \draw[thickline] (-4.3,2.3) -- (4.3,2.3);
  \foreach \x in {-4.2,-3.9,...,4.2} {\draw[thinline] (\x,2.3) -- (\x+0.22,2.52);}
  \draw[thickline] (-4.3,-1.7) -- (4.3,-1.7);
  \foreach \x in {-4.2,-3.9,...,4.2} {\draw[thinline] (\x,-1.7) -- (\x+0.22,-1.92);}
  % axes
  \draw[axis] (0,-1.7) -- (0,2.95) node[above] {$y$};
  \draw[axis] (-4.3,0) -- (4.35,0) node[right] {$x$};
  \fill (0,0) circle (1.4pt) node[below right,inner sep=1pt] {$O$};
  \fill (0,2.3) circle (1.4pt) node[right] {$O_1$};
  \fill (0,-1.7) circle (1.4pt) node[right] {$O_2$};
  % trajectories to the right
  \draw[dashedline] plot[domain=0:1.13,samples=40] (\x,{2.6*\x-0.5*\x*\x});
  \draw[dashedline] plot[domain=1.13:2.85,samples=60]
     (\x,{2.3-1.47*(\x-1.13)-0.5*(\x-1.13)*(\x-1.13)});
  \draw[dashedline] plot[domain=0:3.6,samples=80] (\x,{1.33*\x-0.5*\x*\x});
  % trajectories to the left
  \draw[dashedline] plot[domain=-1.13:0,samples=40] (\x,{-2.6*\x-0.5*\x*\x});
  \draw[dashedline] plot[domain=-2.85:-1.13,samples=60]
     (\x,{2.3+1.47*(\x+1.13)-0.5*(\x+1.13)*(\x+1.13)});
  \draw[dashedline] plot[domain=-3.6:0,samples=80] (\x,{-1.33*\x-0.5*\x*\x});
  % dimensions
  \draw[dim] (-0.55,0) -- (-0.55,2.3) node[midway,left] {$h_1$};
  \draw[dim] (0.55,0) -- (0.55,-1.7) node[midway,right] {$h_2$};
\end{tikzpicture}
\end{figbox}

From (6) the maximum height to which the farthest-flying fragment can rise is
\begin{equation}
h = \frac{v_0^2\sin^2\theta_0}{2g}
= \frac{v_0^2(1-\cos^2\theta_0)}{2g}
= \frac{v_0^4}{4g(v_0^2+gh_2)}. \tag{7}
\end{equation}

Discussion:

(1) If $h_1\geqslant h$, i.e. the farthest-flying fragment does not strike the ceiling during its
flight, then the radius $r$ of the region in which the fragments land is just $R_{\max}$:
\begin{equation*}
r = R_{\max} = \frac{v_0}{g}\sqrt{v_0^2+2gh_2}.
\end{equation*}
The condition for this case can also be expressed as
\begin{equation*}
v_0^2\leqslant 2g\left(h_1+\sqrt{h_1(h_1+h_2)}\right).
\end{equation*}

(2) If $h_1<h$, i.e. $v_0^2>2g\left(h_1+\sqrt{h_1(h_1+h_2)}\right)$, then the farthest-flying
fragment will strike the ceiling during its flight and thereby change its trajectory. For this
case it can be proved that among all the fragments, the fragment whose trajectory is tangent to
the ceiling can fly farthest. The proof is as follows.

(i) From the law of conservation of mechanical energy, whether a fragment is tangent to the
ceiling or collides with it, its speed $v$ on reaching the ceiling is the same. After an elastic
collision with the ceiling, the fragment moves as a projectile thrown obliquely downward with
the original speed $v$, whereas the fragment tangent to the ceiling is equivalent to being
projected horizontally with the same speed $v$ at the point of tangency. From the knowledge of
projectile motion it is known that, at the same height, a body projected horizontally can fly
farther than a body projected obliquely downward in any direction.

(ii) Compared with the various fragments that collide with the ceiling, the fragment tangent to
the ceiling reaches the ceiling at a position farther from $O_1$, the point on the ceiling
directly above the lamp.

From the above two points it is known that the fragment tangent to the ceiling will fly farthest.

Let the angle of projection of the fragment tangent to the ceiling at the moment it flies out be
$\theta'$; then
\begin{equation*}
\frac{v_0^2\sin^2\theta'}{2g} = h_1,
\end{equation*}
i.e. $v_0\sin\theta'=\sqrt{2gh_1}$. Substituting this into (2), we obtain
\begin{equation*}
t^2-2\sqrt{\frac{2h_1}{g}}\,t-\frac{2h_2}{g}=0.
\end{equation*}
The solution of this equation is
\begin{equation*}
t = \sqrt{\frac{2h_1}{g}}\pm\sqrt{\frac{2(h_1+h_2)}{g}}.
\end{equation*}

Discarding the solution with the ``$-$'' sign and substituting the solution with the ``$+$''
sign into (1), we obtain the radius of the region in which the fragments land in this case:
\begin{equation*}
r = \frac{1}{g}\sqrt{v_0^2-2gh_1}
\left[\sqrt{2gh_1}+\sqrt{2g(h_1+h_2)}\right].
\end{equation*}
\end{solution}

\begin{solution}{1.05}
(1) The coordinate system $O$-$xyz$ is as shown in the figure, with the origin $O$ at the point
from which the bullet is fired, the wall being parallel to the $xOz$ plane and at distance $d$
from $O$.

Let the trajectory of a certain bullet lie in the $yOz$ plane, its angle of elevation be
$\theta$, and the aiming point on the wall be $M$; the bullet hits point $P$ on the wall at
height $h$ above the ground. For this bullet,
\begin{align}
v_0\sin\theta\,t-\frac{1}{2}gt^2 &= h, \tag{1}\\
v_0\cos\theta\,t &= d. \tag{2}
\end{align}

From (1) and (2), eliminating $t$, we obtain
\begin{equation*}
h = d\tan\theta-\frac{gd^2}{2v_0^2\cos^2\theta},
\end{equation*}
i.e.
\begin{equation}
\tan\theta = \frac{v_0^2\pm\sqrt{v_0^4-2ghv_0^2-g^2d^2}}{gd}. \tag{3}
\end{equation}

Through $O$ draw another coordinate axis $y'$, as shown in the figure. This axis intersects the
wall at $K$; take $KO'=l$, then $KO=\sqrt{d^2+l^2}$. Suppose that the trajectory of another
bullet lies in the $y'Oz$ plane, its angle of elevation is $\theta'$, the aiming point on the
wall is $N$, and the bullet hits point $Q$ on the wall at the same height $h$ above the ground.
For this bullet a relation similar to (3) is obtained:
\begin{equation}
\tan\theta' = \frac{v_0^2\pm\sqrt{v_0^4-2ghv_0^2-g^2(d^2+l^2)}}{g\sqrt{d^2+l^2}}. \tag{4}
\end{equation}

Through $O'$ take another plane coordinate system $O'x'z'$, with $O'K$ as the $x'$-axis and
$O'M$ as the $z'$-axis. In this coordinate system the coordinates of point $N$ are denoted by
$(x',z')$. From the geometric relations shown in the figure,
\begin{equation*}
x' = l,\qquad z' = \sqrt{d^2+l^2}\,\tan\theta'.
\end{equation*}
Substituting (4), we obtain
\begin{equation*}
z' = \frac{v_0^2}{g}\pm\frac{\sqrt{v_0^4-2ghv_0^2-g^2d^2-g^2x'^2}}{g}.
\end{equation*}
From this it follows that
\begin{equation}
\left(z'-\frac{v_0^2}{g}\right)^2+x'^2
= \left(\frac{\sqrt{v_0^4-2ghv_0^2-g^2d^2}}{g}\right)^2. \tag{5}
\end{equation}

Equation (5) is the equation, in the $O'x'z'$ coordinate system, of a circle whose centre is at
the point $\left(0,\dfrac{v_0^2}{g}\right)$ and whose radius is
$R=\dfrac{1}{g}\sqrt{v_0^4-2ghv_0^2-g^2d^2}$; this circle is precisely the locus of the aiming
points for the bullet to hit the horizontal line at height $h$ on the wall.

\begin{figbox}{1.05}
\begin{tikzpicture}[scale=1.0]
  % wall
  \fill[black!10] (4.95,0.822) -- (3.775,0.023) -- (3.775,2.423) -- (4.95,3.222) -- cycle;
  \draw[thickline] (4.95,0.822) -- (3.775,0.023) -- (3.775,2.423) -- (4.95,3.222) -- cycle;
  % axes
  \draw[axis] (1.25,0.55) -- (1.25,3.15) node[above] {$z$};
  \draw[axis] (1.25,0.55) -- (5.15,0.55) node[right] {$y$};
  \draw[axis] (1.25,0.55) -- (0.55,0.07) node[below left] {$x$};
  \fill (1.25,0.55) circle (1.5pt) node[left] {$O$};
  % O'
  \fill (4.55,0.55) circle (1.4pt) node[right] {$O'$};
  \draw[dashedline] (1.25,0.55) -- (4.55,0.55);
  % y' axis through K
  \draw[dashedline] (1.25,0.55) -- (4.75,0.33);
  \fill (4.275,0.363) circle (1.4pt) node[below right,inner sep=1pt] {$K$};
  % aiming points and impact points
  \fill (4.55,2.50) circle (1.4pt) node[right] {$M$};
  \fill (4.55,1.60) circle (1.4pt) node[right] {$P$};
  \fill (4.275,2.313) circle (1.4pt) node[left] {$N$};
  \fill (4.275,1.413) circle (1.4pt) node[left] {$Q$};
  \draw[dashedline] (4.55,2.50) -- (1.25,0.55);
  \draw[dashedline] (4.275,2.313) -- (1.25,0.55);
  \draw[thinline] (4.55,1.60) -- (4.275,1.413);
  % dimensions
  \draw[dim] (4.275,0.363) -- (4.55,0.55) node[midway,below,inner sep=1pt] {$l$};
  \draw[dim] (1.25,0.15) -- (4.55,0.15) node[midway,below] {$d$};
  \draw[dim] (4.72,0.55) -- (4.72,1.60) node[midway,right] {$h$};
\end{tikzpicture}
\end{figbox}

(2) If the bullet is to hit a point on the wall above the horizontal line of height $h$, it is
only necessary to substitute $h'\ (>h)$ for $h$ in (5). It is known that the equation then
becomes a circle of radius $R'\ (<R)$. Hence the aiming points lie inside the circle shown in
(5).

\begin{commentary}
This problem involves three-dimensional space; choosing a suitable coordinate system is the key
to solving it.
\end{commentary}
\end{solution}

\begin{solution}{1.06}
(1) After the bomb is released it moves as a horizontally projected body, with the speed $v_0$ of the aeroplane as its initial speed; hence
\begin{align}
h&=\frac{1}{2}gt^2 \tag{1}\\
L&=v_0t \tag{2}
\end{align}
Combining (1) and (2), we obtain the horizontal distance $L$ between the point at which the bomb is released and the target:
\[
L=v_0\sqrt{\frac{2h}{g}}.
\]

(2) Take the instant at which the bomb is released as $t=0$. At the time $t$ the velocity $\vec{v}_1$ of the bomb is
\begin{equation}
\vec{v}_1=\vec{v}_0+\vec{g}t \tag{3}
\end{equation}
and at this instant the velocity $\vec{v}_2$ of a shell fired from the gun on the ground with initial speed $\vec{v}$ is
\begin{equation}
\vec{v}_2=\vec{v}+\vec{g}t. \tag{4}
\end{equation}
Take the bomb as the frame of reference; in this frame the relative velocity $\vec{v}'$ of the shell is
\begin{equation}
\vec{v}'=\vec{v}_2-\vec{v}_1=\vec{v}-\vec{v}_0. \tag{5}
\end{equation}
The vector relation (5) is shown in Fig.~1.06(a).

\begin{figbox}{1.06(a)}
\begin{tikzpicture}[scale=0.8,>=latex]
  % ground
  \draw[thickline] (-1.0,0)--(9.6,0);
  \fill[hatch] (-1.0,-0.22) rectangle (9.6,0);
  % aeroplane
  \draw[thickline,fill=black!20] (0.10,3.72)--(0.95,3.72)--(1.22,3.60)--(0.95,3.48)--(0.10,3.48)--cycle;
  \draw[thickline,fill=black!20] (0.35,3.72)--(0.75,3.72)--(0.55,4.05)--cycle;
  \draw[thickline,fill=black!20] (0.40,3.48)--(0.85,3.48)--(0.62,3.15)--cycle;
  \draw[force] (1.35,3.60)--(2.85,3.60) node[right] {$\vec{v}_0$};
  % release vertical and height h
  \draw[dashedline] (0,0)--(0,3.60);
  \draw[dim] (-0.75,0)--(-0.75,3.60) node[midway,left] {$h$};
  % relative displacement line: release point -> gun
  \draw[dashedline] (0,3.60)--(5.3,0);
  \draw[thickline] (0.95,3.60) arc (0:-34.2:0.95);
  \node at (1.30,3.32) {$\alpha$};
  % trajectory of the bomb
  \draw[dashedline,domain=0:8,samples=60] plot (\x,{3.60-0.05625*\x*\x});
  % gun
  \draw[thickline,fill=black!30] (5.05,0)--(5.30,0.35)--(5.55,0)--cycle;
  \node[below] at (5.15,-0.30) {gun};
  % target
  \fill[black!70] (7.85,0) rectangle (8.15,0.42);
  \node[above left] at (8.05,0.42) {target};
  % velocity triangle at the point of impact
  \draw[force] (6.60,1.15)--(8.60,1.15) node[midway,above] {$\vec{v}_0$};
  \draw[force] (8.60,1.15)--(7.49,1.90) node[above right] {$\vec{v}'$};
  \draw[force] (6.60,1.15)--(7.49,1.90) node[left] {$\vec{v}$};
  \draw[thickline] (7.40,1.15) arc (0:40.2:0.80);
  \node at (7.62,1.38) {$\theta$};
  % dimensions
  \draw[dim] (0,-0.85)--(8,-0.85) node[midway,below] {$L$};
  \draw[dim] (5.3,-1.25)--(8,-1.25) node[midway,below] {$d$};
\end{tikzpicture}
\end{figbox}

Suppose that after a time $t_1$ the shell hits the bomb in the air; then
\begin{equation}
v't_1=\sqrt{(L-d)^2+h^2}. \tag{6}
\end{equation}
Since the point at which the shell hits the bomb must be in the air, while the longest time for which the bomb can fly in the air is $t=\sqrt{\dfrac{2h}{g}}$, it follows that necessarily
\begin{equation}
t_1\leq t=\sqrt{\frac{2h}{g}}. \tag{7}
\end{equation}
Combining (6) and (7) gives
\begin{equation}
v'\geq\sqrt{(L-d)^2+h^2}\,\sqrt{\frac{g}{2h}}. \tag{8}
\end{equation}
From (8) it is seen that the minimum relative speed required for the shell to hit the bomb in the air is
\[
v_m'=\sqrt{(L-d)^2+h^2}\,\sqrt{\frac{g}{2h}}.
\]
To find the minimum initial speed of the shell relative to the ground, the following discussion may be made.

(i) If $v_0\cos\alpha<v_m'=\sqrt{(L-d)^2+h^2}\sqrt{\dfrac{g}{2h}}$, where $\cos\alpha=\dfrac{L-d}{\sqrt{(L-d)^2+h^2}}$, this condition can be expressed as
\[
v_0<\frac{d^2+h^2}{d}\sqrt{\frac{g}{2h}}.
\]
Then from the vector triangle of Fig.~1.06(a) it is seen that for the value of $v$ to be a minimum the value of $v'$ must be as small as possible; hence we take $v_m'$ and obtain
\begin{align}
v_{\min}^2&=v_m'^2+v_0^2-2v_m'v_0\cos\alpha\notag\\
&=(L-d)^2\frac{g}{2h}+\frac{1}{2}gh+v_0^2-2v_0(L-d)\sqrt{\frac{g}{2h}}\notag
\end{align}
Substituting $v_0=L\sqrt{\dfrac{g}{2h}}$ into the above expression, we obtain the minimum launching speed of the shell in this case:
\[
v_{\min}=\sqrt{\frac{g}{2h}(d^2+h^2)}.
\]
From the geometrical relations shown in the figure, the launching angle $\theta$ of the shell corresponding to $v_{\min}$ satisfies
\[
v_{\min}\sin\theta=v_m'\sin\alpha.
\]
Hence
\[
\sin\theta=\frac{v_m'}{v_{\min}}\sin\alpha=\frac{h}{\sqrt{d^2+h^2}},
\]
or, equivalently,
\[
\tan\theta=\frac{h}{d},\qquad\text{i.e.}\qquad \theta=\arctan\frac{h}{d}.
\]

(ii) If $v_0\cos\alpha\geq v_m'$, the condition for this case can also be expressed as
\[
v_0\geq\frac{d^2+h^2}{d}\sqrt{\frac{g}{2h}}.
\]
The velocity vector triangle for this case is shown in Fig.~1.06(a). From the figure it is seen that we need only make $v\perp v'$ for $v$ to attain a minimum; note that in this case $v'>v_m'$. Hence the minimum initial speed $v_{\min}$ is
\[
v_{\min}=v_0\sin\alpha=\frac{v_0h}{\sqrt{(L-d)^2+h^2}}
=\frac{v_0h}{\sqrt{\left(v_0\sqrt{\dfrac{2h}{g}}-d\right)^2+h^2}}.
\]
The launching angle $\theta$ corresponding to this minimum initial speed is
\[
\theta=\frac{\pi}{2}-\alpha,\qquad\text{or}\qquad
\theta=\arctan\left(\frac{v_0\sqrt{\dfrac{2h}{g}}-d}{h}\right).
\]
The discussion above is all for the case $d<L$, i.e. the case shown in Fig.~1.06(a). If $d\geq L$, the velocity vector triangle is as shown in Fig.~1.06(b). From the geometrical relations shown in the figure it is seen that for $v$ to be a minimum, $v'$ should be as small as possible, so we take $v'=v_m'$, and the law of cosines then gives
\[
v_{\min}=\sqrt{\frac{g}{2h}(h^2+d^2)}.
\]
This value is the same as in case (i), and the corresponding launching angle is also the same as in case (i), that is,
\[
\theta=\arctan\left(\frac{h}{d}\right).
\]

\begin{figbox}{1.06(b)}
\begin{tikzpicture}[scale=1.25,>=latex]
  \coordinate (A) at (0,0);
  \coordinate (B) at (1.60,0);
  \coordinate (C) at (3.14,2.20);
  \draw[force] (A)--(B) node[midway,below] {$\vec{v}_0$};
  \draw[force] (A)--(C) node[below right] {$\vec{v}$};
  \draw[dashedline] (B)--(C) node[midway,above right] {$\vec{v}'$};
  \draw[dashedline] (B)--(2.95,0);
  \draw[thickline] ($(B)+(0.55,0)$) arc (0:55:0.55);
  \node at ($(B)+(0.78,0.42)$) {$\alpha$};
  \draw[thickline] (0.60,0) arc (0:35:0.60);
  \node at (0.78,0.20) {$\theta$};
\end{tikzpicture}
\end{figbox}
\end{solution}

\begin{solution}{1.07}
First find the velocity $v_P$ with which the position of point $P$ changes.

Take the circle as the frame of reference. In this frame the velocity of the straight line in the direction perpendicular to the line is $v_2$, and the direction of the velocity $v_P'$ of point $P$ (its relative velocity) is along the tangent to the circle. As a point on the straight line, $P$ must satisfy
\begin{equation}
v_P'\cos\theta=v_2 \tag{1}
\end{equation}
From the relative motion formula, the velocity $v_P$ with which the position of point $P$ changes relative to the plane of the paper (its absolute velocity) is
\begin{equation}
\vec{v}_P=\vec{v}_P'+\vec{v}_1 \tag{2}
\end{equation}
The vector relation is shown in Fig.~1.07(a). From the geometrical relations shown in the figure we obtain
\begin{align}
v_P&=\sqrt{(v_1-v_P'\sin\theta)^2+v_2^2}\notag\\
&=\sqrt{(v_1-v_2\tan\theta)^2+v_2^2}\notag
\end{align}

\begin{figbox}{1.07(a)}
\begin{tikzpicture}[scale=1.0,>=latex]
  % circle
  \draw[thickline] (0,0) circle (1.5);
  \fill (0,0) circle (1.3pt) node[below left] {$O$};
  % straight line
  \draw[thickline] (-2.0,1.0)--(2.45,1.0);
  \coordinate (P) at (1.118,1.0);
  \fill (P) circle (1.4pt) node[below right] {$P$};
  % radius
  \draw[thickline] (0,0)--(P) node[midway,below right] {$R$};
  % angle theta at P
  \draw[thickline] ($(P)+(-0.55,0)$) arc (180:221.8:0.55);
  \node at ($(P)+(-0.86,-0.22)$) {$\theta$};
  % velocity of the straight line
  \draw[force] (-1.90,1.15)--(-1.90,2.15) node[left] {$\vec{v}_2$};
  % vector triangle at P
  \coordinate (Q) at (0.224,1.998);
  \coordinate (R) at (1.824,1.998);
  \draw[force] (P)--(Q) node[above left] {$\vec{v}_P'$};
  \draw[force] (Q)--(R) node[midway,above] {$\vec{v}_1$};
  \draw[force] (P)--(R) node[above right] {$\vec{v}_P$};
  \draw[thickline] ($(P)+(0.40,0)$) arc (0:54.7:0.40);
  \node at ($(P)+(0.58,0.20)$) {$\alpha$};
  % velocity of the circle
  \draw[force] (2.55,1.0)--(3.35,1.0) node[right] {$v_1$};
\end{tikzpicture}
\end{figbox}

The angle $\alpha$ between the direction of $v_P$ and the straight line is
\[
\cot\alpha=\frac{v_1-v_2\tan\theta}{v_2}=\frac{v_1}{v_2}-\tan\theta,
\]
that is,
\[
\alpha=\operatorname{arccot}\left(\frac{v_1}{v_2}-\tan\theta\right).
\]

Next find the acceleration with which the position of point $P$ changes.

First take the circle as the frame of reference. In this frame, as a point on the circle, $P$ moves in a circle, and indeed in a circle with varying speed, so its relative acceleration $\vec{a}_P'$ can be expressed as
\[
\vec{a}_P'=\vec{a}_n'+\vec{a}_t'.
\]
In the above expression $\vec{a}_n'$ and $\vec{a}_t'$ are the normal and the tangential acceleration respectively. The value of the normal acceleration is
\begin{equation}
a_n'=\frac{v_P'^2}{R}=\frac{v_2^2}{R\cos^2\theta} \tag{3}
\end{equation}
Since the circle moves uniformly relative to the plane of the paper, from the theory of relative motion it is known that the acceleration $\vec{a}_P$ of point $P$ relative to the plane of the paper is equal to $\vec{a}_P'$, that is,
\begin{equation}
\vec{a}_P=\vec{a}_P'=\vec{a}_n'+\vec{a}_t' \tag{4}
\end{equation}
Next take the straight line as the frame of reference. In this frame, as a point on the straight line, $P$ moves along a straight line, so the direction of its relative acceleration $\vec{a}_P''$ is along this straight line. Since the straight line also moves uniformly relative to the plane of the paper, the theory of relative motion likewise gives
\begin{equation}
\vec{a}_P=\vec{a}_P'' \tag{5}
\end{equation}
that is, relative to the plane of the paper the direction of the acceleration $\vec{a}_P$ of point $P$ is along the direction of the straight line. Fig.~1.07(b) shows the vector relation in (4); the value of $a_P$ can then be found:
\begin{equation}
a_P=\frac{a_n'}{\cos\theta} \tag{6}
\end{equation}

\begin{figbox}{1.07(b)}
\begin{tikzpicture}[scale=1.0,>=latex]
  \draw[thickline] (0,0) circle (1.5);
  \fill (0,0) circle (1.3pt) node[below left] {$O$};
  \draw[thickline] (-2.0,1.0)--(2.2,1.0);
  \coordinate (P) at (1.118,1.0);
  \fill (P) circle (1.4pt) node[below right] {$P$};
  \draw[dashedline] (0,0)--(P);
  \draw[thickline] ($(P)+(-0.42,0)$) arc (180:221.8:0.42);
  \node at ($(P)+(-0.66,-0.20)$) {$\theta$};
  % normal acceleration: from P towards O
  \draw[force] (P)--++(221.8:0.70) node[below left] {$\vec{a}_n'$};
  % tangential acceleration: along the tangent, up-left
  \draw[force] (0.60,0.53)--(0.18,1.00) node[above left] {$\vec{a}_t'$};
  % acceleration of P: along the straight line
  \draw[force] (P)--(0.18,1.00) node[above] {$\vec{a}_P$};
\end{tikzpicture}
\end{figbox}

Substituting (3) into (6) gives
\[
a_P=\frac{v_2^2}{R\cos^3\theta}.
\]

\begin{commentary}
The acceleration of the point $P$ in this problem can also be found by first differentiating $v_P'$ with respect to time to obtain $a_t'$, and then using $a_P=a_P'=\sqrt{a_n'^2+a_t'^2}$ to find $a_P$. The preceding method takes the circle and then the straight line as frames of reference in turn and finds by analysis that the direction of $\vec{a}_P$ is along the straight line. Hence one need only find the value of $a_n'$ to obtain the value of $a_P$. This method of solution avoids differentiation and is therefore more convenient.
\end{commentary}
\end{solution}

\begin{solution}{1.08}
(1) Take the point of contact $A$ of the rod with the ground as the frame of reference. In this frame the rod rotates about $A$, so the direction of the velocity $\vec{v}_P'$ (the relative velocity) of the point $P$ of the rod that is in contact with the cylinder is perpendicular to the rod, that is, along the radius of the cylinder towards its centre $O$. Let the angular velocity with which the rod rotates about $A$ at this instant be $\omega$; then
\begin{equation}
v_P'=\omega l \tag{1}
\end{equation}
where $l$ is the length of the rod from the point $P$ to the point $A$. Owing to the constraint by the fixed semi-cylindrical body, the absolute velocity $\vec{v}_P$ of the point $P$ of the rod can have no component perpendicular to the rod and directed obliquely downwards (i.e. towards the point $O$), and of course it can have no component perpendicular to the rod and directed obliquely upwards either, since otherwise the rod would necessarily leave the cylindrical surface; hence at this instant the direction of $\vec{v}_P$ must be along the rod towards the point $A$. From the relative motion formula we obtain
\[
\vec{v}_P=\vec{v}_P'+\vec{v}.
\]
Fig.~1.08(a) shows the vector relation of the above expression. From the geometrical relations shown in the figure we obtain
\begin{equation}
v_P=v\cos\theta \tag{2}
\end{equation}
and at the same time
\begin{equation}
v_P'=v\sin\theta. \tag{3}
\end{equation}

\begin{figbox}{1.08(a)}
\begin{tikzpicture}[scale=1.0,>=latex]
  \draw[thickline] (-1.55,0)--(2.60,0);
  \fill[hatch] (-1.55,-0.20) rectangle (2.60,0);
  \draw[thickline] (1.2,0) arc (0:180:1.2);
  \fill (0,0) circle (1.3pt) node[below left] {$O$};
  \draw[thickline] (0,0)--(0,1.2) node[midway,left] {$R$};
  \draw[thickline] (0,0)--(0.771,0.919);
  \draw[thickline] (1.867,0)--(0.082,1.498);
  \fill (0.771,0.919) circle (1.4pt) node[above right] {$P$};
  \node[below right] at (1.867,0) {$A$};
  \node[above left] at (0.082,1.498) {$B$};
  \draw[thickline] (0,0.50) arc (90:50:0.50);
  \node at (0.23,0.66) {$\theta$};
  \draw[force] (0.771,0.919)--(0.358,0.426) node[midway,below left] {$\vec{v}_P'$};
  \draw[force] (0.358,0.426)--(1.358,0.426) node[midway,below] {$\vec{v}$};
  \draw[force] (0.771,0.919)--(1.358,0.426) node[midway,above right] {$\vec{v}_P$};
  \draw[force] (1.90,0.05)--(2.50,0.05) node[right] {$v$};
\end{tikzpicture}
\end{figbox}

(2) Combining (1) and (3), and noting that $l=R\cot\theta$, we obtain
\[
v\sin\theta=l\omega=R\cot\theta\,\omega,
\]
and hence
\begin{equation}
\omega=\frac{v\sin^2\theta}{R\cos\theta}. \tag{4}
\end{equation}
This $\omega$ is the angular velocity with which the rod rotates about $A$ in the frame of reference attached to $A$. Since this frame translates, $\omega$ is also the angular velocity of the rod relative to the frame of reference of the ground. From Fig.~1.08(b) it is seen that when the rod turns through an angle $\Delta\theta$, the point of intersection $P'$ of the rod and the cylindrical surface also turns through an angle $\Delta\theta$ in its circular motion about $O$, that is, the angular velocity of this circular motion is also $\omega$. Hence
\begin{equation}
v_{P'}=\omega R. \tag{5}
\end{equation}

\begin{figbox}{1.08(b)}
\begin{tikzpicture}[scale=1.0,>=latex]
  \draw[thickline] (-1.55,0)--(2.30,0);
  \fill[hatch] (-1.55,-0.20) rectangle (2.30,0);
  \draw[thickline] (1.2,0) arc (0:180:1.2);
  \fill (0,0) circle (1.3pt) node[below left] {$O$};
  % earlier position of the rod
  \draw[dashedline] (1.867,0)--(0.350,1.273);
  \draw[dashedline] (0,0)--(0.771,0.919);
  % later position of the rod
  \draw[thickline] (1.566,0)--(0.530,1.234);
  \draw[thickline] (0,0)--(0.919,0.771);
  \fill (0.771,0.919) circle (1.3pt) node[above right] {$P$};
  \fill (0.919,0.771) circle (1.3pt) node[right] {$P'$};
  \draw[thickline] (50:0.85) arc (50:40:0.85);
  \node at (0.74,0.78) {$\Delta\theta$};
  \draw[thickline] (0.581,1.175) arc (130:140:0.42);
  \node[fill=white,inner sep=1pt] at (0.53,1.31) {$\Delta\theta$};
\end{tikzpicture}
\end{figbox}

Substituting (4) into (5) gives
\[
v_{P'}=\frac{v\sin^2\theta}{\cos\theta}.
\]
\end{solution}

\begin{solution}{1.09}
(1) When hinge $O_3$ moves along the $x$-axis with speed $v$, the velocities of hinges $O_1$ and
$O_2$ are respectively
\begin{equation*}
v_{O_1} = \frac{1}{3}v,\qquad v_{O_2} = \frac{2}{3}v,
\end{equation*}
and the components along the $x$ direction of the velocities of hinges $A_1$ and $B_3$ are
respectively
\begin{equation*}
v_{A_1,x} = \frac{1}{6}v,\qquad v_{B_3,x} = \frac{5}{6}v.
\end{equation*}

Since point $A_1$ moves in a circle about $O$, the direction of $v_{A_1}$ is perpendicular to the
rod $OA_1$. From the geometric relations it is known that the angle between the direction of
$v_{A_1}$ and the $x$-axis is $\dfrac{\pi}{2}-\theta$. Hence
\begin{equation}
v_{A_1} = \frac{v_{A_1,x}}{\cos\left(\dfrac{\pi}{2}-\theta\right)}
= \frac{v}{6\sin\theta}. \tag{1}
\end{equation}

At the same time the value of the component of the velocity of point $A_1$ in the $y$ direction
can be found:
\begin{equation*}
v_{A_1,y} = v_{A_1}\cos\theta = \frac{v}{6}\cot\theta.
\end{equation*}
Since the value of the component in the $y$ direction of the velocity of point $B_3$ is the same
as that of point $A_1$, i.e.
\begin{equation*}
v_{B_3,y} = v_{A_1,y} = \frac{v}{6}\cot\theta,
\end{equation*}
we therefore obtain
\begin{equation}
v_{B_3} = \sqrt{v_{B_3,x}^2+v_{B_3,y}^2}
= \frac{v}{6}\sqrt{25+\cot^2\theta}. \tag{2}
\end{equation}

In this problem, except for hinge $O$, all the other hinges move uniformly along the $x$
direction. Therefore the accelerations of $A_1$ and $B_3$ are along the $y$-axis; let their
values be $a_{A_1}$ and $a_{B_3}$ respectively. The normal acceleration of hinge $A_1$ in its
circular motion is
\begin{equation}
a_{A_1,n} = \frac{v_{A_1}^2}{l} = \frac{v^2}{36l\sin^2\theta}, \tag{3}
\end{equation}
hence
\begin{equation}
a_{A_1} = \frac{a_{A_1,n}}{\sin\theta} = \frac{v^2}{36l\sin^3\theta}. \tag{4}
\end{equation}
The value of the acceleration of hinge $B_3$ is the same as that of $A_1$, i.e.
\begin{equation*}
a_{B_3} = a_{A_1} = \frac{v^2}{36l\sin^3\theta}.
\end{equation*}

\begin{figbox}{1.09}
\begin{tikzpicture}[scale=0.95]
  \def\lx{0.643}\def\ly{0.766}
  % wall
  \draw[thickline] (0,-1.25) -- (0,1.25);
  \foreach \y in {-1.15,-0.85,...,1.15} {\draw[thinline] (0,\y) -- (-0.18,\y+0.16);}
  % rods
  \draw[thickline] (0,0) -- (\lx,\ly) -- (2*\lx,0);
  \draw[thickline] (0,0) -- (\lx,-\ly) -- (2*\lx,0);
  \draw[thickline] (2*\lx,0) -- (3*\lx,\ly) -- (4*\lx,0);
  \draw[thickline] (2*\lx,0) -- (3*\lx,-\ly) -- (4*\lx,0);
  \draw[thickline] (4*\lx,0) -- (5*\lx,\ly) -- (6*\lx,0);
  \draw[thickline] (4*\lx,0) -- (5*\lx,-\ly) -- (6*\lx,0);
  % x axis
  \draw[axis] (0,0) -- (4.85,0) node[right] {$x$};
  % hinges
  \foreach \p in {(0,0),(\lx,\ly),(2*\lx,0),(\lx,-\ly),(3*\lx,\ly),(4*\lx,0),(3*\lx,-\ly),(5*\lx,\ly),(6*\lx,0),(5*\lx,-\ly)}
     \fill \p circle (1.4pt);
  \node[below left] at (0,0) {$O$};
  \node[below left] at (2*\lx,0) {$O_1$};
  \node[below left] at (4*\lx,0) {$O_2$};
  \node[below right] at (6*\lx,0) {$O_3$};
  \node[above left] at (\lx,\ly) {$A_1$};
  \node[above left] at (3*\lx,\ly) {$A_2$};
  \node[above left] at (5*\lx,\ly) {$A_3$};
  \node[below left] at (\lx,-\ly) {$B_1$};
  \node[below left] at (3*\lx,-\ly) {$B_2$};
  \node[below left] at (5*\lx,-\ly) {$B_3$};
  % angle theta
  \draw (0.55,0) arc (0:50:0.55);
  \node at (0.62,0.28) {$\theta$};
  % velocity decomposition at A1
  \draw[force] (\lx,\ly) -- ++(0.65,0) node[below] {$v_{A_1,x}$};
  \draw[force] (\lx,\ly) -- ++(0,-0.55) node[right] {$v_{A_1,y}$};
  \draw[force] (\lx,\ly) -- ++(0.65,-0.545) node[below right,inner sep=1pt] {$v_{A_1}$};
  \draw (1.05,0.766) arc (0:-40:0.41);
  \node at (1.15,0.62) {$\frac{\pi}{2}-\theta$};
  % velocity of O3
  \draw[force] (6*\lx,0) -- ++(0.62,0) node[below] {$v$};
\end{tikzpicture}
\end{figbox}

(2) Since the angular speeds and the angular accelerations of all the rods are the same, it is
only necessary to find the angular speed $\omega$ and the angular acceleration $\beta$ of rod
$OA_1$:
\begin{equation}
\omega = \frac{v_{A_1}}{l} = \frac{v}{6l\sin\theta}. \tag{5}
\end{equation}
From (3) and (4) the tangential acceleration of $A_1$ can be found:
\begin{equation*}
a_{A_1,t} = \sqrt{a_{A_1}^2-a_{A_1,n}^2}
= \frac{v^2\cos\theta}{36l\sin^3\theta}.
\end{equation*}
From the tangential acceleration, $\beta$ can be found:
\begin{equation}
\beta = \frac{a_{A_1,t}}{l} = \frac{v^2\cos\theta}{36l^2\sin^3\theta}. \tag{6}
\end{equation}
Equations (5) and (6) are the expressions for the angular speed and the angular acceleration of
rod $A_3B_2$.

(3) Hinges $O_1$, $O_2$ and $O_3$ move along the $x$-axis, and their trajectory equations can be
expressed as
\begin{equation*}
y = 0.
\end{equation*}
Hinges $A_1$, $A_2$, $A_3$, $B_1$, $B_2$, $B_3$ can be denoted by the symbols $A_n$, $B_n$
($n=1,2,3$), and their position coordinates are
\begin{align*}
x_n &= (2n-1)l\cos\theta,\\
y_n &= \pm l\sin\theta.
\end{align*}
Hence the trajectory equations are
\begin{equation*}
\frac{x^2}{(2n-1)^2l^2}+\frac{y^2}{l^2} = 1\qquad (n=1,2,3),
\end{equation*}
which are equations of ellipses.
\end{solution}

\begin{solution}{1.10}
Take the disc $A$ as the frame of reference. In this uniformly rotating frame the disc $B$ rolls without slipping along the rim of $A$, and the angular velocity with which $B$ turns about its own centre $O_2$ is
\[
\omega_2'=\omega_2-\omega_1.
\]
Hence the relative velocity $v_{O_2}'$ of the point $O_2$ is
\[
v_{O_2}'=r\omega_2'=r(\omega_2-\omega_1).
\]
From the theory of relative motion it is known that the absolute velocity $v_{O_2}$ of the point $O_2$ is related to the relative velocity $v_{O_2}'$ by
\begin{equation}
\vec{v}_{O_2}=\vec{v}_{O_2}'+\vec{\omega}_1\times\vec{r}_{O_1O_2}, \tag{1}
\end{equation}
where $\vec{r}_{O_1O_2}$ is the vector from $O_1$ to $O_2$; hence
\begin{align}
v_{O_2}&=r(\omega_2-\omega_1)+(R+r)\omega_1\notag\\
&=R\omega_1+r\omega_2.\notag
\end{align}
From the absolute velocity of the point $O_2$ the time $t$ required for the centre $O_2$ of the disc $B$ to make one complete circuit about the centre $O_1$ of the disc $A$ is found to be
\[
t=\frac{2\pi(R+r)}{v_{O_2}}=\frac{2\pi(R+r)}{R\omega_1+r\omega_2}.
\]

\begin{figbox}{1.10}
\begin{tikzpicture}[scale=1.1,>=latex]
  \coordinate (O1) at (0,0);
  \coordinate (O2) at (1.645,0.95);
  \draw[thickline] (O1) circle (1.3);
  \draw[thickline] (O2) circle (0.6);
  \draw[dashedline] (O1)--++(0,1.85);
  \draw[thickline] (O1)--(O2);
  \draw[thickline] (0,0.75) arc (90:60:0.75);
  \node at (0.30,0.90) {$\theta$};
  \fill (O1) circle (1.3pt) node[below left] {$O_1$};
  \fill (O2) circle (1.3pt) node[above right] {$O_2$};
  \node[below right] at (1.126,0.65) {$P$};
  \node at (-0.62,0.62) {$A$};
  \node at (1.80,1.62) {$B$};
  \draw[thickline] (O1)--++(140:1.3) node[midway,above left] {$R$};
  \draw[thickline] (O2)--++(40:0.6) node[midway,above right] {$r$};
  \draw[force] (O1) ++(200:0.42) arc (200:290:0.42);
  \node at (0.52,-0.62) {$\omega_1$};
  \draw[force] (O2) ++(20:0.30) arc (20:-70:0.30);
  \node at ($(O2)+(0.55,-0.12)$) {$\omega_2$};
\end{tikzpicture}
\end{figbox}

\begin{commentary}
This problem concerns relative motion in a frame of reference that rotates uniformly. Equation (1) is the velocity transformation formula, in which $\vec{\omega}_1\times\vec{r}_{O_1O_2}$ is the transport velocity. This expression shows that the transport velocity depends on the position of the particle. This is different from the velocity transformation for relative motion in a frame of reference that translates, where the transport velocity is independent of the position of the particle.
\end{commentary}
\end{solution}

\begin{solution}{1.11}
Take the large steel ring as the frame of reference. In this uniformly rotating frame, let the angular velocity of the small steel ring be $\omega'$ when it rolls without slipping along the inner side of the large ring; then the relative velocity of its centre $C$ is
\begin{equation}
v_C'=\frac{R}{3}\omega'. \tag{1}
\end{equation}

\begin{figbox}{1.11}
\begin{tikzpicture}[scale=1.0,>=latex]
  \def\R{2.0}
  \coordinate (O) at (0,0);
  \coordinate (P) at (-\R,0);
  \coordinate (A) at (0,-\R);
  \coordinate (B) at (0,-\R/3);
  \coordinate (C) at (0,-2*\R/3);
  \draw[thickline] (O) circle (\R);
  \draw[thickline] (C) circle (\R/3);
  \draw[axis] (-\R-0.45,0)--(\R+0.62,0) node[right] {$x$};
  \draw[axis] (0,-\R-0.30)--(0,\R+0.55) node[above] {$y$};
  \draw[dashedline] (P)--(O) node[midway,above] {$R$};
  \draw[dashedline] (C)--(A) node[midway,right] {$\frac{R}{3}$};
  \fill (O) circle (1.3pt) node[above right] {$O$};
  \fill (C) circle (1.3pt) node[right] {$C$};
  \fill (P) circle (1.3pt) node[left] {$P$};
  \node[below] at (A) {$A$};
  \node[right] at (B) {$B$};
  % rotation of the large ring about P
  \draw[force] (P) ++(55:0.62) arc (55:160:0.62);
  \node at ($(P)+(105:0.85)$) {$\omega$};
  % rotation of the small ring
  \draw[force] (C) ++(20:0.38) arc (20:-70:0.38);
  \node at ($(C)+(0.58,-0.33)$) {$\omega'$};
  % velocity of the centre C
  \draw[force] (C)--++(0.85,0) node[above] {$v_C'$};
\end{tikzpicture}
\end{figbox}

According to the conditions given in the problem: during the time $t'$ in which the small steel ring turns once relative to the large steel ring, the large steel ring turns through half a revolution about the point $P$; hence
\[
t'=\frac{2\pi\times\dfrac{2}{3}R}{v_C'}=\frac{4\pi}{\omega'}=\frac{1}{2}\cdot\frac{2\pi}{\omega}.
\]
From the above expression we obtain
\begin{equation}
\omega'=4\omega \tag{2}
\end{equation}
Substituting (2) into (1) gives
\begin{equation}
v_C'=\frac{4}{3}\omega R. \tag{3}
\end{equation}
In order to find the relative accelerations of the points $A$ and $B$ of the small steel ring, we may take the large steel ring as a frame of reference at rest and take the centre $C$ of the small steel ring as a moving (instantaneously translating) frame of reference; in this frame the points $A$ and $B$ both move uniformly in circles, and the value of their relative acceleration is
\[
a_A''=a_B''=\omega'^2\frac{R}{3}=\frac{16}{3}\omega^2R.
\]
The transport acceleration is the acceleration $\vec{a}_C'$ of the point $C$ relative to the large steel ring:
\[
a_C'=\frac{v_C'^2}{\dfrac{2}{3}R}=\frac{8}{3}\omega^2R.
\]
Hence the accelerations of the points $A$ and $B$ relative to the large steel ring are respectively
\[
\vec{a}_A'=\vec{a}_A''+\vec{a}_C'=(a_A''+a_C')\hat{\jmath}=8\omega^2R\,\hat{\jmath},
\]
\[
\vec{a}_B'=\vec{a}_B''+\vec{a}_C'=(a_C'-a_B'')\hat{\jmath}=-\frac{8}{3}\omega^2R\,\hat{\jmath}.
\]
Having found the accelerations of the points $A$ and $B$ relative to the large steel ring, we can use the relative motion formula to find the absolute accelerations of the points $A$ and $B$ respectively:
\begin{align}
\vec{a}_A&=\vec{a}_A'+\vec{\omega}\times(\vec{\omega}\times\vec{r}_A)\notag\\
&=\vec{a}_A'+\vec{\omega}\times[\vec{\omega}\times(R\hat{\imath}-R\hat{\jmath})]\notag\\
&=8\omega^2R\,\hat{\jmath}+\omega^2R\,\hat{\jmath}-\omega^2R\,\hat{\imath}\notag\\
&=-\omega^2R\,\hat{\imath}+9\omega^2R\,\hat{\jmath}\notag\\
&=\omega^2R(-\hat{\imath}+9\hat{\jmath}) \tag{4}
\end{align}
\begin{equation}
\vec{a}_B=\vec{a}_B'+\vec{\omega}\times(\vec{\omega}\times\vec{r}_B)+2\vec{\omega}\times\vec{v}_B' \tag{5}
\end{equation}
In the above expression $\vec{v}_B'$ is the velocity of the point $B$ relative to the large steel ring:
\begin{equation}
\vec{v}_B'=2v_C'=\frac{8}{3}\omega R\,\hat{\imath}. \tag{6}
\end{equation}
Substituting (6) into (5) gives
\begin{align}
\vec{a}_B&=\vec{a}_B'+\vec{\omega}\times\left[\vec{\omega}\times\left(R\hat{\imath}-\frac{R}{3}\hat{\jmath}\right)\right]+2\vec{\omega}\times\frac{8}{3}\omega R\,\hat{\imath}\notag\\
&=-\frac{8}{3}\omega^2R\,\hat{\jmath}+\frac{\omega^2R}{3}\hat{\jmath}-\omega^2R\,\hat{\imath}+\frac{16}{3}\omega^2R\,\hat{\jmath}\notag\\
&=\omega^2R(-\hat{\imath}+3\hat{\jmath}).\notag
\end{align}

\begin{commentary}
Equations (4) and (5) are respectively the relations between the absolute accelerations of the points $A$ and $B$ and their relative accelerations in the uniformly rotating frame of reference; here $\vec{\omega}\times(\vec{\omega}\times\vec{r}_A)$ and $\vec{\omega}\times(\vec{\omega}\times\vec{r}_B)$ are the transport accelerations. Unlike the point $A$, which is at rest in the rotating frame of reference, the point $B$ moves in this frame with velocity $\vec{v}_B'$. Therefore, compared with (4), equation (5) contains one more term $2\vec{\omega}\times\vec{v}_B'$, and this term is the Coriolis acceleration.
\end{commentary}
\end{solution}

\begin{solution}{1.12}
(1) Let the velocity of the centre $O$ of the small cylinder relative to the ground be $v_0$, and
its angular speed of rotation be $\omega$. Since there is no relative slipping between the small
cylinder and cylinders $A$ and $B$ at the points of contact, take $O$ as the moving (translating)
frame of reference. From the relative-motion velocity formula the absolute velocities $v_C$ and
$v_D$ of points $C$ and $D$ on the small cylinder can be obtained, and $v_D$ and $v_C$ are
precisely the velocities of the points of contact of the walls of cylinders $A$ and $B$ with the
small cylinder, so
\begin{align}
v_C &= v_0+\frac{R}{2}\omega = 2R\omega_2, \tag{1}\\
v_D &= -v_0+\frac{R}{2}\omega = R\omega_1. \tag{2}
\end{align}

From (1) and (2) we can solve
\begin{align}
\omega &= \omega_1+2\omega_2, \tag{3}\\
v_0 &= \frac{R}{2}\left(2\omega_2-\omega_1\right). \tag{4}
\end{align}

Since the centre $O$ of the small cylinder moves in a circle of radius $\dfrac{3}{2}R$ about the
axis $O'$ of cylinders $A$ and $B$, the time $t$ required for it to make one circuit about $O'$
relative to the ground is
\begin{equation*}
t = \frac{2\pi\times\dfrac{3}{2}R}{v_0} = \frac{6\pi}{2\omega_2-\omega_1}.
\end{equation*}

Take cylinder $B$ as the moving (uniformly rotating) frame of reference; let the velocity of
point $O$ in this frame be $v_0'$. From the relative-motion velocity formula,
\begin{equation}
v_0' = v_0-\omega_2\times r_0', \tag{5}
\end{equation}
where $r_0'$ is the vector from $O'$ to $O$, $|r_0'|=\dfrac{3}{2}R$. Substituting into (5), we
obtain
\begin{align}
v_0' &= \frac{R}{2}\left(2\omega_2-\omega_1\right)-\omega_2\times\frac{3}{2}R\notag\\
&= -\frac{\omega_1+\omega_2}{2}R. \tag{6}
\end{align}

Hence the time $t'$ for the small cylinder to make one circuit relative to cylinder $B$ is
\begin{equation*}
t' = \frac{2\pi\times\dfrac{3}{2}R}{|v_0'|} = \frac{6\pi}{\omega_1+\omega_2}.
\end{equation*}

(2) First find the acceleration $a_C$ of point $C$ of the small cylinder relative to the ground.

Take $O$ as the moving (translating) frame of reference. Since $O$ moves uniformly in a circle
about $O'$, the value of the entrainment acceleration $a_0$ is
\begin{equation}
a_0 = \frac{v_0^2}{\dfrac{3}{2}R} = \frac{R}{6}\left(2\omega_2-\omega_1\right)^2, \tag{7}
\end{equation}
and since point $C$ moves uniformly in a circle relative to $O$, the relative acceleration is
\begin{equation}
a_C' = \omega^2\cdot\frac{R}{2} = \frac{R}{2}\left(\omega_1+2\omega_2\right)^2. \tag{8}
\end{equation}

From the relative-motion formula,
\begin{equation*}
a_C = a_C'+a_0.
\end{equation*}
From the above and (7), (8), we obtain
\begin{align}
a_C &= a_C'+a_0
= \frac{R}{2}\left(\omega_1+2\omega_2\right)^2+\frac{R}{6}\left(2\omega_2-\omega_1\right)^2
\notag\\
&= \frac{2}{3}R\left[\left(\omega_1+\omega_2\right)^2+3\omega_2^2\right]. \tag{9}
\end{align}
The direction of $a_C$ is from $C$ towards $O'$.

\begin{figbox}{1.12}
\begin{tikzpicture}[scale=1.15]
  \def\R{1.5}
  \coordinate (Op) at (0,0);
  \coordinate (O) at (0,-1.5*\R);
  \draw[thickline] (Op) circle (\R);
  \draw[thickline] (Op) circle (2*\R);
  \draw[thickline] (O) circle (0.5*\R);
  \draw[thinline] (Op) -- (35:\R);
  \draw[thinline] (Op) -- (0,2*\R);
  \node[above right] at (35:0.75*\R) {$R$};
  \node[above] at (1.05,0.12) {$2R$};
  \fill (Op) circle (1.3pt) node[above right] {$O'$};
  \fill (O) circle (1.3pt);
  \fill (0,-\R) circle (1.3pt) node[right] {$D$};
  \fill (0,-2*\R) circle (1.3pt) node[below] {$C$};
  \node[right] at (O) {$O$};
  \node at (-0.95,0.35) {$A$};
  \node at (-2.35,0.25) {$B$};
  % v0
  \draw[force] (O) -- ++(-0.75,0) node[above] {$v_0$};
  % omega of the small cylinder
  \draw[force] (0.42,-1.5*\R) arc (0:-120:0.42);
  \node at (0.78,-2.52) {$\omega$};
  % omega1 of cylinder A
  \draw[force] (200:0.95) arc (200:340:0.95);
  \node at (0,-1.05) {$\omega_1$};
  % omega2 of cylinder B
  \draw[force] ({3.25*cos(-20)},{3.25*sin(-20)}) arc (-20:-60:3.25);
  \node at (2.78,-2.45) {$\omega_2$};
\end{tikzpicture}
\end{figbox}

Next find the acceleration $a_C''$ of point $C$ relative to cylinder $A$.

\textbf{Method 1.}

Take cylinder $A$ as the moving (rotating) frame of reference; in this frame let the relative
velocity of point $O$ be $v_0''$; then from the relative-motion velocity formula
\begin{equation*}
v_0'' = v_0-\omega_1\times r'.
\end{equation*}
Hence
\begin{align}
v_0'' &= \frac{R}{2}\left(2\omega_2-\omega_1\right)
-\left(-\omega_1\times\frac{3}{2}R\right)\notag\\
&= \left(\omega_1+\omega_2\right)R. \tag{10}
\end{align}

In this frame the angular speed of rotation of cylinder $B$ is $\omega_1+\omega_2$. Let the
angular speed of the small cylinder about $O$ be $\omega'$; then the relative velocity $v_C''$ of
point $C$ is
\begin{equation}
v_C'' = \left(\omega_1+\omega_2\right)\cdot 2R = v_0''+\omega'\frac{R}{2}. \tag{11}
\end{equation}

Substituting (10) into (11), we can solve
\begin{equation}
\omega' = 2\left(\omega_1+\omega_2\right). \tag{12}
\end{equation}

Take cylinder $A$ as a fixed frame and take point $O$ as the moving (translating) frame; the
relative acceleration of point $C$ is $\omega'\times\left(\omega'\times r_C''\right)$, so its
absolute acceleration $a_C''$ relative to cylinder $A$ is
\begin{equation}
a_C'' = a_0''+\omega'\times\left(\omega'\times r_C''\right). \tag{13}
\end{equation}
In (13), $r_C''$ is the vector from $O$ to $C$, $|r_C''|=\dfrac{R}{2}$, and
$a_0''=\dfrac{v_0''^2}{\dfrac{3}{2}R}=\dfrac{2}{3}\left(\omega_1+\omega_2\right)^2R$.
Substituting this and (12) into (13), we obtain
\begin{align*}
a_C'' &= \frac{2}{3}\left(\omega_1+\omega_2\right)^2R
+4\left(\omega_1+\omega_2\right)^2\times\frac{R}{2}\\
&= \frac{8}{3}\left(\omega_1+\omega_2\right)^2R.
\end{align*}
The direction of $a_C''$ is from $C$ towards $O'$.

\textbf{Method 2.}

Take cylinder $A$ as the moving (rotating) frame of reference, and find $a_C''$ directly from
the relative-acceleration formula for a rotating frame of reference:
\begin{equation}
a_C'' = a_C-\omega_1\times\left(\omega_1\times r_C'\right)-2\omega_1\times v_C''. \tag{14}
\end{equation}

In (14), the value of $v_C''$ can be found by substituting (10) and (12) into (11):
\begin{align}
v_C'' &= v_0''+\omega'\frac{R}{2}
= \left(\omega_1+\omega_2\right)R+2\left(\omega_1+\omega_2\right)\times\frac{R}{2}\notag\\
&= 2\left(\omega_1+\omega_2\right)R. \tag{15}
\end{align}

Substituting (9) and (15) into (14), we obtain
\begin{align*}
a_C'' &= \frac{2}{3}R\left[\left(\omega_1+\omega_2\right)^2+3\omega_2^2\right]
-\omega_1^2\times 2R+2\omega_1\times 2\left(\omega_1+\omega_2\right)R\\
&= \frac{8}{3}\left(\omega_1+\omega_2\right)^2R.
\end{align*}
\end{solution}

\begin{solution}{1.13}
When the direction of the relative velocity $\vec{v}'=\vec{v}_A-\vec{v}_B$ of the two particles $A$ and $B$ is parallel to the line joining the two particles, the vector relation between the velocities $\vec{v}_A$, $\vec{v}_B$ of $A$ and $B$ and $\vec{v}'$ is as shown in the figure. Applying the sine rule to $\triangle AOB$, we obtain
\[
\frac{a}{\sin\phi}=\frac{b}{\sin(\theta+\phi)}.
\]
The above expression may also be written as
\begin{equation}
b\sin\phi=a\sin(\theta+\phi). \tag{1}
\end{equation}

\begin{figbox}{1.13}
\begin{tikzpicture}[scale=1.30,>=latex]
  \def\a{1.30}
  \def\b{2.60}
  \coordinate (O) at (0,0);
  \coordinate (A) at (36.87:\a);
  \coordinate (B) at (\b,0);
  \draw[thickline] (O) circle (\a);
  \draw[thickline] (O) circle (\b);
  \draw[thinline] (O)--(A);
  \draw[thinline] (O)--(B);
  \draw[dashedline] (A)--(B);
  \draw[thinline] (0.42,0) arc (0:36.87:0.42);
  \node at (0.56,0.14) {$\theta$};
  \draw[thinline] ($(B)+(-0.52,0)$) arc (180:153.43:0.52);
  \node at (1.90,0.20) {$\phi$};
  \node[below left] at (O) {$O$};
  \node[above left] at (A) {$A$};
  \node[right] at (B) {$B$};
  \node[above right] at (0.60,0.42) {$a$};
  \node[below] at (1.35,0) {$b$};
  % tangent direction at B
  \draw[dashedline] (\b,-0.60)--(\b,1.55);
  \draw[force] (B)--++(90:0.80) node[right] {$\vec{v}_B$};
  % vector triangle at A: v_A = v_B + v'
  \coordinate (P) at ($(A)+(0,0.85)$);
  \coordinate (Q) at ($(A)+(-1.02,1.36)$);
  \draw[force] (A)--(Q) node[midway,below left] {$\vec{v}_A$};
  \draw[force] (A)--(P) node[midway,right] {$\vec{v}_B$};
  \draw[force] (P)--(Q) node[midway,above] {$\vec{v}'$};
  \draw[thinline] ($(Q)+(0.36,-0.48)$) arc (-53.13:-26.57:0.60);
  \node at (0.62,1.70) {$\alpha$};
\end{tikzpicture}
\end{figbox}

Applying the sine rule to the vector triangle formed by $\vec{v}_A$, $\vec{v}_B$ and $\vec{v}'$, we obtain
\begin{equation}
\frac{v_A}{\sin[\pi-(\theta+\alpha)]}=\frac{v_B}{\sin\alpha}. \tag{2}
\end{equation}
Since at this instant the vector line of $\vec{v}'$ is parallel to the line joining $A$ and $B$, we have
\[
\pi-(\theta+\alpha)=\pi-\left(\frac{\pi}{2}-\phi\right),
\]
hence
\begin{equation}
\alpha=\frac{\pi}{2}-(\theta+\phi). \tag{3}
\end{equation}
Substituting (3) into (2), we obtain
\begin{equation}
\frac{v_A}{\cos\phi}=\frac{v_B}{\cos(\theta+\phi)}. \tag{4}
\end{equation}
From the condition given in the problem, that the speeds of the two particles are inversely proportional to the radii, we have
\begin{equation}
\frac{v_A}{v_B}=\frac{b}{a}. \tag{5}
\end{equation}
Substituting (5) into (4), we obtain
\begin{equation}
a\cos\phi=b\cos(\theta+\phi). \tag{6}
\end{equation}
Adding the two sides of (1) and (6), we obtain
\[
b\sin\phi+a\cos\phi=a\sin(\theta+\phi)+b\cos(\theta+\phi).
\]
Multiplying both sides of the above expression by $\dfrac{1}{\sqrt{a^2+b^2}}$ and setting $\tan\beta=\dfrac{b}{a}$, the above expression becomes
\begin{equation}
\cos(\beta-\phi)=\sin(\theta+\phi+\beta). \tag{7}
\end{equation}
From (7) we obtain
\[
\theta+\phi+\beta=\frac{\pi}{2}-\beta+\phi,
\]
that is,
\[
\theta=\frac{\pi}{2}-2\beta.
\]
Therefore
\[
\cos\theta=\sin2\beta=\frac{2ab}{a^2+b^2}.
\]
Hence $\theta$ is
\[
\theta=\arccos\frac{2ab}{a^2+b^2}.
\]
\end{solution}

\begin{solution}{1.14}
Let the distance of the small ring from the end $A$ of the rod be $x$, as shown in the figure. In the rotating reference frame that moves together with the rod, among the external forces acting on the small ring two have components along the rod: the weight $mg$ and the centrifugal force of inertia $f_c=m\omega^2x\sin\alpha$. Suppose that when the small ring moves to the position $x=x_0$, the resultant of the forces acting on it along the rod is zero, that is,
\[
mg\cos\alpha=m\omega^2x_0\sin^2\alpha.
\]
Solving, we obtain
\begin{equation}
x_0=\frac{g\cos\alpha}{\omega^2\sin^2\alpha}. \tag{1}
\end{equation}

\begin{figbox}{1.14}
\begin{tikzpicture}[scale=1.15,>=latex]
  \def\al{24}
  \coordinate (A) at (0,0.30);
  \coordinate (B) at ({0.30+3.3*sin(\al)},{0.30+3.3*cos(\al)});
  \coordinate (M) at ({0.30+1.85*sin(\al)},{0.30+1.85*cos(\al)});
  \draw[thinline] (0,-0.15)--(0,4.25);
  \draw[force] (0,3.95) ++(150:0.32) arc (150:-150:0.32);
  \node[right] at (0.30,3.95) {$\omega$};
  \draw[thickline] (A)--(B);
  \node[above right] at (B) {$B$};
  \node[below left] at (A) {$A$};
  \draw[thinline] (0,0.98) arc (90:{90-\al}:0.68);
  \node at (0.20,1.12) {$\alpha$};
  % ring
  \draw[thickline,fill=white] (M) circle (0.09);
  % forces on the ring
  \draw[force] (M)--++(0:1.05) node[right] {$f_c$};
  \draw[force] (M)--++(-90:0.90) node[below] {$mg$};
  % distance x along the rod
  \draw[dim] ($(A)+(-0.16,0.07)$)--($(M)+(-0.16,0.07)$) node[midway,left] {$x$};
  \draw[thinline] ($(A)+(-0.22,0.08)$)--($(A)+(-0.10,0.03)$);
  \draw[thinline] ($(M)+(-0.22,0.08)$)--($(M)+(-0.10,0.03)$);
\end{tikzpicture}
\end{figbox}

(1) If $x_0<l$, that is, $\omega>\sqrt{\dfrac{g\cos\alpha}{l\sin^2\alpha}}$, then in this case, when the small ring moves in the region $x<x_0$, the resultant external force along the rod is directed towards $A$, while when it moves in the region $x>x_0$, the resultant external force along the rod is directed towards $B$. It is seen that the small ring need only reach the point $x=x_0$, and it can then be accelerated to the end $B$. In the rotating reference frame, the equation of motion of the small ring is
\begin{equation}
m\frac{\dd v}{\dd t}=m\omega^2x\sin^2\alpha-mg\cos\alpha. \tag{2}
\end{equation}
Equation (2) can be written as
\[
\frac{\dd v}{\dd x}\frac{\dd x}{\dd t}=\omega^2\sin^2\alpha\,x-g\cos\alpha,
\]
that is,
\[
v\,\dd v=\left(\omega^2\sin^2\alpha\,x-g\cos\alpha\right)\dd x.
\]
Integrating both sides between definite limits, we obtain
\begin{equation}
\int_{v_0}^{0}v\,\dd v=\int_{0}^{x_0}\left(\omega^2\sin^2\alpha\,x-g\cos\alpha\right)\dd x. \tag{3}
\end{equation}
Carrying out the integration, we obtain
\begin{equation}
-\frac{1}{2}v_0^2=\frac{1}{2}\omega^2\sin^2\alpha\,x_0^2-g\cos\alpha\,x_0. \tag{4}
\end{equation}
Substituting (1) into (4), we obtain
\[
v_0=\sqrt{2g\cos\alpha\,\frac{g\cos\alpha}{\omega^2\sin^2\alpha}-\omega^2\sin^2\alpha\,\frac{g^2\cos^2\alpha}{\omega^4\sin^4\alpha}}
=\frac{g}{\omega}\cot\alpha.
\]
Therefore it is only necessary that $v_0\geq\dfrac{g}{\omega}\cot\alpha$, and the small ring can move to the end $B$.

(2) If $x_0\geq l$, that is, $\omega\leq\sqrt{\dfrac{g\cos\alpha}{l\sin^2\alpha}}$, then during the motion of the small ring from $A$ to $B$ the direction of the resultant force along the rod is always towards the end $A$. For this case the integral in (3) should be written as
\begin{equation}
\int_{v_0}^{0}v\,\dd v=\int_{0}^{l}\left(\omega^2\sin^2\alpha\,x-g\cos\alpha\right)\dd x. \tag{5}
\end{equation}
Integrating (5), we obtain
\[
-\frac{1}{2}v_0^2=\frac{1}{2}\omega^2\sin^2\alpha\,l^2-g\cos\alpha\,l,
\]
\[
\therefore\quad v_0=\sqrt{2gl\cos\alpha-\omega^2\sin^2\alpha\,l^2}.
\]
That is, provided $v_0\geq\sqrt{2gl\cos\alpha-\omega^2\sin^2\alpha\,l^2}$, the small ring can move to the end $B$.

\begin{commentary}
In the rotating reference frame, two of the external forces acting on the small ring have components along the rod: the first, the weight, is a constant force; the second, the centrifugal force of inertia, is a variable force. As the position of the small ring rises, this force increases accordingly. Hence the resultant of the forces acting on the small ring along the rod gradually decreases from a comparatively large value directed towards the end $A$ to zero, and then its direction is towards the end $B$ and it gradually increases. In this problem it is only necessary to analyse whether or not there is a point on the rod at which the direction of this resultant changes, and the upper limit of the integral on the right-hand side of the equals sign in (3) or (5) can then be determined.
\end{commentary}
\end{solution}

\begin{solution}{1.15}
(1) From the equation of the shape of the bent rod, $y=\dfrac{1}{2a}x^2$, it follows that the values of $y$ of the small ring at $x=2a$ and at $x=a$ are respectively
\[
y(2a)=2a,\qquad y(a)=\frac{1}{2}a.
\]
The small ring is released from rest at $x=2a$. Let the speed with which it arrives at $x=a$ be $v$; then from the law of conservation of mechanical energy we obtain
\[
\frac{1}{2}mv^2=mg\left[y(2a)-y(a)\right],
\]
\begin{equation}
v=\sqrt{2g\left[y(2a)-y(a)\right]}=\sqrt{3ga}. \tag{1}
\end{equation}
In what follows we find the radius of curvature of the rod at $x=a$.

\textit{Method 1.}\quad The curve of the shape of the rod is compared with the parabolic trajectory of a body projected horizontally. That is, a body is projected from point $O$ with initial speed $v_0$ (along the $x$ direction), and the direction of the gravitational acceleration $g$ of this body is taken to be along the $y$ direction; then the equations of motion of this body are
\begin{align}
x&=v_0t \tag{2}\\
y&=\frac{1}{2}gt^2 \tag{3}
\end{align}

\begin{figbox}{1.15(1)}
\begin{tikzpicture}[scale=1.35,>=latex]
  \draw[axis] (-0.45,0)--(3.10,0) node[right] {$x$};
  \draw[axis] (0,-0.45)--(0,3.10) node[above] {$y$};
  \draw[thickline,domain=-1.95:2.45,samples=80] plot (\x,{0.5*\x*\x});
  \node[below left] at (0,0) {$O$};
  % point at x=a
  \coordinate (P) at (1,0.5);
  \fill (P) circle (1.3pt);
  % tangent at P: slope 1, crosses the x-axis at x=0.5
  \draw[dashedline] (0.30,-0.20)--(1.55,1.05);
  \draw[thinline] (0.78,0) arc (0:45:0.28);
  \node at (0.92,0.20) {$\theta$};
  \node[below] at (1,0) {$a$};
  \node[below] at (2,0) {$2a$};
\end{tikzpicture}
\end{figbox}

Eliminating $t$ from (2) and (3), we obtain
\begin{equation}
y=\frac{g}{2v_0^2}x^2. \tag{4}
\end{equation}
Comparing (4) with the equation of the shape of the bent rod, $y=\dfrac{1}{2a}x^2$, we obtain
\begin{equation}
a=\frac{v_0^2}{g}. \tag{5}
\end{equation}
At $x=a$, $x=a=v_0t$, so $t=\dfrac{a}{v_0}$; substituting (5) into this gives $t=\dfrac{v_0}{g}$. Hence at this instant
\[
v_y=gt=v_0=v_x.
\]
It is seen that here the tangent to the trajectory makes an angle $\theta=\dfrac{\pi}{4}$ with the $x$-axis, so the normal acceleration here is
\[
a_n=g\cos\theta=\frac{\sqrt{2}}{2}g.
\]
From this the radius of curvature $\rho(a)$ here can be found:
\[
\rho(a)=\frac{v^2}{a_n}=\frac{v_x^2+v_y^2}{a_n}
=\frac{2v_0^2}{\dfrac{\sqrt{2}}{2}g}=2\sqrt{2}\,\frac{v_0^2}{g}=2\sqrt{2}\,a.
\]

\textit{Method 2.}\quad Use the formula $\rho=\left|\dfrac{\left[1+(y')^2\right]^{3/2}}{y''}\right|$ to solve the problem.

$y'=\dfrac{\dd y}{\dd x}=\dfrac{x}{a}$, $y''=\dfrac{\dd^2y}{\dd x^2}=\dfrac{1}{a}$. Substituting these two expressions into the formula, we obtain
\[
\rho=\left|\frac{\left[1+\left(\dfrac{x}{a}\right)^2\right]^{3/2}}{\dfrac{1}{a}}\right|_{x=a}=2\sqrt{2}\,a.
\]

\begin{figbox}{1.15(2)}
\begin{tikzpicture}[scale=1.45,>=latex]
  % rod piece
  \draw[thickline,domain=0.40:2.60,samples=50]
     plot (\x,{1.0+0.5774*(\x-1.5)+0.15*(\x-1.5)*(\x-1.5)});
  % tangent at the ring
  \draw[dashedline] (0.70,0.538)--(2.55,1.606);
  \coordinate (M) at (1.5,1.0);
  \draw[thickline,fill=white] (M) circle (0.075);
  % normal force
  \draw[force] (M)--++(120:1.15) node[above left] {$N$};
  % weight
  \draw[force] (M)--++(-90:0.95) node[below] {$mg$};
  % angle theta with the horizontal
  \draw[dashedline] (M)--++(0:1.15);
  \draw[thinline] ($(M)+(0.52,0)$) arc (0:30:0.52);
  \node at ($(M)+(0.62,0.16)$) {$\theta$};
\end{tikzpicture}
\end{figbox}

When the small ring moves to $x=a$, the forces acting on it being as shown in Fig. 2, its equation of motion in the normal direction is
\begin{equation}
N-mg\cos\theta=m\frac{v^2}{\rho}. \tag{6}
\end{equation}
Substituting (1) into (6), we can solve for
\[
N=mg\cos\theta+m\frac{v^2}{\rho}=\frac{5}{4}\sqrt{2}\,mg.
\]

(2) In the rotating reference frame that rotates together with the rod, the forces acting on the ring, which is at rest relative to the rod, are as shown in Fig. 3. Hence
\begin{equation}
mg\sin\theta=f_c\cos\theta=m\omega_0^2x\cos\theta. \tag{7}
\end{equation}
where the slope of the tangent to the track is
\begin{equation}
\frac{\dd y}{\dd x}=\tan\theta=\frac{x}{a}. \tag{8}
\end{equation}

\begin{figbox}{1.15(3)}
\begin{tikzpicture}[scale=1.15,>=latex]
  \draw[axis] (-1.30,0)--(2.75,0) node[right] {$x$};
  \draw[axis] (0,-0.35)--(0,2.75) node[above] {$y$};
  \node[below left] at (0,0) {$O$};
  \draw[thickline,domain=-1.85:2.30,samples=80] plot (\x,{0.5*\x*\x});
  % rotation of the rod
  \draw[force] (0,2.20) ++(150:0.30) arc (150:-150:0.30);
  \node[right] at (0.28,2.20) {$\omega_0$};
  % ring
  \coordinate (M) at (1,0.5);
  \draw[thickline,fill=white] (M) circle (0.075);
  % tangent
  \draw[dashedline] (0.42,-0.08)--(1.62,1.12);
  % forces
  \draw[force] (M)--++(0:1.10) node[right] {$f_c$};
  \draw[force] (M)--++(-90:0.85) node[below] {$mg$};
  \draw[force] (M)--++(135:1.05) node[above left] {$N$};
  \draw[thinline] ($(M)+(0.45,0)$) arc (0:45:0.45);
  \node at ($(M)+(0.60,0.17)$) {$\theta$};
\end{tikzpicture}
\end{figbox}

Substituting (8) into (7), we obtain
\[
\omega_0=\sqrt{\frac{g}{a}}.
\]
That is, the bent rod need only rotate about the $y$-axis with angular velocity $\omega_0=\sqrt{\dfrac{g}{a}}$, and the small ring can then remain at rest relative to the rod at any position on the rod.

(3) If $\omega=\dfrac{\sqrt{2}}{2}\omega_0=\sqrt{\dfrac{g}{2a}}$, then in the reference frame rotating together with the rod, during the process in which the small ring slides from $x=2a$ to $x=a$ the work done by gravity is
\[
W_g=mg\left[y(2a)-y(a)\right]=mg\left(2a-\frac{a}{2}\right)=\frac{3}{2}mga.
\]
In addition, the work done by the centrifugal force of inertia $f=m\omega^2x$ can be found by analogy with the elastic force $f=kx$, and is
\[
W_{f_c}=\frac{1}{2}m\omega^2\left[a^2-(2a)^2\right]=-\frac{3}{2}m\omega^2a^2.
\]
From the work-energy theorem we obtain
\[
\frac{1}{2}mv'^2=W_g+W_{f_c}=\frac{3}{2}mga-\frac{3}{2}m\omega^2a^2,
\]
that is,
\begin{equation}
v'^2=3ga-3\,\frac{g}{2a}a^2=\frac{3}{2}ga. \tag{9}
\end{equation}

At $x=a$, besides the weight $mg$, the centrifugal force of inertia $f_c=m\omega^2x$ and the supporting force $N$, the small ring, since it has a velocity $v'$ along the rod, is also acted upon by the Coriolis force $f_C=2mv'\omega\cos\theta$. The direction of this force is that of $\vec{v}'\times\vec{\omega}$, that is, it is perpendicular to the plane of the page and points inwards (it is not marked in Fig. 4). The supporting force $N$ exerted by the rod on the small ring should have components in two directions, that is,
\begin{equation}
\vec{N}=\vec{N}_1+\vec{N}_2. \tag{10}
\end{equation}
Here $N_1$ lies in the plane of the page, and its value is determined by the equation of motion of the small ring in the normal direction within the plane of the page,
\begin{equation}
N_1-mg\cos\theta-m\omega^2a\sin\theta=m\frac{v'^2}{\rho}. \tag{11}
\end{equation}
Substituting (9) and $\rho=2\sqrt{2}\,a$ into (11), we obtain
\begin{align}
N_1&=\frac{\sqrt{2}}{2}mg+m\times\frac{g}{2a}\times a\times\frac{\sqrt{2}}{2}+m\frac{\dfrac{3}{2}ga}{2\sqrt{2}\,a}\notag\\
&=\frac{9}{8}\sqrt{2}\,mg. \tag{12}
\end{align}
And $N_2$ in (10) is equal in magnitude and opposite in direction to the Coriolis force, that is, its direction is perpendicular to the plane of the page and points outwards, and its value is
\begin{align}
N_2&=2mv'\omega\cos\theta\notag\\
&=2m\sqrt{\frac{3}{2}ga}\times\sqrt{\frac{g}{2a}}\times\frac{\sqrt{2}}{2}=\frac{\sqrt{6}}{2}mg. \tag{13}
\end{align}

\begin{figbox}{1.15(4)}
\begin{tikzpicture}[scale=1.15,>=latex]
  \draw[axis] (-1.30,0)--(2.75,0) node[right] {$x$};
  \draw[axis] (0,-0.35)--(0,2.75) node[above] {$y$};
  \node[below left] at (0,0) {$O$};
  \draw[thickline,domain=-1.85:2.30,samples=80] plot (\x,{0.5*\x*\x});
  \draw[force] (0,2.20) ++(150:0.30) arc (150:-150:0.30);
  \node[right] at (0.28,2.20) {$\omega$};
  \coordinate (M) at (1,0.5);
  \draw[thickline,fill=white] (M) circle (0.075);
  \draw[dashedline] (0.42,-0.08)--(1.62,1.12);
  \draw[force] (M)--++(0:1.10) node[right] {$f_c$};
  \draw[force] (M)--++(-90:0.85) node[below] {$mg$};
  \draw[force] (M)--++(135:1.05) node[above left] {$N_1$};
  \draw[thinline] ($(M)+(0.45,0)$) arc (0:45:0.45);
  \node at ($(M)+(0.60,0.17)$) {$\theta$};
\end{tikzpicture}
\end{figbox}

From (12) and (13) the value of the supporting force $N$ of the rod on the small ring can be obtained; this value is also the value of the force $N'$ exerted by the small ring on the rod,
\[
N'=N=\sqrt{N_1^2+N_2^2}=\frac{\sqrt{258}}{8}mg\approx2.01mg.
\]

\begin{commentary}
In the rotating reference frame, apart from the weight, the centrifugal force of inertia and the supporting force, the small ring, since it has a velocity $v'$, is also acted upon by the Coriolis force, and this force can be expressed as $\vec{f}_C=2m\vec{v}'\times\vec{\omega}$. The Coriolis force is very easily overlooked, and it should be noted that the direction of this force is perpendicular both to $\vec{v}'$ and to $\vec{\omega}$, that is, it is perpendicular to the plane (the plane of the page) formed by these two vectors. Since the small ring has no motion in this direction, the Coriolis force is balanced by the component $N_2$ of the supporting force, while of the supporting force only the value of the other component $N_1$ within the plane of the page has to be found. According to the Pythagorean theorem, the value of the resultant $N$ of $N_1$ and $N_2$ can then be found.
\end{commentary}
\end{solution}

\begin{solution}{1.16}
(1) On the small disc take an arbitrary diameter, and on this diameter take two small segments that are symmetric with respect to the centre $O'$ of the disc; the centres of these two small segments are $A_1$ and $A_2$, as shown in Fig. 2.

\begin{figbox}{1.16(1)}
\begin{tikzpicture}[scale=1.15,>=latex]
  \def\Rx{3.2}
  \def\Ry{1.15}
  \draw[thickline] (0,0) ellipse ({\Rx} and {\Ry});
  \draw[thickline] (0,-0.10) ellipse ({\Rx} and {\Ry});
  \draw[thickline] (0,0)--(0,2.95);
  \draw[force] (0,2.95) ++(-150:0.35) arc (-150:150:0.35);
  \node[right] at (0.32,2.95) {$\Omega$};
  \draw[thinline] (0,0)--(-2.62,-0.66);
  \node[above right] at (-1.32,-0.32) {$R$};
  \coordinate (S) at (1.75,0.30);
  \draw[thickline] (S) ellipse (0.55 and 0.27);
  \draw[thickline] (S)--++(0,0.42);
  \node[right] at ($(S)+(0.28,0.12)$) {$r$};
  \draw[dim] (0,0.80)--(1.75,0.80);
  \node[above] at (0.87,0.84) {$d$};
  \node[below left] at (0,0) {$O$};
  \node[below right] at (S) {$O'$};
\end{tikzpicture}
\end{figbox}

\begin{figbox}{1.16(2)}
\begin{tikzpicture}[scale=1.25,>=latex]
  \coordinate (O) at (0,0);
  \coordinate (Op) at (2.4,0);
  \coordinate (A1) at (1.98,0.74);
  \coordinate (A2) at (2.83,-0.74);
  \draw[thinline] (O)--(A1) node[midway,above left] {$r_1$};
  \draw[thinline] (O)--(A2) node[midway,below left] {$r_2$};
  \draw[thinline] (A1)--(A2);
  % lines of action of the friction forces
  \draw[dashedline] (2.102,0.412)--(1.579,1.818);
  \draw[dashedline] (2.742,-1.079)--(3.120,0.373);
  % lever arms
  \draw[thinline] (Op)--(1.687,1.527);
  \draw[thinline] (Op)--(2.677,-1.329);
  \node[fill=white,inner sep=0.6pt] at (2.16,0.62) {$l_1$};
  \node[fill=white,inner sep=0.6pt] at (2.70,-0.62) {$l_2$};
  % velocities and friction forces
  \draw[force] (A1)--++(110.4:0.90) node[left] {$v_1$};
  \draw[force] ($(A1)+0.13*(0.937,0.349)$)--++(110.4:0.55) node[right] {$f_1$};
  \draw[force] (A2)--++(75.4:0.95) node[above right] {$v_2$};
  \draw[force] ($(A2)+0.13*(0.968,-0.252)$)--++(75.4:0.60) node[below right] {$f_2$};
  % points
  \fill (O) circle (1.4pt);
  \fill (Op) circle (1.4pt);
  \fill (A1) circle (1.4pt);
  \fill (A2) circle (1.4pt);
  \node[below left] at (O) {$O$};
  \node[above left] at (A1) {$A_1$};
  \node[right] at ($(A2)+(0.10,-0.06)$) {$A_2$};
  \node[above right] at (Op) {$O'$};
\end{tikzpicture}
\end{figbox}

When the small disc has just been put down, the velocities of the two points on the large disc that are in contact with the points $A_1$ and $A_2$ of the small disc are respectively $v_1$ and $v_2$, their directions being respectively perpendicular to $r_1$ and $r_2$. At this time the directions of the friction forces $f_1$ and $f_2$ acting at the points $A_1$ and $A_2$ are respectively along the directions of $v_1$ and $v_2$. From Fig. 2 it is seen that the lever arm $l_1$ of $f_1$ with respect to $O'$ is smaller than the lever arm $l_2$ of $f_2$ with respect to $O'$, so the resultant torque of $f_1$ and $f_2$ with respect to $O'$ makes the small disc rotate with acceleration about $O'$ in the same sense of rotation as the large disc.

Suppose that at some instant the angular velocity of the rotation of the small disc is $\omega(<\Omega)$; the velocities of the points $A_1$ and $A_2$ on the small disc at this instant are respectively $u_1=\vec{\omega}\times\vec{r}_1'$ and $u_2=\vec{\omega}\times\vec{r}_2'$, while the velocities of the two points on the large disc that are in contact with $A_1$ and $A_2$ are respectively $v_1=\vec{\Omega}\times\vec{r}_1$ and $v_2=\vec{\Omega}\times\vec{r}_2$. Hence the relative velocities at the two points of contact are respectively
\[
v_1'=v_1-u_1,\qquad v_2'=v_2-u_2.
\]
Since $v_1-v_2=\vec{\Omega}\times(\vec{r}_1-\vec{r}_2)$, while $\vec{r}_1-\vec{r}_2=\vec{r}_1'-\vec{r}_2'$, that is, its direction is the direction of the line joining $A_1$ and $A_2$. From this it is known that the direction of $v_1-v_2$ is perpendicular to the line joining $A_1$ and $A_2$, and since the directions of $u_1$ and $u_2$ are respectively perpendicular to $r_1'$ and $r_2'$, that is, both are perpendicular to the line joining $A_1$ and $A_2$, the directions of $u_1$ and $u_2$ lie along the vector line of $v_1-v_2$.

\begin{figbox}{1.16(3)}
\begin{tikzpicture}[scale=1.30,>=latex]
  \coordinate (O) at (0,0);
  \coordinate (Op) at (1.6,0);
  \coordinate (A1) at (1.32,0.49);
  \coordinate (A2) at (1.88,-0.49);
  \draw[thickline,dashed] (Op) circle (0.822);
  \draw[thinline] (O)--(A1);
  \draw[thinline] (O)--(A2);
  \draw[thinline] (A1)--(A2);
  \draw[thinline] (O)--(Op) node[midway,below] {$R$};
  \node[below left] at (O) {$O$};
  \node[below right] at (Op) {$O'$};
  \node[below left] at (A1) {$A_1$};
  \node[right] at (A2) {$A_2$};
  \node[above left] at (0.66,0.24) {$r_1$};
  \node[below left] at (0.94,-0.24) {$r_2$};
  % fan of vectors at A1
  \draw[force] (A1)--++(106:1.15) node[left] {$v_1$};
  \draw[force] (A1)--++(70:1.15) node[right] {$v_2$};
  \draw[force] (A1)--++(25:0.75) node[right] {$-u_1$};
  \draw[force] (A1)--++(205:0.75) node[below left] {$-u_2$};
  \draw[force] (A1)--++(75.7:1.47) node[above right] {$v_1'$};
  \draw[force] (A1)--++(110.6:0.82) node[above left] {$v_2'$};
  \draw[force] ($(A1)+0.06*(0.969,-0.247)$)--++(75.7:0.55) node[right] {$f_1$};
  \draw[force] ($(A1)-0.06*(0.936,0.351)$)--++(110.6:0.50) node[below left] {$f_2$};
\end{tikzpicture}
\end{figbox}

The vector diagram of $v_2'=v_2-u_2$ at the point $A_2$ is drawn over at the point $A_1$ for convenience of comparison; then the directions of the friction forces $f_1$ and $f_2$ are respectively along the directions of their own relative velocities $v_1'$ and $v_2'$, as shown in Fig. 3. Since
\[
u_1-u_2=\vec{\omega}\times(\vec{r}_1'-\vec{r}_2')=\vec{\omega}\times(\vec{r}_1-\vec{r}_2),
\]
\[
v_1-v_2=\vec{\Omega}\times(\vec{r}_1-\vec{r}_2),
\]
and $\omega<\Omega$, from the above two expressions it is known that
\[
|u_1-u_2|<|v_1-v_2|.
\]
Therefore there is an angle between $f_1$ and $f_2$, that is, the resultant torque of $f_1$ and $f_2$ with respect to the point $O'$ is still not zero, so that the small disc continues to rotate with acceleration.

When $\omega=\Omega$, we have
\[
|u_1-u_2|=|v_1-v_2|.
\]
At this time the directions of $f_1$ and $f_2$ are both along the direction $f$, since
\[
v_1'=v_1-u_1=\vec{\Omega}\times(\vec{r}_1-\vec{r}_1')=\vec{\Omega}\times\vec{R},
\]
\[
v_2'=v_2-u_2=\vec{\Omega}\times(\vec{r}_2-\vec{r}_2')=\vec{\Omega}\times\vec{R}.
\]
It is seen that $f_1=f_2=f$, and the direction is perpendicular to $\vec{R}$; therefore the resultant torque of $f_1$ and $f_2$ with respect to $O'$ is zero. The rotation of the small disc then becomes steady, and the steady angular velocity is
\[
\omega=\Omega.
\]

(2) When $\omega=\Omega$, the directions of the friction forces acting on any two symmetric small segments on any diameter of the small disc are all perpendicular to $\vec{R}$, that is, perpendicular to the direction of the line $OO'$; hence the total friction force acting on the small disc is $f=\mu mg$, and its equivalent point of application is $O'$. From Newton's second law it is known that the torque of the friction force acting on the large disc with respect to the point $O$ is
\[
M'=\mu mgd.
\]
Therefore, in order to maintain the rotational speed of the large disc unchanged, the external torque that must be applied is $M=\mu mgd$.

\begin{commentary}
In the solution of this problem, the vector diagram of $v_2'=v_2-u_2$ at the point $A_2$ is drawn at the point $A_1$ and combined with the vector diagram of $v_1'=v_1-u_1$ at the point $A_1$ into a single geometrical figure. From this it can be seen very clearly that, in the case $\omega<\Omega$, because $|u_1-u_2|<|v_1-v_2|$ the directions of the friction forces $f_1$ and $f_2$ are different, and the correct conclusion is thereby obtained that under the action of this friction torque the small disc continues to rotate with acceleration until $\omega=\Omega$.
\end{commentary}
\end{solution}

\begin{solution}{1.17}
Since the weights acting on the two balls $A$ and $B$ pass through their respective centres, while the surfaces of the three balls are all smooth, the force $N$ exerted by ball $C$ on balls $A$ and $B$ also passes through the centres of the two balls. Therefore, in equilibrium the tension $T$ of the string also passes through the centres of the two balls.

\begin{figbox}{1.17}
\begin{tikzpicture}[scale=1.05,>=latex]
  \coordinate (O)  at (0,3.60);
  \coordinate (OA) at (-1.362,0.927);
  \coordinate (OB) at (1.362,0.927);
  \coordinate (OC) at (0,1.556);
  \coordinate (PA) at (-0.863,1.907);
  \coordinate (PB) at (0.863,1.907);
  % strings
  \draw[thickline] (O)--(PA);
  \draw[thickline] (O)--(PB);
  % balls
  \draw[thickline] (OA) circle (1.10);
  \draw[thickline] (OB) circle (1.10);
  \draw[thickline] (OC) circle (0.55);
  % vertical axis
  \draw[dashedline] (O)--(0,0.85);
  % line of centres A-C
  \draw[thinline] (OA)--(OC);
  % angle theta at O
  \draw[thinline] ($(O)+(243:0.72)$) arc (243:270:0.72);
  \node at ($(O)+(256:0.97)$) {$\theta$};
  % angle phi at OA
  \draw[thinline] ($(OA)+(24.8:0.50)$) arc (24.8:63:0.50);
  \node at ($(OA)+(44:0.86)$) {$\phi$};
  % R and r on the line of centres
  \node[fill=white,inner sep=0.6pt] at ($(OA)+0.55*(0.908,0.419)$) {$R$};
  \node[fill=white,inner sep=0.6pt] at ($(OC)+0.30*(-0.908,-0.419)$) {$r$};
  % forces on C
  \draw[force] (OC)--++(24.8:1.15) node[above right] {$N$};
  \draw[force] (OC)--++(155.2:1.15) node[above left] {$N$};
  \draw[force] (OC)--++(-90:1.10) node[below] {$mg$};
  % forces on A
  \draw[force] (OA)--++(63:1.05) node[left] {$T$};
  \draw[force] (OA)--++(-90:0.90) node[left] {$Mg$};
  \draw[force] (OA)--++(204.8:1.30) node[below left] {$N$};
  % forces on B
  \draw[force] (OB)--++(117:1.05) node[right] {$T$};
  \draw[force] (OB)--++(-90:0.90) node[right] {$Mg$};
  \draw[force] (OB)--++(335.2:1.30) node[below right] {$N$};
  % labels
  \fill (O) circle (1.5pt);
  \node[above right] at (O) {$O$};
  \node[below] at ($(OA)+(0,-1.32)$) {$A$};
  \node[below] at ($(OB)+(0,-1.32)$) {$B$};
  \node[right] at ($(OC)+(0.42,0.28)$) {$C$};
  \node[above left] at ($(OA)+(-0.10,0.12)$) {$O_A$};
  \node[above right] at ($(OB)+(0.10,0.12)$) {$O_B$};
  \node[right] at ($(OC)+(0.06,0.10)$) {$O_C$};
\end{tikzpicture}
\end{figbox}

In the figure it is known that $OO_A=3R$, $O_AO_C=\dfrac{3}{2}R$; from the sine rule we obtain
\[
\frac{\dfrac{3}{2}R}{\sin\theta}=\frac{3R}{\sin(\theta+\phi)},
\]
that is,
\begin{equation}
\sin(\theta+\phi)=2\sin\theta. \tag{1}
\end{equation}
From the equilibrium of the forces acting on ball $C$ we obtain
\begin{equation}
2N\cos(\theta+\phi)=mg=nMg. \tag{2}
\end{equation}
From the equilibrium of the forces acting on ball $A$ (or $B$) in the horizontal direction and in the vertical direction, the following two expressions are obtained respectively:
\begin{align}
T\sin\theta&=N\sin(\theta+\phi) \tag{3}\\
T\cos\theta-N\cos(\theta+\phi)-Mg&=0 \tag{4}
\end{align}
From (1), (2), (3) and (4) we can solve and obtain
\begin{equation}
\sin\theta=\frac{1}{4}\sqrt{\frac{4+4n-3n^2}{1+n}}. \tag{5}
\end{equation}
When balls $A$ and $B$ are in contact with each other, $\theta$ has a minimum value $\theta_0$; from the geometrical relation we obtain
\begin{equation}
\sin\theta\geq\sin\theta_0=\frac{1}{3}. \tag{6}
\end{equation}
From (5) and (6) we obtain
\[
\frac{1}{4}\sqrt{\frac{4+4n-3n^2}{1+n}}\geq\frac{1}{3}.
\]
Solving, we obtain
\[
n\leq\frac{10+8\sqrt{10}}{27}\approx1.3.
\]
Therefore the following analysis can be made:

(1) If $n>1.3$, the system cannot be in equilibrium.

(2) If $n=1.3$, then when the system is in equilibrium $\theta=\theta_0=\arcsin\dfrac{1}{3}$, that is, balls $A$ and $B$ lean against each other, but there is no interaction between the two balls.

(3) If $n<1.3$, the system has two possible states of equilibrium:

(i) $\theta=\theta_0=\arcsin\dfrac{1}{3}$, that is, balls $A$ and $B$ are pressed tightly against each other, and there is an interaction between the two balls; the magnitude of this force $F$ is related to $n$: the larger $n$ is, the smaller $F$ is, and when $n=1.3$, $F=0$.

(ii) $\theta=\arcsin\left[\dfrac{1}{4}\sqrt{\dfrac{4+4n-3n^2}{1+n}}\right]$, that is, balls $A$ and $B$ are separated; the smaller $n$ is, the farther apart they are. If $n\to0$, then $\theta\to\arcsin\dfrac{1}{2}$.
\end{solution}

\begin{solution}{1.18}
Let the radius of the large ring be $R$. The nail is located at the point $P$, and the small ball is fixed at the point $A$ of the large ring, as shown in the figure. The large ring is acted upon by the supporting force $N$ and the friction force $f$ of the nail. The direction of the resultant $F$ of these two forces must be vertically upward. The point $B$ is the intersection of the vertical line through $P$ with $OA$. Let $OB=r$.

\begin{figbox}{1.18}
\begin{tikzpicture}[scale=1.15,>=latex]
  \coordinate (O) at (0,0);
  \coordinate (P) at (0.704,1.324);
  \coordinate (B) at (0.704,-0.429);
  \coordinate (A) at (1.280,-0.780);
  \draw[thickline] (O) circle (1.5);
  \draw[dashedline] (0.704,2.48)--(0.704,-1.70);
  \draw[thinline] (O)--(A);
  \fill (O) circle (1.4pt);
  \fill (P) circle (1.5pt);
  \fill (A) circle (2.2pt);
  \fill (B) circle (1.2pt);
  \node[below left] at (O) {$O$};
  \node[right] at ($(P)+(0.06,0.03)$) {$P$};
  \node[right] at ($(B)+(0.06,0.02)$) {$B$};
  \node[right] at ($(A)+(0.10,-0.03)$) {$A$};
  \node[above left] at (0.34,0.72) {$R$};
  \node[right] at (0.80,-0.20) {$r$};
  % forces at the nail
  \draw[force] (P)--++(62:0.95) node[right] {$N$};
  \draw[force] (P)--++(152:0.50) node[left] {$f$};
  \draw[force] (P)--++(90:1.10) node[right] {$F$};
  % weights
  \draw[force] (O)--++(-90:0.65) node[below] {$Mg$};
  \draw[force] (A)--++(-90:0.75) node[below] {$mg$};
  % angles
  \draw[thinline] ($(P)+(242:0.48)$) arc (242:270:0.48);
  \node at ($(P)+(257:0.70)$) {$\theta$};
  \draw[thinline] ($(B)+(90:0.50)$) arc (90:148.6:0.50);
  \node at ($(B)+(-0.30,0.66)$) {$\alpha$};
\end{tikzpicture}
\end{figbox}

For $\triangle OPB$, from the sine rule we obtain
\[
\frac{R}{\sin\alpha}=\frac{r}{\sin\theta},
\]
that is,
\[
\sin\theta=\frac{r}{R}\sin\alpha.
\]
It is seen that when $\alpha=\dfrac{\pi}{2}$, the value of $\theta$ is a maximum, and its maximum value $\theta_m$ is
\begin{equation}
\theta_m=\arcsin\frac{r}{R}. \tag{1}
\end{equation}
When the large ring is in equilibrium, the moments of the external forces acting on it with respect to the point $P$ balance, and we obtain
\[
Mgr\sin\alpha=mg(R-r)\sin\alpha,
\]
that is,
\begin{equation}
\frac{r}{R}=\frac{m}{M+m}. \tag{2}
\end{equation}
Substituting (2) into (1), we obtain
\begin{equation}
\theta_m=\arcsin\frac{m}{M+m}. \tag{3}
\end{equation}
Since the large ring does not slip relative to the nail, the angle $\theta$ must be smaller than or equal to the angle of friction $\phi_{\max}(=\arctan\mu)$, that is,
\begin{equation}
\theta\leq\phi_{\max}. \tag{4}
\end{equation}
Expression (3) represents the maximum value of $\theta$ for equilibrium when different points of the large ring are hung on the nail; as long as this angle is smaller than or equal to the angle of friction, the equilibrium of the large ring hung on the nail at any position can be ensured. In addition, from $\tan\phi_{\max}=\mu$ we obtain $\sin\phi_{\max}=\dfrac{\mu}{\sqrt{\mu^2+1}}$; hence we obtain
\[
\sin\theta_m=\frac{m}{M+m}\leq\sin\phi_m=\frac{\mu}{\sqrt{\mu^2+1}}.
\]
From the above expression we can solve and obtain
\[
\mu\geq\frac{m}{\sqrt{M(M+2m)}}.
\]
\end{solution}

\begin{solution}{1.19}
(1) Let $\alpha$ be the maximum value $\alpha_m$ at which the system can be in equilibrium. At this time the tendency of relative motion between ball 1 and the inclined plane is greatest, that is, the friction force acting on ball 1 is the maximum static friction force, $f_1=\mu_1mg\cos\theta$, while the tendency of relative motion between ball 2 and the inclined plane has not reached its maximum, that is, $f_2$ is an ordinary static friction force. Since the direction of the motion of ball 1 relative to the inclined plane is perpendicular to the string, it is seen that $f_1$ is also perpendicular to the string. Let the tension in the string be $T$, and let the component of the weight of each of the two balls along the inclined plane be $F$, that is, $F=mg\sin\theta$. Fig. 2 shows the vector relations of the external forces acting on the two balls; drawing the two vector diagrams together will be helpful for comparing the magnitudes and directions of the various vectors. From the geometrical relations shown in the figure we obtain
\begin{equation}
T=F\cos\alpha_m=mg\sin\theta\cos\alpha_m \tag{1}
\end{equation}
\begin{equation}
f_1=\mu_1mg\cos\theta=F\sin\alpha_m=mg\sin\theta\sin\alpha_m \tag{2}
\end{equation}

\begin{figbox}{1.19(1)}
\begin{tikzpicture}[scale=1.15,>=latex]
  \coordinate (V) at (0,0);
  \coordinate (Va) at (3.60,0.55);
  \coordinate (Vb) at (1.30,2.50);
  \draw[thickline] (V)--(Va)--($(Va)+(Vb)$)--(Vb)--cycle;
  \coordinate (P1) at (2.14,1.248);
  \coordinate (P2) at (3.116,2.041);
  \draw[thickline] (P1)--(P2);
  \draw[thickline,fill=white] (P1) circle (0.10);
  \draw[thickline,fill=white] (P2) circle (0.10);
  \draw[thinline] (0.80,0.122) arc (8.68:62.5:0.80);
  \node at (0.92,0.58) {$\theta$};
  % line of greatest slope through ball 1
  \draw[dashedline] (P1)--++(62.5:1.15);
  \draw[thinline] ($(P1)+(39.1:0.52)$) arc (39.1:62.5:0.52);
  \node at ($(P1)+(52:0.72)$) {$\alpha$};
  % forces
  \draw[force] (P1)--++(242.5:0.80) node[left] {$F$};
  \draw[force] (P2)--++(242.5:0.80) node[left] {$F$};
  \draw[force] (P1)--++(129.1:0.62) node[above left] {$f_1$};
  \node[below right] at ($(P1)+(0.04,-0.06)$) {$1$};
  \node[above right] at ($(P2)+(0.04,0.06)$) {$2$};
\end{tikzpicture}
\end{figbox}

\begin{figbox}{1.19(2)}
\begin{tikzpicture}[scale=1.02,>=latex]
  \coordinate (B) at (0,0);
  \coordinate (A) at (0,-1.90);
  \coordinate (C) at (0.692,-1.601);
  \coordinate (D) at (0.692,-3.501);
  \draw[force] (B)--(A) node[midway,left] {$F$};
  \draw[force] (B)--(C) node[pos=0.55,left] {$T$};
  \draw[force] (C)--(A) node[midway,below] {$f_1$};
  \draw[force] (B)--(D) node[pos=0.72,right] {$f_2$};
  \draw[thinline] ($(B)+(0,-0.60)$) arc (-90:-66.6:0.60);
  \node at (0.18,-0.76) {$\alpha_m$};
\end{tikzpicture}
\end{figbox}

From (2) we can solve and obtain
\[
\sin\alpha_m=\mu_1\cot\theta=\sqrt{\frac{\mu_1}{\mu_2}}.
\]
At the same time, we also have
\begin{align}
f_2^2&=F^2+T^2+2FT\cos\alpha_m\notag\\
&=(mg\sin\theta)^2+(mg\sin\theta\cos\alpha_m)^2+2(mg\sin\theta)(mg\sin\theta\cos\alpha_m)\cos\alpha_m\notag\\
&=m^2g^2\sin^2\theta\left(1+3\cos^2\alpha_m\right). \tag{3}
\end{align}
Since $f_2$ is an ordinary static friction force, that is, $f_2\leq\mu_2mg\cos\theta$; substituting this relation into (3), we obtain
\[
f_2^2=m^2g^2\sin^2\theta\left(1+3\cos^2\alpha_m\right)\leq\left(\mu_2mg\cos\theta\right)^2,
\]
that is,
\begin{equation}
\tan^2\theta\left(1+3\cos^2\alpha_m\right)\leq\mu_2^2. \tag{4}
\end{equation}
Substituting $\tan\theta=\sqrt{\mu_1\mu_2}$ and $\cos\alpha_m=\sqrt{\dfrac{\mu_2-\mu_1}{\mu_2}}$ into (4), we obtain
\begin{equation}
\mu_1\mu_2\left[1+\frac{3(\mu_2-\mu_1)}{\mu_2}\right]\leq\mu_2^2. \tag{5}
\end{equation}
Let $x=\mu_1/\mu_2$; then (5) becomes
\begin{equation}
3x^2-4x+1\geq0. \tag{6}
\end{equation}
Solving this inequality, we obtain two sets of solutions:
\[
x_1\geq1,\quad\text{which does not agree with the given condition }\mu_1<\mu_2\text{, and is discarded;}
\]
\[
x_2\leq\frac{1}{3},\quad\text{that is, }\mu_2\geq3\mu_1.
\]
From the above analysis it is known that if $\mu_2\geq3\mu_1$, then the maximum value $\alpha_m$ of $\alpha$ satisfies
\[
\sin\alpha_m=\sqrt{\frac{\mu_1}{\mu_2}}.
\]

(2) If $\mu_1<\mu_2<3\mu_1$, then when $\alpha$ takes the maximum value $\alpha_m$ at which the system can be in equilibrium, the tendency of relative motion between ball 2 and the inclined plane reaches its maximum, that is, the friction force acting on ball 2 is the maximum static friction force, so that $f_2=\mu_2mg\cos\theta$, while $f_1$ is an ordinary static friction force, that is, $f_1\leq\mu_1mg\cos\theta$. In this case the vector relations of the external forces acting on the two balls are as shown in Fig. 3. From the geometrical relations shown in the figure we obtain
\begin{equation}
f_1^2=T^2+(mg\sin\theta)^2-2Tmg\sin\theta\cos\alpha_m\leq(\mu_1mg\cos\theta)^2 \tag{7}
\end{equation}
\begin{equation}
f_2^2=T^2+(mg\sin\theta)^2+2Tmg\sin\theta\cos\alpha_m=(\mu_2mg\cos\theta)^2 \tag{8}
\end{equation}

\begin{figbox}{1.19(3)}
\begin{tikzpicture}[scale=1.02,>=latex]
  \coordinate (B) at (0,0);
  \coordinate (A) at (0,-1.90);
  \coordinate (C) at (0.760,-1.750);
  \coordinate (D) at (0.760,-3.650);
  \draw[force] (B)--(A) node[midway,left] {$F$};
  \draw[force] (B)--(C) node[pos=0.60,left] {$T$};
  \draw[force] (C)--(A) node[midway,above] {$f_1$};
  \draw[force] (B)--(D) node[pos=0.78,right] {$f_2$};
  \draw[thinline] ($(B)+(0,-0.62)$) arc (-90:-67.8:0.62);
  \node at (0.19,-0.78) {$\alpha_m$};
\end{tikzpicture}
\end{figbox}

Adding the two sides of (7) and (8), we obtain
\begin{equation}
T\leq\sqrt{\left(\frac{\mu_1^2+\mu_2^2}{2}\right)(mg\cos\theta)^2-(mg\sin\theta)^2}. \tag{9}
\end{equation}
Subtracting the two sides of (7) and (8), we obtain
\begin{equation}
4Tmg\sin\theta\cos\alpha_m\geq(\mu_2^2-\mu_1^2)(mg\cos\theta)^2. \tag{10}
\end{equation}
Substituting (9) into (10), we obtain
\[
\cos\alpha_m\geq\frac{\sqrt{2}(\mu_1+\mu_2)}{4\sqrt{\mu_1\mu_2}}.
\]

\begin{commentary}
The given condition of this problem is $\mu_2>\mu_1$, and for this reason it is very easy to conclude that, when $\alpha$ takes the value $\alpha_m$, it is the friction force $f_1$ acting on ball 1 that reaches the maximum static friction force while the friction force $f_2$ acting on ball 2 does not reach the maximum static friction force; such a conclusion would lead to the omission of the discussion of case (2) in the above analysis. From the above analysis it is known that when $\mu_2\geq3\mu_1$ it is indeed true that when $\alpha=\alpha_m$, $f_1=\mu_1mg\cos\theta$ and $f_2\leq\mu_2mg\cos\theta$; but if $\mu_1<\mu_2<3\mu_1$, then $f_2=\mu_2mg\cos\theta$ while $f_1\leq\mu_1mg\cos\theta$. By calculation the value of $\alpha_m$ or its range of values can be found in the two cases separately.
\end{commentary}
\end{solution}

\begin{solution}{1.20}
\textit{Method 1.}\quad Let the components along the horizontal and vertical directions of the force acting on the end $A_1$ of rod $O_1A_1$ be $F_{1x}$ and $F_{1y}$ respectively, their directions being as shown in Fig. 1. From symmetry it is known that $F_{1y}$ should be equal to half of the weight of the rods below the hinges $A_1$ and $B_1$. Since the mass of the rod between any two hinges is $m$, we have
\begin{equation}
F_{1y}=5mg. \tag{1}
\end{equation}

\begin{figbox}{1.20(1)}
\begin{tikzpicture}[scale=1.00,>=latex]
  \coordinate (O1) at (0,4.80);
  \coordinate (A1) at (-0.85,3.95);
  \coordinate (B1) at (0.85,3.95);
  \coordinate (O2) at (0,3.10);
  \coordinate (A2) at (-0.85,2.25);
  \coordinate (B2) at (0.85,2.25);
  \coordinate (O3) at (0,1.40);
  \coordinate (A3) at (-0.85,0.55);
  \coordinate (B3) at (0.85,0.55);
  \coordinate (O4) at (0,-0.30);
  % ceiling
  \draw[thickline] (-1.8,4.95)--(1.8,4.95);
  \foreach \xx in {-1.7,-1.45,...,1.7}
    \draw[thinline] (\xx,4.95)--++(-0.10,0.17);
  % rods
  \draw[thickline] (O1)--(A1);
  \draw[thickline] (O1)--(B1);
  \draw[thickline] (A1)--(B2);
  \draw[thickline] (B1)--(A2);
  \draw[thickline] (A2)--(B3);
  \draw[thickline] (B2)--(A3);
  \draw[thickline] (A3)--(O4);
  \draw[thickline] (B3)--(O4);
  % string O3-O4
  \draw[thinline] (O3)--(O4);
  \foreach \p in {O1,A1,B1,O2,A2,B2,O3,A3,B3,O4}
    \draw[fill=white,line width=0.5pt] (\p) circle (1.5pt);
  % forces at A1
  \draw[force] (A1)--++(0:0.80) node[above right] {$F_{1x}$};
  \draw[force] (A1)--++(90:0.55) node[above left] {$F_{1y}$};
  \draw[force] (-0.425,4.375)--++(-90:0.55) node[pos=0.55,right] {$mg$};
  % labels
  \node[right] at ($(O1)+(0.08,-0.06)$) {$O_1$};
  \node[below left] at ($(A1)+(-0.06,-0.04)$) {$A_1$};
  \node[right] at ($(B1)+(0.09,0)$) {$B_1$};
  \node[right] at ($(O2)+(0.08,-0.06)$) {$O_2$};
  \node[left] at ($(A2)+(-0.09,0)$) {$A_2$};
  \node[right] at ($(B2)+(0.09,0)$) {$B_2$};
  \node[right] at ($(O3)+(0.08,-0.06)$) {$O_3$};
  \node[left] at ($(A3)+(-0.09,0)$) {$A_3$};
  \node[right] at ($(B3)+(0.09,0)$) {$B_3$};
  \node[right] at ($(O4)+(0.08,-0.06)$) {$O_4$};
\end{tikzpicture}
\end{figbox}

From the equilibrium of the moments of the external forces acting on rod $O_1A_1$ with respect to the point $O_1$, we obtain
\begin{equation}
F_{1x}l\sin\theta=F_{1y}l\cos\theta+mg\,\frac{l}{2}\cos\theta. \tag{2}
\end{equation}
Substituting $\theta=45\degree$ and (1) into (2), we obtain
\begin{equation}
F_{1x}=\frac{11}{2}mg. \tag{3}
\end{equation}

The external forces acting on rod $A_1B_2$ that have moments with respect to the point $O_2$ are as shown in Fig. 2, where $F_{2x}$ and $F_{2y}$ are the components in the horizontal and vertical directions of the force acting on the end $B_2$. Similarly, we have
\begin{equation}
F_{2y}=3mg. \tag{4}
\end{equation}

\begin{figbox}{1.20(2)}
\begin{tikzpicture}[scale=1.10,>=latex]
  \coordinate (A1) at (0,2.0);
  \coordinate (O2) at (1.0,1.0);
  \coordinate (B2) at (2.0,0);
  \draw[thickline] (A1)--(B2);
  \draw[thickline,fill=white] (A1) circle (0.05);
  \draw[thickline,fill=white] (O2) circle (0.05);
  \draw[thickline,fill=white] (B2) circle (0.05);
  \draw[force] (A1)--++(0:0.85) node[below] {$F_{1x}$};
  \draw[force] (A1)--++(90:0.75) node[left] {$F_{1y}$};
  \draw[force] (B2)--++(0:0.85) node[right] {$F_{2x}$};
  \draw[force] (B2)--++(-90:0.75) node[right] {$F_{2y}$};
  \draw[force] (O2)--++(-90:0.60) node[right] {$2mg$};
  \node[left] at (A1) {$A_1$};
  \node[above left] at (O2) {$O_2$};
  \node[above right] at (B2) {$B_2$};
  \draw[thinline] (1.0,0.55) arc (-90:-45:0.45);
  \node at (1.22,0.62) {$\theta$};
\end{tikzpicture}
\end{figbox}

From the equilibrium of the moments of the external forces acting on rod $A_1B_2$ with respect to the point $O_2$, we obtain
\begin{equation}
F_{1y}l\cos\theta+F_{1x}l\sin\theta+F_{2y}l\cos\theta=F_{2x}l\sin\theta. \tag{5}
\end{equation}
Substituting (1), (3) and (4) into (5), we can solve and obtain
\begin{equation}
F_{2x}=\frac{27}{2}mg. \tag{6}
\end{equation}

Taking rod $B_2A_3$ and rod $A_3O_4$ as a single whole, the external forces acting on it that have moments with respect to the point $O_3$ are as shown in Fig. 3. In the figure, $N$ is the external force acting on the end $O_4$; from the symmetry of the whole hinge structure it is known that $N$ must be horizontal. Hence, from the equilibrium of the moments of this whole with respect to the point $O_3$, we obtain
\begin{equation}
F_{2x}l\sin\theta+F_{2y}l\cos\theta+mg\,\frac{l}{2}\cos\theta=N\sqrt{2}\,l. \tag{7}
\end{equation}
Substituting (4) and (6) into (7), we can solve and obtain
\begin{equation}
N=\frac{17}{2}mg. \tag{8}
\end{equation}

\begin{figbox}{1.20(3)}
\begin{tikzpicture}[scale=1.05,>=latex]
  \coordinate (B2) at (2.4,3.2);
  \coordinate (A3) at (0.4,1.2);
  \coordinate (O3) at (1.4,2.2);
  \coordinate (O4) at (1.4,0.2);
  \coordinate (M4) at (0.9,0.7);
  \draw[thickline] (B2)--(A3);
  \draw[thickline] (A3)--(O4);
  \draw[thickline,fill=white] (B2) circle (0.05);
  \draw[thickline,fill=white] (A3) circle (0.05);
  \draw[thickline,fill=white] (O3) circle (0.05);
  \draw[thickline,fill=white] (O4) circle (0.05);
  \draw[force] (B2)--++(0:0.85) node[right] {$F_{2x}$};
  \draw[force] (B2)--++(-90:0.70) node[right] {$F_{2y}$};
  \draw[force] (O4)--++(180:0.85) node[left] {$N$};
  \draw[force] (O4)--++(90:0.90) node[right] {$\dfrac{T}{2}$};
  \draw[force] (O3)--++(-90:0.55) node[right] {$2mg$};
  \draw[force] (M4)--++(-90:0.55) node[left] {$mg$};
  \node[right] at ($(B2)+(0.08,-0.05)$) {$B_2$};
  \node[left] at ($(A3)+(-0.08,0)$) {$A_3$};
  \node[above right] at (O3) {$O_3$};
  \node[below] at (O4) {$O_4$};
\end{tikzpicture}
\end{figbox}

From the equilibrium of the moments of rod $A_3O_4$ with respect to the point $A_3$, we obtain
\begin{equation}
mg\,\frac{l}{2}\cos\theta+Nl\sin\theta=\frac{T}{2}l\cos\theta. \tag{9}
\end{equation}
In (9), $T$ is the tension in the string; again because of the symmetry of the whole hinge structure, the force acting on the end $O_4$ of rod $A_3O_4$ should be $\dfrac{T}{2}$. Substituting (8) into (9), we can solve and obtain
\[
T=18mg.
\]

\textit{Method 2.}\quad This problem can also be solved by means of the principle of virtual work.

Let the point $O_2$ undergo a very small virtual displacement $\delta x$ upwards; then the points $O_3$ and $O_4$ undergo the corresponding virtual displacements $2\delta x$ and $3\delta x$, and the virtual work done by the tension in the string is
\begin{equation}
\delta W=T\cdot3\delta x-T\cdot2\delta x=T\delta x. \tag{10}
\end{equation}
At this time the gravitational potential energy of the whole hinge structure increases by $\delta E_p$:
\begin{align}
\delta E_p&=4mg\,\frac{\delta x}{2}+4mg\left(\delta x+\frac{\delta x}{2}\right)+4mg\left(2\delta x+\frac{\delta x}{2}\right)\notag\\
&=18mg\,\delta x. \tag{11}
\end{align}
According to the principle of virtual work, $\delta W=\delta E_p$; substituting (10) and (11) into it, we immediately obtain
\[
T=18mg.
\]

\begin{commentary}
The principle of virtual work used in Method 2 of this problem may be described roughly as follows.

When a mechanical system is in equilibrium, then for any particle $P_i$ of the system the forces acting on it are the active force $\vec{F}_i$ and the passive force (that is, the constraint reaction) $\vec{N}_i$, and we have
\[
\vec{F}_i+\vec{N}_i=0\qquad(i=1,2,\cdots,N).
\]
If the particle $P_i$ undergoes a virtual displacement (that is, any displacement consistent with the constraints) $\delta\vec{r}_i$, then
\[
\vec{F}_i\cdot\delta\vec{r}_i+\vec{N}_i\cdot\delta\vec{r}_i=0\qquad(i=1,2,\cdots,N).
\]
Adding the equations for all the particles of the system, we obtain
\[
\sum_{i=1}^{N}\vec{F}_i\cdot\delta\vec{r}_i+\sum_{i=1}^{N}\vec{N}_i\cdot\delta\vec{r}_i=0.
\]
For a system with ideal constraints (that is, $\sum\limits_{i=1}^{N}\vec{N}_i\cdot\delta\vec{r}_i=0$) in equilibrium, the equilibrium condition is
\[
\delta W=\sum\vec{F}_i\cdot\delta\vec{r}_i=0,
\]
that is, the equilibrium condition of a mechanical system with ideal constraints is that, for any virtual displacement, the sum of the virtual work done by all the active forces is equal to zero.

For the mechanical system of this problem — the hinge structure — when it is in equilibrium, the active forces acting on it are gravity and the tension of the string, and the total virtual work done by them during a virtual displacement is zero. Here the virtual work done by gravity is equal to the change of its gravitational potential energy. Hence the conclusion is obtained that the virtual work done by the tension of the string is equal to the change of the gravitational potential energy.
\end{commentary}
\end{solution}

\begin{solution}{1.21}
When the sphere is in equilibrium, the external forces acting on it are as shown in the figure. From the equilibrium of the moments of these external forces about the centre $C$ of the sphere we obtain
\begin{equation}
f_1=f_2=f. \tag{1}
\end{equation}
From the vanishing of the resultant of the external forces in the horizontal and in the vertical direction we obtain respectively
\begin{align}
N_1-N_2\sin\theta-f_2\cos\theta&=0, \tag{2}\\
N_2\cos\theta-mg-f_1-f_2\sin\theta&=0. \tag{3}
\end{align}
Solving (1), (2) and (3) simultaneously, we obtain
\begin{equation}
N_1=\frac{1+\sin\theta}{\cos\theta}f+mg\tan\theta. \tag{4}
\end{equation}

\begin{figbox}{1.21}
\begin{tikzpicture}[scale=0.95,>=latex]
  % vertical board hinged at O
  \draw[thickline] (0,0)--(0,6.0);
  \fill (0,0) circle (1.3pt);
  \node[below left] at (0,0) {$O$};
  % inclined plane through O
  \draw[thickline] (0,0)--(5.2,3.64);
  % horizontal reference at O
  \draw[dashedline] (0,0)--(2.3,0);
  % angle theta
  \draw[thickline] (0.95,0) arc (0:35:0.95);
  \node at (1.0,0.33) {$\theta$};
  % sphere
  \draw[thickline] (1.6,3.074) circle (1.6);
  \fill (1.6,3.074) circle (1.3pt);
  \node[above left] at (1.62,3.10) {$C$};
  % contact points
  \fill (0,3.074) circle (1.3pt) node[above left] {$A$};
  \fill (2.518,1.759) circle (1.3pt) node[below right] {$B$};
  % normal N1 from the board
  \draw[force] (0,3.074)--(1.30,3.074) node[above] {$N_1$};
  % friction f1 at A (downward along the board)
  \draw[force] (0,3.074)--(0,1.95) node[left] {$f_1$};
  % normal N2 at B (perpendicular to the plane, into the sphere)
  \draw[force] (2.518,1.759)--(1.83,2.74) node[right] {$N_2$};
  % friction f2 at B (down the plane)
  \draw[force] (2.518,1.759)--(1.21,0.84) node[below left] {$f_2$};
  % weight
  \draw[force] (1.6,3.074)--(1.6,1.57) node[below] {$mg$};
  % radius
  \draw[dashedline] (1.6,3.074)--(2.518,1.759);
  \node[below right] at (2.70,1.60) {$R$};
  % angle alpha between board and OC
  \draw[dashedline] (0,0)--(1.6,3.074);
  \draw[thickline] (0,1.15) arc (90:62.5:1.15);
  \node at (0.30,1.22) {$\alpha$};
  % force F at the top of the board
  \draw[force] (0,6.0)--(1.35,6.0) node[above] {$F$};
  % length l
  \draw[dim] (-0.45,0)--(-0.45,6.0) node[midway,left] {$l$};
\end{tikzpicture}
\end{figbox}

Also, from (1) and (3) we obtain
\begin{equation}
N_2=\frac{1+\sin\theta}{\cos\theta}f+\frac{mg}{\cos\theta}. \tag{5}
\end{equation}

(i) If $N_1\mu_1\geq N_2\mu_2$: when $F$ increases until the equilibrium of the system is just broken, relative sliding begins to occur at the contact point $B$ between the sphere and the inclined plane, while at the contact point $A$ between the sphere and the board the maximum tendency to relative sliding has not yet been reached; that is, at this moment $f=N_2\mu_2$. Substituting this relation into (4) and (5), we obtain
\begin{align}
N_2&=\frac{mg}{\cos\theta-\mu_2(1+\sin\theta)}, \tag{6}\\
N_1&=\frac{(\sin\theta+\mu_2\cos\theta)mg}{\cos\theta-\mu_2(1+\sin\theta)}. \tag{7}
\end{align}
Substituting (6) and (7) into $N_1\mu_1\geq N_2\mu_2$, we obtain
\[
\frac{(\sin\theta+\mu_2\cos\theta)mg\mu_1}{\cos\theta-\mu_2(1+\sin\theta)}
\geq\frac{mg\mu_2}{\cos\theta-\mu_2(1+\sin\theta)},
\]
from which the condition for this case follows:
\begin{equation}
\mu_1\geq\frac{\mu_2}{\sin\theta+\mu_2\cos\theta}. \tag{8}
\end{equation}

(ii) If $N_1\mu_1<N_2\mu_2$, then when the equilibrium of the system is just broken it is at the contact point $A$ between the sphere and the board that relative sliding occurs, while at the contact point $B$ the maximum tendency to relative sliding has not yet been reached; at this moment $f=N_1\mu_1$. Substituting this relation into (4) and (5), we obtain
\begin{align}
N_1&=\frac{mg\sin\theta}{\cos\theta-\mu_1(1+\sin\theta)}, \tag{9}\\
N_2&=\frac{1-\mu_1\cos\theta}{\cos\theta-\mu_1(1+\sin\theta)}\,mg. \tag{10}
\end{align}
Substituting (9) and (10) into $N_1\mu_1<N_2\mu_2$, we obtain
\[
\frac{mg\sin\theta\,\mu_1}{\cos\theta-\mu_1(1+\sin\theta)}
<\frac{(1-\mu_1\cos\theta)\mu_2\,mg}{\cos\theta-\mu_1(1+\sin\theta)},
\]
from which the condition for this case follows:
\begin{equation}
\mu_1<\frac{\mu_2}{\sin\theta+\mu_2\cos\theta}. \tag{11}
\end{equation}
From the equilibrium of the moments of the external forces acting on the board about point $O$ we obtain
\begin{equation}
Fl=N_1\cdot AO. \tag{12}
\end{equation}
From the geometrical relations shown in the figure we obtain
\begin{equation}
AO=R\cot\alpha=R\cot\left(\frac{\pi}{4}-\frac{\theta}{2}\right). \tag{13}
\end{equation}
Substituting (13) into (12) gives
\begin{equation}
F=\frac{R}{l}\cot\left(\frac{\pi}{4}-\frac{\theta}{2}\right)N_1. \tag{14}
\end{equation}
It is only necessary to substitute (7) and (9) respectively into (14) to obtain the value of $F$ in the two cases. The results may therefore be summarized as follows.

(i) If $\mu_1\geq\dfrac{\mu_2}{\sin\theta+\mu_2\cos\theta}$, then
\[
F=\frac{R(\sin\theta+\mu_2\cos\theta)}{l\left[\cos\theta-\mu_2(1+\sin\theta)\right]}\,mg\cot\left(\frac{\pi}{4}-\frac{\theta}{2}\right).
\]
(ii) If $\mu_1<\dfrac{\mu_2}{\sin\theta+\mu_2\cos\theta}$, then
\[
F=\frac{R}{l}\cdot\frac{\sin\theta}{\cos\theta-\mu_1(1+\sin\theta)}\,mg\cot\left(\frac{\pi}{4}-\frac{\theta}{2}\right).
\]

\begin{commentary}
In this problem the friction coefficients $\mu_1$ and $\mu_2$ at the two contact points $A$ and $B$ are different, so when $F$ increases until the equilibrium of the system is just broken, which of the contact points begins to slide relatively is determined by the relative magnitudes of $\mu_1$ and $\mu_2$. Therefore two cases must be discussed in the solution, and the condition and the corresponding value of $N_1$ must be found separately for each case.
\end{commentary}

\end{solution}

\begin{solution}{1.22}
(1) The upper ends of the three rods lean against one another, as shown in Fig. 1. By symmetry, the resultant $F$ of the forces exerted by the other two rods on the upper end of any one rod (for instance rod $OA$) must be in the horizontal direction, as shown in Fig. 2. In Fig. 1, $D$ is the midpoint of $BC$, and
\[
AD=OD=\frac{\sqrt{3}}{2}l,
\]
\begin{equation}
\cos\alpha=\frac{\sqrt{3}}{3}. \tag{1}
\end{equation}
From the equilibrium of the moments of the external forces acting on rod $OA$ about point $A$ we obtain
\begin{equation}
Fl\sin\alpha-mg\,\frac{l}{2}\cos\alpha=0. \tag{2}
\end{equation}
Substituting (1) into (2), we obtain
\[
F=\frac{\sqrt{2}}{4}mg.
\]

\begin{figbox}{1.22(a)}
\begin{tikzpicture}[scale=0.95,>=latex]
  % base triangle (dashed) and auxiliary lines
  \draw[dashedline] (-2.2,0)--(1.9,-0.45);
  \draw[dashedline] (1.9,-0.45)--(2.2,0.75);
  \draw[dashedline] (-2.2,0)--(2.2,0.75);
  \draw[dashedline] (-2.2,0)--(2.05,0.15);
  \draw[dashedline] (0.6,3.4)--(2.05,0.15);
  % rods
  \draw[thickline] (0.6,3.4)--(-2.2,0);
  \draw[thickline] (0.6,3.4)--(1.9,-0.45);
  \draw[thickline] (0.6,3.4)--(2.2,0.75);
  % labels
  \node[above] at (0.6,3.42) {$O$};
  \node[left] at (-2.2,0) {$A$};
  \node[below] at (1.9,-0.47) {$B$};
  \node[right] at (2.22,0.78) {$C$};
  \node[right] at (2.05,0.15) {$D$};
\end{tikzpicture}
\end{figbox}

\begin{figbox}{1.22(b)}
\begin{tikzpicture}[scale=0.95,>=latex]
  % rod OA
  \draw[thickline] (0,0)--(2.8,3.2);
  \node[above right] at (2.8,3.2) {$O$};
  \node[below left] at (0,0) {$A$};
  % horizontal line through A
  \draw[thickline] (0,0)--(4.6,0);
  \node[right] at (4.6,0) {$D$};
  % force F at O (horizontal)
  \draw[force] (2.8,3.2)--(1.6,3.2) node[above] {$F$};
  % weight at midpoint
  \draw[force] (1.4,1.6)--(1.4,0.55) node[right] {$mg$};
  \node[left] at (1.35,1.62) {$\dfrac{l}{2}$};
  % angle alpha at A
  \draw[thickline] (1.15,0) arc (0:48.8:1.15);
  \node at (1.05,0.42) {$\alpha$};
\end{tikzpicture}
\end{figbox}

(2) After a small ball of mass $m$ has been fixed at the midpoint of rod $OA$, the forces on the three rods all change, and the configuration is no longer symmetric; however, rods $OB$ and $OC$ are in the same situation. The force at the upper end of these two rods may be regarded as, besides the original force $F$, an additional force $T_B$ and $T_C$ respectively, with $T_B=T_C$. Since these two rods are still in equilibrium, $T_B$ and $T_C$ must be along the directions of their respective rods; hence the resultant of these two forces is along $OD$, and its reaction $T$ acts on the upper end of rod $OA$, as shown in Fig. 3. From the vanishing of the total moment of $T$ and of the weight of the small ball about point $A$, we obtain
\[
Tl\sin\alpha-mg\,\frac{l}{2}\cos\alpha=0.
\]
Solving, we obtain
\[
T=\frac{\sqrt{2}}{4}mg.
\]

From the vector relation of $F$ and $T$ shown in Fig. 3 we obtain the force $F_A$ acting on the upper end of rod $OA$ as
\begin{align*}
F_A&=2F\cos\left(\frac{\pi}{2}-\alpha\right)=2F\sin\alpha
=2\times\frac{\sqrt{2}}{4}mg\sqrt{\frac{2}{3}}\\
&=\frac{\sqrt{3}}{3}mg.
\end{align*}

\begin{figbox}{1.22(c)}
\begin{tikzpicture}[scale=0.9,>=latex]
  % horizontal line AD
  \draw[thickline] (0,0)--(4.4,0);
  % rod OA
  \draw[thickline] (0,0)--(2.933,4.148);
  \node[below left] at (0,0) {$A$};
  \node[right] at (4.4,0) {$D$};
  \node[above right] at (2.933,4.148) {$O$};
  % forces at O
  \draw[force] (2.933,4.148)--(1.33,4.148) node[above] {$F$};
  \draw[force] (2.933,4.148)--(2.40,5.66) node[right] {$T$};
  \draw[force] (2.933,4.148)--(0.80,5.66) node[above left] {$F_A$};
  % dashed parallelogram
  \draw[dashedline] (1.33,4.148)--(0.80,5.66);
  \draw[dashedline] (2.40,5.66)--(0.80,5.66);
  % angle alpha at O
  \draw[thickline] (2.933,4.148)++(234.7:0.95) arc (234.7:289.5:0.95);
  \node at (2.72,3.12) {$\alpha$};
\end{tikzpicture}
\end{figbox}

(3) From the vanishing of the resultant of the forces acting on rod $OA$ in the vertical and in the horizontal direction we obtain respectively
\begin{align}
N_A&=2mg-T\sin(\pi-2\alpha), \tag{1}\\
f_A&=F+T\cos(\pi-2\alpha). \tag{2}
\end{align}
Substituting $T=F=\dfrac{\sqrt{2}}{4}mg$ into (1) and (2), we obtain
\begin{align}
N_A&=\frac{5}{3}mg, \tag{3}\\
f_A&=\frac{\sqrt{2}}{3}mg. \tag{4}
\end{align}
Substituting (3) and (4) into $f_A\leq N_A\mu_A$, we obtain
\begin{equation}
\mu_A\geq\frac{f_A}{N_A}=\frac{\sqrt{2}}{5}. \tag{5}
\end{equation}
The forces acting on rod $OB$ are shown in Fig. 4. From the vanishing of the resultant of the forces on this rod in the vertical and in the horizontal direction we obtain respectively
\begin{align}
N_B&=mg+T_B\sin\alpha, \tag{6}\\
f_B&=F+T_B\cos\alpha. \tag{7}
\end{align}

\begin{figbox}{1.22(d)}
\begin{tikzpicture}[scale=0.9,>=latex]
  % rod OB
  \draw[thickline] (0,0)--(2.933,4.148);
  \node[below left] at (0,0) {$B$};
  \node[above right] at (2.933,4.148) {$O$};
  % horizontal reference at O
  \draw[dashedline] (2.933,4.148)--(0.95,4.148);
  % force F at O
  \draw[force] (2.933,4.148)--(1.33,4.148) node[above] {$F$};
  % force TB along the rod (from O towards B)
  \draw[force] (2.933,4.148)--(1.79,2.55) node[right] {$T_B$};
  % angle alpha at O
  \draw[thickline] (2.933,4.148)++(234.7:0.95) arc (234.7:180:0.95);
  \node at (2.02,3.88) {$\alpha$};
  % forces at B
  \draw[force] (0,0)--(0,1.5) node[left] {$N_B$};
  \draw[force] (0,0)--(1.6,0) node[below] {$f_B$};
  % weight at the midpoint
  \draw[force] (1.467,2.074)--(1.467,1.0) node[right] {$mg$};
\end{tikzpicture}
\end{figbox}

From the vector relations shown in Fig. 5, the relation between $T_B$, $T_C$ and $T$ is
\[
T=2T_B\cos 30\degree,
\]
that is,
\begin{equation}
T_B=\frac{1}{\sqrt{3}}T=\frac{1}{\sqrt{3}}\times\frac{\sqrt{2}}{4}mg=\frac{\sqrt{6}}{12}mg. \tag{8}
\end{equation}

\begin{figbox}{1.22(e)}
\begin{tikzpicture}[scale=0.95,>=latex]
  % base triangle and auxiliary lines
  \draw[dashedline] (-2.2,0)--(1.9,-0.45);
  \draw[dashedline] (1.9,-0.45)--(2.2,0.75);
  \draw[dashedline] (-2.2,0)--(2.2,0.75);
  \draw[dashedline] (-2.2,0)--(2.05,0.15);
  \draw[dashedline] (0.6,3.4)--(2.05,0.15);
  % rods
  \draw[thickline] (0.6,3.4)--(-2.2,0);
  \draw[thickline] (0.6,3.4)--(1.9,-0.45);
  \draw[thickline] (0.6,3.4)--(2.2,0.75);
  % vectors at O
  \draw[force] (0.6,3.4)--(0.02,5.10) node[right] {$T_B$};
  \draw[force] (0.6,3.4)--(-0.33,4.94) node[left] {$T_C$};
  \draw[force] (0.6,3.4)--(-0.66,6.23) node[above] {$T$};
  % labels
  \node[right] at (0.62,3.38) {$O$};
  \node[left] at (-2.2,0) {$A$};
  \node[below] at (1.9,-0.47) {$B$};
  \node[right] at (2.22,0.78) {$C$};
  \node[right] at (2.05,0.15) {$D$};
\end{tikzpicture}
\end{figbox}

Substituting (8) into (6) and (7), we obtain
\begin{align}
N_B&=mg+\frac{\sqrt{6}}{12}mg\times\sqrt{\frac{2}{3}}=\frac{7}{6}mg, \tag{9}\\
f_B&=\frac{\sqrt{2}}{4}mg+\frac{\sqrt{6}}{12}mg\times\frac{1}{\sqrt{3}}=\frac{\sqrt{2}}{3}mg. \tag{10}
\end{align}
Substituting (9) and (10) into $f_B\leq\mu_BN_B$, we obtain
\begin{equation}
\mu_B\geq\frac{f_B}{N_B}=\frac{2\sqrt{2}}{7}. \tag{11}
\end{equation}
Since rods $B$ and $C$ are in exactly the same situation as regards the forces, the minimum coefficient of friction required for the equilibrium of rod $C$ is equal to that of rod $B$. Comparing (5) with (11), we obtain that the coefficient of friction between the rods and the ground must satisfy
\[
\mu\geq\frac{2\sqrt{2}}{7}.
\]

\begin{commentary}
In the solution of part (3) it is very easy to omit the discussion and calculation of the minimum coefficient of friction required for the equilibrium of rod $OB$ or $OC$, and to suppose mistakenly that, since the small ball is fixed at the midpoint of rod $OA$, the system is certainly in equilibrium as long as rod $OA$ can remain in equilibrium. That leads to the erroneous conclusion that it is sufficient to satisfy $\mu\geq\dfrac{\sqrt{2}}{5}$.
\end{commentary}

\end{solution}

\begin{solution}{1.23}
Through the centre $O$ of the ring establish a rectangular coordinate system: the $x$-axis is along the direction $PA$, the $yOz$ plane is the plane of the ring, the $y$-axis is along the horizontal direction, and the $z$-axis is along the vertical direction, as shown in the figure.

Since the direction of the supporting force $\vec{N}$ exerted by the ring on the rod is perpendicular both to the rod and to the ring, its component in the plane of the ring must be along the radial direction; denote it by $N_1$. Its other component $N_2$ is along the negative direction of the $x$-axis. Moreover, the direction of $N$ ($=N_1+N_2$) is perpendicular to the rod, hence the resultant of the component of $N_1$ in the plane $PAB$ and of $N_2$ is perpendicular to the rod; that is,
\begin{equation}
N_1\cos\theta\sin\alpha=N_2\cos\alpha. \tag{1}
\end{equation}
From the geometrical relations shown in the figure,
\begin{equation}
\tan\alpha=\frac{2}{3}. \tag{2}
\end{equation}
From (1) and (2) we obtain
\begin{equation}
N_1=3N_2. \tag{3}
\end{equation}
The direction of the friction force $\vec{f}$ acting on the rod is determined by the following analysis. First, the rod has a tendency to slide downwards relative to point $B$ in the plane of the ring (the $yOz$ plane), so the component of $\vec{f}$ in this plane must be along the upward tangent at point $B$; denote it by $f_1$. Secondly, the ring has a tendency to move in the direction $BA$ relative to the rod, so the other component $f_2$ of $\vec{f}$ must be along the direction $BA$, as shown in the figure. Moreover, the direction of $f$ ($=f_1+f_2$) is perpendicular both to the rod and to the direction of $N$. From this it follows that the component of $f_1$ along the rod must be equal and opposite to $f_2$, so that the two cancel each other; that is,
\begin{equation}
f_1\sin\theta\sin\alpha=f_2. \tag{4}
\end{equation}
Substituting $\sin\alpha=\dfrac{2}{\sqrt{13}}$ and $\sin\theta=\dfrac{\sqrt{3}}{2}$ into (4), we obtain
\begin{equation}
f_1=\sqrt{\frac{13}{3}}\,f_2. \tag{5}
\end{equation}

\begin{figbox}{1.23}
\begin{tikzpicture}[scale=1.35,>=latex]
  % projection: (x,y,z) -> x*(1,0) + y*(0.65,0.30) + z*(0,1)
  % ring in the yOz plane
  \draw[thickline] plot[domain=0:360,samples=120]
     ({0.65*cos(\x)},{0.30*cos(\x)+sin(\x)});
  % axes
  \draw[axis] (0,0)--(2.35,0) node[right] {$x$};
  \draw[axis] (0,0)--(1.27,0.585) node[above right] {$y$};
  \draw[axis] (0,0)--(0,1.75) node[above] {$z$};
  \node[below left] at (0,0) {$O$};
  % radius to B
  \coordinate (B) at (0.5629,0.6579);
  \draw[dashedline] (0,0)--(B);
  \node[left] at (0.30,0.40) {$R$};
  % angle theta between OB and the z axis
  \draw[thickline] (0,0.95) arc (90:49.4:0.95);
  \node at (0.22,0.85) {$\theta$};
  % rod through P, B, A
  \coordinate (A) at (1.70,1.03);
  \coordinate (P) at (-0.786,0.321);
  \draw[thickline] (P)--(A);
  \node[below left] at (P) {$P$};
  \node[right] at (A) {$A$};
  \fill (B) circle (1.2pt);
  \node[below right] at (B) {$B$};
  % length 3R/2 along the rod
  \node[above] at (1.28,1.02) {$\frac{3}{2}R$};
  % forces at B
  \draw[force] (B)--($(B)+(0.62,0.72)$) node[above right] {$N_1$};
  \draw[force] (B)--($(B)+(-0.72,0.02)$) node[below] {$N_2$};
  \draw[force] (B)--($(B)+(-0.34,0.65)$) node[left] {$f_1$};
  \draw[force] (B)--($(B)+(0.42,0.12)$) node[below right] {$f_2$};
  % angle alpha at B between the rod and the x direction
  \draw[thickline] ($(B)+(0.75,0)$) arc (0:15.7:0.75);
  \node at ($(B)+(0.82,0.13)$) {$\alpha$};
  % z' axis through A
  \draw[dashedline] (A)--($(A)+(0,0.95)$) node[above] {$z'$};
\end{tikzpicture}
\end{figbox}

Through point $A$ establish a $z'$ axis parallel to the $z$-axis, as shown in the figure. From the equilibrium of the moments of the external forces acting on the rod about the $z'$ axis (since the weight of the rod is along the vertical direction, its moment about the $z'$ axis is zero), we obtain
\begin{equation}
N_1\sin\theta\times\frac{3}{2}R+N_2R\sin\theta=f_1\cos\theta\times\frac{3}{2}R. \tag{6}
\end{equation}
Substituting (3) into (6), we obtain
\begin{equation}
N_2=\frac{\sqrt{3}}{11}f_1. \tag{7}
\end{equation}
Substituting (3) into $N=\sqrt{N_1^2+N_2^2}$, we obtain
\begin{equation}
N=\sqrt{10}\,N_2. \tag{8}
\end{equation}
From the above analysis the value of the friction force $f$ is obtained as
\begin{align}
f&=\sqrt{(f_1\sin\theta\cos\alpha)^2+(f_1\cos\theta)^2}\notag\\
&=\sqrt{\frac{10}{13}}\,f_1. \tag{9}
\end{align}
Combining (7), (8) and (9), we obtain
\[
f=\frac{11}{\sqrt{39}}N\leq N\mu,
\]
from which we obtain
\[
\mu\geq\frac{11}{\sqrt{39}}.
\]

\begin{commentary}
This problem asks for the minimum coefficient of friction required for the rod to be in equilibrium at the position $\theta=60\degree$. For this purpose the relation between the supporting force $N$ and the friction force $f$ exerted by the ring on the rod at this position must first be found, and the correct analysis and judgement of the directions of $N$ and $f$ is the key to obtaining this relation: since the direction of the supporting force between two contact surfaces is perpendicular to both surfaces, the direction of $N$ is perpendicular both to the rod and to the ring; and the direction of $f$ must be perpendicular to $N$, while, because the rod has a tendency to rotate about point $A$, the direction of $f$ is also perpendicular to the rod. Once the directions of $N$ and $f$ are known, the relations between their components can be found, and the relation between these two forces is obtained by requiring the resultant moment of $N$ and $f$ about the vertical axis through $A$ to vanish.
\end{commentary}

\end{solution}

\begin{solution}{1.24}
The speed $v$ with which the water reaches the inner bottom of the garbage can after spraying out from the fountain nozzle is
\begin{equation}
v=\sqrt{v_0^2-2gh}. \tag{1}
\end{equation}
From the continuity of the water stream it follows that the mass $\dfrac{\Delta m}{\Delta t}$ of water that rushes onto the inner bottom of the can per second should be equal to the mass of water sprayed from the fountain nozzle per second; that is,
\begin{equation}
\frac{\Delta m}{\Delta t}=\rho v_0S_0. \tag{2}
\end{equation}
In (2), $\rho$ is the density of water. Half of the water that rushes onto the inner bottom of the can is adsorbed on the inner bottom and flows down along the wall of the can, while the other half splashes vertically downwards with its original speed; hence the change of momentum per second of the water rushing onto the inner bottom is
\begin{equation}
\frac{\Delta p}{\Delta t}=\frac{1}{2}\frac{\Delta m}{\Delta t}\cdot v+\frac{1}{2}\frac{\Delta m}{\Delta t}\cdot 2v. \tag{3}
\end{equation}
Substituting (2) into (3), we obtain
\begin{equation}
\frac{\Delta p}{\Delta t}=\frac{3}{2}\rho S_0v_0v. \tag{4}
\end{equation}

\begin{figbox}{1.24}
\begin{tikzpicture}[scale=0.95,>=latex]
  % ground
  \draw[thickline] (-3.6,0)--(3.6,0);
  \fill[hatch] (-3.6,-0.22) rectangle (3.6,0);
  % inverted can: closed bottom at the top, open mouth at the bottom
  \draw[thickline] (-1.35,4.55)--(-1.05,1.55);
  \draw[thickline] (1.35,4.55)--(1.05,1.55);
  \draw[thickline] (-1.45,4.55)--(1.45,4.55);
  \draw[thickline] (-1.35,4.70)--(1.35,4.70);
  \draw[thickline] (-1.35,4.55)--(-1.35,4.70);
  \draw[thickline] (1.35,4.55)--(1.35,4.70);
  \draw[thickline] (-1.05,1.55)--(1.05,1.55);
  % nozzle
  \draw[thickline] (-0.22,0)--(-0.22,0.55);
  \draw[thickline] (0.22,0)--(0.22,0.55);
  \fill[hatch] (-0.22,0) rectangle (0.22,0.12);
  % jet
  \draw[force] (0,0.55)--(0,2.2) node[midway,right] {$v_0$};
  \draw[force] (0,2.6)--(0,4.32) node[midway,right] {$v$};
  % splashing water inside the can
  \draw[thickline,domain=0:1.0,samples=40] plot ({0.15+1.05*\x},{4.32-1.1*\x*\x-0.25*\x});
  \draw[thickline,domain=0:1.0,samples=40] plot ({-0.15-1.05*\x},{4.32-1.1*\x*\x-0.25*\x});
  \draw[thickline,domain=0:0.7,samples=30] plot ({0.10+0.5*\x},{4.32-2.6*\x*\x});
  \draw[thickline,domain=0:0.7,samples=30] plot ({-0.10-0.5*\x},{4.32-2.6*\x*\x});
  % dimension h
  \draw[dim] (2.55,0)--(2.55,4.62) node[midway,right] {$h$};
  % label
  \node[below] at (0,-0.05) {$S_0$};
\end{tikzpicture}
\end{figbox}

\end{solution}

\begin{solution}{1.25}
Let the density of the raindrop be $\rho$. Take the instant at which the raindrop begins to
fall as $t=0$. At time $t$ the mass of the raindrop is $m(t)$, its velocity $v(t)$, its radius
$r(t)$ and its volume $V(t)=\frac{4}{3}\pi r^3$. On its way down the raindrop adsorbs all the
mist that it encounters, so that its volume and mass increase continually; hence
\begin{equation*}
\frac{\dd m}{\dd t}=\frac{\dd(\rho V)}{\dd t}
=\frac{\dd}{\dd t}\left(\rho\frac{4}{3}\pi r^3\right)
=4\pi r^2\rho\frac{\dd r}{\dd t}
=\rho_0\pi r^2v,
\end{equation*}
where $\rho_0$ is the density of the mist. From the above equation we obtain
\begin{equation}
\frac{\dd r}{\dd t}=\frac{\rho_0}{4\rho}v. \tag{1}
\end{equation}

Since the viscous drag and the buoyant force acting on the raindrop during its fall are
neglected, the only external force is gravity; hence, from the momentum theorem,
\begin{equation}
\frac{\dd(mv)}{\dd t}=mg. \tag{2}
\end{equation}
Substituting $m=\rho V$ into (2), we have
\begin{equation*}
\frac{\dd(Vv)}{\dd t}=Vg.
\end{equation*}
The above equation can also be written as
\begin{equation}
\frac{\dd(Vv)}{\dd V}\frac{\dd V}{\dd t}=Vg. \tag{3}
\end{equation}
From (1) we have
$\frac{\dd V}{\dd t}=4\pi r^2\frac{\dd r}{\dd t}=\pi r^2\frac{\rho_0}{\rho}v$; substituting
this into (3), we obtain
\begin{equation}
\frac{\dd(Vv)}{\dd V}\cdot\frac{\pi r^2\rho_0}{\rho}v=Vg. \tag{4}
\end{equation}
Substituting $r=\sqrt[3]{\frac{3V}{4\pi}}$ into (4), we obtain
\begin{equation*}
(Vv)\,\dd(Vv)=\sqrt[3]{\frac{16\rho g}{9\pi\rho_0}}\,V^{4/3}\,\dd V.
\end{equation*}
Integrating both sides of the above equation,
\begin{equation*}
\int_0^{Vv}(Vv)\,\dd(Vv)
=\sqrt[3]{\frac{16\rho g}{9\pi\rho_0}}\int_0^V V^{4/3}\,\dd V,
\end{equation*}
we obtain
\begin{equation*}
\frac{1}{2}(Vv)^2=\frac{3}{7}\sqrt[3]{\frac{16}{9\pi}\frac{\rho g}{\rho_0}}\,V^{7/3}.
\end{equation*}
Hence
\begin{equation*}
v^2=\frac{4}{7}\sqrt[3]{\frac{6}{\pi}\frac{\rho g}{\rho_0}}\,V^{1/3}
=\frac{8\rho g}{7\rho_0}r,
\end{equation*}
that is,
\begin{equation}
v=\sqrt{\frac{8\rho g}{7\rho_0}r}. \tag{5}
\end{equation}

Substituting (5) into (1), we obtain
\begin{equation*}
\frac{\dd r}{\dd t}=\frac{\rho_0}{4\rho}\sqrt{\frac{8\rho g}{7\rho_0}r}
=\sqrt{\frac{\rho_0 g}{14\rho}r}.
\end{equation*}
Multiplying both sides of the above equation by $\frac{\dd t}{\sqrt r}$ and then integrating
both sides,
\begin{equation*}
\int_0^r\frac{\dd r}{\sqrt r}=\sqrt{\frac{\rho_0 g}{14\rho}}\int_0^t\dd t,
\end{equation*}
we obtain
\begin{equation*}
2\sqrt r=\sqrt{\frac{\rho_0 g}{14\rho}}\,t,
\end{equation*}
that is,
\begin{equation}
r=\frac{\rho_0 g}{56\rho}t^2. \tag{6}
\end{equation}
Substituting (6) into (5), we obtain
\begin{equation}
v=\sqrt{\frac{8\rho g}{7\rho_0}\times\frac{\rho_0 g}{56\rho}t^2}=\frac{1}{7}gt. \tag{7}
\end{equation}
From (7) we immediately obtain
\begin{equation*}
a=\frac{\dd v}{\dd t}=\frac{1}{7}g.
\end{equation*}

One method of solving this problem has been given above; another method is as follows.

From (2) we obtain
\begin{equation}
3v\frac{\dd r}{\dd t}+r\frac{\dd v}{\dd t}=rg. \tag{8}
\end{equation}
Substituting (1) into (8), we obtain
\begin{equation}
\frac{\dd^2r}{\dd t^2}+\frac{3}{r}\left(\frac{\dd r}{\dd t}\right)^2=\frac{\rho_0}{4\rho}g. \tag{9}
\end{equation}
The particular solution of this non-homogeneous differential equation is
\begin{equation}
r=bt^2. \tag{10}
\end{equation}
Substituting (10) into (9), we can solve for $b$:
\begin{equation}
b=\frac{1}{56}\frac{\rho_0}{\rho}g. \tag{11}
\end{equation}
From (1), (10) and (11) we obtain
\begin{equation}
v=\frac{4\rho}{\rho_0}\frac{\dd r}{\dd t}
=\frac{4\rho}{\rho_0}\times 2bt=\frac{1}{7}gt. \tag{12}
\end{equation}
Hence
\begin{equation*}
a=\frac{\dd v}{\dd t}=\frac{1}{7}g.
\end{equation*}
\end{solution}

\begin{solution}{1.26}
Take the instant at which the block begins to slide down the inclined plane as $t=0$. At time
$t$ let the velocity of the block be $v(t)$, and let the total mass of the block together with
that part of the rope which moves with the block at velocity $v$ (this part having length $x$)
be $M(t)$. Then, from the momentum theorem,
\begin{equation*}
\frac{\dd(Mv)}{\dd t}=Mg\sin\theta.
\end{equation*}
The above equation can be rewritten as
\begin{equation*}
\frac{\dd(Mv)}{\dd x}\frac{\dd x}{\dd t}=Mg\sin\theta.
\end{equation*}
Substituting $v=\frac{\dd x}{\dd t}$ into the above equation and multiplying both sides by
$M\,\dd x$, and noting that $M=m+\lambda x$, we obtain
\begin{equation*}
(Mv)\,\dd(Mv)=(m+\lambda x)^2g\sin\theta\,\dd x.
\end{equation*}

\begin{figbox}{1.26}
\begin{tikzpicture}[scale=0.95]
  % ground
  \draw[thickline] (-0.6,0) -- (5.9,0);
  % wedge with a horizontal shelf at the top
  \draw[thickline] (0,0) -- (4.0,2.2) -- (5.6,2.2) -- (5.6,0);
  \draw[thickline] (0,0) -- (5.6,0);
  % angle of inclination
  \draw (0.85,0) arc (0:28.8:0.85);
  \node at (1.10,0.26) {$\theta$};
  % block on the slope
  \begin{scope}[shift={(2.15,1.1825)},rotate=28.8]
    \draw[thickline,fill=black!12] (-0.40,-0.20) rectangle (0.40,0.20);
  \end{scope}
  \node at (2.15,0.82) {$m$};
  % rope from the block to the top edge and along the shelf
  \draw[thickline] (2.50,1.52) -- (3.98,2.19);
  \draw[thickline] (4.02,2.20) -- (4.45,2.20);
  % coil of rope on the shelf
  \draw[thickline] (4.45,2.20) .. controls (4.45,2.72) and (5.25,2.72) .. (5.25,2.20);
  \draw[thickline] (4.57,2.20) .. controls (4.57,2.60) and (5.13,2.60) .. (5.13,2.20);
  \draw[thickline] (4.69,2.20) .. controls (4.69,2.48) and (5.01,2.48) .. (5.01,2.20);
  \node[right] at (5.28,2.48) {$\lambda$};
\end{tikzpicture}
\end{figbox}

Integrating both sides of the above equation,
\begin{equation*}
\int_0^{Mv}(Mv)\,\dd(Mv)=\int_0^x(m+\lambda x)^2g\sin\theta\,\dd x,
\end{equation*}
we obtain
\begin{equation*}
\frac{1}{2}(Mv)^2=\frac{g}{3\lambda}\sin\theta\left[(m+\lambda x)^3-m^3\right].
\end{equation*}
From this we can solve
\begin{equation*}
v=\sqrt{\frac{2g\sin\theta\left[(m+\lambda x)^3-m^3\right]}{3\lambda(m+\lambda x)^2}}.
\end{equation*}
\end{solution}

\begin{solution}{1.27}
When the extension of the spring is $x_0$, that is, when the spring force is equal to the
weight of the small ball, the ball leaves the table. We have
\begin{equation}
x_0=\frac{mg}{k}. \tag{1}
\end{equation}
After this, take the upper end of the spring, which moves uniformly upward at speed $v_0$, as
the frame of reference. In this inertial frame the ball begins to move downward with initial
speed $v_0$ and does so until it stops; at this instant the extension of the spring reaches its
maximum for the first time. Let this maximum extension be $x$. From the conservation of
mechanical energy of the system,
\begin{equation}
\frac{1}{2}mv_0^2-mgx_0+\frac{1}{2}kx_0^2=\frac{1}{2}kx^2-mgx. \tag{2}
\end{equation}
The zero points of both the gravitational potential energy and the elastic potential energy of
the system are taken at the position of the undeformed spring. Substituting (1) into (2), we
obtain
\begin{equation}
kx^2-2mgx+\frac{(mg)^2}{k}-mv_0^2=0. \tag{3}
\end{equation}

\begin{figbox}{1.27}
\begin{tikzpicture}[scale=1.0]
  % table
  \draw[thickline] (-1.5,0) -- (1.5,0);
  \foreach \x in {-1.4,-1.1,...,1.4} {\draw[thinline] (\x,0) -- (\x-0.16,-0.16);}
  % ball
  \shade[ball color=black!30] (0,0.20) circle (0.20);
  \draw[thickline] (0,0.20) circle (0.20);
  \node[right] at (0.24,0.16) {$m$};
  % spring
  \draw[thickline,decorate,decoration={zigzag,segment length=5pt,amplitude=2.5pt}]
    (0,0.40) -- (0,2.10);
  \node[right] at (0.14,1.25) {$k$};
  % upper end of the spring
  \draw[thickline] (-0.45,2.10) -- (0.45,2.10);
  \draw[thickline] (-0.45,1.95) -- (-0.45,2.10);
  \draw[thickline] (0.45,1.95) -- (0.45,2.10);
  % velocity of the upper end
  \draw[force] (0,2.15) -- (0,2.95) node[right] {$v_0$};
\end{tikzpicture}
\end{figbox}

Solving this quadratic equation, a reasonable solution is
\begin{equation}
x=\frac{mg}{k}+v_0\sqrt{\frac{m}{k}}. \tag{4}
\end{equation}
The time elapsed while the extension of the spring changes from $x_0$ to $x$ is equal to one
quarter of the period of the simple harmonic motion of this oscillator, that is,
\begin{equation}
t=\frac{T}{4}=\frac{\pi}{2}\sqrt{\frac{m}{k}}. \tag{5}
\end{equation}
From (5) the height $h$ of the ball above the table at this instant is found to be
\begin{equation}
h=v_0t-(x-x_0)=\left(\frac{\pi}{2}-1\right)v_0\sqrt{\frac{m}{k}}. \tag{6}
\end{equation}
From the work-energy theorem for the system, the work $W$ done by the force acting on the upper
end of the spring during the whole process is
\begin{equation}
W=\frac{1}{2}mv_0^2+mgh+\frac{1}{2}kx^2. \tag{7}
\end{equation}
Substituting (4) and (6) into (7), we obtain
\begin{equation*}
W=mv_0^2+\frac{\pi}{2}mgv_0\sqrt{\frac{m}{k}}+\frac{1}{2}\frac{m^2g^2}{k}.
\end{equation*}

\begin{commentary}
The key to the solution of this problem is that, after the ball leaves the table, one changes
to the frame of reference moving together with the upper end of the spring. This is an inertial
frame; by the principle of relativity of mechanics, the laws of mechanics have the same form in
every inertial frame. In this frame the force acting on the upper end of the spring does no
work, so after the ball leaves the table the mechanical energy of the system is conserved;
equation (2) is thus obtained, and hence the maximum extension $x$ of the spring. Moreover, in
this frame the ball performs simple harmonic motion, so the time $t$ from the ball leaving the
table until the spring first reaches its maximum extension is easily found. In the table frame,
on the other hand, the mechanical energy of the system is not conserved, and the ball does not
perform simple harmonic motion either, so $x$ and $t$ are not easily found.
\end{commentary}
\end{solution}

\begin{solution}{1.28}
Suppose that under the action of the force $F$, after a time $t$, ball $B$ has moved a distance $s$, and at this moment the rod makes an angle $\theta$ with the horizontal, as shown in Fig.~1.28(a). Let the velocity of ball $B$ be $v$, and let the velocity of ball $A$ relative to ball $B$ be $v'$. From the work-energy principle we have
\begin{equation}
Fs=\frac{1}{2}mv^2+\frac{1}{2}m\left[(v-v'\sin\theta)^2+(v'\cos\theta)^2\right]-mgl(1-\sin\theta) \tag{1}
\end{equation}
\begin{figbox}{1.28(a)}
\begin{tikzpicture}[scale=0.95,>=latex]
  % ground
  \fill[hatch] (-0.4,-0.25) rectangle (6.9,0);
  \draw[thickline] (-0.4,0)--(6.9,0);
  % rod
  \coordinate (B) at (4.30,0.18);
  \coordinate (A) at (1.94,2.03);
  \draw[thickline,line width=1.4pt] (B)--(A);
  % balls
  \draw[thickline] (B) circle (0.18);
  \draw[thickline] (A) circle (0.18);
  \node[above right] at (A) {$A$};
  \node[below] at (B) {$B$};
  % angle theta at B
  \draw[dashedline] (B)--++(180:1.0);
  \draw[thickline] ($(B)+(142:0.75)$) arc (142:180:0.75);
  \node at ($(B)+(161:1.02)$) {$\theta$};
  % force F
  \draw[force] ($(B)+(0.18,0)$)--++(0:1.25) node[right] {$\vec{F}$};
  % displacement s
  \draw[dashedline] (2.00,0)--(2.00,-0.80);
  \draw[dashedline] (4.30,0)--(4.30,-0.80);
  \draw[dim] (2.00,-0.55)--(4.30,-0.55) node[midway,below] {$s$};
\end{tikzpicture}
\end{figbox}

The acceleration $a_C$ of the centre of mass of the system along the horizontal direction is
\begin{equation}
a_C=\frac{F}{2m} \tag{2}
\end{equation}
The displacement $x_C$ of the centre of mass along the horizontal direction during this process is
\begin{equation}
x_C=\frac{1}{2}a_Ct^2 \tag{3}
\end{equation}

From the geometric relation shown in Fig.~1.28(a), the displacement $x_C$ of the centre of mass along the horizontal direction is
\begin{equation}
x_C=s-\frac{l}{2}\cos\theta \tag{4}
\end{equation}
From (2), (3) and (4) we obtain
\begin{equation}
t=\sqrt{\frac{2m}{F}(2s-l\cos\theta)} \tag{5}
\end{equation}
From the momentum theorem,
\begin{equation}
Ft=mv+m(v-v'\sin\theta) \tag{6}
\end{equation}
Substituting (5) into (6) gives
\begin{equation}
v=\frac{1}{2}v'\sin\theta+\sqrt{\frac{F}{2m}(2s-l\cos\theta)} \tag{7}
\end{equation}
Substituting (7) into (1), we obtain
\begin{equation}
v'^2=\frac{2Fl\cos\theta+4mgl(1-\sin\theta)}{m(1+\cos^2\theta)} \tag{8}
\end{equation}
Let the force exerted on ball $B$ by the rod be $T$, directed along the rod, and let the normal force exerted on it by the ground be $N$; the forces acting on ball $B$ are then as shown in Fig.~1.28(b). Accordingly, the equations of motion of ball $B$ along the horizontal and vertical directions are respectively
\begin{align}
F+T\cos\theta&=ma \tag{9}\\
N-T\sin\theta-mg&=0 \tag{10}
\end{align}

\begin{figbox}{1.28(b)}
\begin{tikzpicture}[scale=0.95,>=latex]
  % ground
  \fill[hatch] (-2.9,-0.25) rectangle (2.2,0);
  \draw[thickline] (-2.9,0)--(2.2,0);
  % ball B
  \coordinate (B) at (0,0.18);
  \draw[thickline] (B) circle (0.18);
  \node[above right] at (B) {$B$};
  % rod
  \draw[thickline,line width=1.4pt] (B)--++(142:2.4);
  % angle theta
  \draw[dashedline] (B)--++(180:1.0);
  \draw[thickline] ($(B)+(142:0.72)$) arc (142:180:0.72);
  \node at ($(B)+(161:0.99)$) {$\theta$};
  % forces
  \draw[force] ($(B)+(0,0.18)$)--++(90:1.15) node[above] {$\vec{N}$};
  \draw[force] ($(B)+(0.18,0)$)--++(0:1.30) node[right] {$\vec{F}$};
  \draw[force] ($(B)+(0,-0.18)$)--++(-90:1.05) node[below] {$m\vec{g}$};
  \draw[force] (B)--++(-38:1.35) node[right] {$\vec{T}$};
\end{tikzpicture}
\end{figbox}
Take ball $B$ as the reference frame; then the forces acting on ball $A$ are as shown in Fig.~1.28(c). In the figure the inertial force is $f_i=ma$, where $a$ is the acceleration of ball $B$ at this moment. Hence the normal equation of motion of ball $A$, which moves in a circle about $B$ in this frame, is
\begin{equation}
mg\sin\theta-T-ma\cos\theta=m\frac{v'^2}{l} \tag{11}
\end{equation}

\begin{figbox}{1.28(c)}
\begin{tikzpicture}[scale=0.95,>=latex]
  % ball A
  \coordinate (A) at (0,0);
  \draw[thickline] (A) circle (0.18);
  \node[right] at ($(A)+(0.18,0)$) {$A$};
  % rod towards B (lower right)
  \draw[thickline,line width=1.4pt] (A)--++(-38:1.85);
  % forces on A
  \draw[force] (A)--++(142:2.05) node[above left] {$\vec{T}$};
  \draw[force] ($(A)+(-0.18,0)$)--++(180:1.35) node[left] {$\vec{f}_i$};
  \draw[force] ($(A)+(0,-0.18)$)--++(-90:1.15) node[below] {$m\vec{g}$};
  % angle theta
  \draw[dashedline] (A)--++(180:0.95);
  \draw[thickline] ($(A)+(142:0.70)$) arc (142:180:0.70);
  \node at ($(A)+(161:0.97)$) {$\theta$};
\end{tikzpicture}
\end{figbox}
Substituting (8) and (9) into (11) gives
\begin{equation}
T=\frac{mg(6\sin\theta-\sin^3\theta-4)-F\cos\theta(\cos^2\theta+3)}{(1+\cos^2\theta)^2} \tag{12}
\end{equation}
Substituting (12) into (10), we obtain the force $N$ exerted by the horizontal surface on ball $B$ at this moment:
\begin{align*}
N&=mg+T\sin\theta\\
&=mg\left[1+\frac{\sin\theta(6\sin\theta-\sin^3\theta-4)-\dfrac{F}{mg}\sin\theta\cos\theta(\cos^2\theta+3)}{(1+\cos^2\theta)^2}\right].
\end{align*}

\begin{commentary}
The system in this problem consists of two small balls connected by the two ends of a light rod. Whatever the motion of such a system, the force exerted by the rod on a ball at any instant must be directed along the rod; otherwise its reaction, acting on the two ends of the light rod, would subject the light rod to a torque. This contradicts the conclusion, obtained from the law of rotation, that the resultant torque acting on a rigid body of zero moment of inertia is zero. Moreover, in order to find the force exerted by the rod on ball $B$, one may discuss the motion of ball $A$ in the frame of ball $B$ (a non-inertial frame). The velocity of ball $A$ relative to ball $B$ can then be found by applying, in the inertial frame, such principles as the law of motion of the centre of mass, the work-energy principle and the momentum theorem.
\end{commentary}

\end{solution}

\begin{solution}{1.29}
Suppose that when the angle between the string and the rod is $\theta$, the velocity of the small ring $A$ is $v_A$, and the velocity of the small ball $B$ relative to the small ring is $v'$. Then the velocity $v_B$ of ball $B$ relative to the ground is given by the formula for relative motion, namely
\begin{equation}
v_B=v'+v_A \tag{1}
\end{equation}
The vector relation of (1) is shown in Fig.~1.29(a). From the figure we obtain
\begin{equation}
v_{Bx}=v_A-v'\sin\theta \tag{2}
\end{equation}
\begin{equation}
v_B=\sqrt{(v_A-v'\sin\theta)^2+(v'\cos\theta)^2} \tag{3}
\end{equation}
From the conservation of horizontal momentum and the conservation of mechanical energy of the system we obtain respectively
\begin{equation}
mv_A+mv_{Bx}=0 \tag{4}
\end{equation}
\begin{equation}
\frac{1}{2}mv_A^2+\frac{1}{2}mv_B^2-mgl\sin\theta=0 \tag{5}
\end{equation}
where the zero of gravitational potential energy is taken at the point from which ball $B$ is released. Substituting (2) and (3) into (4) and (5) respectively gives
\begin{equation}
mv_A+m(v_A-v'\sin\theta)=0 \tag{6}
\end{equation}
\begin{equation}
\frac{1}{2}mv_A^2+\frac{1}{2}m\left[(v_A-v'\sin\theta)^2+(v'\cos\theta)^2\right]-mgl\sin\theta=0 \tag{7}
\end{equation}
Solving (6) and (7) simultaneously, we obtain
\begin{equation}
v'^2=\frac{4gl\sin\theta}{1+\cos^2\theta} \tag{8}
\end{equation}
\begin{equation}
v_A=\frac{1}{2}\sin\theta\sqrt{\frac{4gl\sin\theta}{1+\cos^2\theta}} \tag{9}
\end{equation}
Let the angle between $v_B$ and $v'$ be $\alpha$. From the geometric relation shown in Fig.~1.29(a), and using the sine rule, we have
\begin{equation}
\frac{v_A}{\sin\alpha}=\frac{v'}{\sin\left(\frac{\pi}{2}+\theta-\alpha\right)} \tag{10}
\end{equation}
Substituting (8) and (9) into (10), we obtain
\[
\frac{v'}{v_A}=\frac{\cos(\theta-\alpha)}{\sin\alpha}
=\frac{\cos\theta\cos\alpha+\sin\theta\sin\alpha}{\sin\alpha}
=\frac{2}{\sin\theta}.
\]
Hence we can solve for
\[
\cot\alpha=\frac{1+\cos^2\theta}{\sin\theta\cos\theta},
\]
that is,
\begin{equation}
\begin{cases}
\sin\alpha=\dfrac{\sin\theta\cos\theta}{\sqrt{1+3\cos^2\theta}}\\[2ex]
\cos\alpha=\dfrac{1+\cos^2\theta}{\sqrt{1+3\cos^2\theta}}
\end{cases} \tag{11}
\end{equation}

\begin{figbox}{1.29(a)}
\begin{tikzpicture}[scale=0.95,>=latex]
  % rod (rail) drawn as a double line
  \draw[thickline] (0.2,3.30)--(6.6,3.30);
  \draw[thickline] (0.2,3.16)--(6.6,3.16);
  % ring A
  \draw[thickline] (1.75,3.05) rectangle (2.25,3.41);
  \node[above] at (2.00,3.41) {$A$};
  \draw[force] (2.25,3.23)--++(0:1.15) node[above] {$v_A$};
  % string
  \coordinate (A) at (2.00,3.23);
  \coordinate (B) at (4.60,1.72);
  \draw[thickline] (A)--(B);
  % angle theta at A
  \draw[thickline] ($(A)+(0:0.85)$) arc (0:-30:0.85);
  \node at ($(A)+(-15:1.12)$) {$\theta$};
  % ball B
  \draw[thickline] (B) circle (0.10);
  \node[above right] at (B) {$B$};
  % direction of the string extended beyond B
  \draw[dashedline] (B)--++(-30:1.35);
  % velocity v' (perpendicular to the string)
  \draw[force] (B)--++(-120:1.35) node[below left] {$v'$};
  % velocity v_B
  \draw[force] (B)--++(-106:1.90) node[below] {$v_B$};
  % theta at B between the vertical and v'
  \draw[dashedline] (B)--++(-90:1.05);
  \draw[thickline] ($(B)+(-120:0.62)$) arc (-120:-90:0.62);
  \node at ($(B)+(-103:0.88)$) {$\theta$};
  % weight
  \draw[force] ($(B)+(0.10,-0.06)$)--++(-90:0.75) node[right] {$m\vec{g}$};
\end{tikzpicture}
\end{figbox}

Let the tension in the string at this moment be $T$, and let the acceleration of the small ring $A$ be $a$. Then
\begin{equation}
T\cos\theta=ma \tag{12}
\end{equation}
Take the small ring $A$ as the reference frame; in this non-inertial frame the forces acting on ball $B$ are as shown in Fig.~1.29(b). In the figure $f_i=ma$. From the figure, the normal equation of motion of ball $B$, which moves in a circle, is
\begin{equation}
T+ma\cos\theta-mg\sin\theta=m\frac{v'^2}{l} \tag{13}
\end{equation}
Substituting (8) into (13) gives
\begin{equation}
T=\frac{5+\cos^2\theta}{1+\cos^2\theta}mg\sin\theta-ma\cos\theta \tag{14}
\end{equation}
Solving (12) and (14) simultaneously, we obtain
\begin{equation}
T=\frac{5+\cos^2\theta}{(1+\cos^2\theta)^2}mg\sin\theta \tag{15}
\end{equation}

\begin{figbox}{1.29(b)}
\begin{tikzpicture}[scale=0.95,>=latex]
  \coordinate (B) at (0,0);
  \draw[thickline] (B) circle (0.12);
  \node[right] at ($(B)+(0.12,0)$) {$B$};
  % string
  \draw[thickline] (B)--++(150:2.0);
  % tension
  \draw[force] (B)--++(150:2.0) node[above left] {$\vec{T}$};
  % angle theta
  \draw[dashedline] (B)--++(180:1.05);
  \draw[thickline] ($(B)+(150:0.75)$) arc (150:180:0.75);
  \node at ($(B)+(165:1.02)$) {$\theta$};
  % inertial force
  \draw[force] ($(B)+(-0.12,0)$)--++(180:1.30) node[left] {$\vec{f}_i$};
  % weight
  \draw[force] ($(B)+(0,-0.12)$)--++(-90:1.10) node[below] {$m\vec{g}$};
\end{tikzpicture}
\end{figbox}

(1) The direction of the velocity $v_B$ of ball $B$ relative to the ground is just the direction of the tangent to its trajectory at that point. When the speed of the ball reaches its maximum, the resultant force acting on the ball along this tangential direction must be zero. Hence
\begin{equation}
mg\cos(\theta-\alpha)=T\sin\alpha \tag{16}
\end{equation}
Substituting (11) and (15) into (16), we obtain a quadratic equation in $\cos^2\theta$:
\begin{equation}
3\cos^4\theta+8\cos^2\theta-3=0 \tag{17}
\end{equation}
The solution of this equation is
\[
\cos^2\theta=\frac{1}{3},
\]
that is,
\begin{equation}
\begin{cases}
\cos\theta=\dfrac{\sqrt{3}}{3}\\[2ex]
\sin\theta=\sqrt{\dfrac{2}{3}}
\end{cases} \tag{18}
\end{equation}
Substituting (8) and (9) into (3), and then substituting (18), we obtain the maximum speed $v_{B\max}$ of ball $B$:
\begin{align*}
v_{B\max}&=\sqrt{(v_A-v'\sin\theta)^2+(v'\cos\theta)^2}\\
&=\sqrt{\frac{gl(4-3\sin^2\theta)\sin\theta}{1+\cos^2\theta}}=\sqrt{\frac{\sqrt{6}}{2}gl}.
\end{align*}
(2) Substituting (18) into (15), we obtain the tension in the string when ball $B$ has its maximum speed:
\[
T=\frac{5+\cos^2\theta}{(1+\cos^2\theta)^2}mg\sin\theta=\sqrt{6}mg.
\]
The above is one method of solving this problem; of course, the maximum value of $v_B$ can also be found by the method of differentiation.

\end{solution}

\begin{solution}{1.30}
(1) Suppose that when the small ball has turned through an angle $\theta$ relative to the ring, the speed of the ring is $V$ and the speed of the small ball relative to the ring is $v'$; the speed of the small ball relative to the ground is $v$. The relation between these three velocity vectors is given by the formula for relative motion, namely
\begin{equation}
v=v'+V \tag{1}
\end{equation}
From the geometric relation shown in Fig.~1.30 we obtain
\begin{equation}
v_x=v'\cos\theta-V \tag{2}
\end{equation}
\begin{equation}
v=\sqrt{(v'\cos\theta-V)^2+(v'\sin\theta)^2} \tag{3}
\end{equation}

\begin{figbox}{1.30}
\begin{tikzpicture}[scale=0.9,>=latex]
  % ground
  \fill[hatch] (-2.9,-2.62) rectangle (2.9,-2.28);
  \draw[thickline] (-2.9,-2.28)--(2.9,-2.28);
  % ring (tube)
  \draw[thickline,line width=1pt] (0,0) circle (2.28);
  \draw[thickline,line width=1pt] (0,0) circle (1.95);
  \node[below] at (0,-1.02) {$M$};
  % centre O and x axis
  \fill (0,0) circle (1.1pt);
  \node[below left] at (-0.08,-0.06) {$O$};
  \draw[axis] (0,0)--(1.05,0) node[right] {$x$};
  % radius R
  \draw[dim] (0,0)--(150:2.11);
  \node at (150:2.45) {$R$};
  % ball
  \coordinate (P) at (45:2.11);
  \fill (P) circle (0.13);
  \node[above right] at ($(P)+(0.10,0.10)$) {$m$};
  % radius to the ball and angle theta
  \draw[dashedline] (0,0)--(P);
  \draw[thickline] (0,0.85) arc (90:45:0.85);
  \node at (67:1.12) {$\theta$};
  % velocity v' of the ball relative to the ring
  \draw[force] (P)--++(-45:1.05) node[right] {$v'$};
  % weight
  \draw[force] (P)--++(-90:1.00) node[right] {$m\vec{g}$};
  % velocity V of the ring
  \draw[force] (135:2.11)--++(180:1.05) node[left] {$V$};
\end{tikzpicture}
\end{figbox}

From the conservation of momentum of the system along the horizontal ($x$) direction, and using (2), we obtain
\begin{equation}
m(v'\cos\theta-V)-MV=0 \tag{4}
\end{equation}
Take the centre of the ring as the zero of gravitational potential energy; from the conservation of mechanical energy of the system, and using (3), we obtain
\begin{equation}
mgR=mgR\cos\theta+\frac{1}{2}MV^2+\frac{1}{2}m\left[(v'\cos\theta-V)^2+(v'\sin\theta)^2\right] \tag{5}
\end{equation}
From (4) and (5) we can solve for
\begin{equation}
v'=\sqrt{\frac{2(M+m)gR(1-\cos\theta)}{M+m\sin^2\theta}} \tag{6}
\end{equation}
\begin{equation}
V=\sqrt{\frac{2m^2gR\cos^2\theta(1-\cos\theta)}{(M+m)(M+m\sin^2\theta)}} \tag{7}
\end{equation}
When the leftward speed of the ring reaches its maximum, the small ball separates from the tube of the ring, that is, there is no longer any interaction between the two; at this moment the acceleration of the ring is zero. Now take the ring as the reference frame; at this moment this frame is an inertial frame. In this frame the normal equation of motion of the small ball, which moves in a circle along the tube, is
\begin{equation}
mg\cos\theta=m\frac{v'^2}{R} \tag{8}
\end{equation}
Substituting (6) into (8) gives
\begin{equation}
\frac{M}{m}=\frac{2(1-\cos\theta)-\sin^2\theta\cos\theta}{3\cos\theta-2} \tag{9}
\end{equation}
Substituting $\theta=45\degree$ into (9), we obtain
\begin{equation}
\frac{M}{m}=\sqrt{2}+\frac{1}{2}\approx 1.914 \tag{10}
\end{equation}
Substituting (10) and $\theta=45\degree$ into (7), we obtain the maximum leftward speed $V_{\max}$ of the ring:
\begin{equation*}
V_{\max}=\sqrt{(17\sqrt{2}-24)gR}\approx 0.2\sqrt{gR}
\end{equation*}
(2) Substituting (6) and (7) into (3), we obtain the absolute speed $v$ of the small ball:
\begin{equation}
v=\sqrt{\frac{2gR(1-\cos\theta)\left[M^2+m(2M+m)\sin^2\theta\right]}{(M+m)(M+m\sin^2\theta)}} \tag{11}
\end{equation}
Let the angle between $v$ and the vertical direction be $\alpha$. Then from the geometric relation shown in Fig.~1.30, and from (6) and (7), we obtain
\[
\tan\alpha=\frac{v'\cos\theta-V}{v'\sin\theta}=\cot\theta-\frac{V}{v'\sin\theta}=\frac{M}{M+m}\cot\theta,
\]
that is,
\begin{equation}
\sin\alpha=\frac{M\cos\theta}{\sqrt{M^2+m(m+2M)\sin^2\theta}} \tag{12}
\end{equation}
Since at this moment there is no interaction force between the small ball and the ring, the acceleration of the small ball is equal to the acceleration of gravity, that is, $a=g$; the component of the acceleration along the normal to the trajectory is
\begin{equation}
a_n=a\sin\alpha=g\sin\alpha \tag{13}
\end{equation}
The radius of curvature $\rho$ of the trajectory of the small ball at this point is
\begin{equation}
\rho=\frac{v^2}{a_n}=\frac{v^2}{g\sin\alpha} \tag{14}
\end{equation}
Substituting (11) and (12) into (14), we obtain
\begin{equation}
\rho=\frac{2(1-\cos\theta)\left[M^2+m(2M+m)\sin^2\theta\right]^{\frac{3}{2}}}{M\cos\theta(M+m)(M+m\sin^2\theta)}R \tag{15}
\end{equation}
Substituting $\theta=45\degree$ and (10) into (15), we obtain
\[
\rho=\frac{8}{7}\left(46\sqrt{2}-65\right)\left(\frac{13}{4}+2\sqrt{2}\right)^{\frac{3}{2}}R\approx 0.92R.
\]

\begin{commentary}
The radius of curvature of the trajectory of the small ball at a certain point can be found in this problem either by first finding the equation of the trajectory of the small ball and then calculating the radius of curvature, or, as in the solution above, by first finding the normal acceleration $a_n$ of the small ball at this point and then using the formula $a_n=\dfrac{v^2}{\rho}$ to find the radius of curvature $\rho$. Since the moment when the leftward speed of the ring reaches its maximum is precisely the moment when the small ball and the ring separate from each other, that is, when there is no interaction force, the acceleration of the small ball at this moment is equal to the acceleration of gravity, while the normal to the trajectory at this point is the direction perpendicular to the tangent, i.e. the direction perpendicular to the velocity of the small ball at this moment. Hence, by using the conservation of momentum of the system, the conservation of energy and the principle of relative motion to find this velocity and its direction, the radius of curvature of the trajectory at this point can be obtained.
\end{commentary}

\end{solution}

\begin{solution}{1.31}
(1) From the moment of release until the string connected to $C$ makes an angle $\alpha$ with the vertical line, the height $H$ through which object $C$ descends is
\begin{equation}
H=l\cot\alpha \tag{1}
\end{equation}
while objects $A$ and $B$ rise by $h$:
\begin{equation}
h=\frac{l}{\sin\alpha}-l \tag{2}
\end{equation}
Let the speeds of $A$ (or $B$) and $C$ at this moment be $v$ and $V$ respectively; from the law of conservation of mechanical energy of the system we have
\begin{equation}
MgH-2mgh=\frac{1}{2}MV^2+2\times\frac{1}{2}mv^2 \tag{3}
\end{equation}
and the constraint relation between $v$ and $V$ is
\begin{equation}
v=V\cos\alpha \tag{4}
\end{equation}
Solving (1), (2), (3) and (4) simultaneously, we obtain
\begin{equation}
V=\sqrt{\frac{M\cos\alpha+2m(\sin\alpha-1)}{(M+2m\cos^2\alpha)\sin\alpha}\cdot 2gl} \tag{5}
\end{equation}

\begin{figbox}{1.31}
\begin{tikzpicture}[scale=0.92,>=latex]
  % support line through the pulleys
  \draw[thickline] (0.30,3.40)--(7.70,3.40);
  % pulleys
  \draw[thickline] (1.00,3.40) circle (0.13);
  \draw[thickline] (7.00,3.40) circle (0.13);
  \node[above left] at (0.92,3.55) {$a$};
  \node[above right] at (7.08,3.55) {$b$};
  % dimensions l, l
  \draw[dim] (1.15,3.80)--(3.95,3.80) node[midway,above] {$l$};
  \draw[dim] (4.05,3.80)--(6.85,3.80) node[midway,above] {$l$};
  % initial position of C
  \draw[dashedline] (3.72,3.22) rectangle (4.28,3.58);
  \node at (4.00,3.40) {$M$};
  \node[above left] at (3.74,3.55) {$C$};
  % strings of A and B
  \draw[thickline] (0.87,3.40)--(0.87,1.92);
  \draw[thickline] (7.13,3.40)--(7.13,1.92);
  % A and B
  \draw[thickline] (0.56,1.92) rectangle (1.18,1.32);
  \node at (0.87,1.62) {$m$};
  \node[left] at (0.56,1.62) {$A$};
  \draw[thickline] (6.82,1.92) rectangle (7.44,1.32);
  \node at (7.13,1.62) {$m$};
  \node[right] at (7.44,1.62) {$B$};
  % C at its present position
  \draw[thickline] (3.68,1.87) rectangle (4.32,1.27);
  \node at (4.00,1.57) {$M$};
  \node[below] at (4.00,1.25) {$C$};
  % strings to C
  \draw[thickline] (1.00,3.40)--(3.68,1.87);
  \draw[thickline] (7.00,3.40)--(4.32,1.87);
  % angle alpha at C
  \draw[dashedline] (4.00,1.87)--(4.00,3.05);
  \draw[thickline] (4.00,2.40) arc (90:150:0.55);
  \node at (4.12,2.68) {$\alpha$};
  % height H
  \draw[dim] (4.00,3.40)--(4.00,1.89) node[midway,right] {$H$};
  % weights
  \draw[force] (0.87,1.32)--++(-90:0.55) node[below] {$m\vec{g}$};
  \draw[force] (7.13,1.32)--++(-90:0.55) node[below] {$m\vec{g}$};
  \draw[force] (4.00,1.27)--++(-90:0.50) node[below] {$M\vec{g}$};
  % tensions
  \draw[force] (0.87,2.55)--++(90:0.60) node[right] {$\vec{T}$};
  \draw[force] ($(1.00,3.40)!0.45!(3.68,1.87)$)--++(150:0.60) node[left] {$\vec{T}$};
\end{tikzpicture}
\end{figbox}

When the speed of object $C$ reaches its maximum, its acceleration is zero. Hence, from the force analysis of object $C$, we obtain
\begin{equation}
2T\cos\alpha=Mg \tag{6}
\end{equation}
At this moment the equation of motion of object $A$ (or $B$) is
\begin{equation}
T-mg=ma \tag{7}
\end{equation}
In a plane polar coordinate system with $a$ (or $b$) as the pole, since the acceleration of object $C$ is zero, its radial acceleration is zero. That is,
\begin{equation}
a_r=\ddot{r}-r\dot{\theta}^2=0 \tag{8}
\end{equation}
In (8), $\ddot{r}$ is just the $a$ of (7), while $r=\dfrac{l}{\sin\alpha}$ and $\dot{\theta}=\dfrac{V\sin\alpha}{l/\sin\alpha}=\dfrac{V\sin^2\alpha}{l}$. Substituting these relations into (8), we obtain
\begin{equation}
a=\frac{V^2\sin^3\alpha}{l} \tag{9}
\end{equation}
Solving (5), (6), (7) and (9) simultaneously, and substituting $\alpha=60\degree$, we obtain
\begin{equation}
4\left(\frac{M}{m}\right)^2-5\left(\frac{M}{m}\right)+(10-6\sqrt{3})=0 \tag{10}
\end{equation}
Solving this quadratic equation in $\dfrac{M}{m}$, we obtain
\begin{equation}
\frac{M}{m}=1.324 \tag{11}
\end{equation}
Substituting (11) and $\alpha=60\degree$ into (5), we obtain the maximum speed of object $C$:
\[
V_{\max}=0.706\sqrt{gl}.
\]
(2) If $\dfrac{M}{m}=1$ is taken, then (5) becomes
\begin{equation}
V=\sqrt{\frac{\cos\alpha+2(\sin\alpha-1)}{(1+2\cos^2\alpha)\sin\alpha}\cdot 2gl} \tag{12}
\end{equation}
The radial equation of motion of object $C$ is
\begin{equation}
T+T\cos2\alpha-mg\cos\alpha=m\left[\frac{(V\sin\alpha)^2}{l/\sin\alpha}-a\right] \tag{13}
\end{equation}
The equation of motion of object $A$ (or $B$) is
\begin{equation}
T-mg=ma \tag{14}
\end{equation}
Let the acceleration of object $C$ be $a_C$, directed along the vertical direction; the equation of motion of object $C$ along this direction is
\begin{equation}
2T\cos\alpha-mg=ma_C \tag{15}
\end{equation}
Solving (12), (13), (14) and (15) simultaneously, and substituting $\alpha=60\degree$, we obtain
\[
a_C=\left(\frac{2}{\sqrt{3}}-1\right)g\approx 0.155g.
\]
The direction of this acceleration is vertically upward.

\begin{commentary}
A key point in solving this problem is that the moment when the speed of object $C$ reaches its maximum is not the moment when the speeds of objects $A$ and $B$ reach their maxima. In the plane polar coordinate system with the bearing $a$ (or $b$) as the pole, when the speed of object $C$ reaches its maximum both its radial acceleration and its transverse acceleration are zero, and from the vanishing of the radial acceleration it follows that $\ddot{r}=r\dot{\theta}^2$, where $\ddot{r}$ is precisely the acceleration of the motion of $A$ and $B$ at this moment. This shows that, as long as the angular velocity $\dot{\theta}$ of the rotation of the string connected to object $C$ is not zero at this moment, the acceleration of $A$ and $B$ is not zero, that is, the speeds of $A$ and $B$ are not at their maxima at this moment. This point is very easily overlooked.
\end{commentary}

\end{solution}

\begin{solution}{1.32}
(1) After the two bodies are released, let the acceleration of the centre of mass of the system be $a_C$. Then
\begin{align*}
-Mg\mu+mg\mu&=(M+m)a_C,\\
a_C&=-\frac{M-m}{M+m}g\mu=-1\unitv{m/s^2}.
\end{align*}
Let the distances from the centre of mass to $A$ and $B$ at this moment be $l_A$ and $l_B$; then
\[
\begin{cases}
l_A+l_B=l\\
Ml_A=ml_B
\end{cases}
\]
From these two equations we obtain
\[
l_A=\frac{m}{M+m}l=\frac{1}{3}l,
\]
\[
l_B=\frac{M}{M+m}l=\frac{2}{3}l.
\]

\begin{figbox}{1.32(a)}
\begin{tikzpicture}[scale=0.92,>=latex]
  % ground
  \fill[hatch] (-0.8,-0.24) rectangle (8.2,0);
  \draw[thickline] (-0.8,0)--(8.2,0);
  % block A
  \draw[thickline] (0.70,0) rectangle (1.90,0.95);
  \node at (1.30,0.47) {$M$};
  \node[above] at (1.30,1.02) {$A$};
  % block B
  \draw[thickline] (5.90,0) rectangle (7.10,0.95);
  \node at (6.50,0.47) {$m$};
  \node[above] at (6.50,1.02) {$B$};
  % elastic string (spring) between the blocks
  \draw[thickline] (1.90,0.47)--(2.15,0.47);
  \draw[thickline] (2.15,0.47)--(2.35,0.63)--(2.55,0.31)--(2.75,0.63)--(2.95,0.31)
    --(3.15,0.63)--(3.35,0.31)--(3.55,0.63)--(3.75,0.31)--(3.95,0.63)--(4.15,0.31)
    --(4.35,0.63)--(4.55,0.31)--(4.75,0.63)--(4.95,0.31)--(5.15,0.63)--(5.35,0.47);
  \draw[thickline] (5.35,0.47)--(5.90,0.47);
\end{tikzpicture}
\end{figbox}

When the string is at its natural length $l_0$, by the same method we can find the lengths $l_{0A}$ and $l_{0B}$ of the two segments of the elastic string from $A$ and $B$ to the centre of mass:
\[
l_{0A}=\frac{m}{M+m}l_0=\frac{1}{3}l_0,
\]
\[
l_{0B}=\frac{M}{M+m}l_0=\frac{2}{3}l_0.
\]
The stiffness coefficients $k_A$ and $k_B$ of these two segments of elastic string are respectively
\[
k_A=\frac{l_0}{l_{0A}}k=\frac{M+m}{m}k=3k,
\]
\[
k_B=\frac{l_0}{l_{0B}}k=\frac{M+m}{M}k=\frac{3}{2}k.
\]
Hence the angular frequency of the simple harmonic motion performed by the two bodies in the centre-of-mass frame is
\[
\omega_A=\omega_B=\omega=\sqrt{\frac{k(M+m)}{Mm}}=\sqrt{20}\unitv{s^{-1}}.
\]
Let the distance between the equilibrium position $O'$ of block $B$ and the position $O$ at which the string is at its natural length be $b_B$; the forces acting on $B$ are as shown in Fig.~1.32(b), where the inertial force is $f_{iB}=ma_C$. Then
\[
k_Bb_B=f_{iB}+f_B=m|a_C|+mg\mu,
\]
\[
b_B=\frac{m(g\mu+|a_C|)}{k_B}=0.2\unitv{m}.
\]
Hence the amplitude $A_B$ of the simple harmonic motion performed by $B$ is
\begin{align*}
A_B&=(l_B-l_{0B})-b_B=\frac{2}{3}(l-l_0)-b_B\\
&=0.4\unitv{m}.
\end{align*}

\begin{figbox}{1.32(b)}
\begin{tikzpicture}[scale=0.95,>=latex]
  % ground
  \fill[hatch] (-2.6,-0.24) rectangle (2.6,0);
  \draw[thickline] (-2.6,0)--(2.6,0);
  % block B
  \draw[thickline] (0.20,0) rectangle (1.20,0.90);
  \node at (0.70,0.45) {$m$};
  \node[above] at (0.70,0.96) {$B$};
  % spring force
  \draw[force] (0.20,0.45)--++(180:1.10) node[left] {$F_B$};
  % friction
  \draw[force] (1.20,0.62)--++(0:1.00) node[right] {$f_B$};
  % inertial force
  \draw[force] (1.20,0.26)--++(0:1.00) node[right] {$f_{iB}$};
  % positions O' and O
  \draw[dashedline] (0.20,-0.24)--(0.20,-0.62);
  \draw[dashedline] (1.20,-0.24)--(1.20,-0.62);
  \draw[dim] (0.20,-0.48)--(1.20,-0.48) node[midway,above] {$b_B$};
  \node[below left] at (0.20,-0.55) {$O'$};
  \node[below right] at (1.20,-0.55) {$O$};
\end{tikzpicture}
\end{figbox}

When the elastic string returns to its natural length, $B$ reaches point $O$, that is, the point $\frac{1}{2}A_B$ to the left of $O'$. Let the time from the release of $B$ to its arrival at point $O$ be $t_1$; from the amplitude vector diagram shown in Fig.~1.32(c) we obtain
\[
\omega t_1=\frac{2\pi}{3},
\]
that is,
\[
t_1=\frac{2\pi}{3\omega}=\frac{\sqrt{5}}{15}\pi\unitv{s}.
\]

\begin{figbox}{1.32(c)}
\begin{tikzpicture}[scale=0.95,>=latex]
  \draw[thickline] (0,0) circle (1.55);
  \draw[axis] (-1.95,0)--(1.95,0);
  \draw[axis] (0,-1.95)--(0,1.95);
  \node[below right] at (0,0) {$O$};
  % radius at 120 degrees
  \draw[thickline] (0,0)--(120:1.55);
  \node[above left] at (120:1.55) {$A_B$};
  \draw[dashedline] (120:1.55)--(-0.775,0);
  % A_B/2
  \draw[dim] (0,0)--(-0.775,0);
  \node[below] at (-0.39,0) {$\dfrac{A_B}{2}$};
  % phase angle
  \draw[thickline] (0.65,0) arc (0:120:0.65);
  \node at (52:0.98) {$\omega t_1$};
\end{tikzpicture}
\end{figbox}

At this moment the velocity $v_B$ of $B$ is
\[
v'_B=-A_B\omega\sin\frac{2\pi}{3}=-\frac{2}{5}\sqrt{15}\unitv{m/s}.
\]
From the properties of the centre-of-mass frame it is known that block $A$, also after the time $t_1$ from release, has the segment of the string between $A$ and the centre of mass restored to its natural length. At this moment the velocity $v'_A$ of $A$ is
\[
v'_A=-\frac{m}{M}v'_B=\frac{1}{2}\cdot\frac{2}{5}\sqrt{15}=\frac{1}{5}\sqrt{15}\unitv{m/s}.
\]
At this moment the velocity $v_C$ of the centre of mass is
\[
v_C=a_Ct_1=-\frac{\sqrt{5}}{15}\pi\unitv{m/s}.
\]
Hence at this moment the velocities of $A$ and $B$ relative to the ground are respectively
\[
v_A=v'_A+v_C=\frac{1}{5}\sqrt{15}-\frac{\sqrt{5}}{15}\pi=\frac{\sqrt{5}}{5}\left(\sqrt{3}-\frac{\pi}{3}\right)\unitv{m/s},
\]
\[
v_B=v'_B+v_C=-\frac{2}{5}\sqrt{15}-\frac{\sqrt{5}}{15}\pi=-\frac{\sqrt{5}}{5}\left(2\sqrt{3}+\frac{\pi}{3}\right)\unitv{m/s}.
\]
After this both bodies move with uniform deceleration; the acceleration $a$ is
\[
a=-g\mu=-3\unitv{m/s^2}.
\]
When the elastic string is restored to its natural length, the velocity $v_A$ of block $A$ is the smaller one; when it comes to rest after a time $t'$, the displacement it has covered is $\Delta l_A$, and we have
\[
t'=\frac{v_A}{|a|}=\frac{\sqrt{5}}{15}\left(\sqrt{3}-\frac{\pi}{3}\right)\unitv{s},
\]
\begin{align*}
\Delta l_A&=\frac{1}{2}|a|t'^2=\frac{1}{2}\times 3\times\frac{5}{15^2}\left(\sqrt{3}-\frac{\pi}{3}\right)^2\\
&\approx 0.016\unitv{m}.
\end{align*}
The displacement $\Delta l_B$ covered by block $B$ during this process is
\begin{align*}
\Delta l_B&=v_Bt'+\frac{1}{2}at'^2=-\frac{\sqrt{5}}{5}\left(2\sqrt{3}+\frac{\pi}{3}\right)\times\frac{\sqrt{5}}{15}\left(\sqrt{3}-\frac{\pi}{3}\right)-\frac{3}{2}\frac{5}{15^2}\left(\sqrt{3}-\frac{\pi}{3}\right)^2\\
&=-\frac{1}{30}\left(9-\frac{\pi^2}{3}\right)\approx -0.190\unitv{m}.
\end{align*}
Obviously, the distance covered by the two bodies $A$ and $B$ during this process is
\[
\Delta l_A-\Delta l_B\approx 0.206\unitv{m}<l_0=0.30\unitv{m}.
\]
It is thus seen that when $A$ has stopped the two bodies have not yet collided, and $B$ must cover a distance $\Delta l'_B$ before it collides with $A$:
\[
\Delta l'_B=l_0-\Delta l_A=0.284\unitv{m}.
\]
Let the value of the velocity of $B$ when the two bodies collide be $V_B$; then
\[
V_B^2-v_B^2=2a\Delta l'_B,
\]
\begin{align*}
\therefore\quad V_B&=\sqrt{v_B^2+2a\Delta l'_B}=\sqrt{\frac{1}{5}\left(2\sqrt{3}+\frac{\pi}{3}\right)^2-2\times 3\times 0.284}\\
&\approx 1.54\unitv{m/s}.
\end{align*}
(2) Let the time from the restoration of the elastic string to its natural length until the two bodies collide be $t_2$; then
\[
\Delta l'_B=|v_B|t_2+\frac{1}{2}at_2^2.
\]
Substituting the data into the above expression gives
\[
0.284=\frac{\sqrt{5}}{5}\left(2\sqrt{3}+\frac{\pi}{3}\right)t_2-\frac{3}{2}t_2^2,
\]
that is,
\[
1.5t_2^2-2.02t_2+0.284=0.
\]
We obtain
\[
t_2=\frac{1}{3}(2.02-1.54)=0.16\unitv{s}.
\]
The other solution of the equation, $t_2'=\dfrac{1}{3}(2.02+1.54)$, is discarded. Hence the total time from the release of the two bodies until their collision is
\[
t=t_1+t_2=\frac{\sqrt{5}}{15}\pi+0.16=0.63\unitv{s}.
\]
\end{solution}

\begin{solution}{1.33}
Let the distances of the two stars $M$ and $m$ from the centre of mass before the explosion be $r_M$ and $r_m$ respectively. Then
\begin{align}
r_M+r_m &= d \tag{1}\\
Mr_M &= mr_m \tag{2}
\end{align}
From (1) and (2) we obtain
\begin{align}
r_M &= \frac{m}{M+m}\,d \tag{3}\\
r_m &= \frac{M}{M+m}\,d \tag{4}
\end{align}
The two stars revolve about the centre of mass with the same angular velocity. Let it be $\omega$; then
\[
\frac{GMm}{d^2}=M\omega^2r_M,
\]
that is,
\begin{equation}
Gm=r_M\omega^2d^2. \tag{5}
\end{equation}
Since the explosion of the star of mass $M$ is instantaneous and spherically symmetric, the velocity of the remaining part of this star does not change at that instant. In this way, for the system formed by the remaining part of the star and the star of mass $m$, the centre of mass has velocity $v_C$ given by
\[
v_C=\frac{m\omega r_m-(M-\Delta M)\omega r_M}{M-\Delta M+m}
=\frac{\Delta M\,r_M\omega}{M-\Delta M+m}.
\]
The kinetic energy $E_{kC}$ of the centre of mass is
\begin{equation}
E_{kC}=\frac{1}{2}(M-\Delta M+m)v_C^2
=\frac{\Delta M^2r_M^2\omega^2}{2(M-\Delta M+m)}. \tag{6}
\end{equation}
The total kinetic energy $E_k$ of the system at this time is
\begin{equation}
E_k=\frac{1}{2}(M-\Delta M)(r_M\omega)^2+\frac{1}{2}m(r_m\omega)^2. \tag{7}
\end{equation}
The potential energy $U$ of the system at this time is
\begin{equation}
U=-\frac{G(M-\Delta M)m}{d}. \tag{8}
\end{equation}
Taking the centre of mass as the reference frame, the mechanical energy $E'$ of the system in the centre-of-mass frame is
\begin{equation}
E'=E_k-E_{kC}+U. \tag{9}
\end{equation}
Substituting (6), (7) and (8) into (9), we obtain
\[
E'=\frac{1}{2}(M-\Delta M)(r_M\omega)^2+\frac{1}{2}m(r_m\omega)^2
-\frac{\Delta M^2r_M^2\omega^2}{2(M-\Delta M+m)}-\frac{G(M-\Delta M)m}{d}.
\]
Substituting (5) into the above expression, we obtain
\begin{equation}
E'=-\frac{1}{2}r_M\omega^2\left(Md-2\Delta Md+\Delta Mr_M
+\frac{(\Delta M)^2r_M}{M-\Delta M+m}\right). \tag{10}
\end{equation}
If the mechanical energy of the system in the centre-of-mass frame is less than zero, i.e. $E'<0$, then the binary star will be bound and will not separate from each other; hence
\begin{equation}
Md-2\Delta Md+\Delta Mr_M+\frac{(\Delta M)^2r_M}{M-\Delta M+m}>0. \tag{11}
\end{equation}
Substituting (3) into (11), we obtain
\[
\Delta M<\frac{1}{2}(M+m).
\]

\begin{commentary}
Before the supernova explosion the centre of mass of the system is at rest. Although in the explosion the mass $\Delta M$ lost by the star of mass $M$ flies apart spherically symmetrically, so that the instantaneous velocity of the remaining part of this star does not change, the system formed by the remainder of the star $M$ after the explosion and the star $m$ is not the original system. The centre of mass of this new system is not at rest, but begins to move uniformly with velocity $v_C$; this point is easily overlooked. In addition, in solving this problem use is made of K\"onig's theorem, i.e. that the total kinetic energy of a system equals the sum of the kinetic energy of the centre of mass and the kinetic energy of the system in the centre-of-mass frame.
\end{commentary}

\end{solution}

\begin{solution}{1.34}
Let the remaining nucleus $m_3$ be emitted with velocity $v_3$, and let the angle between the direction of $v_3$ and the direction of $v_1$ be $\theta$, as shown in the figure. Take the direction of $v_1$ as the $x$ direction and the direction of $v_2$ as the $y$ direction. From the conservation of momentum of the system in the $x$ direction and in the $y$ direction we obtain respectively
\begin{align}
m_1v_1 &= m_3v_3\cos\theta \tag{1}\\
m_2v_2 &= m_3v_3\sin\theta \tag{2}
\end{align}

\begin{figbox}{1.34}
\begin{tikzpicture}[scale=1.0,>=latex]
  % axes
  \draw[axis] (-0.5,0)--(4.6,0) node[right] {$x$};
  \draw[axis] (0,-0.5)--(0,3.2) node[above] {$y$};
  % m1
  \fill (0.6,0) circle (2.2pt) node[above left] {$m_1$};
  \draw[force] (0.75,0)--(1.95,0) node[below right] {$v_1$};
  % m2
  \fill (3.9,2.4) circle (2.2pt) node[right] {$m_2$};
  \draw[force] (3.9,2.55)--(3.9,3.05) node[right] {$v_2$};
  \draw[dashedline] (3.9,0)--(3.9,2.4);
  % m3 at the vertex
  \fill (2.6,0.7) circle (2.2pt) node[below left] {$m_3$};
  \draw[force] (2.6,0.7)--(3.85,0.1) node[below right] {$v_3$};
  \draw[thickline] ($(2.6,0.7)+(0:0.6)$) arc (0:-25.6:0.6);
  \node at (3.25,0.92) {$\theta$};
  \draw[dashedline] (2.6,0.7)--(3.4,0.7);
\end{tikzpicture}
\end{figbox}

Squaring both sides of (1) and (2) respectively and then adding the two equations, we obtain
\[
m_1^2v_1^2+m_2^2v_2^2=m_3^2v_3^2.
\]
The above expression may also be rewritten as
\begin{equation}
\frac{1}{2}m_3v_3^2=\frac{1}{2}\frac{m_1^2}{m_3}v_1^2+\frac{1}{2}\frac{m_2^2}{m_3}v_2^2. \tag{3}
\end{equation}
Since in this nuclear reaction a quantity of non-mechanical energy of magnitude $Q$ is converted into mechanical energy, the energy relation of this reaction is
\begin{equation}
\frac{1}{2}m_1v_1^2+Q=\frac{1}{2}m_2v_2^2+\frac{1}{2}m_3v_3^2. \tag{4}
\end{equation}
Substituting (3) into (4), we obtain
\begin{equation}
\frac{1}{2}m_1v_1^2+Q=\frac{1}{2}m_1v_1^2\frac{m_1}{m_3}
+\frac{1}{2}m_2v_2^2\left(1+\frac{m_2}{m_3}\right). \tag{5}
\end{equation}
From (5) we immediately obtain the kinetic energy $E_k$ of the particle $m_2$ as
\[
E_k=\frac{1}{2}m_2v_2^2=\frac{m_3}{m_2+m_3}
\left[Q+\frac{1}{2}m_1v_1^2\left(\frac{m_3-m_1}{m_3}\right)\right].
\]

\end{solution}

\begin{solution}{1.35}
(1) Let the kinetic energies of the three blocks 1, 2 and 3 after the explosion be $T_1$, $T_2$ and $T_3$ respectively, and let their momenta be $p_1$, $p_2$ and $p_3$ respectively. The total kinetic energy of the three blocks is known to be $Q$; hence
\begin{equation}
T_1+T_2+T_3=Q. \tag{1}
\end{equation}
From the conservation of momentum of the system we obtain
\begin{equation}
\vec{p}_1+\vec{p}_2+\vec{p}_3=0. \tag{2}
\end{equation}
The vector relation expressed by (2) is shown in Fig.~1.35(a). From the figure we obtain
\[
p_1^2=p_2^2+p_3^2-2p_2p_3\cos\theta_{23},
\]
that is,
\begin{equation}
\cos\theta_{23}=\frac{p_2^2+p_3^2-p_1^2}{2p_2p_3}. \tag{3}
\end{equation}
Since $|\cos\theta_{23}|\le1$, equation (3) may be written as
\begin{equation}
\left(\frac{p_2^2+p_3^2-p_1^2}{2p_2p_3}\right)^2\le1. \tag{4}
\end{equation}
Substituting the relation between momentum and kinetic energy $T=p^2/2m$ into (4), we obtain
\[
\frac{(T_2+T_3-T_1)^2}{4T_2T_3}\le1,
\]
that is,
\begin{equation}
(T_2+T_3-T_1)^2\le4T_2T_3. \tag{5}
\end{equation}

\begin{figbox}{1.35(a)}
\begin{tikzpicture}[scale=1.0,>=latex]
  \coordinate (TL) at (0.35,2.15);
  \coordinate (TR) at (4.45,1.85);
  \coordinate (B)  at (1.95,-0.15);
  \draw[thickline] (TL)--(TR);
  \draw[thickline] (TR)--(B);
  \draw[thickline] (B)--(TL);
  \draw[force] ($(TL)!0.5!(TR)$)--(TR);
  \draw[force] ($(TR)!0.5!(B)$)--(B);
  \draw[force] ($(B)!0.5!(TL)$)--(TL);
  \node at (1.9,2.45) {$p_2$};
  \node at (3.75,0.75) {$p_3$};
  \node at (0.75,0.55) {$p_1$};
  \draw[thickline] ($(TR)+(175.8:0.7)$) arc (175.8:218.7:0.7);
  \node at ($(TR)+(197:1.05)$) {$\theta_{23}$};
\end{tikzpicture}
\end{figbox}

Any point inside the equilateral triangle shown in Fig.~1.35(b) satisfies (1); the points that can also satisfy (2), i.e. that satisfy (5), are the physically admissible points. Take an arbitrary point $P$ inside this triangle, and denote its position by the polar coordinates $(r,\theta)$, taking the centroid $O$ of the triangle as the pole of the polar coordinates and $OA$ as the polar axis. Since $OA=\dfrac{2}{3}Q$, from the figure we obtain
\begin{align}
T_1 &= \frac{Q}{3}+r\cos\theta \tag{6}\\
T_2 &= \frac{Q}{3}-r\cos(60\degree+\theta) \tag{7}\\
T_3 &= \frac{Q}{3}-r\cos(60\degree-\theta) \tag{8}
\end{align}

\begin{figbox}{1.35(b)}
\begin{tikzpicture}[scale=1.0,>=latex]
  \coordinate (A) at (0,3.0);
  \coordinate (B) at (-1.732,0);
  \coordinate (C) at (1.732,0);
  \coordinate (O) at (0,1.0);
  \draw[thickline] (A)--(B)--(C)--cycle;
  \draw[dashedline] (A)--(0,0);
  \fill (O) circle (1.3pt) node[right] {$O$};
  \node[above left] at (A) {$A$};
  \node[left] at (B) {$B$};
  \node[right] at (C) {$C$};
  \draw[force] (O)--++(270:0.75) node[below] {$T_1$};
  \draw[force] (O)--++(20:0.85) node[right] {$T_2$};
  \draw[force] (O)--++(160:0.85) node[left] {$T_3$};
  \draw[dim] (2.45,0)--(2.45,3.0) node[midway,right] {$Q$};
  \draw[dashedline] (C)--(2.45,0);
  \draw[dashedline] (A)--(2.45,3.0);
\end{tikzpicture}
\end{figbox}

Substituting (6), (7) and (8) into (5), we obtain
\begin{equation}
\left(\frac{Q}{3}-2r\cos\theta\right)^2
\le\frac{4}{9}Q^2-\frac{4}{3}Qr\cos\theta+4r^2\cos^2\theta-3r^2. \tag{9}
\end{equation}
From (9) we obtain
\[
r\le\frac{Q}{3}.
\]
Thus the boundary of the region of physically admissible points inside the equilateral triangle is the circumference of the inscribed circle. The admissible region is the region inside this circle, i.e. the shaded region shown in Fig.~1.35(c).

\begin{figbox}{1.35(c)}
\begin{tikzpicture}[scale=0.9,>=latex]
  \coordinate (A) at (0,3.0);
  \coordinate (B) at (-1.732,0);
  \coordinate (C) at (1.732,0);
  \fill[hatch] (0,1.0) circle (1.0);
  \draw[thickline] (A)--(B)--(C)--cycle;
  \draw[thickline] (0,1.0) circle (1.0);
  \node[above] at (A) {$A$};
  \node[left] at (B) {$B$};
  \node[right] at (C) {$C$};
\end{tikzpicture}
\end{figbox}

(2) For the case in which an electron-positron pair decays into three $\gamma$ photons, equations (1), (2), (3) and (4) above all apply; only the speed of a photon is $c$, and the relation between its energy and its momentum must be the relativistic one. Since a photon has no rest mass, its kinetic energy $T$ is just the energy $E$ of the photon; hence
\begin{equation}
T=E=cp. \tag{10}
\end{equation}
Thus (4) becomes
\[
\left|\frac{T_2^2+T_3^2-T_1^2}{2T_2T_3}\right|\le1,
\]
that is,
\[
T_2^2+T_3^2-T_1^2\le2T_2T_3,
\]
or
\begin{equation}
(T_2-T_3)^2\le T_1^2. \tag{11}
\end{equation}
Equation (11) has two solutions:
\begin{align}
T_2 &\le T_1+T_3 \tag{12}\\
T_3 &\le T_1+T_2 \tag{13}
\end{align}
Similarly we obtain
\begin{equation}
T_1\le T_2+T_3. \tag{14}
\end{equation}
Any one of (12), (13) and (14), for example (12), may also be written as
\[
2T_2\le T_1+T_2+T_3=Q,
\]
so that
\begin{equation}
T_2\le\frac{Q}{2}. \tag{15}
\end{equation}
From (13) and (14) we likewise obtain
\begin{align}
T_3 &\le\frac{Q}{2} \tag{16}\\
T_1 &\le\frac{Q}{2} \tag{17}
\end{align}
Equations (15), (16) and (17) are the conditions that a physically admissible point must satisfy; inside the equilateral triangle the region of physically admissible points is the interior of the triangle joining the midpoints of the sides together with its three sides, i.e. the shaded region shown in Fig.~1.35(d).

\begin{figbox}{1.35(d)}
\begin{tikzpicture}[scale=0.9,>=latex]
  \coordinate (A) at (0,3.0);
  \coordinate (B) at (-1.732,0);
  \coordinate (C) at (1.732,0);
  \coordinate (MAB) at (-0.866,1.5);
  \coordinate (MAC) at (0.866,1.5);
  \coordinate (MBC) at (0,0);
  \fill[hatch] (MAB)--(MAC)--(MBC)--cycle;
  \draw[thickline] (A)--(B)--(C)--cycle;
  \draw[thickline] (MAB)--(MAC)--(MBC)--cycle;
  \node[above] at (A) {$A$};
  \node[left] at (B) {$B$};
  \node[right] at (C) {$C$};
\end{tikzpicture}
\end{figbox}

\end{solution}

\begin{solution}{1.36}
Let the length of the string from the ball to the hole at the moment when the ball leaves the conical surface be $r$, and let the speed of the ball be $v$. Since none of the external forces acting on the ball has a vertical component of moment about point $O$, the angular momentum of the ball about the axis of symmetry of the cone (i.e. the vertical axis through the small hole $O$) is conserved; hence
\[
mv_0r_0\sin\alpha=mvr\sin\alpha.
\]
That is,
\begin{equation}
v=\frac{r_0}{r}v_0. \tag{1}
\end{equation}
When the ball leaves the conical surface, the forces acting on it are as shown in the figure. Hence
\begin{align}
T\cos\alpha &= mg \tag{2}\\
T\sin\alpha &= m\frac{v^2}{r\sin\alpha} \tag{3}
\end{align}

\begin{figbox}{1.36}
\begin{tikzpicture}[scale=1.0,>=latex]
  \coordinate (O) at (0,3.2);
  \coordinate (BL) at (-2.5,0.5);
  \coordinate (BR) at (2.5,0.5);
  \draw[thickline] (O)--(BL);
  \draw[thickline] (O)--(2.9,0.28);
  \draw[dashedline] (O)--(0,-0.1);
  \draw[dashedline] (0,0.5) ellipse (2.5 and 0.55);
  \node[above] at (O) {$O$};
  \draw[thickline] (-0.75,3.2) arc (-90:-113:0.75);
  \node at (-1.05,2.82) {$\alpha$};
  \draw[thickline] (0.75,3.2) arc (-90:-67:0.75);
  \node at (1.05,2.82) {$\alpha$};
  \fill (2.05,1.14) circle (2.4pt);
  \node[right] at (2.05,1.48) {$r_0$};
  \node[right] at (1.30,2.22) {$T$};
  \draw[force] (2.05,1.10)--(2.05,0.22) node[below] {$mg$};
\end{tikzpicture}
\end{figbox}

From (2) and (3) we obtain
\begin{equation}
v=\sin\alpha\sqrt{\frac{gr}{\cos\alpha}}. \tag{4}
\end{equation}
Comparing (1) and (4), we obtain
\begin{equation}
\begin{aligned}
r&=\sqrt[3]{\frac{v_0^2r_0^2\cos\alpha}{g\sin^2\alpha}}
=\sqrt[3]{\frac{\frac{15}{16}\times\frac{1}{4}\times\frac{1}{2}}{9.8\times\frac{3}{4}}}\\
&=0.25\unitv{m},
\end{aligned} \tag{5}
\end{equation}
\begin{equation}
\begin{aligned}
v&=v_0r_0\sqrt[3]{\frac{g\sin^2\alpha}{v_0^2r_0^2\cos\alpha}}
=\frac{\sqrt{15}}{4}\times\frac{1}{2}\times
\sqrt[3]{\frac{9.8\times\frac{3}{4}}{\frac{15}{16}\times\frac{1}{4}\times\frac{1}{2}}}\\
&=1.92\unitv{m/s}.
\end{aligned} \tag{6}
\end{equation}
From the principle of work and energy we can find the work $W$ done by the force pulling the string during this process:
\begin{equation}
W=\frac{1}{2}m(v^2-v_0^2)+mg(r_0-r)\cos\alpha. \tag{7}
\end{equation}
Substituting (5), (6) and the data into (7), we obtain
\[
W=2.61\unitv{J}.
\]

\end{solution}

\begin{solution}{1.37}
After the lead bullet is shot into the small ball $B$ and remains inside it, the spring begins to rotate about the centre of mass with an initial angular velocity $\omega_0$, while the two small balls vibrate along the direction of the spring relative to the centre of mass, and the centre of mass $C$ of the system always moves in a straight line with uniform velocity. Let the distance between the balls $A$, $B$ and the centre of mass $C$ be $l_A$ and $l_B$ respectively when the length of the spring is $l$; then
\begin{align}
ml_A &= 2ml_B \tag{1}\\
l_A+l_B &= l \tag{2}
\end{align}
From (1) and (2) we obtain
\begin{align*}
l_A &= \frac{2}{3}l\\
l_B &= \frac{1}{3}l
\end{align*}

\begin{figbox}{1.37}
\begin{tikzpicture}[scale=1.0,>=latex]
  \coordinate (A) at (0,2.9);
  \coordinate (B) at (0,0.4);
  \fill (A) circle (3.0pt);
  \fill[white] (A) circle (1.4pt);
  \node[right] at (0.22,2.9) {$A$};
  \draw[thickline,spring] (0,2.72)--(0,0.58);
  \node[right] at (0.35,1.65) {$C$};
  \fill (B) circle (3.0pt);
  \fill[white] (B) circle (1.4pt);
  \node[right] at (0.22,0.4) {$B$};
  \draw[force] (-0.28,0.4)--(-1.5,0.4) node[below] {$v_0$};
\end{tikzpicture}
\end{figbox}

It is thus seen that when the spring is at its natural length $a$ and when its length has become $2a$, the distances of the balls $A$ and $B$ from the centre of mass are $\dfrac{2}{3}a$, $\dfrac{1}{3}a$ and $\dfrac{4}{3}a$, $\dfrac{2}{3}a$ respectively. Before and after the lead bullet is shot into the small ball $B$, the angular momentum of the system about the spatial point $C$ at which the centre of mass is located at that instant is conserved; hence
\begin{equation}
mv_0\cdot\frac{1}{3}a=\left[2m\left(\frac{a}{3}\right)^2+m\left(\frac{2a}{3}\right)^2\right]\omega_0. \tag{3}
\end{equation}
From (3) we obtain
\begin{equation}
\omega_0=\frac{v_0}{2a}. \tag{4}
\end{equation}
Take the centre of mass as the reference frame. In the centre-of-mass frame the angular momentum of the system about the centre of mass is conserved and the mechanical energy is conserved. Let the angular velocity of rotation be $\omega$ when the length of the spring reaches its maximum value $2a$; then respectively
\begin{equation}
mv_0\cdot\frac{1}{3}a=\left[2m\left(\frac{2a}{3}\right)^2+m\left(\frac{4a}{3}\right)^2\right]\omega, \tag{5}
\end{equation}
\begin{equation}
\frac{1}{2}\left[2m\left(\frac{a}{3}\right)^2+m\left(\frac{2a}{3}\right)^2\right]\omega_0^2
=\frac{1}{2}\left[2m\left(\frac{2a}{3}\right)^2+m\left(\frac{4a}{3}\right)^2\right]\omega^2+\frac{1}{2}ka^2. \tag{6}
\end{equation}
Solving (4), (5) and (6) simultaneously, we obtain
\[
v_0=2a\sqrt{\frac{2k}{m}}.
\]

\begin{commentary}
This problem discusses the conservation of angular momentum and the conservation of mechanical energy of a system in the centre-of-mass frame. Since the centre of mass of the system always moves in a straight line with uniform velocity, the centre-of-mass frame is an inertial frame, and by the principle of relativity of mechanics its laws of mechanics are the same as in any inertial frame. However, the centre-of-mass frame is a special reference frame: even if the centre-of-mass frame were a non-inertial frame, in solving problems by means of the principle of work and energy, the theorem of angular momentum and so on, one need not take the action of inertial forces into account, i.e. it is equivalent to an ordinary inertial frame.
\end{commentary}

\end{solution}

\begin{solution}{1.38}
The orbits of the two comets are both parabolas; hence the mechanical energy of the system formed by each comet together with the Sun is zero. Let the speeds of the two comets just before their collision, when they have moved to a distance $R$ from the Sun, be $v_1$ and $v_2$ respectively; then
\begin{align}
\frac{1}{2}m_1v_1^2-\frac{GM_sm_1}{R} &= 0 \tag{1}\\
\frac{1}{2}m_2v_2^2-\frac{GM_sm_2}{R} &= 0 \tag{2}
\end{align}
From (1) and (2) we obtain
\begin{equation}
v_1=v_2=\sqrt{\frac{2GM_s}{R}}. \tag{3}
\end{equation}
The momentum of the system formed by the two comets is conserved before and after their collision, that is,
\begin{equation}
m_1v_1+m_2v_2=(m_1+m_2)v. \tag{4}
\end{equation}
In (4), $v$ is the velocity with which the two comets combine into a single composite celestial body after the collision. Since $v_1\perp v_2$, as shown in the figure, we have
\[
(m_1v_1)^2+(m_2v_2)^2=(m_1+m_2)^2v^2.
\]

\begin{figbox}{1.38}
\begin{tikzpicture}[scale=0.8,>=latex]
  % momentum diagram
  \coordinate (O) at (0,2.6);
  \draw[force] (O)--(-0.9,3.5) node[left] {$m_1v_1$};
  \draw[force] (O)--(0.9,3.5) node[right] {$m_2v_2$};
  \draw[force] (O)--(0,4.4) node[above] {$(m_1+m_2)v$};
  \draw[dashedline] (-0.9,3.5)--(0,4.4);
  \draw[dashedline] (0.9,3.5)--(0,4.4);
  % comet trajectories and the Sun
  \draw[dashedline,domain=-1.0:1.4,samples=60]
     plot (\x,{1.35-0.9*\x+0.25*\x*\x});
  \draw[dashedline,domain=-1.4:1.0,samples=60]
     plot (\x,{1.35+0.9*\x+0.25*\x*\x});
  \fill (0,0.5) circle (2.0pt);
  \draw[dotted,thickline] (0,0.5) circle (0.42);
  \node[right] at (0.3,0.35) {Sun};
  \fill (0,1.35) circle (1.5pt);
  \draw[dim] (0.55,0.5)--(0.55,1.33) node[midway,right] {$R$};
\end{tikzpicture}
\end{figbox}

we immediately obtain
\begin{equation}
v^2=\frac{m_1^2+m_2^2}{(m_1+m_2)^2}\,\frac{2GM_s}{R}. \tag{5}
\end{equation}
Since the collision is a perfectly inelastic one, mechanical energy is lost in the collision. Therefore the mechanical energy of the system formed by the composite celestial body and the Sun after the collision must be less than zero. Let the mass of the composite body be $M(=m_1+m_2)$; then
\[
E=\frac{1}{2}Mv^2-\frac{GM_sM}{R}<0.
\]
From this it is seen that the orbit of the composite celestial body must be an elliptical orbit with the Sun as focus. Since the semi-major axis $a$ of an elliptical orbit is determined by the energy, whereas the semi-minor axis $b$ is determined jointly by the energy and the angular momentum, and since from the given conditions of this problem it is impossible to determine the angle between the velocity of the composite body after the collision of the two comets and the radius vector at that place, i.e. it is impossible to determine the angular momentum, the semi-minor axis $b$ of the elliptical orbit cannot be determined. Since the semi-major axis $a$ is determined by the energy $E$ alone, in order to find $a$ we may as well assume that the velocity of the composite celestial body after the collision is perpendicular to the radius vector at that place, i.e. that the point of collision is the perihelion (or aphelion) of the elliptical orbit. Let the radius vector and the velocity of this composite body be $r$ and $v'$ when it moves to the aphelion (or perihelion); then from the conservation of angular momentum of the system about the centre of the Sun and from the conservation of mechanical energy we obtain respectively
\begin{align}
Mv'r &= MvR \tag{6}\\
\frac{1}{2}Mv'^2-\frac{GM_sM}{r} &= \frac{1}{2}Mv^2-\frac{GM_sM}{R} \tag{7}
\end{align}
From (5), (6) and (7) we obtain
\[
r=\frac{m_1^2+m_2^2}{2m_1m_2}R.
\]
From this the semi-major axis $a$ of the elliptical orbit is found to be
\[
a=\frac{1}{2}(r+R)=\frac{(m_1+m_2)^2}{4m_1m_2}R.
\]
That is, the orbit of the composite celestial body after the collision of the two comets is an elliptical orbit with the Sun as focus and semi-major axis $\dfrac{(m_1+m_2)^2}{4m_1m_2}R$.

The semi-major axis $a$ may also be found from the following formula:
\begin{equation}
v^2=k^2\left(\frac{2}{r}-\frac{1}{a}\right), \tag{8}
\end{equation}
where $k^2=GM_s$ and $r=R$. Substituting (5) into (8), we immediately obtain
\[
a=\frac{2k^2r}{k^2-v^2r}=\frac{(m_1+m_2)^2}{4m_1m_2}R.
\]

\end{solution}

\begin{solution}{1.39}
Let the speed with which the space shuttle together with the lunar module revolves about the Moon in a circular orbit be $v_0$; then
\begin{equation}
\frac{GM_m(M+m)}{(3R_m)^2}=(M+m)\frac{v_0^2}{3R_m}. \tag{1}
\end{equation}
In (1), $M_m$ is the mass of the Moon. From (1) we obtain
\begin{equation}
v_0=\sqrt{\frac{GM_m}{3R_m}}. \tag{2}
\end{equation}
Its period $T_0$ of revolution about the Moon is
\begin{equation}
T_0=\frac{2\pi(3R_m)}{v_0}=6\pi R_m\sqrt{\frac{3R_m}{GM_m}}. \tag{3}
\end{equation}
Since the gravitational acceleration at the surface of the Moon $g_m$ is
\begin{equation}
g_m=\frac{GM_m}{R_m^2}, \tag{4}
\end{equation}
substituting (4) into (3) and inserting the data, we obtain
\begin{equation}
\begin{aligned}
T_0 &= 6\pi R_m\sqrt{\frac{3}{g_mR_m}}
=6\times3.14\times1.74\times10^6\times\sqrt{\frac{3}{1.62\times1.74\times10^6}}\\
&=33812\unitv{s}\approx9.4\unitv{h}.
\end{aligned} \tag{5}
\end{equation}
Let the velocities of the lunar module and the space shuttle immediately after their separation be $v_1$ and $v_2$ respectively; from the conservation of momentum,
\begin{equation}
(M+m)v_0=mv_1+Mv_2. \tag{6}
\end{equation}
After the lunar module and the space shuttle have separated, the two move along their respective elliptical orbits, as shown in the figure. Let the speed of the lunar module when it reaches the surface of the Moon be $v_1'$; then from the conservation of angular momentum about the centre of the Moon and the conservation of mechanical energy we obtain respectively
\begin{align}
3R_mmv_1 &= R_mmv_1' \tag{7}\\
\frac{1}{2}mv_1^2-\frac{GM_mm}{3R_m} &= \frac{1}{2}mv_1'^2-\frac{GM_mm}{R_m}. \tag{8}
\end{align}
Solving (7) and (8) simultaneously, we obtain
\begin{equation}
v_1=\sqrt{\frac{GM_m}{6R_m}}=\frac{\sqrt{2}}{2}v_0. \tag{9}
\end{equation}
Substituting (9) into (6), we obtain
\begin{equation}
v_2=v_0+\frac{m}{M}\left(1-\frac{\sqrt{2}}{2}\right)v_0. \tag{10}
\end{equation}

\begin{figbox}{1.39}
\begin{tikzpicture}[scale=0.85,>=latex]
  % initial circular orbit of radius 3R_m, drawn in perspective
  \draw[thickline] (0,0) ellipse (1.75 and 1.35);
  % orbit of the space shuttle: perilune 3R_m, apolune kR_m
  \draw[dashedline] (0,-0.625) ellipse (1.874 and 1.975);
  % the Moon
  \fill[gray!30] (0,0) circle (0.45);
  \draw[thickline] (0,0) circle (0.45);
  \node[above right] at (0.06,0.16) {$R_m$};
  \draw[dim] (0,0)--(145:1.36) node[midway,below left] {$3R_m$};
  \node[below] at (0,-2.65) {$kR_m$};
\end{tikzpicture}
\end{figbox}

Let the speed of the space shuttle be $v_2'$ when it moves to its farthest point from the Moon, and let the distance of this point from the centre of the Moon be $kR_m$. Likewise, from the conservation of angular momentum of the space shuttle about the centre of the Moon and the conservation of mechanical energy we obtain respectively
\begin{align}
3R_mMv_2 &= kR_mMv_2' \tag{11}\\
\frac{1}{2}Mv_2^2-\frac{GM_mM}{3R_m} &= \frac{1}{2}Mv_2'^2-\frac{GM_mM}{kR_m}. \tag{12}
\end{align}
Solving (11) and (12) simultaneously, we obtain
\begin{equation}
k=\frac{3}{2\left(\dfrac{v_0}{v_2}\right)^2-1}. \tag{13}
\end{equation}
Substituting (10) into (13) and inserting $\dfrac{M}{m}=2$, we obtain
\[
k\approx5.75.
\]
Thus the semi-major axis $a_M$ of the elliptical orbit of the space shuttle is
\[
a_M=\frac{1}{2}(k+3)R_m\approx4.38R_m,
\]
while the semi-major axis $a_m$ of the elliptical orbit of the lunar module is
\[
a_m=\frac{1}{2}(3+1)R_m=2R_m.
\]
According to Kepler's third law, the ratios of the period $T_m$ of the lunar module and the period $T_M$ of the space shuttle to $T_0$ are found respectively to be
\[
\frac{T_m}{T_0}=\left(\frac{a_m}{3R_m}\right)^{\frac{3}{2}}
=\left(\frac{2}{3}\right)^{\frac{3}{2}}=0.54,
\]
\[
\frac{T_M}{T_0}=\left(\frac{a_M}{3R_m}\right)^{\frac{3}{2}}
=\left(\frac{4.38}{3}\right)^{\frac{3}{2}}=1.76.
\]
That is,
\begin{align}
T_m &= 0.54T_0 \tag{14}\\
T_M &= 1.76T_0 \tag{15}
\end{align}
In order that the lunar module may return along its original elliptical orbit to the point of separation and dock with the space shuttle, the time $\Delta t$ during which it may stay on the Moon is
\begin{equation}
\Delta t=(n+1)T_M-T_m\qquad(n=0,1,2,\dots). \tag{16}
\end{equation}
Substituting (14), (15) and (5) into (16), we obtain
\[
\Delta t=(1.76n+1.22)\times9.4\unitv{h}\qquad(n=0,1,2,\dots).
\]
From this expression it is seen that the shortest time during which the lunar module may stay on the surface of the Moon is
\[
\Delta t_{\min}=11.5\unitv{h}.
\]

\end{solution}

\begin{solution}{1.40}
(1) The gravitational force exerted on the bullet by the gas cloud is
\begin{equation}
\vec{F}(r)=
\begin{cases}
-\dfrac{GMmr}{R^3}\,\hat{r}, & (r\le R)\\[2mm]
-\dfrac{GMm}{r^2}\,\hat{r}, & (r>R)
\end{cases} \tag{1}
\end{equation}
If the gravitational potential energy is taken to be zero at infinity, the gravitational potential energy is found from the following integral:
\begin{equation}
U(r)=\int_r^\infty\vec{F}(r)\cdot\dd\vec{l}=\int_r^\infty F(r)\dd r. \tag{2}
\end{equation}
In the region $r\le R$,
\begin{equation}
\begin{aligned}
U(r)&=\int_r^R\left(-\frac{GMmr}{R^3}\right)\dd r
+\int_R^\infty\left(-\frac{GMm}{r^2}\right)\dd r\\
&=-\frac{GMm}{2R^3}\left(3R^2-r^2\right).
\end{aligned} \tag{3}
\end{equation}
In the region $r>R$,
\begin{equation}
U(r)=\int_R^\infty\left(-\frac{GMm}{r^2}\right)\dd r=-\frac{GMm}{r}. \tag{4}
\end{equation}

\begin{figbox}{1.40(a)}
\begin{tikzpicture}[scale=1.0,>=latex,xscale=1.02,yscale=1.25]
  \draw[axis] (-0.2,0)--(5.7,0) node[right] {$r$};
  \draw[axis] (0,-1.8)--(0,0.55) node[above] {$U(r)$};
  \draw[thickline,domain=0:1,samples=40] plot (\x,{-1.5+0.5*\x*\x});
  \draw[thickline,domain=1:5.4,samples=60] plot (\x,{-1/\x});
  \draw[dashedline] (0,-1.25)--(5.4,-1.25);
  \draw[dashedline] (0.707,-1.25)--(0.707,0);
  \node[left] at (-0.08,-1.25) {$E$};
  \node[left] at (-0.08,-1.0) {$-\dfrac{GMm}{R}$};
  \node[left] at (-0.08,-1.52) {$-\dfrac{3GMm}{2R}$};
  \node[below left,font=\scriptsize] at (0.707,-0.02) {$r_0$};
  \node[below right,font=\scriptsize] at (1.0,-0.02) {$R$};
  \node[below left] at (0,0) {$O$};
\end{tikzpicture}
\end{figbox}

Fig.~1.40(a) is the potential-energy curve. From the statement of the problem, the total energy of the bullet is $E=-\dfrac{5GMm}{4R}$, and the region in which the bullet can appear must satisfy $E\ge U$. From the figure it is seen that the bullet can appear only in the region $r\le r_0$, and $r_0<R$; that is to say, the bullet is always inside the cloud.

(2) The effective potential energy $U_{\mathrm{eff}}(r)$ of the bullet is
\begin{equation}
U_{\mathrm{eff}}(r)=U(r)+\frac{L^2}{2mr^2}. \tag{5}
\end{equation}
In (5), $L$ is the angular momentum of the bullet about the centre of the gas cloud, $L=m\sqrt{\dfrac{GMR}{32}}$. Substituting this expression together with (3) into (5), we obtain
\begin{equation}
U_{\mathrm{eff}}(r)=-\frac{GMm}{2R^3}\left(3R^2-r^2\right)+\frac{GMmR}{64r^2}. \tag{6}
\end{equation}
The mechanical energy $E$ of the bullet is
\begin{equation}
E=\frac{1}{2}m\dot{r}^2+U_{\mathrm{eff}}(r). \tag{7}
\end{equation}
At the turning points of the orbit, $\dot{r}=0$, hence
\begin{equation}
E=-\frac{5GMm}{4R}=U_{\mathrm{eff}}(r)
=-\frac{GMm\left(3R^2-r^2\right)}{2R^3}+\frac{GMmR}{64r^2}. \tag{8}
\end{equation}
From (8) we obtain the equation
\begin{equation}
32\left(\frac{r}{R}\right)^4-16\left(\frac{r}{R}\right)^2+1=0. \tag{9}
\end{equation}
The solutions of this quadratic equation are
\[
\left(\frac{r}{R}\right)^2=\frac{2\pm\sqrt{2}}{8},
\qquad\text{that is,}\qquad
r_1=\sqrt{\frac{2+\sqrt{2}}{8}}\,R,\quad
r_2=\sqrt{\frac{2-\sqrt{2}}{8}}\,R.
\]
At a turning point of the orbit the distance of the bullet from the centre of the gas cloud is $r_1=\sqrt{\dfrac{2+\sqrt{2}}{8}}\,R$ or $r_2=\sqrt{\dfrac{2-\sqrt{2}}{8}}\,R$.

(3) From the discussion and calculation of part (2) it is known that equation (9) has two solutions $r_1$ and $r_2$, i.e. in the effective-potential-energy curve of the bullet shown in Fig.~1.40(b) the mechanical energy $E=-\dfrac{5GMm}{4R}$ of the bullet has two points of intersection with the effective potential-energy curve. Therefore the motion of the bullet is an elliptical orbit, whose semi-major axis $a$ and semi-minor axis $b$ are respectively
\[
a=r_1=\frac{1}{2}\sqrt{1+\frac{\sqrt{2}}{2}}\,R,
\]

\begin{figbox}{1.40(b)}
\begin{tikzpicture}[scale=1.0,>=latex,xscale=1.05,yscale=1.35]
  \draw[axis] (-0.2,0)--(5.7,0) node[right] {$r$};
  \draw[axis] (0,-1.65)--(0,0.75) node[above] {$U(r)$};
  \draw[thickline,domain=0.16:1,samples=70]
     plot (\x,{-1.5+0.5*\x*\x+0.015625/(\x*\x)});
  \draw[thickline,domain=1:5.4,samples=70]
     plot (\x,{-1/\x+0.015625/(\x*\x)});
  \draw[dashedline,domain=0.22:3.6,samples=70]
     plot (\x,{0.015625/(\x*\x)});
  \draw[dashedline] (0,-1.25)--(5.4,-1.25);
  \draw[dashedline] (0.2706,-1.25)--(0.2706,0.05);
  \draw[dashedline] (0.6533,-1.25)--(0.6533,0.05);
  \draw[dashedline] (1.0,-1.25)--(1.0,0.05);
  \node[left] at (-0.08,-1.25) {$E$};
  \node[below left,font=\scriptsize] at (0.2706,-0.02) {$r_1$};
  \node[below right,font=\scriptsize] at (0.6533,-0.02) {$r_2$};
  \node[below right,font=\scriptsize] at (1.0,-0.02) {$R$};
  \node[right] at (0.15,-1.48) {$U_{\mathrm{eff}}$};
  \node[right] at (1.35,0.62) {$\dfrac{L^2}{2mr^2}$};
\end{tikzpicture}
\end{figbox}

and
\[
b=r_2=\frac{1}{2}\sqrt{1-\frac{\sqrt{2}}{2}}\,R.
\]
The shape is roughly as shown in Fig.~1.40(c). It is seen that the orbit has four turning points in all. Since the area $S$ of the ellipse is
\[
S=\pi ab=\frac{\sqrt{2}}{8}\pi R^2,
\]
and since the angular momentum of the bullet about the centre of the gas cloud is conserved, its areal velocity $k$ is constant:
\[
k=\frac{\dd S}{\dd t}=\frac{L}{2m}=\frac{1}{2}\sqrt{\frac{GMR}{32}}.
\]
Therefore the period $T$ of the motion of the bullet in its orbit is
\[
T=\frac{S}{k}=\frac{2\pi R^{3/2}}{\sqrt{GM}}.
\]

\begin{figbox}{1.40(c)}
\begin{tikzpicture}[scale=1.0,>=latex]
  \draw[axis] (-2.3,0)--(2.3,0) node[right] {$x$};
  \draw[axis] (0,-1.8)--(0,1.8) node[above] {$y$};
  \draw[thickline,rotate=22] (0,0) ellipse (1.85 and 1.0);
  \draw[dashedline,rotate=22] (-1.85,0)--(1.85,0);
  \draw[dashedline] (0,0)--(22:1.85);
  \draw[dashedline] (0,0)--(112:1.0);
  \node at (22:1.05) {$r_1$};
  \node at (112:0.62) {$r_2$};
  \node[below left] at (0,0) {$O$};
\end{tikzpicture}
\end{figbox}

\begin{commentary}
When a particle moves in a one-dimensional potential field, the potential-energy curve is often used to discuss its motion; and when a particle moves in two dimensions (i.e. with the two variables $r$ and $\theta$) in the gravitational field of a fixed centre of force, the effective potential energy $U_{\mathrm{eff}}(r)$ is often introduced in order to discuss conveniently the range of variation of its radius vector $r$. The effective potential energy consists of two parts: one is the gravitational potential energy $U(r)$, the other is the transverse kinetic energy $\frac{1}{2}m(r\dot{\theta})^2$ contained in its mechanical energy. Since the angular momentum of the particle about the fixed centre of force is conserved, i.e. $L=mr^2\dot{\theta}=C$, the $\dot{\theta}$ in the expression for the transverse kinetic energy may be replaced by $\dfrac{L}{mr^2}$; hence
\[
U_{\mathrm{eff}}(r)=U(r)+\frac{L^2}{2mr^2}.
\]
It is seen that the effective potential energy is a function of $r$ alone, and the mechanical energy $E$ of the system is also a function of $r$ alone, namely
\[
E=\frac{1}{2}m\dot{r}^2+U_{\mathrm{eff}}(r).
\]
This is analogous to the motion of a particle in one dimension (a single variable $r$); once the effective potential-energy curve is drawn, the range of the radial motion of the particle can be discussed very conveniently.
\end{commentary}
\end{solution}

\begin{solution}{1.41}
(1) According to scheme 1, the process in which a spaceship launched from the ground escapes the solar system directly may be divided into two stages. The first stage is the stage of first escaping the Earth's gravitational field. In order that the launch speed be the smallest, the Earth's orbital velocity should be used as far as possible in the solar system; hence the spaceship should be launched in the direction of the Earth's orbital velocity. When the spaceship has reached a place sufficiently far from the Earth, that is, when it has already escaped the Earth's gravitational field, it may be considered that the spaceship is still approximately in the Earth's orbit around the Sun, and at this time the spaceship still has a certain speed relative to the Earth; during this stage the gravitational potential energy due to the Sun may be regarded as unchanged. The second stage is the process in which the spaceship, after escaping the Earth's gravity, continues to escape from the Sun until it flies to an infinite distance from the Sun.

Taking the Earth as the frame of reference, let the speed of the spaceship relative to the Earth after it has escaped the Earth's gravity be $v_{a1}$; then from the conservation of mechanical energy we have
\begin{equation}
\frac{1}{2}mv_a^2-\frac{GM_em}{R_e}=\frac{1}{2}mv_{a1}^2 \tag{1}
\end{equation}
In (1), $M_e$ and $R_e$ are respectively the mass and the radius of the Earth. Now take the Sun as the frame of reference. By the statement of the problem the Earth moves in a circular orbit around the Sun; let the orbital speed of the Earth be $v_e$, the mass of the Sun be $M_s$, and the radius of the Earth's orbit be $r_{se}$; then
\[
\frac{GM_sM_e}{r_{se}^2}=M_e\frac{v_e^2}{r_{se}},
\]
that is,
\begin{equation}
v_e=\sqrt{\frac{GM_s}{r_{se}}} \tag{2}
\end{equation}

In the frame of the Sun, the velocity $v_{a2}$ with which the spaceship leaves the Earth's gravitational field is
\begin{equation}
v_{a2}=v_{a1}+v_e \tag{3}
\end{equation}
For the spaceship on the Earth's orbit to continue to escape from the Sun, the mechanical energy of the system formed by it and the Sun must be greater than or equal to zero, that is,
\[
\frac{1}{2}mv_{a2}^2-\frac{GM_sm}{r_{se}}\ge 0.
\]
The minimum $v_{a2}$ corresponds to taking the equality sign in the above expression; hence
\begin{equation}
v_{a2}=\sqrt{\frac{2GM_s}{r_{se}}}=\sqrt{2}\,v_e \tag{4}
\end{equation}
Substituting (4) into (3), we obtain
\begin{equation}
v_{a1}=(\sqrt{2}-1)v_e=12.3\ (\mathrm{km/s}) \tag{5}
\end{equation}
Substituting (5) into (1), we obtain
\[
v_a=\sqrt{v_{a1}^2+2G\frac{M_e}{R_e}}=16.4\ (\mathrm{km/s}).
\]

(2) In the frame of the Earth, let the launch speed of the spaceship at the Earth's surface be $v_b$, and let its speed when it leaves the Earth's gravity be $v_{b1}$; then from the conservation of mechanical energy,
\begin{equation}
\frac{1}{2}mv_b^2-\frac{GM_em}{R_e}=\frac{1}{2}mv_{b1}^2 \tag{6}
\end{equation}
In the frame of the Sun, the velocity $v_{b2}$ of the spaceship at this time is
\begin{equation}
v_{b2}=v_{b1}+v_e \tag{7}
\end{equation}
Thereafter, while the spaceship moves from the Earth's orbit to the orbit of Mars, it is acted on by the Sun's gravitation alone. When the spaceship reaches the orbit of Mars, let its velocity be $v$ ($=v_r+v_\theta$), as shown in the figure. Let the radius of the orbit of Mars be $r_{sm}$; then from the conservation of the angular momentum of the spaceship about the centre of the Sun we obtain
\begin{equation}
mv_{b2}r_{se}=mv_\theta r_{sm} \tag{8}
\end{equation}
From the conservation of mechanical energy,
\begin{equation}
\frac{1}{2}mv_{b2}^2-\frac{GM_sm}{r_{se}}=\frac{1}{2}m\left(v_r^2+v_\theta^2\right)-\frac{GM_sm}{r_{sm}} \tag{9}
\end{equation}
Substituting (6), (7) and (8) into (9), and using $v_e=\sqrt{GM_s/r_{se}}$ and $g=GM_e/R_e^2$, we obtain

\begin{figbox}{1.41(a)}
\begin{tikzpicture}[scale=1.0,>=latex]
  % orbits
  \draw[thickline] (0,0) circle (1.72);
  \draw[thickline] (0,0) circle (2.92);
  \fill (0,0) circle (1.6pt);
  \node[below] at (0,-0.14) {Sun};
  % radii
  \draw[dashedline] (0,0)--(-16:1.72);
  \node[below right] at (0.55,-0.30) {$r_{se}$};
  \draw[dashedline] (0,0)--(28:2.92);
  \node[right] at (1.30,0.72) {$r_{sm}$};
  % departure and arrival points
  \coordinate (P1) at (165:1.72);
  \coordinate (P2) at (90:2.92);
  \fill (P1) circle (1.5pt);
  \fill (P2) circle (1.5pt);
  % transfer trajectory
  \draw[dashedline] (P1) .. controls (-1.10,1.55) and (-0.45,2.35) .. (P2);
  % velocity at the Earth's orbit
  \draw[force] (P1)--++(74:0.85) node[left] {$v_{b2}$};
  % velocity components at the orbit of Mars
  \draw[force] (P2)--++(90:0.95) node[right] {$v_r$};
  \draw[force] (P2)--++(0:0.95) node[below] {$v_\theta$};
  \draw[force,dashedline] (P2)--++(0.95,0.95) node[right] {$v$};
\end{tikzpicture}
\end{figbox}

\[
\begin{aligned}
v_r^2&=v_{b2}^2-v_\theta^2-\frac{2GM_s}{r_{se}}+\frac{2GM_s}{r_{sm}}\\
&=(v_{b1}+v_e)^2\left[1-\left(\frac{r_{se}}{r_{sm}}\right)^2\right]-2v_e^2\left(1-\frac{r_{se}}{r_{sm}}\right)
\end{aligned}
\]
Hence
\begin{equation}
v_r=\left[\left(\sqrt{v_b^2-2R_eg}+v_e\right)^2\left(1-\frac{r_{se}^2}{r_{sm}^2}\right)-2v_e^2\left(1-\frac{r_{se}}{r_{sm}}\right)\right]^{\frac{1}{2}} \tag{10}
\end{equation}
\begin{equation}
v_\theta=\frac{r_{se}}{r_{sm}}\left(\sqrt{v_b^2-2R_eg}+v_e\right) \tag{11}
\end{equation}

(3) According to scheme 2, the motion of the spaceship may be divided into three stages. The first stage is the same as the first stage of scheme 1, being the stage in which the spaceship escapes the Earth's gravity. The second stage is the stage in which, after escaping the Earth's gravity, the spaceship moves from the Earth's orbit to the orbit of Mars, while the third stage is the stage in which the spaceship enters the gravitational field of Mars somewhere on the orbit of Mars, changes its direction of motion under the action of Mars's gravity, and leaves the gravitational field of Mars at another point of the orbit of Mars along the tangential direction of that orbit, so that it can make use of the orbital velocity of Mars and continue to escape from the solar system with a greater speed relative to the Sun.

The discussion of the first and second stages is exactly the same as in part (2). When the spaceship flies to a certain point of the orbit of Mars, its speed relative to the Sun is $v$ and its speed relative to Mars is $v'$; from the formula for relative motion we have
\begin{equation}
v'=v-v_m \tag{12}
\end{equation}
In (12), $v_m$ is the orbital velocity of Mars; (12) may also be expressed as
\begin{equation}
v'^2=(v_\theta-v_m)^2+v_r^2 \tag{13}
\end{equation}
The orbital velocity $v_m$ of Mars is obtained from
\begin{align}
\frac{GM_sM_m}{r_{sm}^2}&=M_m\frac{v_m^2}{r_{sm}}\notag\\
v_m&=\sqrt{\frac{GM_s}{r_{sm}}} \tag{14}
\end{align}
By the statement of the problem, when the spaceship leaves the gravitational field of Mars at another point of the orbit of Mars, the direction of its velocity is the tangential direction of the orbit, as shown in the figure. Taking Mars as the frame of reference, it is known from the conservation of mechanical energy that the speed with which it leaves the gravitational field of Mars is still $v'$, that is,
Hence, in the frame of the Sun, the velocity $v_{b3}$ of the spaceship is
\begin{equation}
v_{b3}=v'+v_m \tag{15}
\end{equation}
At the orbit of Mars the spaceship escapes from the solar system under the action of the Sun's gravitation alone, so the mechanical energy of the system formed by the spaceship and the Sun must be at least zero; hence
\[
\frac{1}{2}mv_{b3}^2-\frac{GM_sm}{r_{sm}}=0,
\]
that is,
\begin{equation}
v_{b3}=\sqrt{\frac{2GM_s}{r_{sm}}}=\sqrt{2}\,v_m \tag{16}
\end{equation}

\begin{figbox}{1.41(b)}
\begin{tikzpicture}[scale=1.0,>=latex]
  % orbits
  \draw[thickline] (0,0) circle (1.85);
  \draw[thickline] (0,0) circle (3.05);
  \fill (0,0) circle (1.6pt);
  \node[below] at (0,-0.12) {Sun};
  % radii
  \draw[dashedline] (0,0)--(20:3.05);
  \node at (20:2.45) {$r_{sm}$};
  \draw[dashedline] (0,0)--(-32:1.85);
  \node at (-32:1.25) {$r_{se}$};
  % departure point and transfer trajectory
  \coordinate (P1) at (200:1.85);
  \coordinate (P2) at (40:3.05);
  \fill (P1) circle (1.5pt);
  \fill (P2) circle (1.5pt);
  \draw[dashedline] (P1) .. controls (-0.55,-0.20) and (0.95,1.35) .. (P2);
  % velocity at the Earth's orbit (tangential)
  \draw[force] (P1)--++(110:0.85) node[above left] {$v_{b2}$};
  % velocities at the orbit of Mars
  \draw[force] (P2)--++(15:1.15) node[right] {$v$};
  \draw[force] (P2)--++(-50:0.95) node[below right] {$v_m$};
  \draw[force,dashedline] (P2)--++(-50:1.55) node[right] {$v_{b3}$};
\end{tikzpicture}
\end{figbox}
\end{solution}

\begin{solution}{1.42}
\begin{note}
The page on which this solution begins is missing from the scan; the text below is its continuation. Equations (1)--(5) are not reproduced.
\end{note}
where $\delta(t)\ll r_0$. Adopting plane polar coordinates, the radial equation of motion of the planet is
\begin{equation}
-\frac{GMm}{r^2}-kmr=m\left(\ddot r-r\dot\theta^2\right) \tag{6}
\end{equation}
Since the angular momentum of the planet about the centre of the star, $L(=mr^2\dot\theta)$, is conserved,
\begin{equation}
\dot\theta=\frac{L}{mr^2} \tag{7}
\end{equation}
Substituting (5) and (7) into (6), we obtain
\[
-\frac{GM}{(r_0+\delta)^2}-k(r_0+\delta)=\ddot\delta-\frac{L^2}{m^2(r_0+\delta)^3}.
\]
Since $\delta\ll r_0$, the term $(\delta/r_0)^2$ and the higher-order terms in the above expression may all be neglected, and the above expression becomes
\begin{equation}
\ddot\delta=-\frac{GM}{r_0^2}\left(1-\frac{2\delta}{r_0}\right)-kr_0\left(1+\frac{\delta}{r_0}\right)+\frac{L^2}{m^2r_0^3}\left(1-\frac{3\delta}{r_0}\right) \tag{8}
\end{equation}
Substituting (4) into (8), we obtain
\begin{equation}
\ddot\delta+\left(\frac{L^2}{m^2r_0^4}+3k\right)\delta=0 \tag{9}
\end{equation}
This is the equation of motion of simple harmonic motion, and the angular frequency of the simple harmonic motion is
\begin{equation}
\omega_r=\sqrt{\frac{L^2}{m^2r_0^4}+3k} \tag{10}
\end{equation}
When the planet moves in a circle of radius $r_0$ about the star, its angular velocity is $\omega_0=\dfrac{L}{mr_0^2}$; if there is no cosmic dust in space, that is, if $\rho=0$ and hence $k=0$, then $\omega_r=\omega_0$, which shows that when the planet moves once around the star, its radial vibration in $r$ also goes through exactly one period. In this process the value of $r$ has two maxima and two minima, so the orbit of the planet is a stable ellipse. When cosmic dust is present, that is, when $\rho\neq0$ and hence $k\neq0$, it is known from (10) that $\omega_r>\omega_0$. This shows that the orbit of the planet is no longer a strictly elliptical one; but since $k$ is very small, $\omega_r$ is only slightly larger than $\omega_0$, so the shape of the orbit is still roughly an ellipse, only the ellipse does not close: the radial vibration of the planet is ahead of the circular motion by one phase, and the ellipse will precess, as shown in the figure. That is, if the planet moves anticlockwise about the star, the direction of precession will be clockwise. Let the angular velocity of precession be $\omega_p$; then the time required for the axis of the ellipse to turn through one revolution is $T_p$, where
\begin{equation}
T_p=\frac{2\pi}{\omega_p} \tag{11}
\end{equation}
and when the ellipse has precessed through one revolution, the phase difference between the radial vibration of the planet and its circular motion is $2\pi$, that is,
\begin{equation}
(\omega_r-\omega_0)T_p=2\pi \tag{12}
\end{equation}

\begin{figbox}{1.42}
\begin{tikzpicture}[scale=1.0,>=latex]
  % precessing ellipse: the same focus (the star) at the origin
  \foreach \a in {0,40,80,120}{
    \draw[thickline,rotate=\a] (1.32,0) ellipse (1.60 and 0.90);
  }
  \fill (0,0) circle (1.6pt);
  \node[below left] at (-0.06,-0.06) {star};
  % precession arrow
  \draw[force] (1.80,2.30) .. controls (2.30,2.75) and (2.75,2.30) .. (2.35,1.85);
  \node[right] at (2.80,2.35) {$\omega_p$};
\end{tikzpicture}
\end{figbox}

From (11) and (12) we obtain
\begin{align}
\omega_p=\omega_r-\omega_0&=\sqrt{\frac{L^2}{m^2r_0^4}+3k}-\frac{L}{mr_0^2}\notag\\
&=\frac{L}{mr_0^2}\left[\left(1+\frac{3km^2r_0^4}{L^2}\right)^{\frac{1}{2}}-1\right]\notag\\
&\approx\frac{L}{mr_0^2}\cdot\frac{3km^2r_0^4}{2L^2}=\frac{3kmr_0^2}{2L} \tag{13}
\end{align}
Substituting $k=\dfrac{4}{3}\pi G\rho$ and $L=mr_0^2\omega_0=m\sqrt{GMr_0}$ into (13) gives
\[
\omega_p=2\pi\rho\sqrt{\frac{Gr_0^3}{M}}.
\]

\begin{commentary}
The motion of a planet about a star is a two-dimensional problem, and in plane polar coordinates the two variables $r$ and $\theta$ are used to describe its motion. Using the conservation of the angular momentum of the planet about the centre of the star, $\dot\theta$ in its radial equation of motion is replaced by $\dfrac{L}{mr^2}$, so that the equation becomes a one-dimensional equation of motion with the single variable $r$. From the equation it is known that the small radial vibration of the planet is a simple harmonic vibration; comparing the angular frequency of the radial vibration with the angular frequency of the orbital motion, the conclusion that the elliptical orbit of the planet precesses can be obtained, and both the direction of the precession and the angular frequency of precession can be found.
\end{commentary}
\end{solution}

\begin{solution}{1.43}
(1) Under the action of the central attractive force $F=kr^N$ the particle moves in a circle of radius $r_0$; let its speed of motion be $v_0$; then
\[
kr_0^N=\frac{mv_0^2}{r_0},
\]
that is,
\begin{equation}
v_0=\sqrt{\frac{kr_0^{N+1}}{m}} \tag{1}
\end{equation}
When it suffers a small radial disturbance its motion will deviate slightly from the circular orbit, and its position vector $r$ may be written as
\begin{equation}
r=r_0+\delta \tag{2}
\end{equation}
In (2), $\delta\ll r_0$. Adopting plane polar coordinates, its radial equation of motion is
\begin{equation}
-kr^N=m\left(\ddot r-r\dot\theta^2\right) \tag{3}
\end{equation}
Since the force acting on the particle is a central force, the angular momentum $L$ of the particle about the centre of force is conserved, $L$ being
\begin{equation}
L=mr^2\dot\theta=mv_0r_0 \tag{4}
\end{equation}
that is,
\begin{equation}
\dot\theta=\frac{L}{mr^2} \tag{5}
\end{equation}
Substituting (2) and (5) into (3), we obtain
\[
-k(r_0+\delta)^N=m\left[\ddot\delta-\frac{L^2}{m^2(r_0+\delta)^3}\right].
\]
Since $\delta\ll r_0$, neglecting $\left(\dfrac{\delta}{r_0}\right)^2$ and higher-order terms in the above expression, we obtain
\begin{equation}
\ddot\delta=-\frac{k}{m}r_0^N+\frac{L^2}{m^2r_0^3}-\delta\left(\frac{k}{m}Nr_0^{N-1}+\frac{3L^2}{m^2r_0^4}\right) \tag{6}
\end{equation}
Substituting (1) into (4), we obtain
\[
L=mv_0r_0=mr_0\sqrt{\frac{kr_0^{N+1}}{m}}.
\]
That is, we have the equality
\begin{equation}
\frac{L^2}{m^2r_0^3}-\frac{k}{m}r_0^N=0 \tag{7}
\end{equation}
Substituting (7) into (6), (6) becomes
\begin{equation}
\ddot\delta+\left(\frac{k}{m}Nr_0^{N-1}+\frac{3L^2}{m^2r_0^4}\right)\delta=0 \tag{8}
\end{equation}
From (8) it is known that if the small radial disturbance is to be stable, then we must have
\begin{equation}
\frac{k}{m}Nr_0^{N-1}+\frac{3L^2}{m^2r_0^4}>0 \tag{9}
\end{equation}
Substituting (4) into (9), we obtain
\begin{equation}
N>\frac{-\dfrac{3v_0^2}{r_0^2}}{\dfrac{k}{m}r_0^{N-1}}=-\frac{3mv_0^2}{kr_0^{N+1}} \tag{10}
\end{equation}
Substituting (1) into (10), we find
\[
N>-3.
\]

(2) Equation (8) is the equation of motion of the small radial vibration of the particle, and from this equation it is known that the angular frequency $\omega$ of this vibration is
\begin{equation}
\omega=\sqrt{\frac{k}{m}Nr_0^{N-1}+\frac{3L^2}{m^2r_0^4}} \tag{11}
\end{equation}
Substituting $L=mr_0v_0=mr_0\sqrt{\dfrac{kr_0^{N+1}}{m}}$ into (11), we obtain
\begin{equation}
\omega=\sqrt{(N+3)\frac{k}{m}r_0^{N-1}} \tag{12}
\end{equation}
Hence the period $T$ is
\begin{equation}
T=\frac{2\pi}{\omega}=2\pi\sqrt{\frac{m}{(N+3)kr_0^{N-1}}} \tag{12}
\end{equation}

(3) The period $T_0$ of the motion of the particle in a circular orbit of radius $r_0$ is
\begin{equation}
T_0=\frac{2\pi r_0}{v_0}=2\pi\sqrt{\frac{m}{kr_0^{N-1}}} \tag{13}
\end{equation}
If the period of the stable circular motion is to be an integral multiple of the period of the small radial vibration, then we have
\begin{equation}
T_0=nT\qquad(n=1,2,3,\cdots) \tag{14}
\end{equation}
Substituting (12) and (13) into (14), we obtain
\[
2\pi\sqrt{\frac{m}{kr_0^{N-1}}}=n\cdot2\pi\sqrt{\frac{m}{(N+3)kr_0^{N-1}}}.
\]
From the above expression we obtain
\[
N=n^2-3\qquad(n=1,2,3,\cdots),
\]
that is, when $N=-2,1,6,\cdots$, the period of the stable circular motion is $1,2,3,\cdots$ times the period of the small radial vibration.
\end{solution}

\begin{solution}{1.44}
(1) In the rotating frame which rotates about the axis with angular velocity $\omega_0$, the forces acting on one of the balls $m$ are as shown in the figure. The inertial centrifugal force is $f_c=m\omega_0^2l\sin\theta$; hence
\begin{equation}
(T_1-T_2)\cos\theta=mg \tag{1}
\end{equation}
\begin{equation}
(T_1+T_2)\sin\theta=f_c=m\omega_0^2l\sin\theta \tag{2}
\end{equation}
From (1) and (2) we obtain
\begin{equation}
T_2=\frac{1}{2}m\left(\omega_0^2l-\frac{mg}{\cos\theta}\right) \tag{3}
\end{equation}
From the equilibrium of the forces on sleeve $M$ we obtain
\begin{equation}
2T_2\cos\theta=Mg \tag{4}
\end{equation}
Substituting (3) into (4), we obtain
\begin{equation}
\cos\theta=\frac{(M+m)g}{m\omega_0^2l} \tag{5}
\end{equation}
Hence, when sleeve $M$ is in equilibrium, the distance $L$ between the upper and lower sleeves is
\[
L=2l\cos\theta=\frac{2(M+m)g}{m\omega_0^2}.
\]
Of course, $\theta=0$ should also be a possible equilibrium state of the system; in this case $L=2l$.

\begin{figbox}{1.44(b)}
\begin{tikzpicture}[scale=1.0,>=latex]
  % shaft
  \draw[thickline] (0,0.62)--(0,-3.62);
  % fixed upper sleeve
  \draw[thickline,fill=black!75] (-0.13,0.02) rectangle (0.13,0.30);
  % lower sleeve
  \draw[thickline,fill=white] (-0.16,-3.02) rectangle (0.16,-2.70);
  \node[right] at (0.24,-2.86) {$M$};
  % balls and arms
  \coordinate (L) at (-1.06,-1.56);
  \coordinate (R) at (1.06,-1.56);
  \draw[thickline] (0,0.02)--(L);
  \draw[thickline] (0,0.02)--(R);
  \draw[thickline] (L)--(0,-2.86);
  \draw[thickline] (R)--(0,-2.86);
  \node at (-0.72,-0.62) {$l$};
  \node at (0.72,-0.62) {$l$};
  \node at (-0.62,-2.20) {$l$};
  \node at (0.62,-2.20) {$l$};
  \draw[thickline,fill=white] (L) circle (0.13);
  \draw[thickline,fill=white] (R) circle (0.13);
  \node[left] at (-1.26,-1.56) {$m$};
  \node at (0.80,-1.18) {$T_1$};
  \node at (0.72,-1.96) {$T_2$};
  % axes
  \draw[axis] (0,0.02)--(2.10,0.02) node[right] {$x$};
  \draw[axis] (0,0.02)--(0,1.25) node[above] {$y$};
  \node[right] at (0.04,0.34) {$O$};
  % angle
  \draw[thickline] (0,-0.80) arc (270:305:0.80);
  \node at (0.24,-0.94) {$\theta$};
  % forces on the right ball
  \draw[force] (R)--++(0:0.62) node[above right] {$f_c$};
  \draw[force] (R)--++(270:0.62) node[below] {$mg$};
  % rotation
  \draw[force] (0.46,-3.42) arc (0:150:0.46 and 0.17);
  \node[right] at (0.50,-3.56) {$\omega_0$};
\end{tikzpicture}
\end{figbox}

(2) \textit{Method 1.}\quad Set up the coordinate system shown in the figure in the rotating frame; then the coordinates $x$, $y$ of the ball $m$ and the coordinate $Y$ of sleeve $M$ are respectively
\[
x=l\sin\theta,\qquad y=-l\cos\theta,
\]
\[
Y=-2l\cos\theta.
\]
Hence
\[
\dot x=l\cos\theta\,\dot\theta,\qquad \dot y=l\sin\theta\,\dot\theta,\qquad \dot Y=2l\sin\theta\,\dot\theta,
\]
\[
\ddot x=-l\sin\theta\,\dot\theta^2+l\cos\theta\,\ddot\theta,
\]
\[
\ddot y=l\cos\theta\,\dot\theta^2+l\sin\theta\,\ddot\theta,
\]
\[
\ddot Y=2l\cos\theta\,\dot\theta^2+2l\sin\theta\,\ddot\theta.
\]
Then the equations of motion of the ball $m$ and of sleeve $M$ are respectively
\begin{equation}
m\omega_0^2l\sin\theta-(T_1+T_2)\sin\theta=m\ddot x \tag{6}
\end{equation}
\begin{equation}
(T_1-T_2)\cos\theta-mg=m\ddot y \tag{7}
\end{equation}
\begin{equation}
2T_2\cos\theta-Mg=M\ddot Y \tag{8}
\end{equation}
From (6) and (7) we obtain
\begin{equation}
2T_2\cos\theta=m\cot\theta\left(\omega_0^2l\sin\theta-\ddot x\right)-m(\ddot y+g) \tag{9}
\end{equation}
From (8) and (9) we obtain
\[
M(\ddot Y+g)=m\cot\theta\left(\omega_0^2l\sin\theta-\ddot x\right)-m(\ddot y+g).
\]
Substituting the expressions for $\ddot x$, $\ddot y$ and $\ddot Y$ into the above expression, we obtain
\[
2Ml\cos\theta\,\dot\theta^2+\frac{(2M\sin^2\theta+m)l\ddot\theta}{\sin\theta}=m\omega_0^2l\cos\theta-(M+m)g.
\]
Let the system be in equilibrium at $\theta=\theta_0$; then after a small disturbance let $\theta=\theta_0+\theta'$, with $\theta'\ll\theta_0$, so that $\dot\theta=\dot\theta'$ and $\ddot\theta=\ddot\theta'$; since $\dot\theta^2=\dot\theta'^2$ is a second-order small quantity and is neglected, the above expression becomes
\[
(2M\sin^2\theta+m)l\ddot\theta'=m\omega_0^2l\cos\theta\sin\theta-(M+m)g\sin\theta.
\]
Since $\sin\theta=\sin(\theta_0+\theta')\approx\sin\theta_0+\cos\theta_0\cdot\theta'$ and $\cos\theta=\cos(\theta_0+\theta')\approx\cos\theta_0-\sin\theta_0\cdot\theta'$, substitution into the above expression and neglect of second-order and higher-order small quantities give
\[
(2M\sin^2\theta_0+m)l\ddot\theta'-\left[m\omega_0^2l\cos2\theta_0-(M+m)g\cos\theta_0\right]\theta'=0.
\]
Substituting (5) into the above expression, we obtain
\[
(2M\sin^2\theta_0+m)l\ddot\theta'+ml\omega_0^2\sin^2\theta_0\,\theta'=0.
\]
Hence the frequency of the small oscillation is
\[
\nu=\frac{\omega_0}{2\pi}\sqrt{\frac{m\sin^2\theta_0}{2M\sin^2\theta_0+m}}.
\]

\textit{Method 2.}\quad Introduce the centrifugal potential energy in the rotating frame (the conservative force being the inertial centrifugal force). Take the axis of rotation as the zero of potential energy; when the ball is at a distance $x$ from the axis, the centrifugal potential energy $U_{\mathrm{cf}}$ is
\[
U_{\mathrm{cf}}=\int_x^0 m\omega_0^2x\,\dd x=-\frac{1}{2}m\omega_0^2x^2.
\]
Take the upper sleeve as the zero of gravitational potential energy; then when the hinged arms make an angle $\theta$ with the axis, the potential energy $U$ of the system is
\begin{equation}
U=U_{\mathrm{cf}}+U_{\mathrm{g}}=-2\times\frac{1}{2}m\omega_0^2l^2\sin^2\theta-2(m+M)gl\cos\theta \tag{10}
\end{equation}
Let the speed of the ball be $v(=l\dot\theta)$ and the speed of the sleeve be $V$; the constraint relation between them is
\[
V\cos\theta=v\cos\left(\frac{\pi}{2}-2\theta\right).
\]
That is,
\[
V=2v\sin\theta.
\]
The kinetic energy of the system is
\begin{align}
E_k&=2\times\frac{1}{2}mv^2+\frac{1}{2}MV^2\notag\\
&=mv^2+2M\sin^2\theta\,v^2 \tag{11}
\end{align}
From (10) and (11) the mechanical energy of the system is
\begin{equation}
E=E_k+U=(2M\sin^2\theta+m)l^2\dot\theta^2-m\omega_0^2l^2\sin^2\theta-2(M+m)gl\cos\theta \tag{12}
\end{equation}
Since the mechanical energy of the system is conserved,
\[
\frac{\dd E}{\dd t}=4M\sin\theta\cos\theta\,l^2\dot\theta^3+(2M\sin^2\theta+m)l^2\cdot2\dot\theta\ddot\theta-2m\omega_0^2l^2\sin\theta\cos\theta\,\dot\theta+2(M+m)gl\sin\theta\,\dot\theta=0.
\]
The ball executes small oscillations about the equilibrium position. Let $\theta=\theta_0+\theta'$, $\theta'\ll\theta_0$; then $\dot\theta=\dot\theta'$ and $\ddot\theta=\ddot\theta'$, $\sin\theta\approx\sin\theta_0+\cos\theta_0\cdot\theta'$ and $\cos\theta\approx\cos\theta_0-\sin\theta_0\cdot\theta'$. Neglecting second-order and higher-order small quantities, the above expression becomes
\[
(2M\sin^2\theta+m)l\ddot\theta'-m\omega_0^2l\left(\sin\theta_0\cos\theta_0+\cos2\theta_0\,\theta'\right)+(M+m)g\left(\sin\theta_0+\cos\theta_0\cdot\theta'\right)=0.
\]
Substituting (5) into the above expression, we obtain
\[
(2M\sin^2\theta+m)l\ddot\theta'+m\omega_0^2l\sin^2\theta_0\,\theta'=0.
\]
From this the frequency of the small oscillation is obtained as
\[
\nu=\frac{\omega_0}{2\pi}\sqrt{\frac{m\sin^2\theta_0}{2M\sin^2\theta_0+m}}.
\]

(3) If the axis can rotate freely, then when the system suffers a small disturbance the angular velocity of rotation $\omega$ will change; let $\omega=\dot\phi$. Since the system is not acted on by an external torque about the axis, the angular momentum of the system about the axis is conserved; let this angular momentum be $L$,
\begin{equation}
L=2ml^2\sin^2\theta\,\dot\phi \tag{13}
\end{equation}
Taking a point at an infinite distance from the axis as the zero of centrifugal potential energy, the centrifugal potential energy of a ball is
\[
U_{\mathrm{cf}}=\int_x^\infty m\dot\phi^2x\,\dd x=\int_x^\infty m\frac{L^2}{4m^2x^3}\,\dd x=\frac{L^2}{8mx^2}.
\]
Hence the mechanical energy of the system is
\[
E=(2M\sin^2\theta+m)l^2\dot\theta^2+\frac{2\times L^2}{8ml^2\sin^2\theta}-2(M+m)gl\cos\theta.
\]
Since the mechanical energy of the system is conserved,
\[
\frac{\dd E}{\dd t}=4M\sin\theta\cos\theta\,l^2\dot\theta^3+2(2M\sin^2\theta+m)l^2\dot\theta\ddot\theta-\frac{L^2\cos\theta}{2ml^2\sin^3\theta}\dot\theta+2(M+m)gl\sin\theta\,\dot\theta=0.
\]
Let $\theta=\theta_0+\theta'$, $\theta'\ll\theta_0$, and $\dot\theta=\dot\theta'$, $\ddot\theta=\ddot\theta'$; substituting into the above expression and neglecting second-order and higher-order small quantities, the above expression becomes
\[
(2M\sin^2\theta+m)l\ddot\theta'-\frac{L^2\cos\theta_0}{2ml^2\sin^3\theta_0}\left(1-\frac{1+2\cos^2\theta_0}{\sin\theta_0\cos\theta_0}\theta'\right)+2(M+m)g\left(\sin\theta_0+\cos\theta_0\theta'\right)=0.
\]
Substituting $L=2ml^2\sin^2\theta_0\,\omega_0$ and $\cos\theta_0=\dfrac{(M+m)g}{m\omega_0^2l}$ into the above expression, we obtain
\[
(2M\sin^2\theta_0+m)l\ddot\theta'+m\omega_0^2\left(1+3\cos^2\theta_0\right)\theta'=0.
\]
From this the frequency of the small oscillation in this case is obtained as
\[
\nu=\frac{\omega_0}{2\pi}\sqrt{\frac{m\left(1+3\cos^2\theta_0\right)}{2M\sin^2\theta_0+m}}.
\]

\begin{commentary}
In a rotating frame the inertial centrifugal force acting on a particle is a conservative force, and it is customary to introduce a centrifugal potential energy. If the centrifugal potential energy is treated as one of the potential energies of the system, just like the potential energies of other real conservative internal forces, then the question of the mechanical energy of the system may be discussed in the same way as in any inertial frame; in this problem, for instance, since no external force does work and no non-conservative internal force does work, the mechanical energy of the system is conserved. It must be pointed out here that in a rotating frame the particle is acted on not only by the inertial centrifugal force but also by the Coriolis force; however, since the Coriolis force is always perpendicular to the direction of the particle's velocity, this force does no work, and hence the presence of the Coriolis force does not affect the discussion of the mechanical-energy question of the system.
\end{commentary}
\end{solution}

\begin{solution}{1.45}
(1) When the string has turned through an angle $\phi$, let the elongation of the string be $l$, as shown in the figure; then the equations of motion of the ball in the $x$ and $y$ directions are respectively
\begin{equation}
-kl\cos\phi=-kx=m\ddot x \tag{1}
\end{equation}
\begin{equation}
-kl\sin\phi-mg=-ky-mg=m\ddot y \tag{2}
\end{equation}
Let $y'=y+\dfrac{mg}{k}=y+a$; then $\ddot y'=\ddot y$, and substitution into (2) gives
\begin{equation}
-ky'=m\ddot y' \tag{3}
\end{equation}
From (1) and (3) the angular frequencies of the oscillations in the $x$ and $y$ directions are obtained as
\[
\omega_x=\omega_y=\omega=\sqrt{\frac{k}{m}}=\sqrt{\frac{g}{a}}.
\]

\begin{figbox}{1.45(b)}
\begin{tikzpicture}[scale=1.0,>=latex]
  % board
  \draw[thickline,fill=black!3] (-1.75,-2.45) rectangle (3.05,2.35);
  % axes
  \coordinate (O) at (0,0);
  \draw[axis] (O)--(2.95,0) node[right] {$x$};
  \draw[axis] (O)--(0,2.10) node[above] {$y$};
  \node[below left] at (O) {$O$};
  % elastic part below the hole
  \draw[thickline,decorate,decoration={coil,aspect=0.5,segment length=2.6pt,amplitude=1.3pt}]
    (0,-0.12)--(0,-2.42);
  % ball and string
  \coordinate (B) at (2.10,1.15);
  \draw[thickline] (O)--(B);
  \draw[thickline,fill=white] (B) circle (0.13);
  \node[above right] at ($(B)+(0.06,0.06)$) {$m$};
  % angle phi at O
  \draw[thickline] (0.75,0) arc (0:28.7:0.75);
  \node at (0.92,0.22) {$\phi$};
  % initial velocity
  \draw[force] (B)--++(60:0.85) node[above right] {$\vec v_0$};
  \draw[dashedline] (B)--(2.10,0);
  \draw[dim] (0,-0.55)--(2.10,-0.55) node[midway,below] {$a$};
  \draw[thickline] ($(B)+(0.50,0)$) arc (0:60:0.50);
  \node at (2.78,0.62) {$60\degree$};
\end{tikzpicture}
\end{figbox}

From the initial conditions $x_0=a$, $\dot x_0=v_0\cos\theta=\dfrac{1}{2}v_0$, $y_0'=a$, $\dot y_0'=v_0\sin\theta=\dfrac{\sqrt3}{2}v_0$, the amplitudes and initial phases of the oscillations in the $x$ and $y$ directions are obtained as
\[
A_x=\sqrt{x_0^2+\left(\frac{\dot x_0}{\omega_x}\right)^2}=\sqrt{a^2+\frac{ga}{g}\cdot a}=\sqrt2\,a,
\]
\[
A_y=\sqrt{y_0'^2+\left(\frac{\dot y_0'}{\omega_y}\right)^2}=\sqrt{a^2+\frac{ga\cdot3a}{g}}=2a,
\]
\[
\tan\phi_x=\frac{-\dot x_0}{\omega_xx_0}=\frac{-\dfrac{1}{2}\times2\sqrt{ga}}{\sqrt{\dfrac{g}{a}}\cdot a}=-1.
\]
\[
\therefore\ \phi_x=-\frac{\pi}{4},
\]
\[
\tan\phi_y=\frac{-\dot y_0'}{\omega_yy_0'}=\frac{-\dfrac{\sqrt3}{2}\times2\sqrt{ga}}{\sqrt{\dfrac{g}{a}}\cdot a}=-\sqrt3,
\]
\[
\therefore\ \phi_y=-\frac{\pi}{3}.
\]
Hence the equations of the oscillations in the $x$ and $y$ directions are respectively
\begin{equation}
x=\sqrt2\,a\cos\left(\sqrt{\frac{g}{a}}\,t-\frac{\pi}{4}\right) \tag{1}
\end{equation}
\begin{equation}
y'=2a\cos\left(\sqrt{\frac{g}{a}}\,t-\frac{\pi}{3}\right) \tag{2}
\end{equation}
When the ball reaches its highest position, we have
\[
y'=A_y=2a=y+a,
\]
\[
\therefore\ y=a.
\]
Let the time elapsed be $t_1$; then from (2) we obtain
\[
y'=2a\cos\left(\sqrt{\frac{g}{a}}\,t_1-\frac{\pi}{3}\right)=2a,
\]
that is,
\[
\sqrt{\frac{g}{a}}\,t_1-\frac{\pi}{3}=0,
\]
\[
\therefore\ t_1=\frac{\pi}{3}\sqrt{\frac{a}{g}}.
\]

(2) When the ball passes through the $x$-axis again, $y=0$, that is, $y'=a$. Let the time elapsed be $t_2$; then
\[
y'=2a\cos\left(\sqrt{\frac{g}{a}}\,t_2-\frac{\pi}{3}\right)=a,
\]
that is,
\[
\cos\left(\sqrt{\frac{g}{a}}\,t_2-\frac{\pi}{3}\right)=\frac{1}{2}.
\]
\[
\therefore\ \sqrt{\frac{g}{a}}\,t_2-\frac{\pi}{3}=\frac{\pi}{3},
\]
and hence
\[
t_2=\frac{2\pi}{3}\sqrt{\frac{a}{g}}.
\]
The $x$-coordinate at this time is
\begin{align}
x&=\sqrt2\,a\cos\left(\sqrt{\frac{g}{a}}\,t_2-\frac{\pi}{4}\right)\notag\\
&=\sqrt2\,a\cos\left(\sqrt{\frac{g}{a}}\times\frac{2\pi}{3}\sqrt{\frac{a}{g}}-\frac{\pi}{4}\right)\notag\\
&=\sqrt2\,a\cos\frac{5}{12}\pi=0.366a
\end{align}
At this time the velocities $v_x$ and $v_y$ of the ball in the $x$ and $y$ directions are respectively
\begin{align}
v_x&=-A_x\omega_x\sin(\omega_xt+\phi_x)\notag\\
&=-\sqrt{2ga}\sin\frac{5}{12}\pi=-1.366\sqrt{ga},
\end{align}
\begin{align}
v_y&=-A_y\omega_y\sin(\omega_yt+\phi_y)\notag\\
&=-2\sqrt{ga}\sin\frac{\pi}{3}=-\sqrt{3ga}=-1.73\sqrt{ga}.
\end{align}
Let the angle between the direction of the ball's velocity at this time and the $x$-axis be $\alpha$; then
\[
\tan\alpha=\frac{v_y}{v_x}=1.266,
\]
\[
\therefore\ \alpha=51.7\degree.
\]
The magnitude $v$ of the velocity is
\[
v=\sqrt{v_x^2+v_y^2}=2.2\sqrt{ga}.
\]

\begin{commentary}
This is a two-dimensional oscillation problem, that is, the motion of the ball is regarded as the superposition of two oscillations of the same frequency whose directions are mutually perpendicular but whose initial phases and amplitudes are all different. It is only necessary to use the initial conditions to find the initial phases and amplitudes of the oscillations in the two directions, and so to obtain the equations of these two oscillations; then, from the connection between the oscillations in the two directions, the required results can be found.
\end{commentary}
\end{solution}

\begin{solution}{1.46}
(1) After the locking of the small spring is released, the cylinder acquires an upward initial speed $V_0$. From the conservation of momentum, the downward initial speed $v_0$ acquired by the piston is found to be
\[
v_0=\frac{MV_0}{m}=3V_0=\frac{3}{2}g\sqrt{\frac{m}{k}}.
\]
Afterwards the piston moves downwards with uniform acceleration $a$, and the equation of motion is
\[
mg-f=ma.
\]
Since $f=\dfrac{1}{2}mg$, its acceleration $a$ is
\[
a=\frac{1}{2}g.
\]
Take the positive direction of the $x$-axis to be vertically downwards, and take the origin of the coordinate axis at the equilibrium position of the cylinder after the locking of the small spring is released and before the piston leaves the cylinder, as shown in Fig. 1.46(a). Take the instant at which the lock is released as $t=0$. At the instant $t=0$ the position of the cylinder is $x_0=\dfrac{mg-f}{k}$, its speed is $V_0=-\dfrac{g}{2}\sqrt{\dfrac{m}{k}}$, and the angular frequency of the cylinder's oscillation is $\omega=\sqrt{\dfrac{k}{M}}$. Hence its amplitude and initial phase are respectively
\[
A=\sqrt{x_0^2+\left(\frac{V_0}{\omega}\right)^2}
=\sqrt{\left(\frac{mg}{2k}\right)^2+\left(\frac{g}{2}\sqrt{\frac{m}{k}}\cdot\sqrt{\frac{M}{k}}\right)^2}
=\frac{mg}{k}.
\]

\begin{figbox}{1.46(a)}
\begin{tikzpicture}[scale=0.95,>=latex]
  % ceiling
  \draw[thickline] (-1.2,6.2)--(1.2,6.2);
  \fill[hatch] (-1.2,6.2) rectangle (1.2,6.5);
  % big spring
  \draw[thickline] (0,6.2) .. controls (0.35,5.9) and (0.35,5.6) .. (0,5.3)
    .. controls (-0.35,5.0) and (-0.35,4.7) .. (0,4.4)
    .. controls (0.35,4.1) and (0.35,3.8) .. (0,3.5);
  \node[right] at (0.45,4.9) {$k$};
  % cylinder
  \draw[thickline] (-0.75,3.5)--(0.75,3.5)--(0.75,0.0)--(-0.75,0.0)--cycle;
  % piston
  \fill[gray!70] (-0.75,3.0) rectangle (0.75,3.3);
  \draw[thickline] (-0.75,3.3)--(0.75,3.3);
  \node[right] at (0.85,3.15) {$m$};
  % origin O and x axis
  \node[left] at (-1.25,3.5) {$O$};
  \draw[thickline] (-1.18,3.5)--(-0.78,3.5);
  \draw[dim] (1.35,3.5)--(1.35,0.0) node[midway,right] {$x$};
\end{tikzpicture}
\end{figbox}

\[
\tan\phi=\frac{-V_0}{x_0\omega}
=\frac{\dfrac{g}{2}\sqrt{\dfrac{m}{k}}}{\dfrac{mg}{2k}\sqrt{\dfrac{k}{M}}}
=\sqrt{3},
\]
hence
\[
\phi=\frac{\pi}{3}.
\]
Therefore, before the piston leaves the cylinder, the motion of the cylinder can be expressed as
\begin{equation}
x=\frac{mg}{k}\cos\left(\sqrt{\frac{k}{M}}\,t+\frac{\pi}{3}\right) \tag{1}
\end{equation}
In this case the lowest position that the cylinder can reach is $x=A=\dfrac{mg}{k}$.

Suppose that at some instant before the cylinder reaches the above lowest position the piston has already left the cylinder. Can the cylinder afterwards move to a position lower than the above lowest position? Let the piston leave the cylinder when the cylinder has moved to the position $x'$ ($<A$), and let its speed at this instant be $V'$. From the conservation of mechanical energy,
\begin{equation}
\frac{1}{2}kx'^2+\frac{1}{2}MV'^2=\frac{1}{2}kA^2. \tag{2}
\end{equation}
Since the piston has left the cylinder, the equilibrium position of the cylinder has changed: compared with the situation before the separation it has risen by $x_0'=\dfrac{f}{k}=\dfrac{mg}{2k}$. If the new equilibrium position is taken as the origin of the new coordinate axis $x'$, and the amplitude of the subsequent oscillation of the cylinder is denoted by $A'$, then again from the conservation of mechanical energy,
\begin{equation}
\frac{1}{2}k\left(x'+x_0'\right)^2+\frac{1}{2}MV'^2=\frac{1}{2}kA'^2. \tag{3}
\end{equation}
Rewriting (2) and (3) gives
\begin{equation}
\left(x'\right)^2+\left(\sqrt{\frac{M}{k}}V'\right)^2=A^2, \tag{4}
\end{equation}
\begin{equation}
\left(x'+x_0'\right)^2+\left(\sqrt{\frac{M}{k}}V'\right)^2=A'^2. \tag{5}
\end{equation}
The relation among the magnitudes of the quantities in (4) and (5) is similar to the relation of the Pythagorean theorem in a right triangle, and can be represented by the geometrical relation shown in Fig. 1.46(b). From the figure it is known that
\begin{equation}
A+x_0'>A'. \tag{6}
\end{equation}
Equation (6) shows that, if the piston leaves the cylinder at any position before the cylinder has moved to the position $x=A$, then although the amplitude $A'$ of the subsequent motion is greater than $A$, its equilibrium position has risen by $x_0'$, so its lowest position is still higher than the position $x=A$. Therefore the lowest position that the cylinder can reach is $x=A=\dfrac{mg}{k}$, and the elapsed time $t$ can be found from (1):
\[
x=\frac{mg}{k}\cos\left(\sqrt{\frac{k}{M}}\,t+\frac{\pi}{3}\right)=\frac{mg}{k}.
\]

\begin{figbox}{1.46(b)}
\begin{tikzpicture}[scale=0.9,>=latex]
  % right angle vertex B and horizontal leg
  \coordinate (B) at (0,0);
  \coordinate (R) at (4.2,0);
  \coordinate (M) at (0,2.2);
  \coordinate (T) at (0,3.2);
  \draw[thickline] (T)--(B);
  \draw[thickline] (B)--(R);
  \draw[thickline] (T)--(R) node[pos=0.35,above left] {$A'$};
  \draw[thickline] (M)--(R) node[pos=0.72,below right] {$A$};
  % right angle mark
  \draw[thickline] (0.22,0)--(0.22,0.22)--(0,0.22);
  % dimension for x_0'
  \draw[dim] (-0.75,2.2)--(-0.75,3.2) node[midway,left] {$x_0'$};
  \draw[dashedline] (-0.75,2.2)--(0,2.2);
  \draw[dashedline] (-0.75,3.2)--(0,3.2);
  % label x'
  \node[left] at (-0.2,1.1) {$x'$};
  % bottom leg label
  \draw[dim] (0,-0.55)--(4.2,-0.55) node[midway,below] {$\sqrt{\dfrac{M}{k}}\,V'$};
  \draw[dashedline] (0,0)--(0,-0.55);
  \draw[dashedline] (4.2,0)--(4.2,-0.55);
\end{tikzpicture}
\end{figbox}

That is,
\[
\sqrt{\frac{k}{M}}\,t+\frac{\pi}{3}=2\pi,
\]
and hence
\[
t=\frac{5\pi}{3}\sqrt{\frac{M}{k}}.
\]

(2) For the cylinder to be able to reach the above lowest position, the piston must not leave the cylinder before the cylinder reaches this position; hence the length $l$ of the cylinder must satisfy
\[
\begin{aligned}
l&\ge v_0t+\frac{1}{2}at^2-(A-x_0)\\
&=\frac{3}{2}g\sqrt{\frac{m}{k}}\times\frac{5\pi}{3}\sqrt{\frac{M}{k}}
+\frac{1}{4}g\left(\frac{5\pi}{3}\sqrt{\frac{M}{k}}\right)^2-\frac{mg}{2k}\\
&=33.6\,\frac{mg}{k}.
\end{aligned}
\]
\end{solution}

\begin{solution}{1.47}
When the system is in a steady state, let the distance from the small ball moving on the table top to the hole $O$ be $r_0$; then
\[
m\omega_0^2r_0=Mg.
\]
Hence
\begin{equation}
r_0=\frac{Mg}{m\omega_0^2}. \tag{1}
\end{equation}
After the system is slightly disturbed, the motion of the small ball on the table top can be described in a polar coordinate system with origin $O$, and its radial equation of motion is
\begin{equation}
-T=m\left(\ddot{r}-r\dot{\theta}^2\right). \tag{2}
\end{equation}
In the above expression $T$ is the tension in the string. The equation of motion of the block $M$ is
\begin{equation}
T-Mg=Ma. \tag{3}
\end{equation}
The quantity $a$ in (3) should be equal to $\ddot{r}$ in (2), that is,
\begin{equation}
a=\ddot{r}. \tag{4}
\end{equation}
From (2), (3) and (4) we obtain
\begin{equation}
(M+m)\ddot{r}-mr\dot{\theta}^2+Mg=0. \tag{5}
\end{equation}
Since the angular momentum $L$ of the small ball $m$ about the point $O$ is conserved, we have
\[
L=m\omega_0r_0^2=m\dot{\theta}r^2.
\]

\begin{figbox}{1.47}
\begin{tikzpicture}[scale=0.9,>=latex]
  % table top in perspective
  \draw[thickline] (-3.4,-0.9)--(2.9,-1.45)--(4.4,1.05)--(-1.9,1.6)--cycle;
  % hole O
  \coordinate (O) at (0.55,0.05);
  \fill (O) circle (1.4pt) node[above left] {$O$};
  % dashed circular path
  \draw[dashedline] (O) ellipse (1.75 and 0.62);
  % ball m
  \fill (2.3,-0.05) circle (2.2pt) node[right] {$m$};
  % radius
  \draw[force] (O)--(2.3,-0.05) node[midway,above] {$r_0$};
  % string through the hole
  \draw[thickline] (O)--(0.55,-3.15);
  % block M
  \draw[thickline] (0.2,-3.15) rectangle (0.9,-3.85);
  \node at (0.55,-3.5) {$M$};
\end{tikzpicture}
\end{figbox}

That is,
\begin{equation}
\dot{\theta}=\frac{L}{mr^2}. \tag{6}
\end{equation}
Substituting (6) into (5) gives
\begin{equation}
(M+m)\ddot{r}-\frac{L^2}{mr^3}+Mg=0. \tag{7}
\end{equation}
Let $r=r_0+\delta$, where $\delta\ll r_0$; substituting this into (7) gives
\[
(M+m)\ddot{\delta}-\frac{L^2}{mr_0^3}\left(1+\frac{\delta}{r_0}\right)^{-3}+Mg=0.
\]
Since $\dfrac{\delta}{r_0}\ll1$, we take the approximation in the above expression and neglect small quantities of second and higher order. Substituting (1) and $L=m\omega_0r_0^2$, we obtain
\[
(M+m)\ddot{\delta}+\frac{3L^2}{mr_0^4}\delta=0,
\]
or
\begin{equation}
\ddot{\delta}+\frac{3m\omega_0^2}{M+m}\delta=0. \tag{8}
\end{equation}
From (8) the period of the small oscillation of the block is obtained at once:
\[
T=\frac{2\pi}{\omega_0}\sqrt{\frac{M+m}{3m}}.
\]
This problem may also be solved by a method similar to Method 2 of Problem 1.44: by taking the first derivative of the mechanical energy of the system the equation of oscillation is obtained, and hence the period of oscillation is found. This is not repeated here; the reader is asked to work it out independently.
\end{solution}

\begin{solution}{1.48}
Suppose that the circle rolls uniformly without slipping along a straight line; then the coordinates of the point $P$ on it at time $t$ can be expressed as
\begin{align}
x&=R\left(\omega t+\sin\omega t\right) \tag{1}\\
y&=R\left(1-\cos\omega t\right) \tag{2}
\end{align}
where $\omega$ is the angular speed of rotation about the centre of the circle when the circle rolls without slipping.

\begin{figbox}{1.48}
\begin{tikzpicture}[scale=0.8,>=latex]
  % axes
  \draw[axis] (-4.6,0)--(5.0,0) node[right] {$x$};
  \draw[axis] (0,0)--(0,3.9) node[above] {$y$};
  \node[below left] at (0,0) {$O$};
  % rolling line y = 2R
  \draw[thickline] (-4.1,2)--(4.8,2);
  \draw[dashedline] (-3.1416,2)--(-3.75,2);
  \draw[dim] (-3.75,0)--(-3.75,2) node[midway,left] {$2R$};
  % cycloid
  \draw[thickline,domain=-180:180,samples=180]
     plot ({(\x*3.14159/180+sin(\x))},{1-cos(\x)});
  % generating circle at the right cusp
  \draw[thickline] (3.1416,1) circle (1);
  \draw[thickline] (3.1416,1)--(3.1416,2);
  \node[right] at (3.16,1.5) {$R$};
  \fill (3.1416,2) circle (1.2pt);
  % point P on the left branch
  \coordinate (P) at (-1.571,1.0);
  \fill (P) circle (1.7pt) node[above left] {$P$};
  % tangent and horizontal at P
  \draw[dashedline] ($(P)+(-0.85,0.85)$)--($(P)+(0.85,-0.85)$);
  \draw[dashedline] ($(P)+(-1.1,0)$)--($(P)+(1.1,0)$);
  \draw[thickline] ($(P)+(0.7,0)$) arc (0:-45:0.7);
  \node at ($(P)+(0.85,-0.28)$) {$\theta$};
  % weight
  \draw[force] (P)--++(0,-0.95) node[below] {$m\vec{g}$};
\end{tikzpicture}
\end{figbox}

Differentiating both sides of (1) and (2) gives
\begin{align}
\dd x&=R\omega\left(1+\cos\omega t\right)\dd t \tag{3}\\
\dd y&=R\omega\sin\omega t\,\dd t \tag{4}
\end{align}
From (3) and (4) the element of arc length $\dd s$ traversed by the point $P$ along the cycloid during the time $\dd t$ is
\begin{align}
\dd s&=\sqrt{(\dd x)^2+(\dd y)^2}=R\omega\sqrt{2+2\cos\omega t}\,\dd t\notag\\
&=R\omega\sqrt{4\cos^2\left(\frac{\omega t}{2}\right)}\,\dd t=2R\omega\cos\frac{\omega t}{2}\,\dd t. \tag{5}
\end{align}
Integrating both sides of (5),
\[
\int_0^s\dd s=\int_0^t2R\omega\cos\frac{\omega t}{2}\,\dd t,
\]
we obtain
\begin{equation}
s=4R\sin\frac{\omega t}{2}. \tag{6}
\end{equation}
When a small ring is threaded on a smooth steel wire of the same shape as this cycloid and slides along the wire, its tangential acceleration $\ddot{s}$ is
\begin{equation}
\ddot{s}=-g\sin\theta. \tag{7}
\end{equation}
Since the slope of the tangent is
\[
\frac{\dd y}{\dd x}=\tan\theta,
\]
we have
\begin{equation}
\sin\theta=\frac{\dd y}{\sqrt{(\dd x)^2+(\dd y)^2}}=\frac{\dd y}{\dd s}. \tag{8}
\end{equation}
Substituting (8) into (7) gives
\begin{equation}
\ddot{s}=-g\frac{\dd y}{\dd s}. \tag{9}
\end{equation}
Substituting (4) and (5) into (9) gives
\begin{equation}
\ddot{s}=-g\,\frac{\sin\omega t}{2\cos\dfrac{\omega t}{2}}=-g\sin\frac{\omega t}{2}. \tag{10}
\end{equation}
Substituting (6) into (10) gives
\begin{equation}
\ddot{s}+\frac{g}{4R}s=0. \tag{11}
\end{equation}
Equation (11) shows that the small ring performs a motion similar to simple harmonic motion along this steel wire, and its period is
\[
T=2\pi\sqrt{\frac{4R}{g}}=4\pi\sqrt{\frac{R}{g}}.
\]
Hence the time required for the small ring to slide from the highest point of the steel wire to the lowest point is
\[
t=\frac{1}{4}T=\pi\sqrt{\frac{R}{g}}.
\]
\end{solution}

\begin{solution}{1.49}
(1) The rectangular coordinate system is chosen as shown in the figure. The source of sound is located at the origin $O$ of the coordinate system. The wind that blows over the ground is divided along the vertical ($y$) direction into many thin layers, and the wind speed within each thin layer $\dd y$ is uniform. The sound wave is refracted successively at the interfaces of the thin layers. Suppose that a sound wave is incident on an interface at an angle of incidence $\theta$ in the layer where the wind speed is $v$, and enters the layer where the wind speed is $v+\dd v$ at an angle of refraction $\theta+\dd\theta$; then from the law of refraction,
\[
\frac{\sin(\theta+\dd\theta)}{\sin\theta}
=\frac{V_s+(v+\dd v)\sin(\theta+\dd\theta)}{V_s+v\sin\theta}.
\]
Expanding the above expression, neglecting small quantities of higher order and rearranging, we obtain
\[
\frac{\dd v}{V_s}=\frac{\dd\left(\sin\theta\right)}{\sin^2\theta}.
\]

\begin{figbox}{1.49}
\begin{tikzpicture}[scale=0.9,>=latex]
  % axes
  \draw[axis] (0,0)--(7.8,0) node[right] {$x$};
  \draw[axis] (0,0)--(0,4.8) node[above] {$y$};
  \node[below left] at (0,0) {$O$};
  % wind layers
  \foreach \yy in {0.6,1.1,...,4.3}{\draw[thin,gray] (0.1,\yy)--(7.3,\yy);}
  % ray path
  \draw[thickline] (0,0) .. controls (2.6,1.05) and (4.3,2.1) .. (5.3,2.6)
     .. controls (6.1,2.95) and (6.7,3.55) .. (7.15,4.25);
  % angle at the origin
  \draw[thickline] (1.0,0) arc (0:22:1.0);
  \node at (1.25,0.28) {$\theta_0$};
  % normal at the interface point
  \coordinate (Q) at (5.3,2.6);
  \draw[dashedline] (Q)--++(0,1.6);
  \draw[dashedline] (Q)--++(0,-1.15);
  \fill (Q) circle (1.3pt);
  % angles theta and theta+dtheta
  \draw[thickline] ($(Q)+(0,0.85)$) arc (90:29.4:0.85);
  \node at ($(Q)+(0.44,0.98)$) {$\theta+\dd\theta$};
  \draw[thickline] ($(Q)+(0,-0.85)$) arc (270:206.6:0.85);
  \node at ($(Q)+(-0.55,-0.85)$) {$\theta$};
  % layer thickness
  \draw[dim] (4.35,2.1)--(4.35,2.6) node[midway,left] {$\dd y$};
  \draw[dashedline] (4.35,2.1)--(6.2,2.1);
  % wind velocities
  \draw[force] (5.9,2.1)--++(1.0,0) node[right] {$v$};
  \draw[force] (5.9,2.62)--++(1.15,0) node[right] {$v+\dd v$};
\end{tikzpicture}
\end{figbox}

Substituting $v=v_0+ky^2$ into the above expression and integrating both sides,
\[
\int_0^y\frac{2ky}{V_s}\,\dd y=\int_{\theta_0}^{\theta}\frac{\dd\left(\sin\theta\right)}{\sin^2\theta},
\]
we obtain
\begin{equation}
\frac{k}{V_s}y^2=\frac{1}{\sin\theta_0}-\frac{1}{\sin\theta}. \tag{1}
\end{equation}
Since
\[
\sin\theta=\frac{\dd x}{\sqrt{(\dd x)^2+(\dd y)^2}}=\frac{1}{\sqrt{1+\left(\dfrac{\dd y}{\dd x}\right)^2}},
\]
substitution of this relation into (1) gives
\[
\frac{\dd y}{\dd x}=\pm\sqrt{\frac{1}{\sin^2\theta_0}\left(1-\frac{k\sin\theta_0}{V_s}y^2\right)^2-1}.
\]
By the condition given in the problem, $\dfrac{ky^2}{V_s}\ll1$, that is, $\dfrac{k\sin\theta_0}{V_s}y^2\ll1$. Hence, expanding the above expression and neglecting small quantities of higher order, we obtain
\[
\begin{aligned}
\frac{\dd y}{\dd x}&=\pm\sqrt{\frac{1}{\sin^2\theta_0}\left(1-\frac{2k\sin\theta_0}{V_s}y^2\right)-1}\\
&=\pm\sqrt{\cot^2\theta_0-\frac{2k}{V_s\sin\theta_0}y^2},
\end{aligned}
\]
or
\[
\frac{\dd y}{\sqrt{\cot^2\theta_0-\dfrac{2k}{V_s\sin\theta_0}y^2}}=\pm\dd x.
\]
Integrating both sides of the above expression gives
\[
\sqrt{\frac{V_s\sin\theta_0}{2k}}\arcsin\left(\sqrt{\frac{2k\sin\theta_0}{V_s}}\cdot\frac{y}{\cos\theta_0}\right)=\pm x+C.
\]
From the boundary conditions $x_0=0$, $y_0=0$ the constant of integration is found to be $C=0$; dropping the ``$-$'' sign in the above expression, we obtain the trajectory along which the sound wave propagates:
\begin{equation}
y=\cos\theta_0\sqrt{\frac{V_s}{2k\sin\theta_0}}\sin\left(\sqrt{\frac{2k}{V_s\sin\theta_0}}\,x\right). \tag{2}
\end{equation}

(2) When the sound wave reaches the ground, $y=0$, that is,
\[
y=\cos\theta_0\sqrt{\frac{V_s}{2k\sin\theta_0}}\sin\left(\sqrt{\frac{2k}{V_s\sin\theta_0}}\,x\right)=0.
\]
Hence
\[
\sqrt{\frac{2k}{V_s\sin\theta_0}}\,x=0\quad\text{or}\quad\pi.
\]
The value $0$ corresponds to the position of the source of sound; therefore the distance $x$ from the source to the position at which the sound can be heard clearly is
\[
x=\pi\sqrt{\frac{V_s\sin\theta_0}{2k}}.
\]
Since $k$ is very small, the value of $x$ is comparatively large, which shows that the sound can be transmitted rather far.
\end{solution}

\begin{solution}{1.50}
(1) Let the equation of the incident wave be
\[
y_1=A\cos\left(\omega t-kx+\phi_1\right).
\]
In the above expression $k=\dfrac{2\pi}{\lambda}$ is the wave number. The equation of the vibration excited by this wave at $x=-a$ is
\[
y_1(x=-a)=A\cos\left(\omega t+ka+\phi_1\right).
\]
According to the problem this vibration is given by $y=A\cos(\omega t+\phi_0)$; hence
\[
ka+\phi_1=\phi_0,
\]
and we obtain
\[
\phi_1=\phi_0-ka.
\]
Thus
\[
y_1=A\cos\left(\omega t-kx+\phi_0-ka\right).
\]
Let the equation of the reflected wave be
\[
y_2=B\cos\left(\omega t+kx+\phi_2\right).
\]
Since $x=a$ is a fixed end, the reflection there is total, that is,
\[
B=A.
\]
At this fixed end the reflected wave suffers a half-wave loss; hence
\[
ka+\phi_2=-ka+\phi_0-ka+\pi,
\]
and we obtain
\[
\phi_2=\phi_0-3ka+\pi.
\]
Thus
\[
y_2=A\cos\left(\omega t+kx+\phi_0-3ka+\pi\right).
\]

(2) The incident wave and the reflected wave superpose to form a standing wave:
\[
\begin{aligned}
y_{\text{st}}&=y_1+y_2
=2A\cos\left(kx-ka+\frac{\pi}{2}\right)\cos\left(\omega t+\phi_0-2ka+\frac{\pi}{2}\right)\\
&=2A\sin k(x-a)\sin\left(\omega t+\phi_0-2ka\right).
\end{aligned}
\]
From the standing-wave equation shown above, the positions of the nodes are obtained as follows:
\[
k(x-a)=n\pi\qquad(n=0,-1,-2,\cdots),
\]
that is,
\[
x=a+\frac{n\pi}{k}\qquad(n=0,-1,-2,\cdots).
\]
Since the wave number $k$ is
\[
k=\frac{2\pi}{\lambda}=\frac{\omega}{v}=\omega\sqrt{\frac{\eta}{T}},
\]
the positions of the nodes are
\[
x=a+\frac{n\pi}{\omega}\sqrt{\frac{T}{\eta}}\qquad(n=0,-1,-2,\cdots),
\]
while the positions of the antinodes are
\[
k(x-a)=\left(n+\frac{1}{2}\right)\pi\qquad(n=-1,-2,\cdots),
\]
that is,
\[
x=a+\frac{(2n+1)\pi}{2\omega}\sqrt{\frac{T}{\eta}}\qquad(n=-1,-2,\cdots).
\]

(3) Take an arbitrary small segment $\dd x$ of the standing wave; the kinetic energy of this small segment is
\[
\begin{aligned}
\dd E_k&=\frac{1}{2}\eta\,\dd x\left(\frac{\partial y}{\partial t}\right)^2\\
&=2\eta A^2\omega^2\sin^2k(x-a)\cos^2\left(\omega t+\phi_0-2ka\right)\dd x.
\end{aligned}
\]
From the properties of a standing wave it is known that the total mechanical energy between two adjacent nodes is conserved, and that all points between the two nodes are in phase, that is, they pass through the equilibrium position simultaneously; at that instant the kinetic energy of every mass element reaches its maximum, while the potential energy is zero. The kinetic energy of the segment $\dd x$ at that instant is
\[
\dd E_{k\max}=2\eta A^2\omega^2\sin^2k(x-a)\,\dd x.
\]
Therefore the total energy between two adjacent nodes is
\[
\begin{aligned}
E&=\int\dd E_{k\max}=2\eta A^2\omega^2\int_{a-\frac{\pi}{k}}^{a}\sin^2k(x-a)\,\dd x\\
&=2\eta A^2\omega^2\cdot\frac{\pi}{2k}=\eta A^2\omega^2\frac{\pi}{k}\\
&=\pi\eta A^2\omega\sqrt{\frac{T}{\eta}}.
\end{aligned}
\]
\end{solution}

\begin{solution}{1.51}
The oscillating system shown in the figure has two normal modes: in mode 1, $B$ does not move, while $A$ and $C$ move towards each other. The eigenfrequency of this mode is $\omega_1=\sqrt{\dfrac{k}{m}}$. In mode 2, $A$ and $C$ move in the same direction, while $B$ moves in the opposite direction. The eigenfrequency of this mode is $\omega_2=\sqrt{\dfrac{k(M+2m)}{Mm}}$. A general motion is a superposition of the oscillations of these two normal modes.

\begin{figbox}{1.51}
\begin{tikzpicture}[scale=0.95,>=latex]
  % ground
  \draw[thickline] (-0.3,0)--(7.0,0);
  \fill[hatch] (-0.3,-0.22) rectangle (7.0,0);
  % block A
  \draw[thickline] (0,0) rectangle (0.9,0.8);
  \node at (0.45,0.4) {$m$};
  \node[above] at (0.45,0.82) {$A$};
  % spring A-B
  \draw[thickline,decorate,decoration={coil,aspect=0.6,segment length=4pt,amplitude=3pt}]
    (0.9,0.42)--(2.5,0.42);
  \node[above] at (1.7,0.52) {$k$};
  % block B
  \draw[thickline] (2.5,0) rectangle (3.7,0.9);
  \node at (3.1,0.45) {$M$};
  \node[above] at (3.1,0.92) {$B$};
  % spring B-C
  \draw[thickline,decorate,decoration={coil,aspect=0.6,segment length=4pt,amplitude=3pt}]
    (3.7,0.47)--(5.3,0.47);
  \node[above] at (4.5,0.57) {$k$};
  % block C
  \draw[thickline] (5.3,0) rectangle (6.2,0.8);
  \node at (5.75,0.4) {$m$};
  \node[above] at (5.75,0.82) {$C$};
  % initial velocity of A
  \draw[force] (0.15,1.15)--(1.05,1.15) node[right] {$v_0$};
  % displacement labels
  \draw[dim] (0,-0.55)--(0.9,-0.55) node[midway,below] {$x_1$};
  \draw[dim] (2.5,-0.55)--(3.7,-0.55) node[midway,below] {$x_2$};
  \draw[dim] (5.3,-0.55)--(6.2,-0.55) node[midway,below] {$x_3$};
\end{tikzpicture}
\end{figbox}

After $A$ is given the initial speed $v_0$, the centre of mass of the system begins to move uniformly, and the speed of the centre of mass $v_C$ is
\[
v_C=\frac{mv_0}{M+2m}.
\]
In the centre-of-mass frame the displacements of the three bodies can be expressed respectively as
\begin{align}
x_1'&=A_1\sin\omega_1t+A_2\sin\omega_2t+x_{10}' \tag{1}\\
x_2'&=-B_2\sin\omega_2t+x_{20}' \tag{2}\\
x_3'&=-A_1\sin\omega_1t+A_2\sin\omega_2t+x_{30}' \tag{3}
\end{align}
Since $x_1'=x_2'=x_3'=0$ at $t=0$, we have $x_{10}'=x_{20}'=x_{30}'=0$ in the above three expressions. Differentiating the above three expressions with respect to time $t$ gives
\begin{align}
\dot{x}_1'&=A_1\omega_1\cos\omega_1t+A_2\omega_2\cos\omega_2t \tag{4}\\
\dot{x}_2'&=-B_2\omega_2\cos\omega_2t \tag{5}\\
\dot{x}_3'&=-A_1\omega_1\cos\omega_1t+A_2\omega_2\cos\omega_2t \tag{6}
\end{align}
At $t=0$ we have
\begin{align}
\dot{x}_1'&=v_0-v_C=v_0-\frac{mv_0}{M+2m}=\frac{M+m}{M+2m}v_0, \tag{7}\\
\dot{x}_2'&=-v_C=-\frac{mv_0}{M+2m}, \tag{8}\\
\dot{x}_3'&=-v_C=-\frac{mv_0}{M+2m}. \tag{9}
\end{align}
Substituting (7), (8) and (9) respectively into (4), (5) and (6), we obtain
\[
\begin{aligned}
A_1&=\frac{v_0}{2\omega_1},\\
A_2&=\frac{Mv_0}{2\omega_2(M+2m)},\\
B_2&=\frac{mv_0}{\omega_2(M+2m)}.
\end{aligned}
\]
Substituting the above relations into (1), (2) and (3), and using the principle of relative motion, we can obtain the relations for the variation with time of the displacements of the three bodies relative to the laboratory frame of reference.
\end{solution}

\begin{solution}{1.52}
\begin{note}
The first page of this solution (page 124 of the original) is missing from the available copy.
The translation therefore begins with equation (7), the first line preserved on the following
page.
\end{note}
\begin{align}
-k_1A_1-k_2A_2 &= \frac{m}{2}\left(-A_1\omega^2-A_2\omega^2\right), \tag{7}\\
-\frac{l}{2}\left(k_2A_2-k_1A_1\right)
&= \frac{I_0}{l}\left(-A_2\omega^2+A_1\omega^2\right). \tag{8}
\end{align}
The necessary and sufficient condition for the system of equations (7) and (8) to have a
non-zero solution is
\begin{equation*}
\begin{vmatrix}
k_1-\dfrac{m\omega^2}{2} & k_2-\dfrac{m\omega^2}{2}\\[8pt]
\dfrac{I_0\omega^2}{l}-\dfrac{lk_1}{2} & \dfrac{lk_2}{2}-\dfrac{I_0\omega^2}{l}
\end{vmatrix}=0.
\end{equation*}
Hence
\begin{equation*}
m^2\omega^4-4m(k_1+k_2)\omega^2+12k_1k_2=0.
\end{equation*}
Solving, we obtain
\begin{equation*}
\omega^2=\frac{2(k_1+k_2)\pm 2\sqrt{k_1^2+k_2^2-3k_1k_2}}{m}.
\end{equation*}
The values of $\omega$ corresponding to the ``$+$'' and ``$-$'' signs in the above expression
are the eigenfrequencies of the two normal modes of this system.
\end{solution}

\begin{solution}{1.53}
(1) Take the nozzle as the origin of coordinates, take some horizontal direction as the direction of the $x$-axis and the vertical upward direction as the direction of the $y$-axis. The water sprayed from the nozzle is symmetric about the $y$-axis, so it suffices to consider, in one cross-section (the $xOy$ plane), the jets of water sprayed by the nozzle in all directions, to find the envelope of the parabolas formed by these jets, and then to rotate this envelope about the $y$-axis through one complete revolution; the solid figure is thereby obtained.

In the $xOy$ plane, suppose that the jet of water from a certain hole is ejected at an elevation angle $\theta$; then the equations of motion in the $x$ and $y$ directions are
\begin{align}
x &= v_0\cos\theta\,t \tag{1}\\
y &= v_0\sin\theta\,t-\frac{1}{2}gt^2 \tag{2}
\end{align}
Combining (1) and (2) and eliminating $t$, we obtain
\begin{equation}
y=\tan\theta\,x-\frac{g}{2v_0^2\cos^2\theta}x^2. \tag{3}
\end{equation}
Equation (3) is the trajectory equation of the jet. Setting $u=\tan\theta$, equation (3) can be written as
\begin{equation}
\frac{gx^2}{2v_0^2}u^2-xu+\left(y+\frac{gx^2}{2v_0^2}\right)=0. \tag{4}
\end{equation}
If $x$ and $y$ are regarded as fixed values, then (4) is a quadratic equation in $u$. If this equation is to have real solutions, its discriminant must be greater than or equal to zero, that is,
\[
x^2-4\cdot\frac{gx^2}{2v_0^2}\left(y+\frac{gx^2}{2v_0^2}\right)\geq 0.
\]
Solving, we obtain
\begin{equation}
y\leq\frac{v_0^2}{2g}-\frac{g}{2v_0^2}x^2. \tag{5}
\end{equation}
The inequality expressed by (5) divides the $xOy$ plane into two parts by means of a parabola. Water can reach the various points below this parabola, but cannot reach the points on and above it; this parabola is the required envelope, as shown by the dashed line in Fig. 1.53(a). The equation of this parabola is
\begin{equation}
y=\frac{v_0^2}{2g}-\frac{g}{2v_0^2}x^2. \tag{6}
\end{equation}
Therefore the surface of the solid figure traced out by the water in the air is the paraboloid of revolution formed by rotating the parabola given by (6) about the $y$-axis through one complete revolution; the height of the solid figure is $\dfrac{v_0^2}{2g}$, and the radius of the region wetted on the ground is $R=\dfrac{v_0^2}{g}$.

\begin{figbox}{1.53(a)}
\begin{tikzpicture}[scale=0.95,>=latex]
  % axes
  \draw[axis] (-4.0,0)--(4.15,0) node[right] {$x$};
  \draw[axis] (0,-0.25)--(0,2.35) node[above] {$y$};
  \node[below left] at (0,0) {$O$};
  % ground
  \draw[thickline] (-3.7,0)--(3.7,0);
  % envelope (dashed)
  \draw[dashedline,domain=-3.6:3.6,samples=80] plot (\x,{1.8-0.138889*\x*\x});
  % jets
  \foreach \u/\xl in {0.35/2.245,0.55/3.040,0.75/3.456,1.0/3.6,1.35/3.444,1.9/2.967,2.8/2.281}{
    \draw[thickline,domain=0:\xl,samples=40] plot (\x,{\u*\x-0.138889*(1+\u*\u)*\x*\x});
  }
  \foreach \u/\xl in {-0.35/2.245,-0.55/3.040,-0.75/3.456,-1.0/3.6,-1.35/3.444,-1.9/2.967,-2.8/2.281}{
    \draw[thickline,domain=-\xl:0,samples=40] plot (\x,{\u*\x-0.138889*(1+\u*\u)*\x*\x});
  }
  % height and radius
  \draw[dim] (0,-0.45)--(3.6,-0.45) node[midway,below] {$R$};
  \draw[dim] (3.95,0)--(3.95,1.8) node[midway,right] {$\dfrac{v_0^2}{2g}$};
  \draw[dashedline] (0,1.8)--(3.95,1.8);
\end{tikzpicture}
\end{figbox}

(2) Let the radius of the spherical surface of the small nozzle be $r$, as shown in Fig. 1.53(b). The range of the water ejected at the angle $\theta$ in the figure is
\[
R=\frac{v_0^2\sin2\theta}{g}.
\]
Differentiating both sides of the above expression, we obtain
\[
\dd R=\frac{2v_0^2\cos2\theta}{g}\dd\theta.
\]
Therefore the water ejected from the small holes on the narrow annular band (the dashed band in the figure) on the spherical surface of the nozzle, which makes an angle from $\theta$ to $\theta+\dd\theta$ with the vertical axis, lands, at the points where it splashes down, on an annular band of radius $R$ and width $\dd R$, whose area $\dd S$ is
\[
\dd S=2\pi R\,\dd R=\frac{2\pi v_0^4\sin4\theta}{g^2}\dd\theta.
\]
Let the volume of water received per unit area per second within the region of splash-down points be $k$; then the volume of water received per second on the above annular band of splash-down points is
\[
\dd Q=k\,\frac{2\pi v_0^4\sin4\theta}{g^2}\dd\theta.
\]
The area $\dd S'$ of the narrow annular band $\theta\to\theta+\dd\theta$ on the spherical surface of the nozzle is
\[
\dd S'=2\pi r^2\sin\theta\,\dd\theta.
\]
Let the area of a small hole be $S_0$, and let the distribution of the number of small holes per unit area on the spherical surface be $N(\theta)$; then the volume of water sprayed per second from this narrow annular band, $\dd Q'$, is
\[
\dd Q'=N(\theta)S_0v_0\,\dd S'=2\pi r^2\sin\theta\,N(\theta)S_0v_0\,\dd\theta.
\]
From $\dd Q=\dd Q'$ we obtain
\[
k\,\frac{2\pi v_0^4\sin4\theta}{g^2}\dd\theta=2\pi r^2\sin\theta\,N(\theta)S_0v_0\,\dd\theta.
\]
From this expression $N(\theta)$ is obtained as
\[
N(\theta)=\frac{kv_0^3\sin4\theta}{r^2g^2S_0\sin\theta}=C\,\frac{\sin4\theta}{\sin\theta}.
\]
In the above expression the constant $C$ is a quantity related to $k$, $v_0$, $r$, $g$, $S_0$ and other such quantities.

\begin{figbox}{1.53(b)}
\begin{tikzpicture}[scale=1.0,>=latex]
  % spherical nozzle (upper half)
  \draw[thickline] (-2,0) arc (180:0:2);
  \draw[dashedline] (-2.55,0)--(2.55,0);
  % vertical axis
  \draw[dashedline] (0,0)--(0,2.65);
  % narrow band at theta
  \draw[very thick,dashedline] (40:2.02) arc (40:52:2.02);
  % radius to the band
  \draw[dashedline] (0,0)--(46:2.02);
  % angle theta
  \draw[thickline] (0,0.85) arc (90:46:0.85);
  \node at (0.30,1.02) {$\theta$};
  % ejected water direction
  \draw[force] (46:2.02)--++(46:0.85);
  \node at (2.05,1.72) {$\dd\theta$};
  % radius dimension
  \draw[dim] (0,-0.32)--(2,-0.32) node[midway,below] {$r$};
  \node[below left] at (0,0) {$O$};
\end{tikzpicture}
\end{figbox}
\end{solution}

\begin{solution}{1.54}
(1) Take a small volume element of cross-section $\dd S$ and depth $\dd r$ at a depth $h$ near the surface at polar angle $\theta$. Its volume is $\dd V=\dd S\cdot\dd r$ and its mass is $\dd m=\rho\,\dd V$. In the rotating reference frame that rotates together with the planet at angular velocity $\omega$, the inertial centrifugal force $\dd\vec{f}_i$ acting on this small volume element is
\[
\dd\vec{f}_i=\dd m\,\omega^2r\cos\theta=\rho\omega^2r\cos\theta\,\dd V,
\]
and its direction is shown in Fig. 1.54(a). The universal gravitational force $\dd\vec{F}$ acting on the small volume element is
\[
\dd F=\frac{G\,\dfrac{4}{3}\pi r^3\rho\,\dd m}{r^2}=\frac{4}{3}\pi G\rho^2r\,\dd V,
\]
its direction being along the radial direction $r$. The difference between the radial component of $\dd\vec{f}_i$ and $\dd F$ is balanced by the force produced at the two cross-sections of the small volume element, separated by $\dd r$, by the pressure difference $\dd p(r)$, that is,
\[
-\dd p(r)\,\dd S=\frac{4}{3}\pi G\rho^2r\,\dd V-\rho\omega^2r\cos^2\theta\,\dd V.
\]
Since $\rho=\dfrac{M}{\dfrac{4}{3}\pi R^3}$, the above expression can also be written as
\[
-\dd p(r)=\left(\frac{GM\rho}{R^3}-\rho\omega^2\cos^2\theta\right)r\,\dd r.
\]
Integrating both sides of the above expression as definite integrals, and since $p=0$ at $r=R$, we have
\[
-\int_{p(r)}^{0}\dd p(r)=\int_{r}^{R}\left(\frac{GM\rho}{R^3}-\rho\omega^2\cos^2\theta\right)r\,\dd r.
\]
Integrating, we obtain
\begin{equation}
p(r)=\frac{1}{2}\left(\frac{GM\rho}{R^3}-\rho\omega^2\cos^2\theta\right)\left(R^2-r^2\right). \tag{1}
\end{equation}
Substituting $r=R-h$ into the above expression, and since $h/R\ll1$, expanding the above expression and neglecting quantities of higher order in $h/R$, we obtain
\begin{align}
p(h)&=\frac{1}{2}\rho\left(\frac{GM}{R^3}-\omega^2\cos^2\theta\right)\left[R^2-(R-h)^2\right]\notag\\
&=\frac{1}{2}\rho\left(\frac{GM}{R^3}-\omega^2\cos^2\theta\right)R^2\left[1-\left(1-\frac{h}{R}\right)^2\right]\notag\\
&\approx\left(\frac{GM}{R^2}-\omega^2\cos^2\theta\,R\right)\rho h. \tag{2}
\end{align}
It is seen that within the range $h\ll R$ the pressure is approximately proportional to the depth.

\begin{figbox}{1.54(a)}
\begin{tikzpicture}[scale=0.95,>=latex]
  % planet
  \draw[thickline] (0,0) circle (2);
  \draw[axis] (-2.55,0)--(2.85,0) node[right] {$x$};
  \draw[axis] (0,-2.35)--(0,2.65) node[above] {$y$};
  \node[below left] at (0,0) {$O$};
  % rotation
  \draw[thickline,->] (0.6,2.42) arc (-5:185:0.6 and 0.25);
  \node at (1.1,2.58) {$\omega$};
  % radius and angle
  \draw[dashedline] (0,0)--(40:2);
  \draw[thickline] (0.95,0) arc (0:40:0.95);
  \node at (1.3,0.12) {$\theta$};
  \draw[dim] (-2,-0.3)--(0,-0.3) node[midway,below] {$R$};
  % small volume element
  \draw[thickline] (30:1.62) arc (30:50:1.62);
  \draw[thickline] (30:1.80) arc (30:50:1.80);
  \draw[thickline] (30:1.62)--(30:1.80);
  \draw[thickline] (50:1.62)--(50:1.80);
  % h and dr
  \draw[dim] (24:1.80)--(24:2.02) node[midway,left] {$h$};
  \draw[dim] (56:1.62)--(56:1.80) node[midway,right] {$\dd r$};
  % forces
  \draw[force] (40:1.71)--++(-0.38,-0.32) node[left] {$\dd\vec{F}$};
  \draw[force] (40:1.78)--++(0.55,0) node[right] {$\dd\vec{f}_i$};
\end{tikzpicture}
\end{figbox}

(2) We introduce the inertial centrifugal potential energy in the rotating reference frame (hereafter called the centrifugal potential energy). Taking the axis of rotation as the zero point of this potential energy, the centrifugal potential energy per unit volume at the equatorial surface of this planet (i.e. the centrifugal potential) is
\[
U_{\text{cent}}(\theta=0)=-\int_{0}^{R_e}\rho\omega^2r\,\dd r=-\frac{1}{2}\rho\omega^2R_e^2.
\]
The gravitational potential and the total potential there are respectively
\[
U_{\text{grav}}^{(\theta=0)}=-\frac{\rho GM}{R_e},
\]
\[
U_{\text{tot}}^{(\theta=0)}=U_{\text{cent}}+U_{\text{grav}}=-\frac{1}{2}\rho\omega^2R_e^2-\frac{\rho GM}{R_e}.
\]
At the polar surface of the planet, on the other hand, the centrifugal potential is zero, so its total potential is
\[
U_{\text{tot}}\left(\theta=\frac{\pi}{2}\right)=-\frac{\rho GM}{R_p}.
\]
The surface of the planet is an equipotential surface everywhere; hence from
\[
U_{\text{tot}}(\theta=0)=U_{\text{tot}}\left(\theta=\frac{\pi}{2}\right)
\]
it follows that
\[
-\frac{1}{2}\rho\omega^2R_e^2-\frac{\rho GM}{R_e}=-\frac{\rho GM}{R_p},
\]
that is,
\begin{equation}
R_e-R_p=\frac{\omega^2R_e^2R_p}{2GM}\approx\frac{\omega^2R^4}{2GM}. \tag{3}
\end{equation}
If the mass and the radius of this planet are both similar to those of the Earth, and its rotation is also the same as that of the Earth, that is, $R=6.4\times10^6\unitv{m}$, $M=6\times10^{24}\unitv{kg}$ and $\omega=\dfrac{2\pi}{T}=7.3\times10^{-5}\unitv{s^{-1}}$, then substituting these data into (3) gives the estimate
\[
R_e-R_p\approx11\unitv{km}.
\]

This sub-question can also be solved by the following method. Take a small volume element $\dd V=R_p\,\dd\theta\cdot\dd S$ at polar angle $\theta\to\theta+\dd\theta$ and at a distance $R_p$ from the centre $O$, as shown in Fig. 1.54(b). The transverse component of the inertial centrifugal force acting on this volume is balanced by the force caused by the pressure difference $\dd p(\theta)$ on the two transverse cross-sections $\dd S$ of the small volume element, that is,
\[
-\dd p(\theta)\,\dd S=\rho R_p\,\dd\theta\,\dd S\,\omega^2R_p\cos\theta\sin\theta.
\]
Integrating both sides of the above expression as definite integrals, and since $p=0$ at $\theta=\dfrac{\pi}{2}$, we have
\[
-\int_{p(\theta=0)}^{0}\dd p(\theta)=\int_{0}^{\frac{\pi}{2}}\rho\omega^2R_p^2\sin\theta\cos\theta\,\dd\theta.
\]
Integrating, we obtain
\begin{equation}
p(\theta=0)=\frac{1}{2}\rho\omega^2R_p^2. \tag{4}
\end{equation}
From equation (1) of sub-question (1) it is known that the pressure at $\theta=0$, $r=R_p$ is
\begin{equation}
p(\theta=0,\,h)=\frac{1}{2}\left(\frac{GM}{R_e^3}-\omega^2\right)\rho\left(R_e^2-R_p^2\right). \tag{5}
\end{equation}

\begin{figbox}{1.54(b)}
\begin{tikzpicture}[scale=0.95,>=latex]
  \draw[thickline] (0,0) circle (2);
  \draw[axis] (-2.5,0)--(2.75,0) node[right] {$x$};
  \draw[axis] (0,-2.35)--(0,2.65) node[above] {$y$};
  \node[below left] at (0,0) {$O$};
  \draw[thickline,->] (0.6,2.42) arc (-5:185:0.6 and 0.25);
  \node at (1.05,2.55) {$\omega$};
  % radius to the element
  \draw[dashedline] (0,0)--(55:2);
  \draw[thickline] (0.9,0) arc (0:55:0.9);
  \node at (1.2,0.42) {$\theta$};
  \node[above left] at (55:1.18) {$R_p$};
  % element
  \draw[thickline] (46:1.70)--(64:1.70)--(64:1.94)--(46:1.94)--cycle;
  % dS on the outer face
  \draw[force] (55:1.94)--++(55:0.42) node[above right] {$\dd S$};
  % d theta
  \node at (38:1.78) {$\dd\theta$};
\end{tikzpicture}
\end{figbox}

The pressures $p$ in (4) and (5) must be equal, hence
\[
\frac{1}{2}\omega^2R_p^2=\frac{1}{2}\left(\frac{GM}{R_e^3}-\omega^2\right)\left(R_e^2-R_p^2\right).
\]
That is,
\begin{equation}
R_e-R_p=\frac{\omega^2R_e^5}{GM\left(R_e+R_p\right)}\approx\frac{\omega^2R_e^4}{2GM}\approx\frac{\omega^2R^4}{2GM}. \tag{6}
\end{equation}
\end{solution}

\begin{solution}{1.55}
First consider the tide caused by the Moon. Let the mass of the Moon be $m$, the mass of the Earth be $M_e$, the radius of the solid part of the Earth be $r_0$, and the distance from the Moon to the centre of the Earth be $r_{me}$. The distance from the centre of mass $C$ of the Moon--Earth system to the centre of the Earth is $r_{Ce}$. In the cross-section shown in the figure, the position of a point $P$ under consideration on the water surface is expressed by polar coordinates $(r,\theta)$, and the distance from $C$ to $P$ is denoted by $\eta$. Suppose that the Moon--Earth system rotates about an axis through the centre of mass $C$ perpendicular to the plane of the figure, with angular velocity $\omega$, and that the Earth's axis of rotation is also perpendicular to the plane of the figure. This setup makes it convenient to calculate the relation between the tidal variation at the equator and $\theta$ over a 24-hour period. In the above setup, $r_{Ce}$ is
\begin{equation}
r_{Ce}=\frac{m}{M_e+m}r_{me}. \tag{1}
\end{equation}
In addition,
\begin{equation}
\frac{GM_em}{r_{me}^2}=M_e\omega^2r_{Ce}. \tag{2}
\end{equation}
From (1) and (2) we obtain
\begin{equation}
\omega=\sqrt{\frac{G\left(M_e+m\right)}{r_{me}^3}}. \tag{3}
\end{equation}
Moreover, in the above setup $\eta$ is
\begin{equation}
\eta=\left(r_{Ce}^2+r^2+2rr_{Ce}\cos\theta\right)^{\frac{1}{2}}. \tag{4}
\end{equation}
The potential energy per unit mass of water at the point $P(r,\theta)$ on the Earth's surface, that is, the potential $V$, is composed of three potentials. The first is the Earth's gravitational potential $V_e$:
\begin{equation}
V_e=-\frac{GM_e}{r}. \tag{5}
\end{equation}
The second is the Moon's gravitational potential $V_m$:
\begin{equation}
V_m=\frac{-Gm}{\left(r_{me}^2+r^2+2rr_{me}\cos\theta\right)^{\frac{1}{2}}}. \tag{6}
\end{equation}

\begin{figbox}{1.55}
\begin{tikzpicture}[scale=0.95,>=latex]
  % Moon
  \fill (-6,0) circle (2.6pt);
  \node[above right] at (-6,0) {$m$};
  % Earth: water surface and solid body
  \draw[thickline] (0,0) circle (1.8);
  \draw[thickline] (0,0) circle (1.68);
  \node[below left] at (0,0) {$O$};
  % centre of mass C
  \fill (-2.2,0) circle (1.8pt);
  \node[above left] at (-2.2,0.06) {$C$};
  % line from Moon to Earth, and x axis
  \draw[thickline] (-6,0)--(-1.8,0);
  \draw[dashedline] (0,0)--(2.15,0) node[right] {$x$};
  % point P
  \coordinate (P) at (45:1.8);
  \draw[thickline] (0,0)--(P);
  \draw[dashedline] (-2.2,0)--(P);
  \node[above left] at (45:1.2) {$r$};
  \node[above left] at (-0.55,0.88) {$\eta$};
  \node[above right] at (P) {$P$};
  % angle
  \draw[thickline] (0.8,0) arc (0:45:0.8);
  \node at (0.98,0.33) {$\theta$};
  % radius of the solid part
  \draw[->,line width=0.5pt] (0,0)--(-55:1.68);
  \node[below right] at (-55:0.85) {$r_0$};
  % dimensions
  \draw[dim] (-6,-2.45)--(0,-2.45) node[midway,below] {$r_{me}$};
  \draw[dashedline] (-6,-0.12)--(-6,-2.6);
  \draw[dashedline] (0,-1.95)--(0,-2.6);
  \draw[dim] (-2.2,-1.9)--(0,-1.9) node[midway,below] {$r_{Ce}$};
  \draw[dashedline] (-2.2,-0.12)--(-2.2,-2.0);
\end{tikzpicture}
\end{figbox}

Since $\dfrac{r}{r_{me}}\ll1$, we may expand (6) in powers of $\dfrac{r}{r_{me}}$ and neglect terms of higher order than $\left(\dfrac{r}{r_{me}}\right)^2$, obtaining
\begin{align}
V_m&=\frac{-Gm}{r_{me}}\left[1+\left(\frac{r}{r_{me}}\right)^2+2\left(\frac{r}{r_{me}}\right)\cos\theta\right]^{-\frac{1}{2}}\notag\\
&\approx\frac{-Gm}{r_{me}}\left[1-\left(\frac{r}{r_{me}}\right)\cos\theta+\frac{r^2}{2r_{me}^2}\left(3\cos^2\theta-1\right)\right]. \tag{7}
\end{align}
The third is the centrifugal potential $V_{\text{cent}}$ in the reference frame of the centre of mass of the Moon--Earth system:
\begin{align}
V_{\text{cent}}&=-\frac{1}{2}\omega^2\eta^2=-\frac{1}{2}\frac{G\left(M_e+m\right)}{r_{me}^3}\left(r_{Ce}^2+r^2+2r_{Ce}r\cos\theta\right)\notag\\
&=-\frac{Gm}{r_{me}}\left[\frac{m}{2\left(M_e+m\right)}+\frac{M_e+m}{2m}\left(\frac{r}{r_{me}}\right)^2+\left(\frac{r}{r_{me}}\right)\cos\theta\right]. \tag{8}
\end{align}
Hence $V$ is
\begin{equation}
V=V_e+V_m+V_{\text{cent}}. \tag{9}
\end{equation}
Substituting (5), (7) and (8) into (9), and noting that the surface of the Earth is an equipotential surface, we obtain
\begin{equation}
-\frac{GM_e}{r}-\frac{Gm}{r_{me}}\left[\frac{2M_e+3m}{2\left(M_e+m\right)}+\frac{M_e+3m\cos^2\theta}{2m}\left(\frac{r}{r_{me}}\right)^2\right]=C. \tag{10}
\end{equation}
Substituting $r=r_0+h$ into (10) and making reasonable approximations, we obtain
\[
\frac{GM_e}{r_0^2}h-\frac{3Gmr_0^2}{2r_{me}^3}\cos^2\theta=C.
\]
That is,
\begin{equation}
h=\frac{3mr_0^4\cos^2\theta}{2M_er_{me}^3}+C. \tag{11}
\end{equation}
At $\theta=\dfrac{\pi}{2}$ the height of the water surface is a minimum $h_0$, so the $C$ in (11) is $C=h_0$; while at $\theta=0$ the tide is a maximum, which we take as $h_{\max}$. Hence the tidal variation $(\Delta h)_m$ at the equator during a 24-hour period is
\begin{equation}
(\Delta h)_m=h_{\max}-h_0=\frac{3mr_0^4}{2M_er_{me}^3}. \tag{12}
\end{equation}
Similar to the calculation of the tide caused by the Moon, the tidal variation $(\Delta h)_s$ caused by the Sun is
\begin{equation}
(\Delta h)_s=\frac{3M_sr_0^4}{2M_er_{se}^3}. \tag{13}
\end{equation}
In (13) $M_s$ is the mass of the Sun and $r_{se}$ is the distance from the Sun to the Earth; hence from (12) and (13) we obtain
\[
\frac{(\Delta h)_s}{(\Delta h)_m}=\frac{M_sr_{me}^3}{mr_{se}^3}
=\frac{2.0\times10^{30}\times\left(3.8\times10^8\right)^3}{7.3\times10^{22}\times\left(1.5\times10^{11}\right)^3}\approx0.44.
\]
\end{solution}

\begin{solution}{1.56}
Take the planet and the satellite as the system; the angular momentum $L$ of this system about the centre of the planet is
\begin{equation}
L=I\Omega+mr^2\omega. \tag{1}
\end{equation}
Owing to the action of tidal friction, both the angular velocity of rotation $\Omega$ of the planet and the angular velocity $\omega$ of the satellite's circular orbital motion about the planet will change, and the change in $\omega$ will in turn lead to a change in the orbital radius $r$. The relations among these changes can be obtained by differentiating (1):
\begin{equation}
\dd L=I\,\dd\Omega+mr^2\,\dd\omega+2mr\omega\,\dd r. \tag{2}
\end{equation}
Since the satellite moves in a circle about the planet, we have
\[
G\frac{Mm}{r^2}=m\omega^2r.
\]
That is,
\begin{equation}
\omega^2=\frac{GM}{r^3}. \tag{3}
\end{equation}
Differentiating both sides of (3), we obtain
\[
2\omega\,\dd\omega=-\frac{3GM}{r^4}\dd r=-\frac{3\omega^2}{r}\dd r.
\]
Hence the relation between $\dd\omega$ and $\dd r$ is
\begin{equation}
\dd r=-\frac{2r}{3\omega}\dd\omega. \tag{4}
\end{equation}
Substituting (4) into (2), we obtain
\begin{equation}
\dd L=I\,\dd\Omega-\frac{1}{3}mr^2\,\dd\omega. \tag{5}
\end{equation}
Since the angular momentum of the system about the centre of the planet is conserved, that is, $\dd L=0$, from (5) we obtain the relation between $\dd\Omega$ and $\dd\omega$:
\begin{equation}
\dd\Omega=\frac{mr^2}{3I}\dd\omega. \tag{6}
\end{equation}

(2) The mechanical energy $E$ of the system is
\begin{equation}
E=\frac{1}{2}I\Omega^2+\frac{1}{2}mr^2\omega^2-G\frac{Mm}{r}. \tag{7}
\end{equation}
Since tidal friction causes a decrease in the mechanical energy of the system, differentiating (7) gives
\[
\dd E=I\Omega\,\dd\Omega+mr^2\omega\,\dd\omega+mr\omega^2\,\dd r+\frac{GMm}{r^2}\dd r.
\]
Substituting (3) into the above expression, we obtain
\[
\dd E=I\Omega\,\dd\Omega+mr^2\omega\,\dd\omega+2mr\omega^2\,\dd r.
\]
Substituting (4) and (6) into the above expression, we obtain
\begin{equation}
\frac{\dd E}{\dd\omega}=\frac{1}{3}mr^2\left(\Omega-\omega\right). \tag{8}
\end{equation}
When the system finally reaches a stable state, that is, a state in which the energy does not change, $\dfrac{\dd E}{\dd\omega}=0$; hence from (8) we obtain
\[
\Omega=\omega.
\]
From this it is seen that when the orbital angular velocity $\omega$ of the satellite is the same as the angular velocity of rotation $\Omega$ of the planet, the mechanical energy of the system no longer changes; at this time there is no longer any phenomenon of the tide rising and falling on the planet, and tidal friction no longer exists. For the system to be in a stable state, the mechanical energy at this time should take a minimum value, that is,
\[
\frac{\dd^2E}{\dd\omega^2}>0.
\]
Using (8), we obtain
\[
\frac{\dd^2E}{\dd\omega^2}=\frac{1}{3}mr^2\left(\frac{\dd\Omega}{\dd\omega}-1\right)+\frac{2}{3}mr\left(\Omega-\omega\right)\frac{\dd r}{\dd\omega}>0.
\]
Substituting (6) and $\Omega=\omega$ into the above expression, we obtain
\[
\frac{mr^2}{I}>3.
\]
This is the condition that the system should satisfy.
\end{solution}

\begin{solution}{1.57}
The radial component $v_r$ and the transverse component $v_\theta$ of the wind velocity $\vec{v}$ at point $B$ near the hangar are respectively
\[
v_r=-\left.\frac{\partial\phi}{\partial r}\right|_{r=R}
=v_0\left(1-\frac{R^2}{r^2}\right)\cos\theta\,\bigg|_{r=R}=0,
\]
\[
v_\theta=-\frac{1}{r}\left.\frac{\partial\phi}{\partial\theta}\right|_{r=R}
=-v_0\left(1+\frac{R^2}{r^2}\right)\sin\theta\,\bigg|_{r=R}=-2v_0\sin\theta.
\]
It is known from the statement of the problem that the wind speed far away is $v_0$ and the pressure there is $p_0$; hence from Bernoulli's equation the pressure $p(\theta)$ at point $B$ is
\[
p(\theta)+\frac{1}{2}\rho v_\theta^2=p_0+\frac{1}{2}\rho v_0^2.
\]
That is,
\begin{align}
p(\theta)&=p_0+\frac{1}{2}\rho\left(v_0^2-v_\theta^2\right)\notag\\
&=p_0+\frac{1}{2}\rho v_0^2\left(1-4\sin^2\theta\right).
\end{align}
The pressure inside the hangar, that is, the pressure $p_A$ at the door $A$ of the hangar at $\theta=\pi$, is
\[
p_A=p_0+\frac{1}{2}\rho v_0^2.
\]
Hence the force $F$ exerted by the wind on the hangar is
\begin{align}
F&=\int_{0}^{\pi}\left[p_A-p(\theta)\right]\sin\theta\,\dd s\notag\\
&=\frac{1}{2}\rho v_0^2LR\int_{0}^{\pi}4\sin^2\theta\cdot\sin\theta\,\dd\theta\notag\\
&=2\rho v_0^2LR\left(\frac{1}{3}\cos^3\theta-\cos\theta\right)\bigg|_{0}^{\pi}\notag\\
&=\frac{8}{3}\rho v_0^2LR\approx9.0\times10^5\unitv{N}.
\end{align}

\begin{figbox}{1.57}
\begin{tikzpicture}[scale=0.95,>=latex]
  % ground
  \draw[thickline] (-3.3,0)--(4.0,0);
  \draw[->] (3.95,0)--(4.35,0) node[right] {$x$};
  % hangar (semicircular cylinder)
  \draw[thickline] (-2,0) arc (180:0:2);
  \node[below] at (0,0) {$O$};
  \node[below left] at (-2,0) {$A$};
  % point B and angle
  \coordinate (B) at (60:2);
  \draw[dashedline] (0,0)--(B);
  \draw[thickline] (0.9,0) arc (0:60:0.9);
  \node at (1.1,0.4) {$\theta$};
  \fill (B) circle (1.2pt) node[above right] {$B$};
  % radius
  \draw[->,line width=0.5pt] (0,0)--(125:1.97);
  \node[above left] at (125:1.2) {$R$};
  % streamlines
  \draw[thickline] plot[smooth,tension=0.75] coordinates {(-3.2,0.2) (-2.3,0.22) (-1.5,1.38) (-0.7,1.98) (0,2.14) (0.7,1.98) (1.5,1.38) (2.3,0.22) (3.2,0.2)};
  \draw[thickline] plot[smooth,tension=0.75] coordinates {(-3.2,0.62) (-2.3,0.66) (-1.5,1.62) (-0.7,2.28) (0,2.46) (0.7,2.28) (1.5,1.62) (2.3,0.66) (3.2,0.62)};
  \draw[thickline] plot[smooth,tension=0.75] coordinates {(-3.2,1.12) (-2.3,1.18) (-1.5,2.02) (-0.7,2.72) (0,2.92) (0.7,2.72) (1.5,2.02) (2.3,1.18) (3.2,1.12)};
  \draw[thickline] plot[smooth,tension=0.75] coordinates {(-3.2,1.7) (-2.3,1.78) (-1.5,2.5) (-0.7,3.2) (0,3.42) (0.7,3.2) (1.5,2.5) (2.3,1.78) (3.2,1.7)};
\end{tikzpicture}
\end{figbox}
\end{solution}

\begin{solution}{1.58}
(1) At the moment of release the angle between the two rods and the ground is $\theta=\alpha=60\degree$, and the forces acting on one of the rods, $AB$, are shown in Fig. 1.58(a). Let the horizontal ($x$) and vertical ($y$) components of the acceleration of the centre of mass $C$ of this rod at this moment be $a_{Cx}$ and $a_{Cy}$ respectively; then the equations of motion of the centre of mass of this rod along the $x$ and $y$ directions are respectively
\begin{align}
F&=ma_{Cx} \tag{1}\\
mg-N&=ma_{Cy} \tag{2}
\end{align}
The equation of rotation of this rod about its centre of mass $C$ is
\begin{equation}
\frac{1}{2}lN\cos\theta-F\cdot\frac{l}{2}\sin\theta=\frac{1}{12}ml^2\ddot{\theta}. \tag{3}
\end{equation}

\begin{figbox}{1.58(a)}
\begin{tikzpicture}[scale=0.9,>=latex]
  % ground
  \draw[thickline] (-1.5,0)--(5.3,0);
  % rod
  \coordinate (B) at (0,0);
  \coordinate (A) at (2.3,3.98);
  \draw[thickline] (B)--(A);
  \node[below left] at (B) {$B$};
  \node[above right] at (A) {$A$};
  \coordinate (C) at ($(B)!0.5!(A)$);
  \fill (C) circle (1.2pt) node[above left] {$C$};
  % forces
  \draw[force] (A)--++(-1.2,0) node[above] {$F$};
  \draw[force] (B)--++(0,1.1) node[right] {$N$};
  \draw[force] (C)--++(0,-1.1) node[right] {$mg$};
  % angle alpha
  \draw[thickline] (1.1,0) arc (0:60:1.1);
  \node at (1.35,0.45) {$\alpha$};
  % coordinate directions at A
  \draw[dashedline] (A)--++(1.1,0) node[right] {$x$};
  \draw[dashedline] (A)--++(0,-1.1) node[below] {$y$};
\end{tikzpicture}
\end{figbox}

From the symmetry of the system it is known that the acceleration of the hinge $A$ has no horizontal component, that is, the acceleration of $A$ is along the $y$ direction. Hence, from the formula for relative motion, we obtain
\begin{align}
a_{Ax}&=a_{Cx}-\frac{1}{2}l\ddot{\theta}\sin\theta=a_{Cx}-\frac{\sqrt{3}}{4}l\ddot{\theta}=0 \tag{4}\\
a_{Ay}&=a_{Cy}+\frac{1}{2}l\ddot{\theta}\cos\theta=a_{Cy}+\frac{1}{4}l\ddot{\theta} \tag{5}
\end{align}
From (4) we obtain
\begin{equation}
a_{Cx}=\frac{\sqrt{3}}{4}l\ddot{\theta}. \tag{6}
\end{equation}
The acceleration of the end $B$ of the rod is along the $x$ direction, that is, it has no $y$ component; hence
\[
a_{By}=a_{Cy}-\frac{1}{2}l\ddot{\theta}\cos\theta=a_{Cy}-\frac{1}{4}l\ddot{\theta}=0,
\]
and we obtain
\begin{equation}
a_{Cy}=\frac{1}{4}l\ddot{\theta}. \tag{7}
\end{equation}
Substituting (6) into (1), we obtain
\begin{equation}
F=\frac{\sqrt{3}}{4}ml\ddot{\theta}. \tag{8}
\end{equation}
Substituting (7) into (2), we obtain
\begin{equation}
N=mg-\frac{1}{4}ml\ddot{\theta}. \tag{9}
\end{equation}
Substituting (8), (9) and $\theta=60\degree$ into (3), we obtain
\[
\frac{1}{2}\left(mg-\frac{1}{4}ml\ddot{\theta}\right)-\frac{\sqrt{3}}{2}\times\frac{\sqrt{3}}{4}ml\ddot{\theta}=\frac{1}{6}ml\ddot{\theta}.
\]
That is,
\begin{equation}
l\ddot{\theta}=\frac{3}{4}g. \tag{10}
\end{equation}
From (5), (7) and (10) we obtain
\[
a_A=a_{Ay}=\frac{3}{8}g.
\]

(2) When the angle between the two rods is $\alpha'=90\degree$, the angle between rod $AB$ and the ground is $\theta'=45\degree$. At this moment the rod may be regarded as making pure rotation about the instantaneous axis $O$, as shown in Fig. 2. In the figure the distance between point $O$ and the centre of mass $C$ of the rod is $\frac{1}{2}l$.
\end{solution}

\begin{solution}{1.59}
\begin{note}
The first page of this solution (page 134 of the original) is missing from the available copy.
The translation therefore begins with equation (6), the first line preserved on the following
page.
\end{note}
Since the rod now begins to slide relative to the corner of the table, $f$ is the kinetic
friction force, so
\begin{equation}
f=\mu N. \tag{6}
\end{equation}
From (3) we obtain
\begin{equation}
\frac{\dd\omega}{\dd t}
=\sqrt{\frac{3g}{l}}\cdot\frac{\cos\theta}{2\sqrt{\sin\theta}}\cdot\omega
=\frac{3g}{2l}\cos\theta. \tag{7}
\end{equation}
Substituting (7) into (4), we obtain
\begin{equation}
N=\frac{3}{4}mg\cos\theta. \tag{8}
\end{equation}
Substituting (8) into (6), we obtain
\begin{equation}
f=\frac{3}{4}mg\mu\cos\theta. \tag{9}
\end{equation}
Substituting (3) and (9) into (5), we obtain
\begin{equation*}
\tan\theta=\frac{\mu}{2}.
\end{equation*}

\begin{commentary}
In solving this problem, attention should be paid to the fact that when the rod begins to slide
relative to the corner of the table, the resultant force on the centre of mass of the rod in
the direction along the rod is not zero, that is, it is not $f=mg\sin\theta$; it should be
equal to the centripetal force of the centre of mass in its circular motion. The equation of
motion of the centre of mass in this direction is (5). This is because, from the start of the
motion until the rod begins to slide off the corner of the table, the centre of mass always
moves in a circle.
\end{commentary}
\end{solution}

\begin{solution}{1.60}
(1) The water-filled cylinder rolls without slipping on the inclined plane, so besides gravity it is also acted on by a static friction force $f$; hence the equation of motion of its centre of mass is
\begin{equation}
(m+\mu)g\sin\theta-f=(m+\mu)a. \tag{1}
\end{equation}
The equation of rotation of the cylinder about its own axis is
\begin{equation}
fR=I\beta. \tag{2}
\end{equation}
In (2), $I$ is the moment of inertia of the water-filled cylinder about its own axis. Since, while the cylinder rolls, the water inside it does not rotate together with the cylinder but moves translationally, we have
\begin{equation}
I=mR^2. \tag{3}
\end{equation}
In (2), $\beta$ is the angular acceleration of the cylinder; since the cylinder rolls without slipping,
\begin{equation}
\beta=\frac{a}{R}. \tag{4}
\end{equation}
From (1), (2), (3) and (4) we obtain
\[
a=\frac{m+\mu}{2m+\mu}g\sin\theta.
\]

(2) (i) Suppose that when the cylinder rolls down the inclined plane of the wedge-shaped block, the acceleration of the block is $w$. Take the block as the reference frame. In this frame the forces acting on the cylinder are as shown in the figure, where the inertial force is $f_i=(m+\mu)w$; let the acceleration of the centre of mass of the cylinder relative to the block be $a'$, and let its angular acceleration of rotation be $\beta$. The equation of motion of its centre of mass is
\begin{equation}
(m+\mu)g\sin\theta+(m+\mu)w\cos\theta-f=(m+\mu)a'. \tag{5}
\end{equation}

\begin{figbox}{1.60}
\begin{tikzpicture}[scale=1.02,>=latex]
  % ground
  \draw[thickline] (-0.8,0)--(6.2,0);
  % wedge
  \fill[hatch] (0,0)--(5.2,0)--(5.2,2.6)--cycle;
  \draw[thickline] (0,0)--(5.2,2.6);
  \draw[thickline] (5.2,0)--(5.2,2.6);
  % angle theta
  \draw[thickline] (1.5,0) arc (0:26.57:1.0);
  \node at (1.62,0.28) {$\theta$};
  % height h
  \draw[dim] (-0.35,0)--(-0.35,2.6) node[midway,left] {$h$};
  \draw[dashedline] (0,2.6)--(-0.45,2.6);
  % cylinder
  \coordinate (O) at (2.29,1.93);
  \draw[thickline] (O) circle (0.7);
  \fill (O) circle (1.2pt);
  \coordinate (P) at (2.6,1.3);
  % forces on the cylinder
  \draw[force] (O)--++(0,-1.2) node[below right,fill=white,inner sep=1.5pt] {$(m+\mu)g$};
  \draw[force] (O)--++(-0.58,1.16) node[above left] {$N$};
  \draw[force] (P)--++(0.72,0.36) node[right,fill=white,inner sep=1.5pt] {$f$};
  \draw[force] (O)--++(-1.05,0) node[above] {$f_i$};
  % block acceleration
  \draw[force] (5.45,0.1)--++(0.8,0) node[right] {$w$};
\end{tikzpicture}
\end{figbox}

The equation of rotation of the cylinder about its centre-of-mass axis is
\begin{equation}
fR=I\beta. \tag{6}
\end{equation}
The relation between $a'$ and $\beta$ is the same as in (4), namely
\begin{equation}
a'=\beta R. \tag{7}
\end{equation}
In the direction perpendicular to the inclined plane we have
\begin{equation}
N-(m+\mu)g\cos\theta+(m+\mu)w\sin\theta=0, \tag{8}
\end{equation}
and the equation of motion of the block is
\begin{equation}
N\sin\theta-f\cos\theta=Mw. \tag{9}
\end{equation}
From (3), (6) and (7) we obtain $f=ma'$; substituting this into (5), we obtain
\begin{equation}
(m+\mu)g\sin\theta+(m+\mu)w\cos\theta=(2m+\mu)a'. \tag{10}
\end{equation}
Eliminating $N$ from (8) and (9), and using $f=ma'$, we obtain
\begin{equation}
(m+\mu)g\sin\theta\cos\theta-(m+\mu)w\sin^2\theta-ma'\cos\theta=Mw. \tag{11}
\end{equation}
Eliminating $a'$ from (10) and (11), we obtain the acceleration of the block
\begin{equation}
w=\frac{(m+\mu)^2g\sin\theta\cos\theta}{(M+m+\mu)(2m+\mu)-(m+\mu)^2\cos^2\theta}. \tag{12}
\end{equation}
Substituting (12) into (10), we obtain the acceleration $a'$ of the cylinder relative to the block:
\begin{equation}
a'=\frac{(M+m+\mu)(m+\mu)g\sin\theta}{m(m+\mu)+(m+\mu)^2\sin^2\theta+M(2m+\mu)}. \tag{13}
\end{equation}
(ii) Let $t$ be the time taken by the water-filled cylinder to roll from a height $h$ above the table top down to the table top; then
\[
\frac{h}{\sin\theta}=\frac{1}{2}a't^2,
\]
that is,
\begin{equation}
t=\sqrt{\frac{2h}{a'\sin\theta}}. \tag{14}
\end{equation}
Using (12), (13) and (14), the velocity $u$ of the wedge-shaped block at this moment is obtained as
\[
u=wt=\sqrt{\frac{2gh(m+\mu)^3\cos^2\theta}{(M+m+\mu)\left[(M+m+\mu)(2m+\mu)-(m+\mu)^2\cos^2\theta\right]}}.
\]
\end{solution}

\begin{solution}{1.61}
(1) In the rotating reference frame, the potential energy of rod $AB$ is the sum of the gravitational potential energy and the centrifugal potential energy caused by the inertial centrifugal force, in which the gravitational potential energy is $mga\cos\theta$ (taking point $A$ as the zero of gravitational potential energy), while the centrifugal potential energy can be found by the following calculation.

Take an element $\dd l$ on rod $AB$ at a distance $l$ from end $A$, as shown in Fig.~1.61(a). Let the distance of the element from the axis of rotation be $\xi$. If the axis of rotation is taken as the zero of centrifugal potential energy, then the centrifugal potential energy of this element is
\[
\dd U_{\text{cent}}=-\int_0^{b+l\sin\theta}\frac{m}{2a}\,\dd l\,\omega^2\xi\,\dd\xi
=-\frac{m}{4a}\omega^2(b+l\sin\theta)^2\,\dd l.
\]

\begin{figbox}{1.61(a)}
\begin{tikzpicture}[scale=0.82,>=latex]
  % rotation axis (the y axis)
  \draw[thickline,->] (0,-1.0)--(0,4.4) node[right] {$y$};
  \draw[thickline,->] (0.275,2.974) arc (-60:240:0.55);
  \node at (1.15,3.70) {$\omega$};
  % x axis
  \draw[thickline,->] (0,0)--(5.4,0) node[below right] {$x$};
  \node[below left] at (0,0) {$O$};
  % rod AB
  \coordinate (A) at (1.6,0);
  \coordinate (B) at (4.20,-2.19);
  \draw[thickline] (A)--(B);
  \node[above] at (A) {$A$};
  \node[right] at (B) {$B$};
  % element dl on the rod
  \coordinate (E) at (3.13,-1.29);
  \draw[thickline] ($(E)+0.12*(0.643,0.766)$)--($(E)-0.12*(0.643,0.766)$);
  % centrifugal force on the element
  \draw[force] (E)--++(1.15,0) node[right] {$\dd\vec{F}_c$};
  % weight of the rod
  \draw[force] (B)--++(0,-1.0) node[below] {$m\vec{g}$};
  % angle theta at A
  \draw[thickline] ($(A)+(0.95,0)$) arc (0:-40:0.95);
  \node at (2.62,-0.42) {$\theta$};
  % distance b
  \draw[dim] (0,-0.55)--(1.6,-0.55) node[midway,below] {$b$};
  % distance l along the rod
  \draw[dim] ($(A)+0.36*(-0.643,-0.766)$)--($(E)+0.36*(-0.643,-0.766)$)
     node[midway,above left] {$l$};
  \draw[dashedline] ($(E)+0.16*(-0.643,-0.766)$)--($(E)+0.36*(-0.643,-0.766)$);
\end{tikzpicture}
\end{figbox}

Therefore the centrifugal potential energy of rod $AB$ is
\begin{align}
U_{\text{cent}}&=\int\dd U_{\text{cent}}=\int_0^{2a}-\frac{m}{4a}\omega^2(b+l\sin\theta)^2\,\dd l\notag\\
&=-\frac{m\omega^2}{12a\sin\theta}\left[(b+2a\sin\theta)^3-b^3\right]\notag\\
&=-\frac{1}{2}m\omega^2\left[\frac{1}{3}a^2\sin^2\theta+(b+a\sin\theta)^2\right] \tag{1}
\end{align}

Therefore the potential energy of rod $AB$ is
\begin{equation}
U=U_{\text{cent}}+U_{\text{grav}}
=-\frac{1}{2}m\omega^2\left[\frac{1}{3}a^2\sin^2\theta+(b+a\sin\theta)^2\right]+mga\cos\theta \tag{2}
\end{equation}

When the rod is in equilibrium the condition $\dfrac{\dd U}{\dd\theta}=0$ is satisfied, so
\[
\frac{\dd U}{\dd\theta}=-m\omega^2a\cos\theta\left(b+\frac{4}{3}a\sin\theta\right)-mga\sin\theta=0,
\]
that is,
\begin{equation}
\tan\theta=-\frac{a}{g}\omega^2\left(\frac{b}{a}+\frac{4}{3}\sin\theta\right) \tag{3}
\end{equation}
Equation (3) is the relation satisfied by the possible values of $\theta$ when rod $AB$ is in equilibrium.

(2) Each of the terms on the left and right sides of (3) is set equal to a function, that is, we set $y_1=\tan\theta$ and $y_2=-\dfrac{a}{g}\omega^2\left(\dfrac{b}{a}+\dfrac{4}{3}\sin\theta\right)$; the curves of the two functions are shown in Fig.~1.61(b). From the figure it is seen that $\theta$ has one equilibrium position in each of quadrants II and IV, while whether there is an equilibrium position in quadrant III depends on the relation between $a$ and $b$: there may be two equilibrium positions, or there may be one equilibrium position, or there may be no equilibrium position. This is because the curve $y_1$ is fixed, whereas the shape of the curve $y_2$ is determined by the constants $-\dfrac{b\omega^2}{g}$ and $\dfrac{4a\omega^2}{3g}$, that is, the curve $y_2$ oscillates about $y=-\dfrac{b\omega^2}{g}$. When the value of $\omega$ changes, $y=-\dfrac{b\omega^2}{g}$ changes correspondingly, and its amplitude of oscillation $\dfrac{4a\omega^2}{3g}$ changes accordingly as well, so that the two curves $y_1$ and $y_2$ in the third quadrant either have two points of intersection, or are tangent at one point, or have no point of intersection; Fig.~1.61(b) shows the case of no point of intersection, that is, the case in which there is no equilibrium position.

\begin{figbox}{1.61(b)}
\begin{tikzpicture}[>=latex]
\begin{axis}[
  width=8.8cm, height=5.8cm,
  axis lines=middle,
  xmin=0, xmax=6.6, ymin=-3.0, ymax=2.6,
  xtick={1.5708,3.1416,4.7124,6.2832},
  xticklabels={$\frac{\pi}{2}$,$\pi$,$\frac{3\pi}{2}$,$2\pi$},
  ytick=\empty,
  xlabel={$\theta$},
  samples=200,
  clip=true,
]
\addplot[thickline,domain=0:1.45,restrict y to domain=-2.6:2.4] {tan(x)};
\addplot[thickline,domain=1.70:4.55,restrict y to domain=-2.6:2.4] {tan(x)};
\addplot[thickline,domain=4.80:6.28,restrict y to domain=-2.6:2.4] {tan(x)};
\addplot[thickline,domain=0:6.2832,samples=160] {-1.2-sin(x)};
\addplot[dashedline,domain=0:6.2832,samples=2] {-1.2};
\draw[dashedline] (axis cs:1.5708,-2.6)--(axis cs:1.5708,2.4);
\draw[dashedline] (axis cs:4.7124,-2.6)--(axis cs:4.7124,2.4);
\node[anchor=west] at (axis cs:0.75,1.95) {$y_1=\tan\theta$};
\node[anchor=north east] at (axis cs:6.5,-2.30)
  {$y_2=-\dfrac{a}{g}\omega^2\!\left(\dfrac{b}{a}+\dfrac{4}{3}\sin\theta\right)$};
\node[above left] at (axis cs:0,2.45) {$y$};
\node[above right] at (axis cs:-0.08,-1.2) {$-\dfrac{b\omega^2}{g}$};
\end{axis}
\end{tikzpicture}
\end{figbox}

The stability of an equilibrium position depends on whether the value of $\dfrac{\dd^2U}{\dd\theta^2}$ at that equilibrium position is positive or negative, and $\dfrac{\dd^2U}{\dd\theta^2}$ is
\begin{equation}
\frac{\dd^2U}{\dd\theta^2}=-\frac{4}{3}ma^2\omega^2\left(\cos^2\theta-\sin^2\theta\right)+mab\omega^2\sin\theta-mga\cos\theta \tag{4}
\end{equation}
Substituting (3) into the above expression gives
\begin{equation}
\frac{\dd^2U}{\dd\theta^2}=\frac{ma\cos^2\theta}{\sin\theta}\left(b\omega^2-g\tan^3\theta\right) \tag{5}
\end{equation}
Discussion:

(i) For the equilibrium point of $\theta$ in quadrant II, that is, in the interval $\dfrac{\pi}{2}<\theta<\pi$, we have $\sin\theta>0$ and $\tan\theta<0$; therefore from (5) we obtain
\[
\left.\frac{\dd^2U}{\dd\theta^2}\right|_{\text{II}}>0,
\]
that is, the equilibrium property of this equilibrium position is stable equilibrium.

(ii) For the equilibrium point of $\theta$ in quadrant IV, that is, in the interval $\dfrac{3\pi}{2}<\theta<2\pi$, we have $\sin\theta<0$ and $\tan\theta<0$; therefore from (5) we obtain
\[
\left.\frac{\dd^2U}{\dd\theta^2}\right|_{\text{IV}}<0,
\]
that is, the equilibrium property of this equilibrium position is unstable equilibrium.

(iii) For the equilibrium point of $\theta$ in quadrant III, that is, in the interval $\pi<\theta<\dfrac{3\pi}{2}$, we have $\sin\theta<0$ and $\tan\theta>0$, and (4) can also be expressed as
\begin{align}
\frac{\dd^2U}{\dd\theta^2}&=\frac{ma\cos^2\theta}{\sin\theta}\left[\frac{4}{3}a\omega^2\sin\theta\cdot\tan^2\theta+b\omega^2\left(\tan^2\theta+1\right)\right]\notag\\
&=\frac{ma\omega^2}{\sin\theta}\left(\frac{4}{3}a\sin^2\theta+b\right) \tag{6}
\end{align}
From (6) it is seen that when $\dfrac{4}{3}a\sin^2\theta+b<0$, that is, $b<-\dfrac{4}{3}a\sin^2\theta$, we have $\dfrac{\dd^2U}{\dd\theta^2}>0$, so the equilibrium property of this equilibrium position is stable equilibrium; and when $b>-\dfrac{4}{3}a\sin^2\theta$, we have $\dfrac{\dd^2U}{\dd\theta^2}<0$, so the equilibrium property of this equilibrium position is unstable equilibrium.

In this problem the centrifugal potential energy is the negative of the work done by the inertial centrifugal force during the displacement of the rod from the zero of potential energy (that is, the position of the axis of rotation) to the equilibrium position. In fact, this work is just the kinetic energy of this rod rotating about the axis of rotation in the inertial frame, that is,
\begin{equation}
U_{\text{cent}}=-\frac{1}{2}I\omega^2 \tag{7}
\end{equation}
where $I$ is the moment of inertia of the rod about the axis of rotation. $I$ is
\begin{equation}
I=\frac{1}{3}ma^2\sin^2\theta+m\left(b+a\sin\theta\right)^2 \tag{8}
\end{equation}
Substituting (8) into (7) gives
\[
U_{\text{cent}}=-\frac{1}{2}m\omega^2\left[\frac{1}{3}a^2\sin^2\theta+\left(b+a\sin\theta\right)^2\right].
\]
This expression is completely the same as (1).

\end{solution}

\begin{solution}{1.62}
(1) Before the ball strikes the ground, the velocity of its centre of mass is
\begin{equation}
v_0=\sqrt{2gh} \tag{1}
\end{equation}
Before the collision the angular velocity of the ball is $\omega_0$. Let the angular velocity after the collision be $\omega$ and the velocity of the centre of mass after the collision be $v$, as shown in Fig.~1.62(a). The two components of $v$ are $v_x$ and $v_y$. For the process of the collision between the ball and the ground there are the following two possible cases.

\textit{Case 1.}\quad During the time $\Delta t$ of the collision between the ball and the ground, there is always relative sliding between the points of contact of the ball and the ground, so the friction between the two is always sliding friction, that is, $f=\mu N$, where $N$ is the supporting force of the ground. According to the momentum theorem, in the $x$ and $y$ directions respectively,
\begin{align}
\int_0^{\Delta t}f\,\dd t&=\int_0^{\Delta t}N\mu\,\dd t=mv_x \tag{2}\\
\int_0^{\Delta t}N\,\dd t-mg\Delta t&=mv_y-m(-v_0)=m\left(v_y+v_0\right) \tag{3}
\end{align}

\begin{figbox}{1.62(a)}
\begin{tikzpicture}[scale=0.95,>=latex]
  % ground and axes
  \draw[thickline] (-1.2,0)--(6.4,0);
  \draw[axis] (0,0)--(0,4.6) node[above] {$y$};
  \draw[axis] (0,0)--(6.2,0) node[below right] {$x$};
  \node[below left] at (0,0) {$O$};
  % upper ball
  \draw[thickline] (0,3.0) circle (0.80);
  \draw[thickline] (0,3.0)--(0.62,3.50);
  \node[right] at (0.34,3.36) {$R$};
  \draw[thickline,->] (0.95,3.35) arc (60:-60:0.55);
  \node[right] at (1.10,3.10) {$\omega_0$};
  \draw[dashedline] (0,3.0)--(0,4.3);
  % lower ball at the ground
  \draw[thickline] (0,0.80) circle (0.80);
  \draw[thickline] (0.72,0.80) arc (0:-60:0.72);
  \node at (0.52,0.42) {$\omega$};
  % height h
  \draw[dim] (-1.0,0)--(-1.0,3.0) node[midway,left] {$h$};
  \draw[dashedline] (-1.0,3.0)--(0,3.0);
  % velocity after the collision
  \draw[force] (0,0)--(2.05,1.15) node[above right] {$\vec{v}$};
  \draw[thickline] (1.0,0) arc (0:29.3:1.0);
  \node at (1.22,0.24) {$\theta$};
  \draw[dashedline,domain=0:5.6,samples=50] plot (\x,{0.562*\x-0.0417*\x*\x});
\end{tikzpicture}
\end{figbox}

From the definition of the coefficient of restitution $e$,
\begin{equation}
e=\frac{v_y}{v_0} \tag{4}
\end{equation}
From (2), (3) and (4) we obtain
\begin{equation}
v_x=\mu(1+e)v_0 \tag{5}
\end{equation}
Hence the velocity $v$ of the centre of mass of the ball after the collision is
\begin{equation}
v=\sqrt{v_x^2+v_y^2}=v_0\sqrt{\mu^2(1+e)^2+e^2} \tag{6}
\end{equation}
The angle $\theta$ between the direction of $v$ and the $x$-axis is given by
\begin{equation}
\tan\theta=\frac{v_y}{v_x}=\frac{e}{\mu(1+e)} \tag{7}
\end{equation}
In the collision the rotational speed of the ball becomes smaller because of the action of the friction torque; from the angular momentum theorem,
\begin{equation}
\int_0^{\Delta t}-fR\,\dd t=-N\mu R\Delta t=I\left(\omega-\omega_0\right) \tag{8}
\end{equation}
From (2), (5) and (8) we obtain
\[
\omega=\omega_0-\frac{\mu mR(1+e)v_0}{I}
\]
The moment of inertia of the ball about the axis through its centre of mass is $I=\dfrac{2}{5}mR^2$; substituting this and (1) into the above expression gives
\begin{equation}
\omega=\omega_0-\frac{5\mu(1+e)}{2R}\sqrt{2gh} \tag{9}
\end{equation}

This case is the one in which, while the ball is in contact with the ground, there is always relative sliding between the two; the requirement for this is
\[
\omega R\geq v_x .
\]
Substituting (5) and (9) into the above expression gives
\[
\omega_0\geq\frac{7\mu(1+e)}{2R}\sqrt{2gh}=\omega_{0c},
\]
and this is the initial condition corresponding to this case.

\textit{Case 2.}\quad During the time $\Delta t$ of the collision between the ball and the ground, there is relative sliding between the points of contact of the two only during the first part of the time $\Delta t_1$, while during the latter part of the time there is no longer relative sliding, that is, the ball has already reached pure rolling before the collision ends, and the friction will become zero. Similarly, applying the momentum theorem in the $x$ direction,
\begin{equation}
\int_0^{\Delta t_1}f\,\dd t=\int_0^{\Delta t_1}\mu N\,\dd t=mv_x \tag{10}
\end{equation}
From the angular momentum theorem,
\begin{equation}
\int_0^{\Delta t_1}-fR\,\dd t=\int_0^{\Delta t_1}-\mu NR\,\dd t=I\left(\omega-\omega_0\right) \tag{11}
\end{equation}
Since at the end of this stage the ball has already reached pure rolling, we have
\begin{equation}
\omega=\frac{v_x}{R} \tag{12}
\end{equation}
From (10), (11) and (12) we obtain
\[
v_x=\frac{IR}{I+mR^2}\omega_0=\frac{2}{7}\omega_0R,
\]
while $v_y$ is still given by (4); hence the velocity of the centre of mass after the collision is
\begin{equation}
v=\sqrt{v_x^2+v_y^2}=\sqrt{\left(\frac{2}{7}\omega_0R\right)^2+\left(ev_0\right)^2} \tag{13}
\end{equation}
The angle $\theta$ between the direction of $v$ and the $x$-axis satisfies
\begin{equation}
\tan\theta=\frac{v_y}{v_x}=\frac{7ev_0}{2\omega_0R} \tag{14}
\end{equation}
The angular velocity of the ball after the collision is
\[
\omega=\frac{v_x}{R}=\frac{2}{7}\omega_0 .
\]
The initial condition corresponding to this case is
\[
\omega_0<\frac{7\mu(1+e)}{2R}\sqrt{2gh}=\omega_{0c}.
\]

(2) After the collision between the ball and the ground the centre of the ball performs projectile motion. From the range formula for projectile motion, the horizontal distance $s$ between the first and the second point of landing of the ball is
\begin{equation}
s=\frac{2v^2\sin\theta\cos\theta}{g} \tag{15}
\end{equation}
For case 1, from (7) we obtain
\[
\sin\theta=\frac{e}{\sqrt{e^2+\mu^2(1+e)^2}},\qquad
\cos\theta=\frac{\mu(1+e)}{\sqrt{e^2+\mu^2(1+e)^2}}.
\]
Substituting (6) and the above two expressions into (15) gives
\[
\begin{aligned}
s&=\frac{2}{g}\left[e^2+\mu^2(1+e)^2\right]v_0^2\left[\frac{\mu e(1+e)}{e^2+\mu^2(1+e)^2}\right]\\
&=4\mu(1+e)h .
\end{aligned}
\]
For case 2, from (14) we obtain
\[
\sin\theta=\frac{ev_0}{\sqrt{\left(\dfrac{2}{7}\omega_0R\right)^2+\left(ev_0\right)^2}},\qquad
\cos\theta=\frac{\dfrac{2}{7}\omega_0R}{\sqrt{\left(\dfrac{2}{7}\omega_0R\right)^2+\left(ev_0\right)^2}}.
\]
Substituting (13) and the above two expressions into (15) gives
\[
\begin{aligned}
s&=\frac{2}{g}\left[\left(\frac{2}{7}\omega_0R\right)^2+\left(ev_0\right)^2\right]
\left[\frac{\dfrac{2}{7}\omega_0Rev_0}{\left(\dfrac{2}{7}\omega_0R\right)^2+\left(ev_0\right)^2}\right]\\
&=\frac{4}{7}e\omega_0R\sqrt{\frac{2h}{g}} .
\end{aligned}
\]

(3) Equations (7) and (14) respectively give the relation of $\tan\theta$ with the initial angular velocity $\omega_0$ in case 1 and in case 2. For case 1, that is, when $\omega_0\geq\omega_{0c}$, $\tan\theta$ is independent of $\omega_0$ and is the constant $\dfrac{e}{\mu(1+e)}$; for case 2, that is, when $\omega_0<\omega_{0c}$, $\tan\theta$ is inversely proportional to $\omega_0$. Hence the curve of $\tan\theta$ against $\omega_0$ is as shown in Fig.~1.62(b).

\begin{figbox}{1.62(b)}
\begin{tikzpicture}[scale=1.0,>=latex]
  \draw[axis] (0,0)--(6.8,0) node[below right] {$\omega_0$};
  \draw[axis] (0,0)--(0,3.6) node[above] {$\tan\theta$};
  \node[below left] at (0,0) {$O$};
  \draw[thickline,domain=0.62:2.60,samples=80] plot (\x,{1.10/\x});
  \draw[thickline] (2.60,0.423)--(6.35,0.423);
  \draw[dashedline] (0,0.423)--(2.60,0.423);
  \draw[dashedline] (2.60,0)--(2.60,0.423);
  \node[left] at (0,0.423) {$\dfrac{e}{\mu(1+e)}$};
  \node[below] at (2.60,0) {$\omega_{0c}$};
\end{tikzpicture}
\end{figbox}

\end{solution}

\begin{solution}{1.63}
(1) When the top spins rapidly about its axis of spin, the top has a very large spin angular momentum $L$ along the direction of the axis of spin, as shown in Fig.~1.63(a). The torque $\tau_g$ of the weight acting on the top about the point of support $O$ gives the angular momentum $L$, within the time $\dd t$, an increment $\dd L$ in the direction of $\tau_g$. This means that the axis of spin of the top will precess about the vertical axis through $O$. In addition, the axis of spin of the top in Fig.~1.63(a) inclines to the left, and the left-hand part of the lower end of the axis is in contact with the ground; because of the precession, the direction of the friction $f$ acting on the end of the axis is perpendicular to the plane of the figure and inwards, so the direction of the friction torque $\tau_f$ about the centre of mass $C$ is as shown in Fig.~1.63(b). Therefore, within the time $\dd t$ the angular momentum has an increment $\dd L$ along the direction of $\tau_f$. This causes the axis of spin of the top to perform another kind of precession, that is, while precessing about the vertical axis it also moves gradually towards the vertical axis, and finally becomes vertical.

\begin{figbox}{1.63(a)}
\begin{tikzpicture}[scale=1.0,>=latex]
  % ground
  \fill[hatch] (-1.7,-0.18) rectangle (3.4,0);
  \draw[thickline] (-1.7,0)--(3.4,0);
  \node[below left] at (0,0) {$O$};
  % vertical axis
  \draw[dashedline] (0,0)--(0,4.9);
  % axis of spin
  \draw[thickline] (0,0)--(1.37,3.76);
  \draw[thickline] (0,1.35) arc (90:70:1.35);
  \node at (0.30,1.52) {$\theta$};
  % rotor
  \coordinate (C) at (0.684,1.879);
  \draw[thickline] (C) ellipse [x radius=0.95, y radius=0.22, rotate=-20];
  \node[below left] at (C) {$C$};
  \node[right] at (1.05,1.55) {$m$};
  \draw[thickline] (C)--++(-20:0.95);
  \node[right] at (1.20,1.72) {$R$};
  % spin arrow
  \draw[thickline,->] (0.98,2.55) arc (60:-120:0.30 and 0.55);
  \node[right] at (1.20,2.62) {$\omega$};
  % angular momentum
  \draw[force] (1.37,3.76)--++(70:0.85) node[right] {$\vec{L}$};
  % increment of angular momentum
  \draw[force] (1.37,3.76)--++(180:1.05) node[left] {$\dd\vec{L}$};
  % precession circle
  \draw[dashedline] (0,3.76) ellipse [x radius=1.40, y radius=0.40];
  \draw[thickline,->] (0.62,3.86) arc (35:80:0.72 and 0.20);
  \node at (0.72,4.06) {$\dd\phi$};
\end{tikzpicture}
\end{figbox}

\begin{figbox}{1.63(b)}
\begin{tikzpicture}[scale=1.0,>=latex]
  % vertical
  \draw[dashedline] (0,0)--(0,4.1);
  % axis of spin
  \draw[thickline] (0,0)--(1.20,3.40);
  \node[below right] at (0,0) {$O$};
  % d(theta)
  \draw[thickline] (0,1.75) arc (90:70.6:1.75);
  \node at (0.36,2.02) {$\dd\theta$};
  % dL perpendicular to the axis
  \draw[force] (1.20,3.40)--++(160:1.00) node[left] {$\dd\vec{L}$};
  \node[above right] at (1.70,3.28) {$\vec{\tau}_f$};
  % friction
  \draw[force] (0,0)--++(180:0.95) node[left] {$f$};
\end{tikzpicture}
\end{figbox}

The torque $\tau_g$ of the weight acting on the top about point $O$ is
\begin{equation}
\tau_g=mg\frac{l}{2}\sin\theta \tag{1}
\end{equation}
The increment $\dd L$ of the angular momentum within the time $\dd t$ is
\begin{equation}
\dd L=L\sin\theta\,\dd\phi \tag{2}
\end{equation}
From the angular momentum theorem,
\begin{equation}
\tau_g\,\dd t=\dd L \tag{3}
\end{equation}
From (1), (2) and (3) we obtain the angular velocity $\Omega$ of the precession of the top about the vertical axis:
\begin{equation}
\Omega=\frac{\dd\phi}{\dd t}=\frac{mgl}{2L}=\frac{mgl}{2I\omega} \tag{4}
\end{equation}
where $I$ is the moment of inertia of the top about its axis of spin. From (4) it is seen that $\Omega$ is inversely proportional to $\omega$; as $\omega$ becomes smaller and smaller, the angular velocity of precession $\Omega$ becomes larger and larger.

(2) The friction is $f=\mu mg$, and the torque $\tau_f$ of the friction about the centre of mass $C$ is
\begin{equation}
\tau_f=\mu mg\frac{l}{2} \tag{5}
\end{equation}
The increment $\dd L$ of the angular momentum within the time $\dd t$ is
\begin{equation}
\dd L=-L\,\dd\theta \tag{6}
\end{equation}
From the angular momentum theorem,
\begin{equation}
\tau_f\,\dd t=\dd L \tag{7}
\end{equation}
From (5), (6) and (7) we obtain the angular velocity $\Omega'$ of the precession of the axis of spin of the top towards the vertical axis:
\begin{equation}
\Omega'=\frac{\dd\theta}{\dd t}=-\frac{\mu mgl}{2I\omega} \tag{8}
\end{equation}
Hence the time $t$ that elapses from the moment when the axis of spin makes an angle $\theta$ with the vertical line until it turns to the vertical position is
\[
t=-\frac{\theta}{\Omega'}=\frac{2I\omega\theta}{\mu mgl}.
\]

\end{solution}

\begin{solution}{1.64}
(1) In the reference frame $S'$ of the spaceship, the time that elapses from the spaceship's flying past the Earth to its flying past the space station is $\Delta t_1'=30\times60=1800\unitv{s}$. Since this is the proper time, the time $\Delta t_1$ measured in the space-station--Earth reference frame $S$ is
\[
\Delta t_1=\frac{\Delta t_1'}{\sqrt{1-v^2/c^2}}=\frac{1800}{\sqrt{1-(0.8)^2}}=3000\ (\mathrm{s}),
\]
that is, 50 minutes; hence the reading of the clock on the space station is 12:50.

(2) In frame $S$ the distance between the space station and the Earth is
\[
l=0.8c\times50\times60=2400c .
\]
The time $\Delta t_2$ that elapses from the moment the spaceship sends out the report signal until the Earth receives this signal is
\[
\Delta t_2=\frac{l}{c}=2400\ (\mathrm{s}),
\]
that is, 40 minutes; hence the clock on the Earth at this moment indicates 1:30 p.m.

(3) In frame $S'$ the distance between the space station and the Earth is
\[
l'=0.8c\times30\times60=1440c .
\]
Let the time that elapses in frame $S'$ from the moment the spaceship sends out the report signal until it receives the signal sent back by the Earth be $\Delta t'$; then
\[
(c-v)\frac{\Delta t'}{2}=l'
\]
we obtain
\[
\Delta t'=\frac{2l'}{c-v}=\frac{2l'}{0.2c}=\frac{10l'}{c}=14400\ (\mathrm{s}),
\]
that is, 4 hours; hence when the signal sent back by the Earth is received, the clock on the spaceship indicates 4:30 p.m.

This part of the problem can also be solved by the following method.

In frame $S$, when the Earth receives the report signal sent out by the spaceship, the distance between the spaceship and the Earth is $L=l+v\Delta t_2$. Let the time that elapses from the moment the Earth sends out the reply signal until the spaceship receives this signal be $\Delta t_3$; then
\[
(c-v)\Delta t_3=L=l+v\Delta t_2 ,
\]
so
\[
\Delta t_3=\frac{l+v\Delta t_2}{c-v}=21600\ (\mathrm{s}).
\]
Hence the total time $\Delta t$ that elapses from the moment the spaceship sends out the report signal until it receives the reply signal sent back by the Earth is
\[
\Delta t=\Delta t_2+\Delta t_3=24000\ (\mathrm{s}).
\]
Since in frame $S'$ the time $\Delta t'$ is the proper time and $\Delta t$ is the dilated time, $\Delta t'$ is
\[
\Delta t'=\Delta t\sqrt{1-v^2/c^2}=24000\times0.6=14400\ (\mathrm{s}).
\]
\end{solution}

\begin{solution}{1.65}
Take the ground as the reference frame $S$ and the moving thick glass as the reference frame $S'$, the frame $S'$ moving with velocity $v$ in the $x$ direction relative to $S$. In frame $S'$ the speed of propagation of light in the glass is $u_x'=\dfrac{c}{n}$; from the velocity transformation formula, the speed of this propagation $u_x$ in frame $S$ is
\begin{equation}
u_x=\frac{u_x'+v}{1+\dfrac{u_x'v}{c^2}}=\frac{\dfrac{c}{n}+v}{1+\dfrac{v}{cn}}=\frac{c^2+cnv}{cn+v} \tag{1}
\end{equation}

\begin{figbox}{1.65}
\begin{tikzpicture}[scale=1.0,>=latex]
  \draw[thickline] (2.55,0.15) rectangle (4.55,2.85);
  \node at (3.55,1.50) {$n$};
  \draw[dashedline] (1.30,1.50)--(5.95,1.50);
  \node[left] at (1.30,1.50) {$A$};
  \node[right] at (5.95,1.50) {$B$};
  \draw[force] (4.55,3.10)--++(0.85,0) node[right] {$v$};
  \draw[dashedline] (1.30,1.50)--(1.30,-0.05);
  \draw[dashedline] (5.95,1.50)--(5.95,-0.05);
  \draw[dim] (1.30,-0.05)--(5.95,-0.05) node[midway,below] {$L$};
\end{tikzpicture}
\end{figbox}

In frame $S'$ the thickness of the glass is $D_0$; then the thickness $D$ in frame $S$ is
\begin{equation}
D=D_0\sqrt{1-\beta^2} \tag{2}
\end{equation}
As seen by an observer in frame $S$, let the time required for the light to propagate through the thick glass be $\Delta t_1$; during this time the glass has moved a distance $v\Delta t_1$ in the $x$ direction, hence
\begin{equation}
D+v\Delta t_1=u_x\Delta t_1 \tag{3}
\end{equation}
Substituting (1) and (2) into (3) gives
\begin{equation}
\Delta t_1=\frac{D_0\sqrt{1-\beta^2}}{u_x-v}=\frac{D_0(cn+v)}{c^2\sqrt{1-\beta^2}} \tag{4}
\end{equation}
Let the time that the light takes to propagate through the space outside the glass be $\Delta t_2$; then
\begin{equation}
c\Delta t_2=L-u_x\Delta t_1 \tag{5}
\end{equation}
Substituting (1) and (4) into (5) gives
\begin{equation}
\Delta t_2=\frac{L}{c}-\frac{D_0(c+nv)}{c^2\sqrt{1-\beta^2}} \tag{6}
\end{equation}
Hence the total time $\Delta t$ required for the light to propagate from point $A$ to point $B$ is
\begin{equation}
\Delta t=\Delta t_1+\Delta t_2 \tag{7}
\end{equation}
Substituting (4) and (6) into (7) gives
\[
\begin{aligned}
\Delta t&=\frac{D_0(cn+v)}{c^2\sqrt{1-\beta^2}}+\frac{L}{c}
-\frac{D_0(c+nv)}{c^2\sqrt{1-\beta^2}}\\
&=\frac{L}{c}+\frac{D(n-1)(c-v)}{c^2\sqrt{1-\beta^2}} .
\end{aligned}
\]

Equation (4) in this problem can also be obtained by the following method.

In frame $S'$ the speed of propagation of light in the glass is $u_x'=\dfrac{c}{n}$. Hence the time $\Delta t_1'$ taken by the light to propagate through the thick glass is
\[
\Delta t_1'=\frac{D_0}{u_x'}=\frac{nD_0}{c}.
\]
From the Lorentz transformation, the time $\Delta t_1$ measured in frame $S$ for the light to propagate through the glass can be found:
\[
\begin{aligned}
\Delta t_1&=\frac{\Delta t_1'+\dfrac{v}{c^2}\Delta x'}{\sqrt{1-\beta^2}}
=\frac{\dfrac{nD_0}{c}+\dfrac{v}{c^2}D_0}{\sqrt{1-\beta^2}}\\
&=\frac{D_0(cn+v)}{c^2\sqrt{1-\beta^2}} .
\end{aligned}
\]
This is (4).
\end{solution}

\begin{solution}{1.66}
(1) As shown in Fig.~1.66(a), in the inertial frame $S$ spaceship $B$ flies along the $y$ direction with speed $\dfrac{c}{2}$; therefore, in order that the packet thrown from spaceship $A$ can be received by spaceship $B$, the $y$ component $u_y$ of its velocity $u$ must be the same as the flying speed of spaceship $B$, that is,
\[
u_y=\frac{c}{2}.
\]

\begin{figbox}{1.66(a)}
\begin{tikzpicture}[scale=1.0,>=latex]
  % coordinate axes
  \draw[axis] (0,0)--(1.9,0) node[below right] {$x$};
  \draw[axis] (0,0)--(0,1.9) node[above left] {$y$};
  % initial and final positions of the two spaceships
  \draw[dashedline] (0.70,3.55)--(3.30,3.55)--(3.30,0.70)--(0.70,0.70)--cycle;
  % spaceship A
  \draw[thickline,fill=black] (0.62,3.62)--(0.78,3.62)--(0.70,3.45)--cycle;
  \node[below left] at (0.62,3.50) {$A$};
  % spaceship B
  \draw[thickline,fill=black] (3.22,0.62)--(3.38,0.62)--(3.30,0.79)--cycle;
  \node[above right] at (3.38,0.72) {$B$};
  % velocity of B
  \draw[force] (3.55,1.10)--++(0,-0.75) node[below] {$\dfrac{c}{2}$};
  % displacement
  \draw[dim] (0.70,3.80)--(3.30,3.80) node[midway,above] {$d$};
  \draw[force] (0.48,3.55)--++(0,-2.60) node[midway,left] {$\dfrac{c}{2}$};
\end{tikzpicture}
\end{figbox}

In frame $S$ the speed with which the packet is thrown is $u=\dfrac{3}{4}c$, so the $x$ component $u_x$ of the velocity is
\[
u_x=\sqrt{u^2-u_y^2}=\frac{\sqrt{5}}{4}c .
\]
The coordinate system $S'$, which is at rest relative to spaceship $A$, moves relative to frame $S$ with speed $\dfrac{c}{2}$ along the negative $y$ direction, that is, $v=-\dfrac{c}{2}$; from the velocity transformation formula, the two velocity components $u_y'$ and $u_x'$ of the packet at the moment it is thrown, as measured in frame $S'$, are respectively
\[
u_y'=\frac{u_y-v}{1-\dfrac{u_xv}{c^2}}=\frac{\dfrac{c}{2}+\dfrac{c}{2}}{1+\dfrac{1}{4}}=\frac{4}{5}c,
\]
\[
u_x'=\frac{u_x\sqrt{1-\beta^2}}{1-\dfrac{u_yv}{c^2}}
=\frac{\dfrac{\sqrt{5}}{4}c\sqrt{1-\dfrac{1}{4}}}{1+\dfrac{1}{4}}=\frac{\sqrt{15}}{10}c .
\]

\begin{figbox}{1.66(b)}
\begin{tikzpicture}[scale=1.0,>=latex]
  \draw[axis] (0,0)--(3.5,0) node[below right] {$x'$};
  \draw[axis] (0,0)--(0,2.9) node[above left] {$y'$};
  \draw[force] (0,0)--(2.70,2.05) node[above right] {$\vec{u}'$};
  \draw[thickline] (1.0,0) arc (0:37.2:1.0);
  \node at (1.22,0.26) {$\alpha$};
\end{tikzpicture}
\end{figbox}

As shown in Fig.~1.66(b), therefore, as seen by an observer in frame $S'$ (that is, the reference frame of spaceship $A$), the angle $\alpha$ between the direction of the velocity of the packet at the moment it is thrown and the $x'$ axis is given by
\[
\tan\alpha=\frac{u_y'}{u_x'}=\frac{8}{\sqrt{15}},
\]
that is,
\[
\alpha=\arctan\frac{8}{\sqrt{15}}\approx64\degree .
\]

(2) In frame $S'$ the speed with which the packet is thrown is $u'$:
\[
u'=\sqrt{u_x'^2+u_y'^2}=\frac{\sqrt{79}}{10}c\approx0.89c .
\]
\end{solution}

\begin{solution}{1.67}
(1) Let the proper length of the two rulers $AB$ and $A'B'$ be $l_0$. The observer in the $S'$ frame measures the length of ruler $AB$ to be $\sqrt{1-\beta^2}\,l_0$; in his view, $B'$ meets $B$ first, and $A'$ meets $A$ later. The time difference between these two events is
\[
\Delta t'=\frac{l_0-\sqrt{1-\beta^2}\,l_0}{v}.
\]
When $A$ meets $A'$, the light signal emitted by $A'$ takes a time $l_0/c$ to reach $B'$ in the $S'$ frame. Hence when $B'$ receives this light signal the reading $t_{B'}$ of the clock $B'$ is
\begin{equation}
t_{B'}=\Delta t'+\frac{l_0}{c}=\frac{l_0}{c}+\frac{l_0-\sqrt{1-\beta^2}\,l_0}{v}. \tag{1}
\end{equation}
From the statement of the problem,
\begin{equation}
t_{B'}=1. \tag{2}
\end{equation}
From (1) and (2) we obtain
\begin{equation}
\frac{l_0}{c}=\frac{\beta}{\beta+\left(1-\sqrt{1-\beta^2}\right)}. \tag{3}
\end{equation}

\begin{figbox}{1.67}
\begin{tikzpicture}[scale=0.9,>=latex]
  % --- S' frame: rod A'B' ---
  \node[left] at (-0.12,1.5) {$A'$};
  \fill[black!25] (0,1.32) rectangle (5.0,1.68);
  \draw[thickline] (0,1.32)--(5.0,1.32);
  \draw[thickline] (0,1.68)--(5.0,1.68);
  \draw[thickline] (0,1.32)--(0,1.68);
  \draw[thickline] (5.0,1.32)--(5.0,1.68);
  \node[right] at (5.12,1.5) {$B'$};
  \draw[force] (2.1,2.18)--(3.4,2.18) node[midway,above] {$\vec{v}$};
  \draw[axis] (5.7,1.5)--(7.1,1.5) node[right] {$x'$};
  \node[left] at (-0.85,1.5) {$S'$};
  % --- S frame: rod AB ---
  \node[left] at (-0.12,0.3) {$A$};
  \fill[black!25] (0,0.12) rectangle (5.0,0.48);
  \draw[thickline] (0,0.12)--(5.0,0.12);
  \draw[thickline] (0,0.48)--(5.0,0.48);
  \draw[thickline] (0,0.12)--(0,0.48);
  \draw[thickline] (5.0,0.12)--(5.0,0.48);
  \node[right] at (5.12,0.3) {$B$};
  \draw[axis] (5.7,0.3)--(7.1,0.3) node[right] {$x$};
  \node[left] at (-0.85,0.3) {$S$};
\end{tikzpicture}
\end{figbox}

The observer in the $S$ frame measures the length of ruler $A'B'$ to be $\sqrt{1-\beta^2}\,l_0$; in his view, $A$ meets $A'$ first, and $B$ meets $B'$ later. The time difference between these two events is
\[
\Delta t=\frac{l_0-\sqrt{1-\beta^2}\,l_0}{v}.
\]
When $A$ meets $A'$, the light signal emitted by $A'$ takes a time $l_0/c$ to reach $B$ in the $S$ frame, hence when $B$ receives this light signal the reading $t_B$ of the clock $B$ is
\begin{align}
t_B&=\frac{l_0}{c}-\Delta t=\frac{l_0}{c}-\frac{l_0-\sqrt{1-\beta^2}\,l_0}{v}\notag\\
&=\frac{l_0}{c}\left(1-\frac{1-\sqrt{1-\beta^2}}{\beta}\right). \tag{4}
\end{align}
Substituting (3) into (4) gives
\begin{align*}
t_B&=\frac{\beta}{\beta+\left(1-\sqrt{1-\beta^2}\right)}\cdot\frac{\beta-\left(1-\sqrt{1-\beta^2}\right)}{\beta}\\
&=\sqrt{\frac{1-\beta}{1+\beta}}\qquad(\text{time units}).
\end{align*}

(2) (a) In the view of the observer in the $S'$ frame, after $B'$ has received the signal emitted by $A'$ it immediately emits a reply signal; this signal takes a time $l_0/c$ to reach $A'$, hence when $A'$ receives this reply signal the reading $t_{A'}$ of the clock $A'$ is
\begin{equation}
t_{A'}=\frac{2l_0}{c}. \tag{5}
\end{equation}
Substituting (3) into (5) gives
\[
t_{A'}=\frac{2\beta}{\beta+\left(1-\sqrt{1-\beta^2}\right)}
=1+\sqrt{\frac{1-\beta}{1+\beta}}\qquad(\text{time units}).
\]

(b) Let the reading of clock $A$ in the $S$ frame be $t_{A1}$ at the moment when $A'$ receives the reply signal. Since $t_{A'}$ is a proper time, $t_{A1}$ is
\[
t_{A1}=\frac{t_{A'}}{\sqrt{1-\beta^2}}=\frac{2l_0}{c\sqrt{1-\beta^2}}.
\]
At this moment the distance between $A'$ and $A$ is
\[
\Delta l_{AA'}=vt_{A1}=\frac{2vl_0}{c\sqrt{1-\beta^2}}.
\]
The light signal then takes a further time $\dfrac{\Delta l_{AA'}}{c}$ to reach $A$, hence when $A$ receives this reply signal the reading $t_A$ of clock $A$ is
\begin{equation}
t_A=t_{A1}+\frac{\Delta l_{AA'}}{c}=\frac{2(1+\beta)l_0}{c\sqrt{1-\beta^2}}. \tag{6}
\end{equation}
Substituting (3) into (6) gives
\[
t_A=1+\sqrt{\frac{1+\beta}{1-\beta}}\qquad(\text{time units}).
\]
\end{solution}

\begin{solution}{1.68}
(1) Take the Galaxy as the $S$ frame and the traveller as the $S'$ frame; the $S'$ frame moves with velocity $v$ along the $x$ direction relative to the $S$ frame, as shown in Fig.~1.68(a). In the $S$ frame the velocity of the light coming from the edge of the Galaxy makes an angle $\varphi=45\degree$ with the $x$-axis. Hence the components of the light velocity along the $x$ and $y$ directions are respectively $u_x=c\cos\varphi$ and $u_y=c\sin\varphi$; from the velocity transformation formulas the two components of the light velocity in the $S'$ frame are respectively
\[
u_x'=\frac{u_x-v}{1-\dfrac{u_xv}{c^2}}=\frac{c\cos\varphi-v}{1-\dfrac{v}{c}\cos\varphi},
\]
\[
u_y'=\frac{u_y\sqrt{1-\beta^2}}{1-\dfrac{u_xv}{c^2}}=\frac{c\sin\varphi\sqrt{1-\beta^2}}{1-\dfrac{v}{c}\cos\varphi}.
\]

\begin{figbox}{1.68(a)}
\begin{tikzpicture}[scale=0.85,>=latex]
  % Galaxy: a curved band
  \draw[thickline] (0.55,0.15) .. controls (-0.55,1.6) and (-0.55,3.6) .. (0.75,5.1);
  \draw[thickline] (1.05,0.05) .. controls (-0.05,1.55) and (-0.05,3.65) .. (1.25,5.0);
  \draw[thickline] (0.55,0.15)--(1.05,0.05);
  \draw[thickline] (0.75,5.1)--(1.25,5.0);
  \node[rotate=90] at (-0.85,2.6) {Galaxy};
  % line of sight
  \draw[dashedline] (0.9,0.5)--(5.3,3.6);
  \draw[force] (5.3,3.6)--(6.6,3.6);
  % angle phi
  \draw[thickline] (0.9,0.5)--(1.9,0.5);
  \draw[thickline] (1.55,0.5) arc (0:35.2:0.65);
  \node at (1.28,0.88) {$\varphi$};
  % axes
  \draw[axis] (4.1,0.6)--(5.5,0.6) node[right] {$x$};
  \draw[axis] (4.1,0.6)--(4.1,2.0) node[above] {$y$};
  \node[below left] at (4.1,0.6) {$O$};
\end{tikzpicture}
\end{figbox}

Let the angle between this light velocity direction and the $x'$-axis in the $S'$ frame be $\varphi'$; then
\[
\tan\varphi'=\frac{u_y'}{u_x'}=\frac{\sqrt{1-\beta^2}}{1-\frac{\beta}{\cos\varphi}}\cdot\tan\varphi.
\]
Substituting $\beta=\dfrac{v}{c}=0.99$ and $\varphi=45\degree$ into the above expression gives
\[
\tan\varphi'=\frac{\sqrt{1-(0.99)^2}}{1-0.99\times\sqrt{2}}=-0.3529,
\]
that is, $\varphi'=160.58\degree$. Fig.~1.68(b) shows the shape of the Galaxy as it appears to the traveller: just as described by the hero of the novel, it looks like a ``goldfish bowl'', and he is flying towards the mouth of the bowl.

(2) According to the formula for the Doppler effect of light, the frequency $\nu'$ of the light from the edge of the Galaxy received by the traveller is
\[
\nu'=\frac{\nu\sqrt{1-\beta^2}}{1-\frac{v}{c}\cos\theta}.
\]

\begin{figbox}{1.68(b)}
\begin{tikzpicture}[scale=0.85,>=latex]
  % bowl (distorted Galaxy)
  \draw[thickline] (2.356,0.655) arc (25:335:2.6 and 1.55);
  \draw[thickline] (1.949,0.486) arc (25:335:2.15 and 1.15);
  \draw[thickline] (2.356,0.655)--(1.949,0.486);
  \draw[thickline] (2.356,-0.655)--(1.949,-0.486);
  % traveller
  \fill (0,0) circle (1.4pt);
  % direction to the edge of the Galaxy
  \draw[force] (0,0)--(160.58:1.25);
  % x' axis line
  \draw[dashedline] (-1.15,0)--(2.45,0);
  % angles
  \draw[thickline] (0:0.9) arc (0:160.58:0.9);
  \node at (120:1.05) {$\varphi'$};
  \draw[thickline] (160.58:0.42) arc (160.58:180:0.42);
  \node[above left] at (172:0.62) {$\theta$};
\end{tikzpicture}
\end{figbox}

In the above expression $\nu$ is the frequency of the light actually emitted at the edge of the Galaxy. $\theta$ and $\varphi'$ are supplementary to each other, as shown in Fig.~1.68(b), that is, $\cos\theta=0.9431$. Substituting this and $\beta=0.99$ into the above expression gives
\[
\frac{\nu'}{\nu}=\frac{\sqrt{1-(0.99)^2}}{1-0.99\times0.9431}=2.127.
\]
From this ratio it is seen that the light from the edge of the Galaxy observed by the traveller is blue-shifted.

(3) Take $\varphi'=0$, that is, $\theta=\pi$, $\cos\theta=-1$. Then
\[
\frac{\nu'}{\nu}=\frac{\sqrt{1-(0.99)^2}}{1+0.99}=0.0709.
\]
This is the light from the centre of the Galaxy, which appears red-shifted to the traveller.

If $\nu'=\nu$ is taken, then
\[
\cos\theta=\frac{1-\sqrt{1-\beta^2}}{\beta}=\frac{1-\sqrt{1-(0.99)^2}}{0.99}=0.8676,
\]
that is, $\theta=29.8\degree$. At $\varphi'=-150.18\degree$ there is a critical angle: if $\theta$ is greater than this angle the light is red-shifted, whereas if $\theta$ is smaller than this angle the light is blue-shifted. This means that the closer to the centre of the Galaxy the light comes from, the more marked is its red shift, and the closer to the edge of the Galaxy the light comes from, the more marked is its blue shift.

When the traveller is at the centre of the Galaxy and flying outwards, for the light coming from the edge of the Galaxy one has $\varphi=\dfrac{\pi}{2}$, $\varphi'=\arctan\left(\dfrac{\sqrt{1-\beta^2}}{-\beta}\right)=171.89\degree$, that is, $\theta=8.11\degree$; hence
\[
\frac{\nu'}{\nu}=\frac{\sqrt{1-(0.99)^2}}{1-0.99\times0.99}=7.089.
\]
From this ratio it is seen that the closer the traveller is to the centre of the Galaxy, the greater is the degree of blue shift of the light from the edge of the Galaxy.

In summary, what the space hero says in the novel written by Hall --- namely his description of flying inside a ``goldfish bowl'' with a blue rim and a red body, towards the mouth of the bowl --- is basically correct, and Feynman lost his bet.
\end{solution}

\begin{solution}{1.69}
(1) Let the proper wavelength of the light emitted by the source be $\lambda_0$. When the source approaches the Earth with velocity $v_1$, the formula for the Doppler effect of light gives the wavelength $\lambda_1$ of the light received by the person $A$ in the $S$ frame, which is at rest relative to the Earth, as
\begin{equation}
\lambda_1=\lambda_0\sqrt{\frac{c-v_1}{c+v_1}}. \tag{1}
\end{equation}
Another inertial frame $S'$ moves relative to the Earth with velocity $v_2$ along the same straight line as $v_1$. Let the velocity of the source observed in the $S'$ frame be $v_1'$; then the velocity transformation formula gives
\begin{equation}
v_1'=\frac{v_1-v_2}{1-\dfrac{v_1v_2}{c^2}}. \tag{2}
\end{equation}
Let the wavelength of the light emitted by the source as seen by the observer $B$ in the $S'$ frame be $\lambda_2$. By the formula for the Doppler effect of light, $\lambda_2$ is
\begin{equation}
\lambda_2=\lambda_0\sqrt{\frac{c-v_1'}{c+v_1'}}. \tag{3}
\end{equation}
When the source recedes from the Earth with the same speed $v_1$, the wavelength of the light seen by $A$ in the $S$ frame is $\lambda_2$, hence
\begin{equation}
\lambda_2=\lambda_0\sqrt{\frac{c+v_1}{c-v_1}}. \tag{4}
\end{equation}
From (1) and (4) we obtain
\begin{equation}
\lambda_0=\sqrt{\lambda_1\lambda_2}=547.7\ (\mathrm{nm}). \tag{5}
\end{equation}
Substituting (5) into (1) gives
\[
\left(\frac{500}{547.7}\right)^2=\frac{c-v_1}{c+v_1}.
\]
Solving, we obtain
\begin{equation}
v_1=\frac{1}{11}c. \tag{6}
\end{equation}
From (3) and (4) we obtain
\begin{equation}
v_1'=-v_1=-\frac{1}{11}c. \tag{7}
\end{equation}
Substituting (6) and (7) into (2) gives
\[
v_2=0.18c.
\]
$v_2$ being positive means that the $S'$ frame moves in the direction approaching the Earth.

(2) When the source recedes from the Earth with speed $v_1$, the speed $v_1''$ of the source as observed by $B$ in the $S'$ frame is
\[
v_1''=\frac{v_1-v_2}{1-\dfrac{v_1v_2}{c^2}}
=\frac{-\frac{1}{11}c-0.18c}{1+\dfrac{0.18}{11}}=-0.267c.
\]
Hence, by the formula for the Doppler effect of light, the wavelength $\lambda_1'$ of the light seen by $B$ is
\[
\lambda_1'=\lambda_0\sqrt{\frac{c-v_1''}{c+v_1''}}
=\lambda_0\sqrt{\frac{1+0.267}{1-0.267}}=720.1\ (\mathrm{nm}).
\]
\end{solution}

\begin{solution}{1.70}
Let the rest mass of this particle be $m_0$, its velocity before decay be $v$, and let the directions of the two $\gamma$ rays, whose energies after the decay are $E_1$ and $E_2$ respectively, make angles $\theta_1$ and $\theta_2$ respectively with the direction of motion of the particle (the $x$ direction), as shown in Fig.~1.70. From the conservation of momentum of the system in the $x$ and $y$ directions we obtain respectively
\begin{equation}
\frac{m_0v}{\sqrt{1-\beta^2}}=\frac{E_1}{c}\cos\theta_1+\frac{E_2}{c}\cos\theta_2, \tag{1}
\end{equation}
\begin{equation}
0=\frac{E_1}{c}\sin\theta_1-\frac{E_2}{c}\sin\theta_2. \tag{2}
\end{equation}

\begin{figbox}{1.70}
\begin{tikzpicture}[scale=0.95,>=latex]
  % particle
  \fill (0,0) circle (1.6pt);
  \draw[force] (-2.1,0)--(-0.15,0) node[midway,above] {$\vec{v}$};
  % gamma rays
  \draw[force] (0,0)--(2.5,1.35) node[right] {$E_1$};
  \draw[force] (0,0)--(2.3,-1.5) node[right] {$E_2$};
  % dashed reference line
  \draw[dashedline] (0,0)--(3.1,0);
  % angles
  \draw[thickline] (1.05,0) arc (0:28.4:1.05);
  \node at (1.28,0.36) {$\theta_1$};
  \draw[thickline] (1.05,0) arc (0:-33.1:1.05);
  \node at (1.28,-0.40) {$\theta_2$};
  % axes
  \draw[axis] (3.6,-0.6)--(3.6,1.7) node[right] {$y$};
  \draw[axis] (3.6,-0.6)--(5.0,-0.6) node[right] {$x$};
\end{tikzpicture}
\end{figbox}

From the statement of the problem,
\begin{equation}
\theta_1+\theta_2=\theta. \tag{3}
\end{equation}
From the conservation of energy of the system before and after the decay of the particle,
\begin{equation}
\frac{m_0c^2}{\sqrt{1-\beta^2}}=E_1+E_2. \tag{4}
\end{equation}
Squaring both sides of (1) and (2) separately and then adding the two equations, we obtain
\[
\frac{m_0^2v^2}{1-v^2/c^2}=\frac{E_1^2}{c^2}+\frac{E_2^2}{c^2}
+2\frac{E_1E_2}{c^2}\left(\cos\theta_1\cos\theta_2-\sin\theta_1\sin\theta_2\right).
\]
Substituting (3) into the above expression gives
\[
\frac{m_0^2c^4}{1-v^2/c^2}\cdot\frac{v^2}{c^4}
=\frac{1}{c^2}\left(E_1^2+E_2^2+2E_1E_2\cos\theta\right).
\]
Substituting (4) into the above expression gives
\[
(E_1+E_2)^2\cdot\frac{v^2}{c^2}=(E_1+E_2)^2-2E_1E_2(1-\cos\theta),
\]
that is,
\[
\sqrt{1-v^2/c^2}=\frac{\sqrt{2E_1E_2(1-\cos\theta)}}{E_1+E_2}.
\]
Substituting this expression into (4), we obtain the rest mass $m_0$ of this particle:
\[
m_0=\frac{1}{c^2}\sqrt{2E_1E_2(1-\cos\theta)}.
\]
\end{solution}

\begin{solution}{1.71}
(1) Let the velocity of the $\pi$ meson be $v$, its rest mass be $m_0=\dfrac{E_{0\pi}}{c^2}=139.6\ \mathrm{MeV}/c^2$, and its momentum $P_\pi$ be
\[
P_\pi=\frac{m_0v}{\sqrt{1-\beta^2}}=\frac{m_0c\beta}{\sqrt{1-\beta^2}}=2\times10^5\ (\mathrm{MeV}/c),
\]
that is,
\[
\frac{\beta}{\sqrt{1-\beta^2}}=\frac{2\times10^5}{139.6}.
\]
We obtain
\[
\beta=0.999\,999\,756.
\]
In the laboratory frame the lifetime $\tau_\pi$ of the $\pi$ meson is
\[
\tau_\pi=\frac{\tau_{\pi0}}{\sqrt{1-\beta^2}}.
\]
Hence the distance $l$ travelled by the $\pi$ meson before it decays is
\begin{align*}
l&=\tau_\pi v=\tau_{\pi0}\cdot c\cdot\frac{\beta}{\sqrt{1-\beta^2}}\\
&=2.6\times10^{-8}\times3\times10^8\times\frac{2\times10^5}{139.6}\\
&=11\,175\ (\mathrm{m}).
\end{align*}

(2) Take the laboratory frame as the $S$ frame and the centre-of-mass frame of the system as the $S'$ frame. In the $S'$ frame the total energy of the system is just the rest energy of the $\pi$ meson, $E'=139.6\ \mathrm{MeV}$. Let the energy of the $\mu$ meson be $E_\mu'$ and the energy of the neutrino be $E_\nu'$; from the conservation of energy,
\begin{equation}
E'=E_\mu'+E_\nu'. \tag{1}
\end{equation}
In the centre-of-mass frame $S'$, let the momentum of the $\mu$ meson be $p'$; from the properties of the centre-of-mass frame, the momentum of the neutrino is $-p'$. For these two particles respectively we have
\begin{equation}
E_\mu'^2=E_{\mu0}^2+p'^2c^2, \tag{2}
\end{equation}
\begin{equation}
E_\nu'=p'c. \tag{3}
\end{equation}
Substituting (2) and (3) into (1) and solving, we obtain
\begin{align*}
p'&=\frac{E'^2-E_{\mu0}^2}{2E'c}=\frac{(139.6)^2-(105.7)^2}{2\times(139.6)}\ (\mathrm{MeV}/c)\\
&=29.783\,9\ (\mathrm{MeV}/c).
\end{align*}
Substituting the above expression into (2) gives
\[
E_\mu'=\sqrt{(105.7)^2+(29.8)^2}=109.816\ (\mathrm{MeV}).
\]
Take the direction of motion of the $\pi$ meson as the positive direction of the $x$-axis, and let $\theta'$ denote the angle between $p'$ and the $x$-axis. From the relativistic momentum transformation formulas, the components of the momentum of the $\mu$ meson along the $x$ and $y$ directions in the $S$ frame are respectively
\begin{equation}
p_{\mu x}=\frac{p_x'+\dfrac{v}{c^2}E_\mu'}{\sqrt{1-\beta^2}}
=\frac{p'\cos\theta'+\dfrac{v}{c^2}E_\mu'}{\sqrt{1-\beta^2}}, \tag{4}
\end{equation}
\begin{equation}
p_{\mu y}=p_y'=p'\sin\theta'. \tag{5}
\end{equation}
Let the angle between the momentum of the $\mu$ meson and the $x$-axis in the $S$ frame be $\theta$; then from (4) and (5) we obtain
\begin{equation}
\tan\theta=\frac{p_{\mu y}}{p_{\mu x}}
=\frac{p'\sin\theta'\sqrt{1-\beta^2}}{p'\cos\theta'+\dfrac{v}{c^2}E_\mu'}. \tag{6}
\end{equation}
Differentiating both sides of (6) with respect to $\theta'$, when $\theta$ takes its maximum value we have $\dfrac{\dd\theta}{\dd\theta'}=0$, and hence
\[
\cos\theta'=-\frac{p'^2c^2}{E_\mu'v}=-0.271\,2,
\]
that is,
\[
\theta'=105.736\degree.
\]
Substituting this value of $\theta'$ into (6), the tangent of the maximum value $\theta_{\max}$ of $\theta$ is found to be
\[
\tan\theta_{\max}=0.000\,196\,6,
\]
that is,
\[
\theta_{\max}=0.011\,267\degree.
\]

(3) The maximum and minimum momenta of the neutrino in the laboratory frame $S$ occur respectively when its momentum in the centre-of-mass frame $S'$ is along the positive and the negative direction of the $x$-axis. From the momentum transformation formulas for relative motion, the maximum momentum $p_{\nu\max}$ and the minimum momentum $p_{\nu\min}$ of the neutrino are respectively
\[
p_{\nu\max}=\frac{p'+\dfrac{v}{c^2}E_\nu'}{\sqrt{1-\beta^2}}
=\frac{p'(1+\beta)}{\sqrt{1-\beta^2}}=85\,358.5\ (\mathrm{MeV}/c),
\]
\[
p_{\nu\min}=\frac{-p'+\dfrac{v}{c^2}E_\nu'}{\sqrt{1-\beta^2}}
=-\frac{p'(1-\beta)}{\sqrt{1-\beta^2}}=-0.010\,413\,7\ (\mathrm{MeV}/c).
\]
\end{solution}

\begin{solution}{1.72}
During its flight from the Earth to the Galaxy the rocket is always in accelerated motion, so if the rocket is taken as the $S'$ frame of reference, this frame is a non-inertial frame. If the special theory of relativity is to be used to solve this kind of problem, one may set up in succession a series of coordinate systems in the $S'$ frame, the coordinate system at any instant $t$ moving uniformly and rectilinearly relative to the Earth with the velocity $v$ of the rocket at that instant during an extremely short interval of time. Hence during this short interval $S'$ may be regarded as an inertial frame, and the special theory of relativity can then be used; the whole accelerating process of the rocket is thus a combination of this series of infinitely short processes of motion.

In the $S'$ frame the rocket moves with acceleration $a'$. In the time $\dd t'$ the increment of the momentum of the rocket is $\dd p'$; according to the laws of dynamics,
\begin{equation}
\frac{\dd p'}{\dd t'}=Ma', \tag{1}
\end{equation}
where $M$ is approximately the rest mass of the rocket at that instant. From the conservation of momentum it is known that the increment of the momentum of the rocket should be equal to the momentum of the photons radiated in the time $\dd t'$; at the same time as the rocket radiates photons, the energy $E$ of the rocket (and of the fuel in it) changes, so according to the mass-energy formula $E=Mc^2$,
\begin{equation}
\dd p'=-\frac{1}{c}\dd E=-c\,\dd M. \tag{2}
\end{equation}
From (1) and (2) we obtain
\[
Ma'=-c\frac{\dd M}{\dd t'}.
\]
The power $N$ of the rocket engine at any moment is
\begin{equation}
N=-\frac{\dd E}{\dd t'}=-c^2\frac{\dd M}{\dd t'}=cMa'. \tag{3}
\end{equation}
Equation (3) shows that the power of the rocket engine depends on the mass of the rocket at that moment; hence one need only find the initial mass $M_0$ of the rocket in order to find the required initial power. Let the rest mass of the rocket at half the journey (that is, at the end of the accelerating process) be $M_h$, its velocity be $v$, and let the total energy of the photons radiated during this half of the journey be $E_r$; then from the conservation of energy and the conservation of momentum we obtain respectively
\begin{equation}
M_0c^2=\frac{M_hc^2}{\sqrt{1-v^2/c^2}}+E_r, \tag{4}
\end{equation}
\begin{equation}
\frac{M_hv}{\sqrt{1-v^2/c^2}}=\frac{E_r}{c}. \tag{5}
\end{equation}
From (4) and (5) we obtain
\begin{equation}
\frac{M_h}{M_0}=\sqrt{\frac{c-v}{c+v}}. \tag{6}
\end{equation}
The second half of the journey may be treated in the same way, giving
\begin{equation}
\frac{M_0'}{M_h}=\sqrt{\frac{c-v}{c+v}}. \tag{7}
\end{equation}
From (6) and (7) we obtain
\begin{equation}
M_0=\frac{c+v}{c-v}M_0'. \tag{8}
\end{equation}
Equation (8) shows that, in order to find the initial mass $M_0$ of the rocket, one must first find the velocity $v$ of the rocket at half the journey. For this purpose, let the velocity of the rocket at time $t$ in the terrestrial frame $S$ be $u$; after a time $\dd t$ the increment of $u$ is $\dd u$, while the corresponding increment in the $S'$ frame is $\dd u'$. Then the velocity transformation formula gives
\[
u+\dd u=\frac{u+\dd u'}{1+\dfrac{u}{c^2}\dd u'},
\]
that is,
\[
(u+\dd u)\left(1+\frac{u}{c^2}\dd u'\right)=u+\dd u'.
\]
Neglecting infinitesimals of higher order in the above expression, we obtain
\begin{equation}
\dd u'=\frac{1}{1-u^2/c^2}\dd u=\frac{1}{1-\beta^2}\dd u. \tag{9}
\end{equation}
Since $\dd t'$ is a proper time, time dilation gives
\begin{equation}
\dd t=\frac{\dd t'}{\sqrt{1-\beta^2}}. \tag{10}
\end{equation}
From (9) and (10) we obtain
\[
a'=\frac{\dd u'}{\dd t'}=\frac{1}{(1-\beta^2)^{3/2}}\frac{\dd u}{\dd t}
=\frac{1}{(1-\beta^2)^{3/2}}a,
\]
where $a=\dfrac{\dd u}{\dd t}$ is the acceleration of the rocket in the $S$ frame. Since $a'$ is constant, $a$ is not constant; hence the above expression may be rewritten as
\[
\frac{\dd u}{\left(1-\dfrac{u^2}{c^2}\right)^{3/2}}=a'\dd t.
\]
Integrating the above expression and using the initial condition $u=0$ at $t=0$, we obtain
\[
\frac{u}{\sqrt{1-u^2/c^2}}=a't,
\]
that is,
\begin{equation}
u=\frac{a't}{\sqrt{1+\dfrac{a'^2}{c^2}t^2}}. \tag{11}
\end{equation}
Since $u=\dfrac{\dd x}{\dd t}$, (11) may be written as
\[
\dd x=\frac{a't}{\sqrt{1+\dfrac{a'^2}{c^2}t^2}}\dd t.
\]
Integrating the above expression and using the initial condition $x=0$ at $t=0$, we obtain
\[
x=\frac{c^2}{a'}\left[\sqrt{1+\frac{a'^2}{c^2}t^2}-1\right].
\]
When the rocket has flown half the journey, $x=\dfrac{R}{2}$; substituting this into the above expression, the terrestrial time $t$ required to complete the half journey is found to be
\begin{equation}
t=\frac{R}{2c}\sqrt{1+\frac{4c^2}{a'R}}. \tag{12}
\end{equation}
Substituting (12) into (11), the velocity $v$ of the rocket at half the journey is obtained as
\[
v=\frac{c}{\sqrt{1+\dfrac{c^2}{a'^2t^2}}}
=c\left[1+\frac{4c^4}{a'^2R^2}\left(1+\frac{4c^2}{a'R}\right)^{-1}\right]^{-1/2}.
\]
Since $\dfrac{c^2}{a'R}\ll1$, neglecting quantities of higher order in the above expression gives
\begin{equation}
v\approx c\left(1-\frac{2c^4}{a'^2R^2}\right). \tag{13}
\end{equation}
Substituting (13) into (8), the rest mass $M_0$ of the rocket at take-off is
\begin{equation}
M_0=\frac{c+v}{c-v}M_0'=\left(\frac{a'^2R^2}{c^4}-1\right)M_0'
\approx\frac{a'^2R^2}{c^4}M_0'. \tag{14}
\end{equation}
Substituting (14) into (3), the power $N$ required by the engine at the moment the rocket starts is
\[
N=cM_0a'=\frac{a'^3R^2}{c^3}M_0'=3.8\times10^{24}\ (\mathrm{J/s}).
\]
\end{solution}

